% !TEX program = pdflatex
\documentclass[12pt]{article}

\usepackage[utf8]{inputenc}
\usepackage{CJKutf8}
\usepackage{amsmath,amssymb}
\usepackage{amsthm}
\usepackage{geometry}
\geometry{
  a4paper,
  top=25mm,
  bottom=25mm,
  left=20mm,
  right=20mm
}
\usepackage{xcolor}
\usepackage{url}
\usepackage[unicode,colorlinks=true,linkcolor=blue,urlcolor=blue]{hyperref}
\usepackage{xurl}
\Urlmuskip=0mu plus 2mu
\sloppy
\usepackage{tocloft}

% 目录中 section 条目使用点线填充到页码
\renewcommand{\cftsecleader}{\cftdotfill{\cftdotsep}}

% 基本定理环境（工作笔记：形式系统风格；参考 NoteBook/notebook1/tex20_pub）
\newtheorem{definition}{定义}[section]
\newtheorem{axiom}{公理}[section]
\newtheorem{assumption}{假设}[section]
\newtheorem{theorem}{定理}[section]
\newtheorem{proposition}{命题}[section]
\newtheorem{corollary}{推论}[section]
\newtheorem{lemma}{引理}[section]
\newtheorem{remark}{备注}[section]
\newtheorem{Certificate}{证书}[section]

% 记号（仅用于本笔记自洽引用，避免依赖主稿宏定义）
\newcommand{\RR}{\mathbb{R}}
\newcommand{\NN}{\mathbb{N}}
\newcommand{\ZZ}{\mathbb{Z}}
\newcommand{\PP}{\mathcal{P}}
\newcommand{\JacGap}{\operatorname{JacGap}}
\newcommand{\indicator}[1]{\mathbf{1}_{#1}}
\newcommand{\code}[1]{\path{#1}}
\newcommand{\GFramework}{\textsf{G-Framework}}
\newcommand{\ORLDTHBase}{\textsf{ODTHB}}

% 避免少量不可控的换行溢出（尤其是混排 CJK 与等宽字体时）
\setlength{\emergencystretch}{6em}

\begin{document}
\begin{CJK*}{UTF8}{gkai}

\title{幂域纤维扩张与 ORL 规范场：从相空间扩张到同调修复与 OHU 正交子丛\\（工作笔记：形式系统版）}
\author{Gao Zheng（高政）}
\date{2026-01-30}
\maketitle

\section*{前言}
\addcontentsline{toc}{section}{前言}

本文的目标不是把幂域、纤维扩张、ORL 规范场、同调修复与量子规范场分别作为
孤立模块处理，而是在 \GFramework{} 的工程--几何接口中，把它们组织成一条从
相空间扩张到规范化修复的形式链条。若这些对象分散出现，幂域容易只被看成集合
扩张，纤维容易只被看成几何背景，ORL 规范场也容易只被看成外加控制项；本文将
它们压合为“如何把状态空间扩展成可调、可修复、可规范化的纤维系统”的问题。

本文的第一层立场是幂域扩张的结构化。幂域不是简单增加候选集合，而是把原始相
空间提升为可承载多路径、多候选、多层调度的纤维对象。通过幂域嵌入，局部状态
不再只是一个点，而可以成为候选族、子丛或可调分支的索引位置。由此，后续同伦
修复、权重态规范化与回路曲率才有可操作的空间。

本文的第二层立场是 ORL 规范场的接口化。ORL 不是替代同调修复的独立算法，而是
把幂扩展纤维中的选择、调度、权重与回路结构写成可规范化的场。有效一阶调度、
\(M\) 阶覆盖细化与闭包趋近共同说明：宏观控制可以先在低阶骨架上运行，再通过
细化与修复逐步逼近更高阶结构。

本文的第三层立场是微观统计与宏观调参的闭合。权重态、回路曲率、同调修复与
量子规范场共同提供一个跨层接口：微观层给出统计权重与局部态，宏观层给出调参
方向与规范化约束，同调层负责检测闭合失败并触发修复。由此，系统不再依赖单次
手工调参，而转向可审计的纤维扩张与规范场更新。

本文的第四层立场是“降级再扩展”的递归口径。系统在高维或高阶结构中直接求解时，
往往会遇到不可控的组合爆炸；本文先将其降级到有效低阶骨架，再通过幂扩展纤维
逐步恢复候选结构。这不是简单丢失信息，而是把不可治理的高阶自由度拆分为可调度、
可覆盖、可闭包逼近的纤维分支。

本文的第五层立场是卡丘空间/流形与量子规范场的接口化。空间或流形在本文中不是
单纯背景，而是承载有效一阶离散骨架和高阶纤维扩张的几何载体。量子规范场则提供
权重态、回路相位、曲率反馈与同调修复之间的统一语言，使微观统计与宏观调参可以
在同一个规范场框架内交换信息。

本文的第六层立场是同调修复对规范场的约束作用。规范场若只调权而不处理闭合失败，
就可能把局部优化误认为全局一致；同调修复负责检查回路、边界和障碍类是否被正确
吸收。由此，ORL 规范场不是任意控制场，而是必须被同调闭合条件约束的可审计场。

本文的第七层立场是为后续策略簇与命中范式提供纤维底座。后续统计同伦命中需要
在候选路径、候选策略和候选残差状态之间进行调制；本文的幂域纤维扩张正提供这种
候选族的几何承载方式。也就是说，命中不是从单点开始，而是在已扩展、已纤维化、
已规范化的候选空间中发生。

因此，本文在整体 \GFramework{} 中承担的是“幂域纤维扩张到 ORL 规范场”的工程
几何桥接作用。它向前连接同伦修复与纤维丛路径语言，向中间连接幂域候选族、
权重态与回路曲率，向后为同伦命中、策略簇与 PDE 残差命中提供可调纤维底座。
概括地说，本文把“如何扩展可控状态空间”推进为“幂域化、纤维化、规范场化并由
同调修复维持闭合”的形式系统。

更具体地说，本文把“多候选”从算法技巧提升为几何对象。候选不只是列表，而是位于
幂域纤维中的局部截面、分支或权重态；选择不只是排序，而是规范场作用下的可审计
变换。这样，策略空间、路径空间和状态空间可以在同一纤维扩张语言中被比较。

从方法论上看，本文为后续策略簇和同伦命中提供“候选空间的容器”。若没有幂域纤维，
命中只能在单点或单轨道上讨论；有了幂域纤维，命中可以发生在候选族、子丛和加权
路径集合上。ORL 规范场则进一步说明如何对这些候选族进行调制、归一和修复。

从整体谱系看，本文把早期同伦修复与后续统计力学命中之间的空间结构补齐。前者
说明障碍可以修复，后者说明目标可以命中；本文说明承载修复和命中的候选空间如何
被扩张、纤维化和规范化。因此，它是 \GFramework{} 工程几何链条中的关键中介。

\setcounter{tocdepth}{3}
\setcounter{secnumdepth}{3}
\tableofcontents
\clearpage

%%%%%%%%%%%%%%%%%%%%%%%%%%%%%%%%%%%%%%%%%%%%%%%%%%%%%%%%%%%%
\section{相空间扩张与幂域嵌入：10 维格点纤维的幂域解释}
\label{sec:tex27-ch1-powerdomain}
%%%%%%%%%%%%%%%%%%%%%%%%%%%%%%%%%%%%%%%%%%%%%%%%%%%%%%%%%%%%

\subsection{形式系统写作约定（公理降级、证书与谓词因果链）}
\label{subsec:tex27-writing-convention}
本笔记参考 \cite{Tex20SemanticFunctorGaugeField} 的工作笔记体例（源文件见
\code{NoteBook/notebook1/tex20_pub/gframework_semantic_functor_field_gauge_field_normalization_interface_formal_system.t
  ex}），
将自然语言叙述压缩为\textbf{可证书化}的形式逻辑条目，并统一使用
\textbf{基于数学符号推演逻辑的谓词因果链}来组织证明要点。为对齐复用接口，本文默认遵循：
\begin{itemize}
  \item \textbf{公理降级}：除不可再降级的最小底座外，正文默认不单列“公理”环境；
  凡需固定为基础条件者，均写成“公理（降级为定理）”，并配套\textbf{证书}与\textbf{证明}。
  若断言包含“对任意阶/任意长度”的全称迭代结构，默认给出\textbf{数学归纳法证书}。
  \item \textbf{定理/引理/命题/推论}：每条编号断言均同时给出（i）证书（可审计见证/构造/不变量/基步-归纳步等），
  （ii）证明（以谓词链闭合方式把推理写成可复核的符号推演步骤）。
  \item \textbf{备注}：仅用于动机、解释、示例与语境对齐，不承担证明义务。
  \item \textbf{证书调用}：对外复用时建议成对引用“断言 + 证书”，以保证他处调用具备可追溯的最小接口。
\end{itemize}

\subsection{定义：自由度空间、格点扩展、相空间与幂域}
\label{subsec:tex27-ch1-defs}

\begin{definition}[任意维度自由度空间]
\label{def:tex27-ch1-theta}
对任意维度 $d\in\NN$，令当前自由度空间为
\[
  \Theta_d \subseteq X^d,
\]
其中 $X$ 是单维自由度的取值集合（可取 $\RR$、离散集合或混合）。
“任意维度”用并对象表示为
\[
  \Theta := \bigsqcup_{d\in\NN}\Theta_d.
\]
\end{definition}

\begin{definition}[10 维格点扩展]
\label{def:tex27-ch1-g10}
取一个离散格点集合（偏移/档位/微状态集合）
\[
  G=\{g_1,\dots,g_K\},
  \qquad
  K\in\NN,\ K\ge 2.
\]
定义 10 维格点空间
\[
  G^{10}:=\underbrace{G\times\cdots\times G}_{10\text{ 次}}.
\]
\end{definition}

\begin{definition}[相空间的“10 维扩张自由度空间”]
\label{def:tex27-ch1-extended-space}
对每个 $d$，定义扩张后的更大自由度空间（相空间扩张）为
\[
  \widetilde{\Theta}_d := \Theta_d \times G^{10}
  \ \subseteq\ X^{d}\times G^{10}.
\]
整体扩张空间为
\[
  \widetilde{\Theta} := \bigsqcup_{d\in\NN} \widetilde{\Theta}_d.
\]
直观上：把原来任意维 $d$ 的状态，统一嵌入到 $(d+10)$ 维的“宏态 + 10 个微自由度”空间。
\end{definition}

\begin{definition}[幂域 / 幂集]
\label{def:tex27-ch1-powerdomain}
最弱版本：幂集
\[
  \PP(Y):=\{A: A\subseteq Y\}.
\]
更贴近“计算/证书化”的版本：有限非空幂域
\[
  \PP_{\mathrm{fin}}^{+}(Y):=\{A\subseteq Y:\ 1\le |A|<\infty\}.
\]
\end{definition}

\subsection{公理组（降级为定理）：纤维扩张结构}
\label{subsec:tex27-ch1-axioms}

\begin{theorem}[公理（降级为定理）A1（相空间扩张是“纤维扩张”而非任意增维）]
\label{thm:tex27-ch1-a1}
存在自然投影
\[
  \pi_d:\widetilde{\Theta}_d\to\Theta_d,\qquad \pi_d(\theta,z)=\theta,
\]
其纤维为
\[
  \pi_d^{-1}(\theta)=\{\theta\}\times G^{10}.
\]
因此“增 10 维”不是随意添加变量，而是把每个宏态 $\theta$ 扩张成一个固定形状的微态集合（纤维）。
\end{theorem}

\begin{Certificate}[证书 A1（笛卡尔积分解）]
\label{cert:tex27-ch1-a1}
定义~\ref{def:tex27-ch1-extended-space} 已给出 $\widetilde{\Theta}_d=\Theta_d\times G^{10}$；
取第一投影即可得到 $\pi_d$，并可直接读出纤维结构。
\end{Certificate}

\begin{proof}[证明（谓词因果链）]
由定义~\ref{def:tex27-ch1-extended-space}，任意元素写为 $(\theta',z')\in\Theta_d\times G^{10}$，
且 $\pi_d(\theta',z')=\theta'$。
固定 $\theta$，则
\[
  \pi_d(\theta',z')=\theta\ \Longleftrightarrow\ \theta'=\theta\ \wedge\ z'\in G^{10},
\]
故 $\pi_d^{-1}(\theta)=\{\theta\}\times G^{10}$。证毕。
\end{proof}

\subsection{定理链：把“纤维扩张”嵌入到幂域结构}
\label{subsec:tex27-ch1-theorems}

\begin{theorem}[定理 1（把“10 维格点扩张”严格嵌入到幂域）]
\label{thm:tex27-ch1-thm1}
定义映射
\[
  \iota_d:\Theta_d\to \PP(\widetilde{\Theta}_d),\qquad
  \iota_d(\theta):=\{\theta\}\times G^{10}.
\]
则 $\iota_d$ 是单射；并且当 $G$ 有限时，
\[
  \iota_d(\Theta_d)\subseteq \PP_{\mathrm{fin}}^{+}(\widetilde{\Theta}_d).
\]
这给出严格含义：把原自由度空间的点提升为扩张相空间中的一个有限集合（纤维），从而自然落在幂域结构里。
\end{theorem}

\begin{Certificate}[证书 1（单射见证与有限性见证）]
\label{cert:tex27-ch1-thm1}
证书由两部分组成：
\begin{enumerate}
  \item 单射见证：用任意 $z\in G^{10}$ 作为“探针”，从 $\{\theta\}\times G^{10}$ 的任意元素读回 $\theta$；
  \item 有限性见证：若 $|G|=K<\infty$，则 $|G^{10}|=K^{10}<\infty$，且 $G^{10}\neq\varnothing$。
\end{enumerate}
\end{Certificate}

\begin{proof}[证明（谓词因果链）]
\begin{itemize}
  \item（单射）若 $\iota_d(\theta_1)=\iota_d(\theta_2)$，则
  $\{\theta_1\}\times G^{10}=\{\theta_2\}\times G^{10}$。
  取任意 $z\in G^{10}$，则 $(\theta_1,z)$ 属于左边也属右边，
  从而 $(\theta_1,z)=(\theta_2,z)$，故 $\theta_1=\theta_2$。
  \item（有限非空）若 $|G|=K<\infty$，则笛卡尔积基数满足 $|G^{10}|=K^{10}<\infty$。
  且由 $K\ge 2$ 知 $G\neq\varnothing$，从而 $G^{10}\neq\varnothing$，
  故对任意 $\theta$，$\iota_d(\theta)=\{\theta\}\times G^{10}$ 属于 $\PP_{\mathrm{fin}}^{+}(\widetilde{\Theta}_d)$。
\end{itemize}
证毕。
\end{proof}

\begin{theorem}[定理 2（“域的幂域”在本构造下的精确解释：纤维集合族是幂域子对象）]
\label{thm:tex27-ch1-thm2}
令集合族
\[
  \mathcal{F}_d := \bigl\{\{\theta\}\times G^{10}:\theta\in\Theta_d\bigr\}.
\]
则
\[
  \mathcal{F}_d = \iota_d(\Theta_d)\ \subseteq\ \PP_{\mathrm{fin}}^{+}(\widetilde{\Theta}_d)\ \subseteq\
    \PP(\widetilde{\Theta}_d).
\]
因此，“10 维格点扩张的相空间”可被严格视为一个“幂域视角”：每个宏态对应一个微态集合（纤维），整体是幂域中的一个可识别子对象。
\end{theorem}

\begin{Certificate}[证书 2（集合族的显式刻画）]
\label{cert:tex27-ch1-thm2}
由定义 $\mathcal{F}_d$ 与定理~\ref{thm:tex27-ch1-thm1} 中的 $\iota_d$ 的像空间完全一致，故可直接作为证书。
\end{Certificate}

\begin{proof}[证明（谓词因果链）]
由 $\mathcal{F}_d$ 的定义可知 $\mathcal{F}_d=\iota_d(\Theta_d)$。
包含关系 $\iota_d(\Theta_d)\subseteq \PP_{\mathrm{fin}}^{+}(\widetilde{\Theta}_d)\subseteq \PP(\widetilde{\Theta}_d)$
由定理~\ref{thm:tex27-ch1-thm1} 与定义~\ref{def:tex27-ch1-powerdomain} 立得。证毕。
\end{proof}

\subsection{引理、命题与推论}
\label{subsec:tex27-ch1-others}

\begin{lemma}[引理 1（维度与基数的严格增长）]
\label{lem:tex27-ch1-lem1}
若 $\Theta_d\subseteq X^d$，则扩张空间 $\widetilde{\Theta}_d=\Theta_d\times G^{10}\subseteq X^d\times G^{10}$ 的自由度个数由 $d$ 增至
  $d+10$。
并且当 $|G|=K$ 有限时，
\[
  |\widetilde{\Theta}_d| = |\Theta_d|\cdot K^{10}.
\]
\end{lemma}

\begin{Certificate}[证书 1（自由度计数与基数乘法法则）]
\label{cert:tex27-ch1-lem1}
证书由两条标准规则给出：
\begin{enumerate}
  \item 笛卡尔积自由度数相加：$X^d\times G^{10}$ 比 $X^d$ 多 10 个独立坐标；
  \item 有限基数乘法：若 $|A|,|B|$ 有限，则 $|A\times B|=|A||B|$，且 $|G^{10}|=K^{10}$。
\end{enumerate}
\end{Certificate}

\begin{proof}[证明（谓词因果链）]
由定义~\ref{def:tex27-ch1-extended-space}，$\widetilde{\Theta}_d=\Theta_d\times G^{10}$，
其坐标由 $(d)$ 个宏观坐标与 10 个微观坐标组成，故自由度个数为 $d+10$。
若 $|G|=K<\infty$，则 $|G^{10}|=K^{10}$，从而
$
|\widetilde{\Theta}_d|
=|\Theta_d\times G^{10}|
=|\Theta_d|\cdot |G^{10}|
=|\Theta_d|\cdot K^{10}.
$
证毕。
\end{proof}

\begin{proposition}[命题 1（“相空间把任意维自由度视为 10 维格点扩张”是一个函子）]
\label{prop:tex27-ch1-prop1}
把“扩张 10 维格点”视为对集合的作用：
\[
  E_{10}(X):=X\times G^{10}.
\]
则对任意映射 $f:\Theta_d\to\Theta_d'$，定义
\[
  E_{10}(f):\Theta_d\times G^{10}\to\Theta_d'\times G^{10},\quad
  E_{10}(f)(\theta,z)=(f(\theta),z),
\]
满足恒等与复合保持：
\[
  E_{10}(\mathrm{id})=\mathrm{id},\qquad
  E_{10}(g\circ f)=E_{10}(g)\circ E_{10}(f).
\]
因此 $E_{10}$ 是一个严格的结构提升算子（端函子）。
\end{proposition}

\begin{Certificate}[证书 1（函子律的逐点验真）]
\label{cert:tex27-ch1-prop1}
证书为两条逐点恒等式验证：
\begin{enumerate}
  \item 对任意 $(\theta,z)$，$E_{10}(\mathrm{id})(\theta,z)=(\mathrm{id}(\theta),z)=(\theta,z)$；
  \item 对任意 $(\theta,z)$，$E_{10}(g\circ f)(\theta,z)=((g\circ f)(\theta),z)=(g(f(\theta)),z)=E_{10}(g)(E_{10}(f)(\theta,
    z))$。
\end{enumerate}
\end{Certificate}

\begin{proof}[证明（谓词因果链）]
按证书~\ref{cert:tex27-ch1-prop1} 的两条逐点计算即可得到恒等与复合保持。证毕。
\end{proof}

\begin{corollary}[推论 1（“幂域”视角的核心：宏态 $\to$ 微态集合）]
\label{cor:tex27-ch1-cor1}
对任意维度 $d$，宏态空间 $\Theta_d$ 可等价地看作幂域中的纤维族 $\mathcal{F}_d$：
\[
  \Theta_d\ \cong\ \mathcal{F}_d\subseteq \PP_{\mathrm{fin}}^{+}(\widetilde{\Theta}_d),
  \qquad
  \theta\leftrightarrow \{\theta\}\times G^{10}.
\]
因此，“把自由度空间视为幂域”在本构造下的严格含义是：把点提升为固定形状集合。
\end{corollary}

\begin{Certificate}[证书 1（显式双射）]
\label{cert:tex27-ch1-cor1}
证书给出一对互逆映射：
\[
  \Theta_d\to\mathcal{F}_d,\ \theta\mapsto \{\theta\}\times G^{10},
  \qquad
  \mathcal{F}_d\to\Theta_d,\ \{\theta\}\times G^{10}\mapsto \theta.
\]
\end{Certificate}

\begin{proof}[证明（谓词因果链）]
由证书~\ref{cert:tex27-ch1-cor1}，两映射互为逆，故给出双射 $\Theta_d\cong \mathcal{F}_d$。
包含关系 $\mathcal{F}_d\subseteq \PP_{\mathrm{fin}}^{+}(\widetilde{\Theta}_d)$ 由定理~\ref{thm:tex27-ch1-thm1} 得证。证毕。
\end{proof}

\subsection{备注：domain-theoretic 幂域、一般化阶数与 ORL 语境对齐}
\label{subsec:tex27-ch1-remarks}

\begin{remark}[关于 domain-theoretic powerdomain（dcpo 语义）]
若把“幂域”理解为域论 power\allowbreak domain（dcpo），可把 $\Theta_d$ 赋予 dcpo（或至少偏序）结构，
并把 $G^{10}$ 赋予离散偏序，则 $\widetilde{\Theta}_d=\Theta_d\times G^{10}$ 仍是 product domain。
进一步可用 Plotkin/\allowbreak Smyth power\allowbreak domain $P(\widetilde{\Theta}_d)$ 表示“非确定性/集合态”\cite{PlotkinPowerdomain,
  SmythPowerdomain}，
而上述 $\iota_d(\theta)=\{\theta\}\times G^{10}$ 落在其“有限集合态”子集中，语义不变。
\end{remark}

\begin{remark}[“10 维”不是本质，只是固定扩张阶数]
上述结论对任意 $n\in\NN$ 成立：把 $10$ 换成 $n$ 即得 $\Theta_d\times G^n$ 与 $\iota_d(\theta)=\{\theta\}\times G^n$。
\end{remark}

\begin{remark}[为什么称为“相空间”扩张]
它把“宏观状态 $\theta$”与“微观隐藏自由度 $z\in G^{10}$”并列成更大的状态；
同时保留投影 $\pi_d(\theta,z)=\theta$ 与纤维 $\{\theta\}\times G^{10}$ 的结构，
可视为离散纤维丛的最小模型。
\end{remark}

\begin{remark}[与 ORL/同伦修复塔语义的对齐提示]
可把 $\pi:\widetilde{\Theta}\to\Theta$ 视为“宏态—微态”纤维丛：
ORL 的边缘算子在 $\widetilde{\Theta}$ 上做离散变分选择，
而证书（如 $\JacGap$ 总量/分量）在 $\Theta$ 上投影评估；
于是“扩张 10 维”等价于“把可选修复动作空间放大 $|G|^{10}$ 而不破坏宏观口径”。
\end{remark}

\clearpage

%%%%%%%%%%%%%%%%%%%%%%%%%%%%%%%%%%%%%%%%%%%%%%%%%%%%%%%%%%%%
\section{递归“降级再扩展”与幂扩展纤维：通式的严格等价}
\label{sec:tex27-ch2-recursive-fiber}
%%%%%%%%%%%%%%%%%%%%%%%%%%%%%%%%%%%%%%%%%%%%%%%%%%%%%%%%%%%%

\subsection{定义：迭代纤维、忘却映射与闭包口径}
\label{subsec:tex27-ch2-defs}

\begin{definition}[$N$ 个格点]
\label{def:tex27-ch2-g}
取有限集合 $G$，满足 $|G|=N\ge 2$。称 $G$ 为“$N$ 个格点”。
\end{definition}

\begin{definition}[第 $m$ 阶纤维：递归“降级再扩展”]
\label{def:tex27-ch2-fm}
令 $F_0:=\{\ast\}$（单点）。对 $m\ge 0$，定义
\[
  F_{m+1}:=F_m\times G.
\]
并把从 $F_m$ 到“1 个格点”的“降级”解释为忘却映射
\[
  q_m:F_m\to\{\ast\},\quad q_m(x)=\ast.
\]
该映射不删除 $x$ 的存在，而是把 $x$ 在宏观口径上视为“一个点”；真实状态仍保留在 $F_m$ 中，
并在下一次扩展时与新格点 $g\in G$ 组成 $(x,g)\in F_{m+1}$。
\end{definition}

\begin{definition}[幂扩展纤维]
\label{def:tex27-ch2-power-fiber}
称 $F_M$ 为“以 $N$ 个格点不断扩张的 $M$ 阶幂扩展纤维”。
\end{definition}

\begin{definition}[内部张开趋近闭包（两种口径）]
\label{def:tex27-ch2-closure}
给出两种可严格化的“闭包”口径：
\begin{enumerate}
  \item 固定阶闭包：$\mathrm{Cl}_M:=F_M$；
  \item 逐层张开闭包：
  \[
    \mathrm{Cl}_{\le M}:=\bigsqcup_{m=0}^{M}F_m.
  \]
  若允许无限迭代，则
  \[
    \mathrm{Cl}_{<\infty}:=\bigsqcup_{m=0}^{\infty}F_m,
    \qquad
    \mathrm{Cl}_{\infty}:=\prod_{m=1}^{\infty}G\ (=G^{\NN}).
  \]
\end{enumerate}
\end{definition}

\subsection{公理组（降级为定理）：降级不是删除}
\label{subsec:tex27-ch2-axiom}

\begin{theorem}[公理（降级为定理）A2（“降级为 1 个格点”不是删除信息）]
\label{thm:tex27-ch2-a2}
“降级为 1 个格点”指忘却映射 $q_m:F_m\to\{\ast\}$ 的宏观口径；
真实状态仍保留为 $F_m$ 的元素，并在下一步扩展形成 $F_{m+1}=F_m\times G$。
\end{theorem}

\begin{Certificate}[证书 A2（规模增长的可检验反例）]
\label{cert:tex27-ch2-a2}
若把“降级”理解为真实删除，则每一步都会抹去历史，永远只剩 $N$ 个格点，无法形成“不断扩张”的结构；
因此“不断扩张到 $M$ 阶并达到 $N^M$ 规模”的目标本身就是该公理的证书。
\end{Certificate}

\begin{proof}[证明（反证；谓词因果链）]
若“降级”为真实删除，则每轮扩展前的状态总是单点 $\{\ast\}$，扩展后总是 $G$；
故任意阶都只有 $N$ 个格点，不可能出现 $N^M$ 规模的“幂扩展”。
这与定义~\ref{def:tex27-ch2-power-fiber} 所要求的“不断扩张到 $M$ 阶”矛盾。故“降级”只能是宏观忘却而非删除。证毕。
\end{proof}

\subsection{引理与定理：递归构造 \texorpdfstring{$\simeq$}{≃} 幂}
\label{subsec:tex27-ch2-lem-thm}

\begin{lemma}[引理 L1（递归构造等价于幂：$F_M\cong G^{M}$）]
\label{lem:tex27-ch2-l1}
对任意 $M\ge 0$，有同构（双射）
\[
  F_M\ \cong\ G^{M}.
\]
\end{lemma}

\begin{Certificate}[数学归纳法证书 L1(M)：幂同构的基步与归纳步]
\label{cert:tex27-ch2-l1}
登记归纳结构：
\[
  P(0)\ \text{成立},\qquad (\forall M\in\NN)\ (P(M)\Rightarrow P(M+1)),
\]
其中 $P(M)$ 表示“存在显式双射 $F_M\cong G^{M}$”。
\end{Certificate}

\begin{proof}[证明（数学归纳法证书；谓词因果链）]
令谓词 $P(M)$ 为：“存在双射 $F_M\cong G^{M}$”。
\begin{itemize}
  \item 基步：$M=0$ 时，$F_0=\{\ast\}\cong G^0$（约定 $G^0$ 为单点），故 $P(0)$ 成立；
  \item 归纳步：假设 $P(M)$ 成立，即 $F_M\cong G^{M}$。则
  \[
    F_{M+1}=F_M\times G\ \cong\ G^{M}\times G\ =\ G^{M+1},
  \]
  同构由归纳假设与笛卡尔积的结合律给出，故 $P(M+1)$ 成立。
\end{itemize}
由数学归纳法，$\forall M\in\NN$，$P(M)$ 成立。证毕。
\end{proof}

\begin{theorem}[定理 T（您的“通式”与“直接幂扩展纤维”严格等价）]
\label{thm:tex27-ch2-t}
在公理~\ref{thm:tex27-ch2-a2} 的语义下，您描述的
\begin{quote}
``$N$ 个格点 $\to$ 降级为 1 个格点 $\to$ 再扩展 $N$ 个格点，反复迭代到 $M$ 阶''
\end{quote}
严格等价于
\[
  F_M=\underbrace{G\times G\times\cdots\times G}_{M\text{ 次}}=G^{M},
\]
即“以 $N$ 个格点不断扩张的 $M$ 阶幂扩展纤维”。
\end{theorem}

\begin{Certificate}[证书 T（递归定义 + 幂同构）]
\label{cert:tex27-ch2-t}
证书由两段数据构成：
\begin{enumerate}
  \item 真实状态空间由定义~\ref{def:tex27-ch2-fm} 的递归构造给出，即 $F_M$；
  \item 引理~\ref{lem:tex27-ch2-l1} 给出 $F_M\cong G^M$。
\end{enumerate}
\end{Certificate}

\begin{proof}[证明（谓词因果链）]
由定义~\ref{def:tex27-ch2-fm}，“迭代降级（忘却）再扩展”得到的真实状态空间就是 $F_M$。
再由引理~\ref{lem:tex27-ch2-l1}，$F_M\cong G^{M}$，故等价成立。证毕。
\end{proof}

\subsection{命题、推论与备注}
\label{subsec:tex27-ch2-others}

\begin{proposition}[命题 P1（与“10 维格点扩展”完全一致）]
\label{prop:tex27-ch2-p1}
若取 $M=10$，并把宏观自由度空间 $\Theta_d$ 作为基底，则扩展空间
\[
  \widetilde{\Theta}_d^{(10)}:=\Theta_d\times F_{10}
\]
满足
\[
  \widetilde{\Theta}_d^{(10)}\ \cong\ \Theta_d\times G^{10}.
\]
\end{proposition}

\begin{Certificate}[证书 P1（取 $M=10$ 并同乘基底）]
\label{cert:tex27-ch2-p1}
由引理~\ref{lem:tex27-ch2-l1} 取 $M=10$ 得 $F_{10}\cong G^{10}$；
两边同乘 $\Theta_d$ 即得结论。
\end{Certificate}

\begin{proof}[证明（谓词因果链）]
由证书~\ref{cert:tex27-ch2-p1} 直接推出同构 $\Theta_d\times F_{10}\cong \Theta_d\times G^{10}$。证毕。
\end{proof}

\begin{corollary}[推论 C1（规模的幂增长）]
\label{cor:tex27-ch2-c1}
若 $|G|=N$，则
\[
  |F_M|=|G^{M}|=N^{M}.
\]
因此“不断扩张的 $M$ 阶幂扩展纤维”在状态数量上是严格的幂增长。
\end{corollary}

\begin{Certificate}[证书 C1（有限集合幂的基数公式）]
\label{cert:tex27-ch2-c1}
证书为标准计数：有限集合 $G$ 的 $M$ 次笛卡尔积满足 $|G^{M}|=|G|^{M}=N^{M}$。
\end{Certificate}

\begin{proof}[证明（谓词因果链）]
由引理~\ref{lem:tex27-ch2-l1}，$F_M\cong G^{M}$，故 $|F_M|=|G^{M}|$。
又由有限集合计数法则，$|G^{M}|=|G|^{M}=N^{M}$。证毕。
\end{proof}

\begin{remark}[若把“降级”为真实删除]
则不会得到 $N^{M}$，而会退化成“每阶都是 $N$”。因此等价成立的关键在于：
降级是\textbf{宏观口径的商化/忘却}，不是抹除。
\end{remark}

\begin{remark}[关于“趋近闭包”的口径选择]
若只讨论固定 $M$，取 $\mathrm{Cl}_M=F_M$ 即足够严格；
若要表达“层层张开直到无限”，则需区分 $\mathrm{Cl}_{<\infty}$ 与 $\mathrm{Cl}_{\infty}$（它们是不同对象），并明确采用哪一个。
\end{remark}

\clearpage

%%%%%%%%%%%%%%%%%%%%%%%%%%%%%%%%%%%%%%%%%%%%%%%%%%%%%%%%%%%%
\section{一阶有效与 \texorpdfstring{$M$}{M} 阶调度：覆盖细化与闭包趋近}
\label{sec:tex27-ch3-first-order-effective}
%%%%%%%%%%%%%%%%%%%%%%%%%%%%%%%%%%%%%%%%%%%%%%%%%%%%%%%%%%%%

\subsection{定义：一阶口径、\texorpdfstring{$M$}{M} 阶扩展与投影}
\label{subsec:tex27-ch3-defs}

\begin{definition}[一阶口径与 $M$ 阶扩展]
\label{def:tex27-ch3-d1}
令“有效自由度（一阶口径）”为
\[
  \Theta^{(1)} := \Theta \times G,
\]
其中 $\Theta$ 是宏观自由度（任意维并对象亦可），$G$ 是 $N$ 个格点（$|G|=N\ge 2$）。

令 $M$ 阶幂扩展纤维为
\[
  F_M := G^M,\qquad \Theta^{(M)} := \Theta \times F_M=\Theta\times G^M.
\]
并定义投影（“降级口径/忘却口径”）
\[
  \pi_{M\to 1}:\Theta^{(M)}\to\Theta^{(1)},\quad
  \pi_{M\to 1}(\theta,g_1,\dots,g_M)=(\theta,g_1).
\]
\end{definition}

\begin{definition}[“刚性/流变”的扇区调度与覆盖]
\label{def:tex27-ch3-d2}
令扇区为覆盖族（允许交集）
\[
  \mathcal{S}=\{S_i\subseteq \Theta^{(M)}\}_{i\in I},
  \qquad
  \bigcup_{i\in I} S_i=\Theta^{(M)}.
\]
\begin{itemize}
  \item \textbf{刚性}：同一扇区内的边/动作集合不变（即时更新权重）；
  \item \textbf{流变}：周期性调整“更偏好进入哪些扇区/哪些跨扇区边”，但仍不必删除边（只调权重）。
\end{itemize}
\end{definition}

\begin{definition}[ORL 学习对象：边权/优先级测度]
\label{def:tex27-ch3-d3}
在任意状态 $u$ 的出边集合 $E(u)$ 上，定义权重场 $w_t(e)$ 与选择测度
\[
  p_t(e\mid u)
  =(1-\varepsilon)\frac{w_t(e)}{\sum_{e'\in E(u)}w_t(e')}
  +\varepsilon\frac{1}{|E(u)|},
  \qquad \varepsilon>0.
\]
并采用“默认权重为 1，窗口内奖励/抑制，窗口外回归 1”的约束：
\[
  w_t(e)=1+\Delta w_t(e),
  \qquad
  e\notin W_t\Rightarrow \Delta w_t(e)=0\Rightarrow w_t(e)=1.
\]
\end{definition}

\subsection{公理（降级为定理）：一阶有效性设计约束}
\label{subsec:tex27-ch3-axiom}

\begin{theorem}[公理（降级为定理）A3（一阶有效性假设：证书与更新仅依赖一阶口径）]
\label{thm:tex27-ch3-a3}
存在风险/目标证书 $J$（例如 $\JacGap$ 总量或其加权版本）与代价项 $C$，
使得 ORL 的更新只依赖
\[
  (J\circ \pi_{M\to 1})(x),\quad (C\circ \pi_{M\to 1})(x)
\]
以及窗口统计；即：高阶坐标 $(g_2,\dots,g_M)$ 不直接进入证书计算，只进入“扇区/纤维调度”。
\end{theorem}

\begin{Certificate}[证书 A3（可核验依赖限制）]
\label{cert:tex27-ch3-a3}
把“真正有效的是 1 阶口径，$M$ 阶用于调度与整体张开闭包趋近”固化为可验真接口：
\[
  \text{证书/代价的观测口径}=\pi_{M\to 1},
  \qquad
  \text{高阶坐标仅进入调度（覆盖选择）层}.
\]
该依赖限制可被实现层直接检查（禁止证书读取 $(g_2,\dots,g_M)$）。
\end{Certificate}

\begin{proof}[证明（降级逻辑；谓词因果链）]
若高阶坐标直接进入证书，则“一阶口径真正有效”的陈述失真；
而“真正有效的是 1 阶口径，$M$ 阶用于调度与整体张开闭包趋近”在形式化上
等价于把证书依赖限制在 $\pi_{M\to 1}$ 上。
因此该假设可作为形式系统的可核验设计约束，并可降级为定理前提。证毕。
\end{proof}

\subsection{定理链：一阶更新不变性与覆盖细化趋近}
\label{subsec:tex27-ch3-theorems}

\begin{theorem}[定理 T1（强结论：在 A3 下，ORL 更新在可观测层面完全由一阶口径决定）]
\label{thm:tex27-ch3-t1}
在公理~\ref{thm:tex27-ch3-a3} 成立时，对任意 $M\ge 1$，
ORL 的权重更新律（即时短窗或周期长窗）在可观测层面满足：
\[
  w_t^{(M)}(e)\ \text{的更新结果仅由}\ \pi_{M\to 1}\ \text{诱导的一阶数据决定}.
\]
因此，“真正有效”在严格意义上成立：高阶仅改变“调度/张开”，不改变“证书驱动的学习核心”。
\end{theorem}

\begin{Certificate}[证书 T1（更新律的观测量因子分解）]
\label{cert:tex27-ch3-t1}
证书要求：把更新律写成
\[
  w_{t+1}^{(M)}=\mathsf{Update}\Big((J\circ\pi_{M\to 1})(\cdot),\ (C\circ\pi_{M\to 1})(\cdot),\ \mathrm{Stats}(W_t)\Big),
\]
从而更新输出只依赖一阶投影的可观测量与窗口统计。
\end{Certificate}

\begin{proof}[证明（谓词因果链）]
\begin{enumerate}
  \item 由公理~\ref{thm:tex27-ch3-a3}，所有进入更新律的量均可写为 $(J\circ\pi_{M\to 1})$ 与 $(C\circ\pi_{M\to 1})$ 的窗口统计；
  \item 窗口统计是对这些量的代数函数（均值、分位数、计数、置信区间等），仍只依赖 $\pi_{M\to 1}$；
  \item 因此更新律输出的 $w_t(e)$ 仅依赖一阶投影数据；高阶坐标不会改变更新的数值结果，只能改变“接下来更可能进入哪些扇区/哪些微状态”，即调度层。
\end{enumerate}
证毕。
\end{proof}

\begin{theorem}[定理 T2（$M$ 阶幂扩展纤维的整体张开闭包趋近：柱集合细化逼近极限纤维）]
\label{thm:tex27-ch3-t2}
令极限纤维
\[
  F_\infty := G^{\NN}=\prod_{m\ge 1}G,
\]
并定义截断投影
\[
  \pi_{\infty\to M}:F_\infty\to F_M,\quad (g_1,g_2,\dots)\mapsto (g_1,\dots,g_M).
\]
则“增加阶数 $M$”等价于对 $F_\infty$ 的开集基进行细化：
任意“只看前 $M$ 坐标”的柱集合
\[
  U_{(a_1,\dots,a_M)}
  :=\{(g_i)\in F_\infty:\ g_1=a_1,\dots,g_M=a_M\}
  \qquad (a_i\in G)
\]
构成一个随 $M$ 递增而越来越细的覆盖体系；因此 $M$ 阶是对“整体张开闭包”的一组截断近似。
\end{theorem}

\begin{Certificate}[证书 T2（柱集合覆盖与细化关系）]
\label{cert:tex27-ch3-t2}
证书由两条可检验事实组成：
\begin{enumerate}
  \item 对任意 $M$，族 $\{U_{(a_1,\dots,a_M)}\}_{(a_1,\dots,a_M)\in G^M}$ 构成对 $F_\infty$ 的划分/覆盖；
  \item 当 $M\leadsto M+1$ 时，每个 $U_{(a_1,\dots,a_M)}$ 被细分为 $N$ 个子柱集合
  $U_{(a_1,\dots,a_M,a_{M+1})}$。
\end{enumerate}
\end{Certificate}

\begin{proof}[证明（谓词因果链）]
\begin{enumerate}
  \item 对任意 $M$，$U_{(a_1,\dots,a_M)}$ 仅由前 $M$ 坐标决定；对任意 $(g_i)\in F_\infty$，其前 $M$ 坐标唯一确定某个 $(a_1,\dots,a_M)$，
  从而 $(g_i)\in U_{(a_1,\dots,a_M)}$，故这些柱集合覆盖 $F_\infty$ 且两两不交；
  \item 当 $M$ 增加到 $M+1$，固定前 $M$ 坐标后，第 $(M+1)$ 坐标可取 $G$ 的任一元素，
  因而 $U_{(a_1,\dots,a_M)}=\bigsqcup_{a_{M+1}\in G}U_{(a_1,\dots,a_M,a_{M+1})}$。
\end{enumerate}
证毕。
\end{proof}

\subsection{命题、推论与备注}
\label{subsec:tex27-ch3-others}

\begin{proposition}[命题 P2（“刚性即时 + 流变周期”在覆盖--细化框架中的作用分工）]
\label{prop:tex27-ch3-p2}
在 $\Theta^{(M)}=\Theta\times G^M$ 上：
\begin{itemize}
  \item 刚性即时迭代主要在固定 $M$ 下，对局部柱集合（或扇区内子图）进行边权调度；
  \item 流变周期迭代主要通过改变“当前更偏好关注哪一类柱集合/哪一类交集扇区”，实现跨扇区的慢尺度迁移；
  \item 二者合起来实现“细化覆盖上的局部优化 + 覆盖层级上的全局调度”。
\end{itemize}
\end{proposition}

\begin{Certificate}[证书 P2（定理 T1 与定理 T2 的角色拆分）]
\label{cert:tex27-ch3-p2}
证书为两条对应关系：
\begin{enumerate}
  \item 由定理~\ref{thm:tex27-ch3-t1}，证书驱动学习核心只看一阶，因此即时更新可视为“固定覆盖单元内的快速调整”；
  \item 由定理~\ref{thm:tex27-ch3-t2}，$M$ 增加对应覆盖细化，周期更新可视为“覆盖单元间的慢尺度选择/迁移”。
\end{enumerate}
\end{Certificate}

\begin{proof}[证明（谓词因果链）]
由定理~\ref{thm:tex27-ch3-t1}，高阶坐标只改变系统在细化覆盖中“落在哪个柱集合/扇区”；
因此即时更新自然对应柱集合内的快速选择，而周期更新对应柱集合间的慢尺度选择。
两者叠加即得到“局部优化 + 全局调度”的分工。证毕。
\end{proof}

\begin{corollary}[推论 C2（一阶有效性的严格等价刻画）]
\label{cor:tex27-ch3-c2}
\[
  \boxed{
  \text{一阶有效}
  \iff
  (J,C)\ \text{仅依赖}\ \pi_{M\to 1}\ \text{且高阶用于扇区/纤维调度}
  }
\]
\end{corollary}

\begin{Certificate}[证书 C2（等价两端的接口同一性）]
\label{cert:tex27-ch3-c2}
证书由公理~\ref{thm:tex27-ch3-a3} 的依赖限制与定理~\ref{thm:tex27-ch3-t1} 的可观测不变性共同给出。
\end{Certificate}

\begin{proof}[证明（谓词因果链）]
“$\Rightarrow$” 方向：一阶有效即意味着证书更新与学习核心不读取高阶坐标，可写成依赖 $\pi_{M\to 1}$ 的形式，得到右端。

“$\Leftarrow$” 方向：若 $(J,C)$ 仅依赖 $\pi_{M\to 1}$ 且高阶只用于调度，则由定理~\ref{thm:tex27-ch3-t1}，
更新在可观测层面由一阶数据决定，即一阶有效。证毕。
\end{proof}

\begin{remark}[复内积空间的“卡丘张开闭包”可严格化为希尔伯特完备化]
若 $(H_0,\langle\cdot,\cdot\rangle)$ 是不完备的内积空间（pre-Hilbert），其完备化为 Hilbert 空间 $H$，满足 $H_0$ 在 $H$ 中稠密。
在“同伦等价”层面：只要它们作为拓扑向量空间，二者都可用
\[
  \mathcal{H}(x,t)=(1-t)x
\]
收缩到 $0$，因此都可缩（contractible），从而同伦等价。
所以“卡丘空间对复内积空间的张开闭包”这一类比可在“完备化/闭包”与“同伦等价（可缩）”两种意义下同时成立。
\end{remark}

\begin{remark}[四维黎曼流形的“卡丘张开闭包”不总是同伦等价，需要额外条件]
给定黎曼流形 $(M,g)$，其距离诱导度量完备化 $\overline{M}$ 是一个完备度量空间；
但一般情况下 $M\hookrightarrow \overline{M}$ 未必同伦等价：例如去掉一点的空间在完备化中加回该点，会改变基本群/高阶同伦群。
因此若要把“卡丘流形对四维黎曼流形的张开闭包”严格化为同伦等价，需要额外给出“边界附加不改变同伦型”的可检验条件，
例如：存在变形收缩 $r:\overline{M}\to M$；或追加部分存在 collar 并可形变收缩回 $M$。
否则该类比应标注为：未证明外延（可作启发但不能当作定理使用）。
\end{remark}

\begin{remark}[结论（严格命题版）]
在“证书更新只依赖一阶投影 $\pi_{M\to 1}$”这一可核验设计约束下：
\[
  \boxed{\begin{aligned}
    &\text{一阶口径负责有效学习（ORL 更新），}\\
    &\text{$M$ 阶幂扩展负责扇区/纤维调度与覆盖细化，}\\
    &\text{从而实现整体张开闭包的截断趋近。}
  \end{aligned}}
\]
\end{remark}

\clearpage

%%%%%%%%%%%%%%%%%%%%%%%%%%%%%%%%%%%%%%%%%%%%%%%%%%%%%%%%%%%%
\section{卡丘空间/流形与量子规范场：有效 1 阶的离散骨架}
\label{sec:tex27-ch4-cauchy-quantum-gauge}
%%%%%%%%%%%%%%%%%%%%%%%%%%%%%%%%%%%%%%%%%%%%%%%%%%%%%%%%%%%%

\subsection{定义：离散骨架、有效 1 阶与量子规范场}
\label{subsec:tex27-ch4-defs}

\begin{definition}[卡丘空间 / 卡丘流形]
\label{def:tex27-ch4-d1}
\begin{itemize}
  \item 卡丘空间：$(X,d)$ 为度量完备（任意卡丘列收敛）；
  \item 卡丘流形：$(M,g)$ 为测地完备（等价于度量完备）。
\end{itemize}
\end{definition}

\begin{definition}[离散等价骨架：$\varepsilon$-网格与邻接图]
\label{def:tex27-ch4-d2}
取尺度 $\varepsilon>0$。令 $V_\varepsilon\subset X$ 为一个 $\varepsilon$-网格（$\varepsilon$-net），满足：
\begin{enumerate}
  \item 覆盖：$\forall x\in X,\ \exists v\in V_\varepsilon,\ d(x,v)\le \varepsilon$；
  \item 分离：$\forall v\neq v'\in V_\varepsilon,\ d(v,v')>\varepsilon$。
\end{enumerate}
构造有向邻接图
\[
  \Gamma_\varepsilon=(V_\varepsilon,E_\varepsilon),\qquad
  (v\to v')\in E_\varepsilon \iff d(v,v')\le c\varepsilon,
\]
其中常数 $c>1$ 固定（用于保证连通与可达性）。
称 $\Gamma_\varepsilon$ 为 $X$ 的离散骨架。（在流形情形，可令边对应最短测地线段。）
\end{definition}

\begin{definition}[有效 1 阶]
\label{def:tex27-ch4-d3}
把“有效 1 阶”严格定义为：只使用一步邻接（边）的自由度与信息，
即只在 $E_\varepsilon$ 上定义与更新对象；更高阶（多步、长程）只通过复合/同伦类出现。
\end{definition}

\begin{definition}[离散规范场：边上的连接]
\label{def:tex27-ch4-d4}
取规范群 $\mathcal{G}$（有限群或紧李群皆可）。
定义离散连接（gauge field）
\[
  U:E_\varepsilon\to \mathcal{G},\qquad e=(v\to v')\mapsto U(e).
\]
规范变换是顶点函数 $g:V_\varepsilon\to \mathcal{G}$，作用为
\[
  U(e:v\to v')\ \mapsto\ U^g(e):=g(v)\,U(e)\,g(v')^{-1}.
\]
这是离散版“联络/规范势”的标准定义：它只占用 1-胞腔（边），因此是“有效 1 阶对象”。
该口径与格点规范理论中“边变量/holonomy”的离散化一致\cite{Wilson1974}。
\end{definition}

\begin{definition}[量子规范场]
\label{def:tex27-ch4-d5}
把“量子”严格化为：在所有离散连接配置空间
\[
  \mathcal{A}_\varepsilon:=\mathcal{G}^{E_\varepsilon}
  \quad(\text{在边集上的群值函数})
\]
上取一个量子态（或概率测度）：
\begin{itemize}
  \item Hilbert 版：$\mathcal{H}_\varepsilon:=L^2(\mathcal{A}_\varepsilon,\mathrm{d}\mu)$，量子规范场为 $\psi\in\mathcal{H}
    _\varepsilon$；
  \item 概率版：量子规范场为 $\mathbb{P}$（在 $\mathcal{A}_\varepsilon$ 上的概率分布）。
\end{itemize}
二者都只依赖边自由度，因此仍是“有效 1 阶”。
\end{definition}

\subsection{公理（降级为定理）：\texorpdfstring{$\varepsilon$}{ε}-网格存在性}
\label{subsec:tex27-ch4-axiom}

\begin{theorem}[公理（降级为定理）A4（卡丘空间可取离散骨架：$\varepsilon$-网格存在性）]
\label{thm:tex27-ch4-a4}
对任意 $\varepsilon>0$，存在 $\varepsilon$-网格 $V_\varepsilon$。
\end{theorem}

\begin{Certificate}[证书 A4（极大 $\varepsilon$-分离集）]
\label{cert:tex27-ch4-a4}
取 $X$ 中一个极大 $\varepsilon$-分离子集 $V_\varepsilon$（任意两点距离 $>\varepsilon$ 且无法再加入新点保持分离性）。
极大性将推出覆盖性：否则若存在 $x\in X$ 与任意 $v\in V_\varepsilon$ 满足 $d(x,v)>\varepsilon$，则可把 $x$ 加入而不破坏分离性，矛盾。
\end{Certificate}

\begin{proof}[证明（构造；谓词因果链）]
考虑所有 $\varepsilon$-分离子集构成的集合族，按包含关系偏序。任意全序链的并仍为 $\varepsilon$-分离子集，
故由 Zorn 引理存在极大元 $V_\varepsilon$。
由证书~\ref{cert:tex27-ch4-a4} 的“极大性 $\Rightarrow$ 覆盖性”推出 $V_\varepsilon$ 同时满足分离与覆盖两条性质，即为 $\varepsilon$-网格。证毕。
\end{proof}

\subsection{定理链：有效 1 阶 \texorpdfstring{$\Longleftrightarrow$}{↔}（量子）离散规范场}
\label{subsec:tex27-ch4-theorems}

\begin{theorem}[定理 1（卡丘空间的有效 1 阶 = 离散骨架上的量子规范场）]
\label{thm:tex27-ch4-thm1}
给定卡丘空间 $(X,d)$ 与尺度 $\varepsilon$。在“有效 1 阶”（只使用一步邻接）口径下，
所有可表达/可学习的局部传输规则都可等价表示为某个离散规范场 $U:E_\varepsilon\to\mathcal{G}$
在配置空间 $\mathcal{A}_\varepsilon=\mathcal{G}^{E_\varepsilon}$ 上的量子态（或分布）$\psi/\mathbb{P}$。
因此：
\[
  \boxed{\text{卡丘空间的有效 1 阶可严格实现为量子规范场（定义~\ref{def:tex27-ch4-d5}）。}}
\]
\end{theorem}

\begin{Certificate}[证书 1（边规则 $\mapsto$ 群值 1-上链；不确定性 $\mapsto$ 态/分布）]
\label{cert:tex27-ch4-thm1}
证书由两段构造组成：
\begin{enumerate}
  \item 在有效 1 阶下，基本自由度为边 $E_\varepsilon$；任一“可组合的局部传输规则”可编码为对每条有向边赋予一个可复合元素，
  即 $U:E_\varepsilon\to\mathcal{G}$；
  \item 当引入学习噪声/统计不确定性时，把确定的 $U$ 提升为配置空间 $\mathcal{A}_\varepsilon$ 上的分布/态（定义~\ref{def:tex27-ch4-d5}）。
\end{enumerate}
\end{Certificate}

\begin{proof}[证明（谓词因果链）]
\begin{enumerate}
  \item 由定义~\ref{def:tex27-ch4-d3}，“有效 1 阶”只允许使用边作为基本邻接自由度；
  \item 在仅有邻接关系的离散骨架上，“局部可组合传输规则”的最一般形式就是：对每条有向边给一个可复合的群元素，并允许在顶点处变换基底，
  这恰是离散连接 $U:E_\varepsilon\to\mathcal{G}$ 与规范作用（定义~\ref{def:tex27-ch4-d4}）；
  \item 一旦引入量子/统计不确定性，只需把 $U$ 的取值从“确定配置”提升为“在 $\mathcal{A}_\varepsilon$ 上的态/分布”（定义~\ref{def:tex27-ch4-d5}）。
\end{enumerate}
故结论成立。证毕。
\end{proof}

\begin{theorem}[定理 2（卡丘流形的有效 1 阶 = 离散等价骨架的有效 1 阶，从而也是量子规范场）]
\label{thm:tex27-ch4-thm2}
设 $(M,g)$ 为卡丘流形（测地完备）。取 $\varepsilon$-网格 $V_\varepsilon\subset M$ 及其邻接图 $\Gamma_\varepsilon$。
则在有效 1 阶口径下，$M$ 的局部一阶结构（沿无穷小位移的并行运输/连接）
与 $\Gamma_\varepsilon$ 的边连接 $U$ 在 $\varepsilon\to 0$ 意义下互相决定（在光滑性假设下）。
因此：
\[
  \boxed{\text{卡丘流形的有效 1 阶等价于其离散等价骨架上的有效 1 阶，从而也是量子规范场。}}
\]
\end{theorem}

\begin{Certificate}[证书 2（连续连接 $\leftrightarrow$ 离散 holonomy 的一阶展开）]
\label{cert:tex27-ch4-thm2}
证书由两段可复算公式给出：
\begin{enumerate}
  \item（流形 $\Rightarrow$ 离散）给定连接 1-形式 $A$，对每条边（测地线段）$e=(v\to v')$ 定义
  \[
    U_A(e):=\mathcal{P}_{\mathrm{ord}}\exp\!\left(\int_{e}A\right)\in\mathcal{G};
  \]
  \item（离散 $\Rightarrow$ 流形）在足够光滑与有界性条件下，对任意 $x\in M$、单位切向量 $v$，
  有一阶展开
  \[
    U_A(v_\varepsilon\to v'_\varepsilon)=\exp\!\big(\varepsilon\,A_x(v)+o(\varepsilon)\big),
  \]
  因而可在选定分支下恢复
  \[
    A_x(v)=\lim_{\varepsilon\to 0}\frac{1}{\varepsilon}\log U_A(v_\varepsilon\to v'_\varepsilon).
  \]
\end{enumerate}
\end{Certificate}

\begin{proof}[证明（分两步；谓词因果链）]
\textbf{(i) 从流形到离散：连接 $\Rightarrow$ 边 holonomy.}
给定主丛连接 $A$，对每条边（取最短测地线段）$e=(v\to v')$，按证书~\ref{cert:tex27-ch4-thm2} 定义 $U_A(e)$，
即得离散连接 $U_A:E_\varepsilon\to \mathcal{G}$。

\textbf{(ii) 从离散到流形：边 holonomy $\Rightarrow$ 一阶连接（$\varepsilon\to 0$）.}
在光滑与局部有界性条件下，沿任意 $x\in M$ 的任意单位切向量 $v$，
取 $v_\varepsilon\in V_\varepsilon$ 逼近 $x$，并取与 $v$ 对齐的邻接点 $v'_\varepsilon$。
由连接的局部展开，边 holonomy 满足
$
U_A(v_\varepsilon\to v'_\varepsilon)=\exp(\varepsilon\,A_x(v)+o(\varepsilon))
$
从而得到恢复公式
$
A_x(v)=\lim_{\varepsilon\to 0}\frac{1}{\varepsilon}\log U_A(v_\varepsilon\to v'_\varepsilon)
$
（在选定局部支与对数分支下）。

因此“有效 1 阶”在流形与离散骨架之间可互相对应；结合定理~\ref{thm:tex27-ch4-thm1}，结论成立。证毕。
\end{proof}

\subsection{推论与备注}
\label{subsec:tex27-ch4-others}

\begin{corollary}[推论 1（打通链条）]
\label{cor:tex27-ch4-cor1}
把定理~\ref{thm:tex27-ch4-thm1} 与定理~\ref{thm:tex27-ch4-thm2} 串起来得到：
\[
  \boxed{\begin{aligned}
    &\text{卡丘空间的有效 1 阶}=\text{量子规范场},\\
    &\text{卡丘流形的有效 1 阶}=\text{离散等价骨架的有效 1 阶},\\
    &\Rightarrow\ \text{卡丘流形的有效 1 阶}=\text{量子规范场}.
  \end{aligned}}
\]
\end{corollary}

\begin{Certificate}[证书 1（链式调用）]
\label{cert:tex27-ch4-cor1}
证书为“定理~\ref{thm:tex27-ch4-thm2} 给出流形 $\Rightarrow$ 骨架的一阶等价；定理~\ref{thm:tex27-ch4-thm1} 给出骨架有效 1 阶 $\Rightarrow$
  量子规范场”的直接复用。
\end{Certificate}

\begin{proof}[证明（谓词因果链）]
由定理~\ref{thm:tex27-ch4-thm2}，卡丘流形的有效 1 阶与其离散骨架的有效 1 阶等价；
再由定理~\ref{thm:tex27-ch4-thm1}，离散骨架的有效 1 阶可实现为量子规范场。
链式合成即得结论。证毕。
\end{proof}

\begin{remark}[必须明确的边界条件]
\begin{enumerate}
  \item “量子”在此是严格定义（定义~\ref{def:tex27-ch4-d5}）：它指“边连接配置空间上的态/分布”，并不自动声称现实物理中的量子场就是该对象；
  \item “流形 $\Leftarrow$ 离散恢复”需要光滑性/有界性：否则离散 holonomy 未必唯一决定连续连接；
  \item 这套证明与“有效 1 阶最关键、$M$ 阶用于调度与闭包趋近”兼容：这里的 1 阶就是边（1-胞腔）上的连接，高阶（多步）只通过路径复合与回路曲率出现。
\end{enumerate}
\end{remark}

\clearpage

%%%%%%%%%%%%%%%%%%%%%%%%%%%%%%%%%%%%%%%%%%%%%%%%%%%%%%%%%%%%
\section{微观统计与宏观调参：权重态的规范化、回路曲率与同调修复}
\label{sec:tex27-ch5-orl-gauge}
%%%%%%%%%%%%%%%%%%%%%%%%%%%%%%%%%%%%%%%%%%%%%%%%%%%%%%%%%%%%

\subsection{定义：图骨架、微观态（边权/连接）与有效 1 阶}
\label{subsec:tex27-ch5-defs}

\begin{definition}[规范场骨架：图与微观态（边权/连接）]
\label{def:tex27-ch5-d1}
令同伦基座的“相空间骨架”为有向图
\[
  \Gamma=(V,E),
\]
其中 $V$ 为状态节点，$E$ 为一步扭曲/切档边。

把 ORL 的边权/优先级权重场视为在边集 $E$ 上的\textbf{微观态（microstate）}：
\[
  w_t:E\to\RR_{>0},\quad e\mapsto w_t(e),
\]
其演化可由任意微观学习/采样规则产生；本节不固定具体更新律，仅要求满足窗口口径：
默认 $1$、窗口内短期偏置、窗口外回归 $1$：
\[
  w_t(e)=1+\Delta w_t(e),\quad e\notin W_t\Rightarrow \Delta w_t(e)=0\Rightarrow w_t(e)=1.
\]
\end{definition}

\begin{remark}[语义澄清：微观统计 $\Rightarrow$ 宏观调参（而非逐边调权）]
本章中，$w_t$（及其等价的边上连接/能量表述）属于\textbf{微观态}；
ODTHB 的职责是从微观采样过程抽取\textbf{微观统计摘要}并据此调节\textbf{宏观算子参数}（如 $\beta,\theta,\ell,\beta(t)$ 等），
而不是把 $w_t(e)$ 当作“逐边控制量”去手工调节（详见第~\ref{subsec:tex27-ch5-odthb}~节）。
更严格地：ODTHB 对微观态与骨架仅有\textbf{只读语义}，禁止对微观进行直接干预（逐边改写 $w_t(e)$、指定连接 $U(e)$、删边/硬剪裁等）；
其影响只能经由\textbf{宏观调参通道}体现。
\end{remark}

\begin{definition}[有效 1 阶：只用边]
\label{def:tex27-ch5-d2}
“有效 1 阶”指：所有可学习对象与证书更新都只依赖边上的数据（一次跳变），高阶路径仅由边复合产生。
\end{definition}

\begin{definition}[选择测度：把权重变成策略]
\label{def:tex27-ch5-d3}
在当前节点 $u\in V$ 的出边集合 $E(u)$ 上，用权重诱导选择测度
\[
  p_t(e\mid u)= (1-\varepsilon)\frac{w_t(e)}{\sum_{e'\in E(u)}w_t(e')}+\varepsilon\frac{1}{|E(u)|},\quad \varepsilon>0.
\]
其中 $\varepsilon>0$ 是“永不硬剪裁”的证书：每条边始终有正概率被试探。
\end{definition}

\subsection{公理（降级为定理）：权重到群元素的可逆绑定}
\label{subsec:tex27-ch5-axiom}

\begin{theorem}[公理（降级为定理）A5（权重到群元素：$U(e)=\exp(-\beta\log w_e)$ 的可逆性）]
\label{thm:tex27-ch5-a5}
取 $\beta>0$，定义
\[
  U_t:E\to \RR_{>0},\qquad U_t(e):=\exp\!\bigl(-\beta\log w_t(e)\bigr)=w_t(e)^{-\beta}.
\]
则映射 $w_t\leftrightarrow U_t$ 在 $\RR_{>0}$ 上双射（可逆）：
\[
  w_t(e)=U_t(e)^{-1/\beta}.
\]
因此可把规范群选为乘法群 $\mathcal{G}_g=(\RR_{>0},\times)$。
\end{theorem}

\begin{Certificate}[证书 A5（$\exp$ 与 $\log$ 的互逆性）]
\label{cert:tex27-ch5-a5}
证书为标准分析事实：$\exp$ 与 $\log$ 在 $\RR_{>0}$ 上互为逆映射，
从而 $U=w^{-\beta}\Leftrightarrow w=U^{-1/\beta}$。
\end{Certificate}

\begin{proof}[证明（谓词因果链）]
\[
\begin{aligned}
  \log U_t(e)
  &=\log\!\bigl(\exp(-\beta\log w_t(e))\bigr)
  =-\beta\log w_t(e),\\
  \log w_t(e)&=-(1/\beta)\log U_t(e),\\
  w_t(e)&=U_t(e)^{-1/\beta}.
\end{aligned}
\]
证毕。
\end{proof}

\subsection{定理链：连接、回路曲率与同调修复}
\label{subsec:tex27-ch5-theorems}

\begin{theorem}[定理 1（微观权重态等价于离散规范连接：有效 1 阶即规范场）]
\label{thm:tex27-ch5-thm1}
令
\[
  A_t(e):=-\beta\log w_t(e)\in\RR,
\]
则
\[
  U_t(e)=\exp(A_t(e)).
\]
于是 $\{A_t(e)\}_{e\in E}$ 是图 $\Gamma$ 上的 1-上链（离散 1-形式），$\{U_t(e)\}$ 是其指数化后的离散连接（规范场）。
由于 $A_t,U_t$ 都只定义在边上，所以它们就是“有效 1 阶”的严格对象。
\end{theorem}

\begin{Certificate}[证书 1（指数化与 1-上链表述）]
\label{cert:tex27-ch5-thm1}
证书为“定义 $A_t(e):=-\beta\log w_t(e)$ 并指数化得到 $U_t(e)=\exp(A_t(e))$”，
从而把正实权重场转写为 $\RR$ 值 1-上链与其对应的乘法群连接。
\end{Certificate}

\begin{proof}[证明（谓词因果链）]
由定义，$A_t$ 是边函数（1-上链），指数化逐点给出 $U_t$，且两者只依赖边集 $E$；
因此它们落在 1-骨架上，即“有效 1 阶对象”。证毕。
\end{proof}

\begin{theorem}[定理 2（回路曲率：$\prod U(e)$ 偏离单位元 $\Leftrightarrow$ 1-上链的周期非零）]
\label{thm:tex27-ch5-thm2}
对任意有向回路（闭路径）
\[
  c=(v_0\to v_1\to\cdots\to v_{m-1}\to v_0),
\]
定义回路 holonomy（回路乘积）
\[
  \mathrm{Hol}_t(c):=\prod_{e\in c}U_t(e)\in\RR_{>0}.
\]
则
\[
  \log \mathrm{Hol}_t(c)=\sum_{e\in c}A_t(e).
\]
特别地：
\[
  \mathrm{Hol}_t(c)=1\quad \Longleftrightarrow\quad \sum_{e\in c}A_t(e)=0.
\]
\end{theorem}

\begin{Certificate}[证书 2（$\log$ 把乘积变为求和）]
\label{cert:tex27-ch5-thm2}
证书为同一条恒等式：
\[
  \log\!\left(\prod_{e\in c}U_t(e)\right)=\sum_{e\in c}\log U_t(e)=\sum_{e\in c}A_t(e),
\]
并用 $\log 1=0$ 给出“单位元 $\Leftrightarrow$ 周期为 0”。
\end{Certificate}

\begin{proof}[证明（谓词因果链）]
\[
  \log\!\left(\prod_{e\in c}U_t(e)\right)
  =\sum_{e\in c}\log U_t(e)
  =\sum_{e\in c}A_t(e).
\]
再由 $\log 1=0$ 得到等价式。证毕。
\end{proof}

\begin{theorem}[定理 3（同调解释：$A_t$ 的同调类由回路周期决定；修复即把同调类推向零）]
\label{thm:tex27-ch5-thm3}
把有向图 $\Gamma$ 视为 1 维 CW 复形，则 $A_t$ 是 1-上链。定义“梯度（精确）”上链：
\[
  (\delta\phi)(v\to v'):=\phi(v')-\phi(v),\quad \phi:V\to\RR.
\]
则：
\begin{enumerate}
  \item $A_t$ 与 $A_t+\delta\phi$ 属于同一同调类，且它们对任意回路 $c$ 的周期相同；
  \item（关键）若对所有回路 $c$ 都有 $\sum_{e\in c}A_t(e)=0$，则存在 $\phi$ 使 $A_t=\delta\phi$（同调类为零）。
\end{enumerate}
等价地：所有回路 holonomy 都为 $1$ 当且仅当 $A_t$ 是精确上链。
\end{theorem}

\begin{Certificate}[证书 3（望远镜求和 + 势函数的构造）]
\label{cert:tex27-ch5-thm3}
证书包含两部分：
\begin{enumerate}
  \item 对任意回路 $c$，有 $\sum_{e\in c}(\delta\phi)(e)=0$（望远镜求和），故加上 $\delta\phi$ 不改变周期；
  \item 若所有回路周期为 0，可选定基点 $v_\ast$ 并对任意 $v$ 取一条从 $v_\ast$ 到 $v$ 的路径 $p$，定义
  $
    \phi(v):=\sum_{e\in p}A_t(e)
  $；
  回路周期为 0 保证该定义与路径选择无关，从而 $A_t=\delta\phi$。
\end{enumerate}
\end{Certificate}

\begin{proof}[证明（谓词因果链）]
\textbf{(1) 周期不变性.}
对任意回路 $c=(v_0\to v_1\to\cdots\to v_0)$，
\[
  \sum_{e\in c}(\delta\phi)(e)
  =\sum_{i=0}^{m-1}\bigl(\phi(v_{i+1})-\phi(v_i)\bigr)
  =0,
\]
故 $A_t$ 与 $A_t+\delta\phi$ 的回路周期相同，因此属于同一同调类。

\textbf{(2) 零周期 $\Rightarrow$ 存在势函数.}
假设对任意回路 $c$，$\sum_{e\in c}A_t(e)=0$。
选定基点 $v_\ast\in V$。对任意 $v\in V$，取一条从 $v_\ast$ 到 $v$ 的有向路径 $p$，定义
\[
  \phi(v):=\sum_{e\in p}A_t(e).
\]
若 $p,p'$ 是两条不同路径，则 $p\cup\overline{p'}$ 构成回路，零周期假设给出 $\sum_{e\in p}A_t(e)=\sum_{e\in p'}A_t(e)$，
故 $\phi(v)$ 与路径选择无关。
对任意边 $e=(v\to v')$，取到 $v$ 的路径并接上 $e$ 得到到 $v'$ 的路径，
则
$
  \phi(v')-\phi(v)=A_t(e)
$
即 $A_t=\delta\phi$。证毕。
\end{proof}

\begin{theorem}[定理 4（JacGap 证书嵌入：JacGap 可取为回路曲率残差范数；ORL 更新即同调修复）]
\label{thm:tex27-ch5-thm4}
取一组回路基（可取生成树外的基本回路）$\{c_i\}_{i=1}^m$。定义分量
\[
  g_i(t):=\big|\log\mathrm{Hol}_t(c_i)\big|
  =\left|\sum_{e\in c_i}A_t(e)\right|.
\]
定义总量证书（点积压缩）
\[
  \JacGap(t):=\langle \mathbf{w},\mathbf{g}(t)\rangle=\sum_{i=1}^m w_i g_i(t),\quad w_i\ge 0.
\]
则：
\begin{enumerate}
  \item $\JacGap(t)=0$ 当且仅当所有基回路 holonomy 为 $1$（等价于同调类为零）；
  \item 若 ORL 的权重更新在窗口内倾向于降低 $\JacGap$，则等价于把 $A_t$ 的同调残差（回路周期）推向 $0$，即“同调修复”。
\end{enumerate}
\end{theorem}

\begin{Certificate}[证书 4（回路分量证书 + 非负加权合成）]
\label{cert:tex27-ch5-thm4}
证书由三条规则给出：
\begin{enumerate}
  \item 由定理~\ref{thm:tex27-ch5-thm2}，$g_i(t)=0\iff \mathrm{Hol}_t(c_i)=1$；
  \item 由 $w_i\ge 0$，$\JacGap(t)=0\iff \forall i,\ g_i(t)=0$；
  \item 基回路周期全为 $0$ 等价于所有回路周期为 $0$，从而由定理~\ref{thm:tex27-ch5-thm3} 得到同调类为零。
\end{enumerate}
\end{Certificate}

\begin{proof}[证明（谓词因果链）]
\begin{enumerate}
  \item 由定理~\ref{thm:tex27-ch5-thm2}，$g_i(t)=0\iff \sum_{e\in c_i}A_t(e)=0\iff \mathrm{Hol}_t(c_i)=1$；
  \item 由 $w_i\ge 0$，$\JacGap(t)=0$ 当且仅当所有 $g_i(t)=0$；
  \item 若所有基回路周期为 $0$，则任意回路周期为它们的线性组合亦为 $0$，由定理~\ref{thm:tex27-ch5-thm3} 得到 $A_t$ 精确（同调类为零）。
\end{enumerate}
至于“更新倾向降低 $\JacGap$”，其等价于降低回路周期残差的范数，从而是把同调残差推向 $0$ 的修复动作。证毕。
\end{proof}

\begin{remark}[建模备注：JacGap 的分量可替换]
此处把 “$\JacGap=$ 回路曲率残差范数” 作为一种可选建模。
若已有更细的 $\JacGap$ 分量定义，可把 $g_i$ 换成既有分量，只要它们仍可解释为某种“回路/障碍周期”即可。
\end{remark}

\subsection{命题、推论与备注}
\label{subsec:tex27-ch5-others}

\begin{proposition}[命题 1（量子规范场的严格化：窗口使 $U_t$ 成为随机场/态）]
\label{prop:tex27-ch5-p1}
由于窗口进出栈与外部噪声，$w_t(e)$ 可视为随机变量（或条件分布下的估计量）。
由定理~\ref{thm:tex27-ch5-a5} 的可逆映射 $U_t(e)=w_t(e)^{-\beta}$，
可把 $U_t$ 视为连接配置空间
\[
  \mathcal{A}:=(\RR_{>0})^{E}
\]
上的随机态。等价的两种严格化为：
\begin{enumerate}
  \item 概率版：量子规范场为 $\mathbb{P}_t$（$\mathcal{A}$ 上的概率测度）；
  \item Hilbert 版：取 $\mathcal{H}=L^2(\mathcal{A},\mathrm{d}\mu)$，量子态为 $\psi_t\in\mathcal{H}$。
\end{enumerate}
\end{proposition}

\begin{Certificate}[证书 1（可逆变换诱导推前测度）]
\label{cert:tex27-ch5-p1}
证书为：逐边可逆变换 $T:(w(e))_{e\in E}\mapsto (w(e)^{-\beta})_{e\in E}$ 在 $\mathcal{A}$ 上给出双射，
从而可把 $w_t$ 的分布通过 $T$ 推前得到 $U_t$ 的分布，即 $\mathbb{P}_t:=T_\#(\mathbb{P}^w_t)$。
\end{Certificate}

\begin{proof}[证明（谓词因果链）]
由定理~\ref{thm:tex27-ch5-a5}，逐点映射 $w(e)\mapsto U(e)=w(e)^{-\beta}$ 在 $\RR_{>0}$ 上可逆，
故其在 $\mathcal{A}=(\RR_{>0})^{E}$ 上逐坐标扩张仍为可逆可测变换 $T$。
若 $w_t$ 是随机变量或分布估计量，则 $U_t=T(w_t)$ 亦为随机变量；
取推前测度即可得到概率版量子态；进一步取 $L^2$ 空间中的密度或波函数表述即得 Hilbert 版。证毕。
\end{proof}

\begin{proposition}[命题 2（刚性/流变双时间尺度 $\Leftrightarrow$ 连接分解 $A=A^{\mathrm{rig}}+A^{\mathrm{rheo}}$）]
\label{prop:tex27-ch5-p2}
若
\[
  w_t(e)=w^{\mathrm{rig}}_t(e)\cdot w^{\mathrm{rheo}}_{k(t)}(e),
\]
则
\[
  A_t(e)=-\beta\log w_t(e)=A^{\mathrm{rig}}_t(e)+A^{\mathrm{rheo}}_{k(t)}(e),
\]
并且回路周期也分解：
\[
  \log\mathrm{Hol}_t(c)=\sum_{e\in c}A^{\mathrm{rig}}_t(e)+\sum_{e\in c}A^{\mathrm{rheo}}_{k(t)}(e).
\]
因此：即时层主要修复“区块内/局部回路”残差，周期层主要修复“跨区块/全局回路”残差。
\end{proposition}

\begin{Certificate}[证书 2（$\log$ 把乘法分解转成加法分解）]
\label{cert:tex27-ch5-p2}
证书为恒等式 $\log(ab)=\log a+\log b$，从而
$-\beta\log(w^{\mathrm{rig}}w^{\mathrm{rheo}})=-\beta\log w^{\mathrm{rig}}-\beta\log w^{\mathrm{rheo}}$，
并把边上加法分解逐回路求和即得周期分解。
\end{Certificate}

\begin{proof}[证明（谓词因果链）]
由 $w_t(e)=w^{\mathrm{rig}}_t(e)\cdot w^{\mathrm{rheo}}_{k(t)}(e)$ 与 $\log$ 的加法性，
有
$
A_t(e)=-\beta\log w_t(e)=-\beta\log w^{\mathrm{rig}}_t(e)-\beta\log w^{\mathrm{rheo}}_{k(t)}(e)
$
定义右两项为 $A^{\mathrm{rig}}_t,A^{\mathrm{rheo}}_{k(t)}$ 即得边上分解。
对回路取和得到周期分解式。证毕。
\end{proof}

\begin{corollary}[推论 1（窗口外回归 1 $\Leftrightarrow$ 回到平直连接）]
\label{cor:tex27-ch5-c1}
窗口不覆盖边 $e$ 时 $w_t(e)=1$，故
\[
  A_t(e)=-\beta\log 1=0,\quad U_t(e)=1.
\]
也即该边回到“平直/无张力”的默认连接，其曲率贡献消失。
\end{corollary}

\begin{Certificate}[证书 1（$\log 1=0$）]
\label{cert:tex27-ch5-c1}
证书为 $\log 1=0$ 与 $U_t(e)=\exp(A_t(e))$ 的直接代入。
\end{Certificate}

\begin{proof}[证明（谓词因果链）]
由 $w_t(e)=1$ 得 $A_t(e)=-\beta\log 1=0$，再由 $U_t(e)=\exp(A_t(e))$ 得 $U_t(e)=1$。证毕。
\end{proof}

\subsection{\ORLDTHBase：ORL 双时间尺度同伦基座（即时 + 周期）}
\label{subsec:tex27-ch5-odthb}

\begin{remark}[定位与抽象性原则]
\ORLDTHBase（\emph{ORL Dual-Timescale Homotopy Base}，ODTHB）用于把“即时层（fast）+ 周期层（slow）”的 ORL 更新组织为
\textbf{可证书化、可持久化、可回放复核}的同伦基座接口。
本节只依赖三类\textbf{抽象数据}（与任何具体项目字段/配置文件/函数名无关）：
\begin{enumerate}
  \item 一个非负风险证书序列 $J(t)\ge 0$（可取为 $\JacGap$ 或其等价口径）；
  \item 一个可持久化的离散同伦事件流（用于构造烈度向量 $\tau$ 与事件口径差分）；
  \item 一个状态--边骨架图 $\Gamma=(V,E)$ 与边代价 $\mathrm{Cost}$（用于即时采样与周期调度；按 ODTHB 严格语义不做删边/硬剪裁）。
\end{enumerate}
\end{remark}

\begin{remark}[ODTHB 的严格语义（微观统计、宏观调节、禁止微观干预）]
在本文中，ODTHB 被严格定义为：
\begin{enumerate}
  \item \textbf{微观统计}：从微观采样过程/事件流中抽取并持久化可审计的统计摘要；
  \item \textbf{宏观调节}：仅通过更新宏观算子参数 $\Lambda=(\beta,\theta,\ell,\ldots)$ 来调节采样核/接受率/退火计划等宏观机制；
  \item \textbf{禁止微观干预}：ODTHB 不允许直接修改微观态（如逐边重写 $w_t(e)$、手工指定连接 $U(e)$），
  也不允许对骨架边集做删边/硬剪裁或强行注入微观跳变；其影响必须经由宏观参数通道体现。
\end{enumerate}
因此本文默认采用 $\varepsilon>0$ 的“永不硬剪裁”探索保底（定义~\ref{def:tex27-ch5-d3}），
并严格避免结构性删边/连通剪裁，因为硬剪裁/删边会破坏微观隧穿（定理~\ref{thm:tex27-ch5-odthb-microtunnel-no-hardprune}）。
\end{remark}

\subsubsection{强命题视角下的定位重构：ODTHB 作为统计力学调节器}
\label{subsubsec:tex27-ch5-odthb-thermostat}

\begin{remark}[与第 8、10 节强命题的对齐]
第~\ref{sec:tex27-ch8-homotopy-statmech}~节把“同伦扭曲演化”解释为统计力学/量子统计力学的自由能变分优化（定义~\ref{def:tex27-ch8-d5}，
定理~\ref{thm:tex27-ch8-t1} 与定理~\ref{thm:tex27-ch8-t3}）；
第~\ref{sec:tex27-ch10-time-space-tradeoff}~节把 ORL 的动作空间严格提升为路径空间 $\Omega=\mathcal{O}^*$ 并给出
“在路径空间上的变分选择”（定理~\ref{thm:tex27-ch10-t2}）。
因此在这些强命题口径下，ODTHB 的更自然定位不是“硬干预算法”，而是维持采样有效混合与可审计口径的
\textbf{调度器/温控器（scheduler/thermostat）}：
快速层从微观采样过程抽取可审计的\textbf{微观统计摘要}并形成局部提议核；
慢速层以这些微观统计为信号，调节\textbf{宏观算子参数}（温度/阈值/退火/步长等），从而影响长期稳态与自由能/风险口径。
其中边权优先级场 $w_t(e)$（定义~\ref{def:tex27-ch5-d1}）被视为采样过程的\textbf{微观态/微观配置}，
并不是 ODTHB 的直接调节对象：ODTHB 不做“逐边手工调权”，而是做“宏观温控/宏观调参”，且按严格语义禁止针对微观态的直接干预（逐边调权/指定连接/删边/硬剪裁）。
\end{remark}

\begin{theorem}[公理（降级为定理）B1（路径集合由可达骨架决定：无硬剪裁下 ODTHB 只改测度）]
\label{thm:tex27-ch5-odthb-support}
在图骨架 $\Gamma=(V,E)$（定义~\ref{def:tex27-ch5-d1}）上，固定起点 $u_0\in V$。
定义长度为 $L\in\NN$ 的\textbf{可达路径集合}为
\[
  \Omega_L(u_0):=\Bigl\{(e_1,\dots,e_L):\ \exists u_1,\dots,u_L\in V,\ e_i=(u_{i-1}\to u_i)\in E\Bigr\}.
\]
若 ODTHB 采用永不硬剪裁的局部选择测度（定义~\ref{def:tex27-ch5-d3}，$\varepsilon>0$），
则对任意 $L\in\NN$ 与任意路径 $\gamma=(e_1,\dots,e_L)\in\Omega_L(u_0)$，
都有 $\mathbb{P}(\gamma\ \text{被逐步采样到})>0$。
因此在不引入硬剪裁的前提下，ODTHB 改变的仅是 $\Omega_L(u_0)$ 上的概率测度（采样核），
而 $\Omega_L(u_0)$ 作为“可达路径集合”仅由骨架 $(V,E)$ 决定。
\end{theorem}

\begin{Certificate}[证书 B1（数学归纳法证书：按路径长度证明全支撑）]
\label{cert:tex27-ch5-odthb-support}
证书为“基步 + 归纳步”：
\begin{enumerate}
  \item（基步 $L=1$）对任意 $e\in E(u_0)$，由 $\varepsilon>0$ 得 $p_t(e\mid u_0)\ge \varepsilon/|E(u_0)|>0$；
  \item（归纳步）若任意 $\gamma_L=(e_1,\dots,e_L)\in\Omega_L(u_0)$ 都有正概率被采样到，
  则对任意 $\gamma_{L+1}=(e_1,\dots,e_L,e_{L+1})\in\Omega_{L+1}(u_0)$，
  其概率为“前缀概率 $\times$ 下一步条件概率”，两者均为正，故仍为正。
\end{enumerate}
\end{Certificate}

\begin{proof}[证明（数学归纳法；谓词因果链）]
对 $L$ 做归纳。
\begin{enumerate}
  \item（基步 $L=1$）任取 $\gamma=(e_1)\in\Omega_1(u_0)$，则 $e_1\in E(u_0)$。
  由定义~\ref{def:tex27-ch5-d3}，$p_t(e_1\mid u_0)=(1-\varepsilon)\frac{w_t(e_1)}{\sum_{e'\in E(u_0)}w_t(e')}
    +\varepsilon\frac{1}{|E(u_0)|}\ge \varepsilon/|E(u_0)|>0$，
  故 $\gamma$ 被采样到的概率 $>0$。
  \item（归纳步）假设对某个 $L$，任意 $\gamma_L=(e_1,\dots,e_L)\in\Omega_L(u_0)$ 都有 $\mathbb{P}(\gamma_L\ \text{被采样到})>0$。
  任取 $\gamma_{L+1}=(e_1,\dots,e_L,e_{L+1})\in\Omega_{L+1}(u_0)$。
  由 $\gamma_{L+1}$ 的可达性，存在中间点 $u_1,\dots,u_{L+1}$ 使得 $e_i=(u_{i-1}\to u_i)\in E$。
  将 $\gamma_{L+1}$ 分解为前缀 $\gamma_L=(e_1,\dots,e_L)\in\Omega_L(u_0)$ 与最后一步 $e_{L+1}\in E(u_L)$。
  由归纳假设有 $\mathbb{P}(\gamma_L\ \text{被采样到})>0$；
  又由定义~\ref{def:tex27-ch5-d3} 的 $\varepsilon>0$，得条件概率 $p_t(e_{L+1}\mid u_L)\ge \varepsilon/|E(u_L)|>0$。
因此
\[
    \mathbb{P}(\gamma_{L+1}\ \text{被采样到})
    =\mathbb{P}(\gamma_L\ \text{被采样到})\cdot p_t(e_{L+1}\mid u_L)\ >0.
  \]
\end{enumerate}
故对任意 $L$ 与任意 $\gamma\in\Omega_L(u_0)$，采样概率均为正。
因此在无硬剪裁且 $\varepsilon>0$ 的前提下，ODTHB 仅改变同一支持集上的\textbf{采样概率测度（采样核参数化）}，
并不改变可达路径集合 $\Omega_L(u_0)$。证毕。
\end{proof}

\begin{definition}[微观隧穿（micro-tunneling）：全支撑试探口径]
\label{def:tex27-ch5-odthb-microtunneling}
沿用定理~\ref{thm:tex27-ch5-odthb-support} 中对可达路径集合 $\Omega_L(u_0)$ 的定义。
给定（可随时间变化的）局部选择测度 $p_t(e\mid u)$。
若对任意 $L\in\NN$ 与任意可达路径 $\gamma\in\Omega_L(u_0)$，
都有 $\mathbb{P}(\gamma\ \text{被逐步采样到})>0$，
则称该采样核满足\textbf{微观隧穿}：
任何可达的“跨势垒”微观跳变序列都不会被硬性置零（不被硬剪裁/删边从支持集中抹去）。
\end{definition}

\begin{theorem}[定理 B2（硬剪裁/删边破坏微观隧穿：支持集收缩）]
\label{thm:tex27-ch5-odthb-microtunnel-no-hardprune}
固定起点 $u_0$ 与骨架 $\Gamma=(V,E)$。
若存在某个时刻 $t$ 与某个\textbf{可达}节点 $u$（即存在 $L$ 与前缀路径 $\gamma_L\in\Omega_L(u_0)$ 使其终点为 $u$），
以及 $u$ 的一条出边 $e\in E(u)$，使得选择测度满足 $p_t(e\mid u)=0$（对该边做硬剪裁/等价删边），
则微观隧穿（定义~\ref{def:tex27-ch5-odthb-microtunneling}）不成立；
更具体地，存在某条可达路径 $\gamma\in\Omega_{L+1}(u_0)$ 使得
\[
  \mathbb{P}(\gamma\ \text{被逐步采样到})=0,
\]
从而可达路径集合的支持集被收缩。
\end{theorem}

\begin{Certificate}[证书 B2（零概率边 $\Rightarrow$ 零概率路径）]
\label{cert:tex27-ch5-odthb-microtunnel-no-hardprune}
取任意到达 $u$ 的可达前缀 $\gamma_L=(e_1,\dots,e_L)\in\Omega_L(u_0)$，
并令 $\gamma:=(e_1,\dots,e_L,e)\in\Omega_{L+1}(u_0)$。
则 $\mathbb{P}(\gamma\ \text{被逐步采样到})$ 的乘积分解包含因子 $p_t(e\mid u)=0$，故整体为 $0$。
\end{Certificate}

\begin{proof}[证明（谓词因果链）]
按证书~\ref{cert:tex27-ch5-odthb-microtunnel-no-hardprune}，
路径逐步采样概率可分解为若干条件概率的乘积，其中包含因子 $p_t(e\mid u)=0$，
故乘积为 $0$。
因此微观隧穿的“全支撑试探”被破坏，支持集收缩。证毕。
\end{proof}

\begin{proposition}[命题 R1（ODTHB 的退化与重构：从“硬干预者”到“微观统计调节框架”）]
\label{prop:tex27-ch5-odthb-r1}
在演化同伦等价口径（定义~\ref{def:tex27-ch8-d5}）与路径空间变分口径（定理~\ref{thm:tex27-ch10-t2}）下，
ODTHB 可等价理解为一个\textbf{双时间尺度统计调节器}：
\begin{enumerate}
  \item 即时层（fast）：从微观采样过程（事件流、局部选择、证书变化等）抽取可审计的微观统计摘要，
  并在给定宏观参数的条件下形成局部提议核/选择测度；其最小接口可取为边选择测度 $p_t(\cdot\mid u)$（定义~\ref{def:tex27-ch5-d3}）。
  \item 周期层（slow）：以微观统计摘要为训练信号/证书输入，调节宏观算子参数（例如逆温度 $\beta$、触发阈值 $\theta$、滞后 $\ell$、
  退火计划 $\beta(t)$、节流/窗口设定、以及步长/投影/clip 等超参数（仅约束\textbf{宏观参数更新幅度}），以维持有效混合并保持可审计的自由能/风险口径。
\end{enumerate}
因此 ODTHB 不应被解释为“计算最优干预动作的硬干预源”；
真正的“干预/扭曲”来源于允许的变换全集 $\mathcal{O}$ 及其复合词空间 $\mathcal{O}^*$（定义~\ref{def:tex27-ch9-d2} 与定义~\ref{def:tex27-ch10-d1}），
而 ODTHB 的职责是：在统计力学口径下用\textbf{微观统计}指导\textbf{宏观调参}，
从而改变采样核/接受率/退火计划等宏观机制，在统计意义上调节系统的变分选择过程。
\end{proposition}

\begin{Certificate}[证书 R1（干预源与调节器的分离）]
\label{cert:tex27-ch5-odthb-r1}
证书由两条可复核接口组成：
\begin{enumerate}
  \item（干预源）路径空间 $\Omega=\mathcal{O}^*$ 与能量 $E(w;s)=\Phi(w(s))$（定义~\ref{def:tex27-ch10-d1}--\ref{def:tex27-ch10-d3}）
    固定了“允许的路径/允许的扭曲”；
  \item（统计调节）边权优先级场 $w_t(e)$（定义~\ref{def:tex27-ch5-d1}）视作微观态，
  其如何随事件流演化属于微观动力学（可由 ORL 的局部规则或采样器自然产生），并非 ODTHB 的直接控制通道。
  ODTHB 的可审计输出是宏观算子参数（$\beta,\theta,\ell,\beta(t)$ 等）及其更新日志，
  它们决定采样核/接受率/门控阈值等宏观机制，从而改变同一支持集上的概率测度；
  在无硬剪裁且 $\varepsilon>0$ 时，其路径支持集保持不变（定理~\ref{thm:tex27-ch5-odthb-support}）。
\end{enumerate}
\end{Certificate}

\begin{proof}[证明（谓词因果链）]
\begin{enumerate}
  \item（路径空间变分）由定理~\ref{thm:tex27-ch10-t2}，当把动作提升为算子词 $w\in\mathcal{O}^*$ 并按 Gibbs/门控选择时，
  ORL 的本质是“在路径空间上做变分选择”；
  \item（支持集由骨架决定）在本节的边骨架口径下，一步扭曲由边 $e\in E$ 表示，ODTHB 只提供边选择测度 $p_t(\cdot\mid u)$。
  若不引入硬剪裁且 $\varepsilon>0$，则由定理~\ref{thm:tex27-ch5-odthb-support}，任意可达路径 $\gamma\in\Omega_L(u_0)$ 均有正概率被采样到，
  因而路径支持集由骨架决定而与权重数值无关；
  \item（调参而非调权）ODTHB 的可审计控制变量是宏观参数（如 $\beta,\theta,\ell,\beta(t)$ 及相关门控/节流设定），
  它们在统计力学口径下改变采样核的参数化（提议核/接受率/退火计划）并影响长期稳态；
  但 ODTHB 不把 $w_t(e)$ 当作“逐边控制量”去手工调节，因此其作用应理解为\textbf{微观统计 $\Rightarrow$ 宏观调参}的温控器行为。
\end{enumerate}
综上，ODTHB 在强命题口径下等价于统计调节器：它改变测度以改善混合与审计口径，而“干预扭曲”的生成闭包来自 $\mathcal{O}^*$。证毕。
\end{proof}

\begin{theorem}[定理 R2（同伦等价“量子”规范场设计 $\cong$ 自由能变分优化）]
\label{thm:tex27-ch5-odthb-r2}
把动作词视为路径（定义~\ref{def:tex27-ch9-d6} 与定义~\ref{def:tex27-ch10-d2}）。
若原子算子非交换且可引入生成元分解（定义~\ref{def:tex27-ch10-d4}），则其 Trotter 极限同伦等价于虚时间演化 $e^{-\tau\widehat{H}}$
（命题~\ref{prop:tex27-ch8-p2} 或定理~\ref{thm:tex27-ch10-t3}）。
因此对“量子”规范场（定义~\ref{def:tex27-ch4-d5}）的可审计设计，可等价转写为：
\textbf{选择缺陷势能 $\Phi$ 与算子全集 $\mathcal{O}$}，
使得系统在自然统计演化下的稳态量子 Gibbs 态 $\rho_\beta$（定义~\ref{def:tex27-ch8-d4}）满足
\[
  \rho_\beta\in \operatorname*{arg\,min}_{\rho}\ \mathcal{F}^Q_\beta[\rho],
\]
即量子自由能泛函的极小（定理~\ref{thm:tex27-ch8-t3}）。
在该转写下，ODTHB 的职责仅为维持采样核的有效混合并调节逆温度/退火计划，
而非把每条边的连接 $U(e)$ 视作需要手工指定的“硬干预输入”。
\end{theorem}

\begin{Certificate}[证书 R2（Trotter 乘积 + 量子自由能变分极小）]
\label{cert:tex27-ch5-odthb-r2}
证书由两条并列引用给出：
\begin{enumerate}
  \item（Trotter）非交换时间有序复合在极限下逼近 $e^{-\tau\widehat{H}}$（命题~\ref{prop:tex27-ch8-p2} 或定理~\ref{thm:tex27-ch10-t3}）；
  \item（变分）量子 Gibbs 态是量子自由能的唯一极小点（定理~\ref{thm:tex27-ch8-t3}）。
\end{enumerate}
\end{Certificate}

\begin{proof}[证明（谓词因果链）]
按证书~\ref{cert:tex27-ch5-odthb-r2}：
\begin{enumerate}
  \item 由 Trotter 乘积（命题~\ref{prop:tex27-ch8-p2} 或定理~\ref{thm:tex27-ch10-t3}），非交换的算子序列在极限下可用虚时间演化 $e^{-\tau\widehat{H}}$
    描述；
  \item 由定理~\ref{thm:tex27-ch8-t3}，该演化的稳态（量子 Gibbs）等价于量子自由能泛函的极小点；
  \item 因而把“设计量子规范场”转写为“设计 $\Phi$ 与 $\mathcal{O}$（从而设计 $\widehat{H}$）并维持有效混合/退火”，即可把目标固化为自由能变分极小口径。
\end{enumerate}
从而得到结论。证毕。
\end{proof}

\begin{corollary}[推论 R3（从算法控制到自然演化：ODTHB 退居温控/调度）]
\label{cor:tex27-ch5-odthb-r3}
在算子完备性证书输入（公理~\ref{thm:tex27-ch9-a9} 的 $\mathrm{Complete}$）与统计力学采样口径（引理~\ref{lem:tex27-ch8-l1}--\ref{lem:tex27-ch8-l2}）下，

干预不再是“算法强制控制”，而可解释为“在路径空间上的自然统计演化”：
\begin{enumerate}
  \item 环境/对手通过缺陷势能 $\Phi$ 指定能量景观；
  \item 算子全集 $\mathcal{O}$（及其复合 $\mathcal{O}^*$）给出允许的跃迁；
  \item ODTHB 只需调节逆温度 $\beta$ 与退火计划，并维持局部提议核的有效混合；
  \item 系统通过 Metropolis/Langevin 型采样自然趋向自由能极小的修复路径分布。
\end{enumerate}
\end{corollary}

\begin{Certificate}[证书 R3（完备性 + 采样稳态 + 变分口径并列）]
\label{cert:tex27-ch5-odthb-r3}
证书由三条并列引用给出：
\begin{enumerate}
  \item（存在可下降修复词）公理~\ref{thm:tex27-ch9-a9} 给出对任意障碍态存在下降修复词；
  \item（稳态为 Gibbs）引理~\ref{lem:tex27-ch8-l2} 给出含噪/接受率的同伦扭曲在平衡时为 Gibbs；
  \item（自由能解释）定理~\ref{thm:tex27-ch8-t1}（经典）与定理~\ref{thm:tex27-ch8-t3}（量子）把 Gibbs 稳态解释为自由能极小。
\end{enumerate}
\end{Certificate}

\begin{proof}[证明（谓词因果链）]
按证书~\ref{cert:tex27-ch5-odthb-r3}：
\begin{enumerate}
  \item 公理~\ref{thm:tex27-ch9-a9} 保证“存在可下降修复路径/修复词”，给出自然演化的可达性底座；
  \item 引理~\ref{lem:tex27-ch8-l2} 给出在有限温统计采样下的 Gibbs 稳态；
  \item 定理~\ref{thm:tex27-ch8-t1} 与定理~\ref{thm:tex27-ch8-t3} 把该稳态解释为自由能极小点，从而“长期趋稳”可解释为“自由能意义的自然优化”。
\end{enumerate}
因此 ODTHB 若仅负责调温/退火与混合调度，系统也可在统计意义下完成修复路径选择。证毕。
\end{proof}

\begin{remark}[工程备注 R1（实现侧的简化与审计接口）]
该视角直接给出“工程最小闭合面”：
\begin{enumerate}
  \item 完备化算子库：把 $\mathcal{O}$ 的 $\mathrm{Complete}$ 由假设提升为可核验证书输入（证书~\ref{cert:tex27-ch9-a9} 的三类之一）；
  \item 精确化势能/证书：让 $\Phi$ 或 $J=\JacGap$ 真正反映障碍强度与代价结构（避免温控调度失真）；
  \item 持久化统计量：落盘 $(\Delta J,\Delta\tau,\beta,\theta,\ell)$ 等慢变量字段，以支持回放审计与离线再训练。
\end{enumerate}
在此框架下，“每一步为什么这样走”可由能量/概率/证书字段解释，从而更接近白盒可审计实现。
\end{remark}

\subsubsection{对象域：证书、烈度与事件口径}
\label{subsubsec:tex27-ch5-odthb-objects}

\begin{definition}[观测证书：总 gap 与分量 gap]
\label{def:tex27-ch5-odthb-j}
沿用定理~\ref{thm:tex27-ch5-thm4} 的口径：令分量向量 $\mathbf{g}(t)=(g_1(t),\dots,g_m(t))$，
并令总量证书为点积压缩
\[
  J(t):=\JacGap(t)=\langle \mathbf{w},\mathbf{g}(t)\rangle=\sum_{i=1}^m w_i g_i(t),\qquad w_i\ge 0.
\]
\end{definition}

\begin{definition}[同伦烈度向量（事件口径）]
\label{def:tex27-ch5-odthb-tau}
固定一个“同伦事件类型”集合 $\mathcal{K}=\{1,\dots,K\}$（例如：不同子规范场/不同修复模块的离散跳变类型；此处只要求可枚举）。
对每个类型 $k\in\mathcal{K}$，令 $N_k(t)\in\ZZ_{\ge 0}$ 表示截至时刻 $t$ 的\textbf{事件口径计数}（只在真实离散跳变发生时递增）。
定义区间 $[t_0,t_1]$ 的同伦烈度差分向量为
\[
  \tau(t_0\to t_1):=\bigl(N_1(t_1)-N_1(t_0),\dots,N_K(t_1)-N_K(t_0)\bigr)\in\ZZ_{\ge 0}^{K}.
\]
并定义加权范数
\[
  |\tau|_w:=\langle \boldsymbol{\omega},\tau\rangle=\sum_{k=1}^K \omega_k \tau_k,\qquad \omega_k\ge 0.
\]
\end{definition}

\begin{remark}[事件口径 vs 调用口径]
为避免把“每步调用”误当作“离散跳变”，ODTHB 默认只使用事件口径 $N_k(t)$ 构造 $\tau$；
即：\textbf{只有当系统状态确实发生离散同伦跳变}时才计数。
这条约束使得 $\tau$ 可作为可审计的“修复强度/扭曲强度”证书，而不是实现细节计数。
\end{remark}

\begin{definition}[强制修复窗口触发器（阈值 + 节流/滞回）]
\label{def:tex27-ch5-odthb-trigger}
给定阈值参数 $\theta$ 与最小节流间隔 $\Delta_{\min}>0$。
令 $t_{\mathrm{last}}$ 为上一次触发时刻（若从未触发则记为 $-\infty$）。
定义触发谓词（示意）：
\[
  \mathrm{fire}(t):=\indicator{J(t)>\theta}\cdot \indicator{t-t_{\mathrm{last}}\ge \Delta_{\min}}.
\]
当 $\mathrm{fire}(t)=1$ 时，系统进入一次“强制修复窗口”，并将该时刻登记为事件时间 $t_k$。
（$\Delta_{\min}$ 可理解为冷却时间/最小跨度/滞回门槛；其具体实现不影响本文形式推演。）
\end{definition}

\begin{definition}[每次窗口的可持久化样本对：$(\Delta\tau,\Delta J)$（允许滞后）]
\label{def:tex27-ch5-odthb-samples}
对事件时间序列 $\{t_k\}$，定义窗口长度 $\Delta t>0$ 与滞后 $\ell\ge 0$。
对每次触发 $k$，构造样本对：
\[
  \Delta\tau_k:=\tau(t_k\to t_k+\Delta t),\qquad
  \Delta J_k:=J(t_k+\ell)-J(t_k).
\]
其中 $\ell$ 用于表达“修复烈度对证书的作用存在滞后”这一常见现实口径。
\end{definition}

\subsubsection{即时层：微观统计摘要与局部提议核（只读微观态）}
\label{subsubsec:tex27-ch5-odthb-fast}

\begin{definition}[持久化事件流与滑动窗口]
\label{def:tex27-ch5-odthb-window}
令持久化事件流为 $\mathcal{H}=\{h_n\}_{n\in\ZZ}$，其中每个事件
\[
  h_n=(u_n,e_n,\Delta J_n,\mathrm{Cost}_n,\ldots),
  \qquad
  u_n\in V,\ e_n\in E(u_n)\subseteq E.
\]
窗口 $W_n$ 可取为“最近 $N$ 次事件”或“最近 $T$ 时间内事件”的集合。
对每条边 $e$，定义窗口覆盖计数
\[
  n_e(n):=\#\{h\in W_n:\ h\text{ 的边为 }e\}.
\]
\end{definition}

\begin{definition}[窗口微观统计摘要（用于宏观调参，而非逐边调权）]
\label{def:tex27-ch5-odthb-microstat}
对每条边 $e$，用窗口 $W_n$ 定义边条件样本数 $n_e(n)$（定义~\ref{def:tex27-ch5-odthb-window}），并定义窗口内的条件均值：
\[
  \overline{\Delta J}_e(n):=
  \begin{cases}
    \mathbb{E}_{W_n}[\Delta J\mid e], & n_e(n)>0,\\
    0, & n_e(n)=0,
  \end{cases}
  \qquad
  \overline{c}_e(n):=
  \begin{cases}
    \mathbb{E}_{W_n}[\mathrm{Cost}\mid e], & n_e(n)>0,\\
    0, & n_e(n)=0.
  \end{cases}
\]
并用收缩因子刻画“统计置信度”（样本少则靠近 $0$）：
\[
  \rho_e(n):=\frac{n_e(n)}{n_e(n)+\alpha}\in[0,1),\qquad \alpha>0.
\]
把
\[
  \mathsf{Stat}_e(n):=\bigl(n_e(n),\rho_e(n),\overline{\Delta J}_e(n),\overline{c}_e(n)\bigr)
\]
称为边 $e$ 的窗口微观统计摘要；把 $\mathsf{Stat}(W_n):=\{\mathsf{Stat}_e(n)\}_{e\in E}$ 称为窗口微观统计摘要族。
该摘要的用途是为 ODTHB 的慢时标宏观调参提供可审计输入，而不是把 $w_t(e)$ 作为“逐边控制量”去手工调节。
\end{definition}

\begin{definition}[宏观算子参数（慢变量；ODTHB 的直接调节对象）]
\label{def:tex27-ch5-odthb-macroparam}
把 ODTHB 的宏观调节变量抽象为参数向量
\[
  \Lambda:=\bigl(\beta,\theta,\ell,\Delta t,\Delta_{\min},\varepsilon,\ldots\bigr),
\]
其中 $\beta$ 为逆温度/退火参数，$\theta$ 为门控阈值，$\ell$ 为滞后，$(\Delta t,\Delta_{\min})$ 为窗口与节流设定，$\varepsilon>0$ 为“永不硬剪裁”的探索保底。
符号 ``$\ldots$'' 允许加入与算子库实现相关的宏观超参数（例如步长/投影/clip/退火速率等），但这些参数均属于\textbf{宏观调参通道}（仅约束宏观更新幅度而非干预微观态）。
\end{definition}

\begin{theorem}[定理（窗口外默认中性：无样本边不携带永久统计偏置）]
\label{thm:tex27-ch5-odthb-reset}
对任意边 $e$，若从某个时刻起该边不再出现在窗口中，即 $\forall n\ge n_0,\ n_e(n)=0$，
则对所有 $n\ge n_0$，有 $\rho_e(n)=0$ 且 $\mathsf{Stat}_e(n)=(0,0,0,0)$。
因此该边在微观统计摘要层面回到“默认中性”，不会对宏观调参造成永久偏置。
\end{theorem}

\begin{Certificate}[证书（由定义直接核验）]
\label{cert:tex27-ch5-odthb-reset}
由定义~\ref{def:tex27-ch5-odthb-microstat}，当 $n_e(n)=0$ 时按分段定义有 $\overline{\Delta J}_e(n)=\overline{c}_e(n)=0$，
且 $\rho_e(n)=n_e(n)/(n_e(n)+\alpha)=0$；代入 $\mathsf{Stat}_e(n)$ 即得结论。
\end{Certificate}

\begin{proof}[证明（谓词因果链）]
对任意 $n\ge n_0$，由假设 $n_e(n)=0$，套用证书~\ref{cert:tex27-ch5-odthb-reset} 的逐点代入即可得到
$\rho_e(n)=0$ 且 $\mathsf{Stat}_e(n)=(0,0,0,0)$。证毕。
\end{proof}

\begin{theorem}[定理（有限记忆：微观统计摘要只依赖窗口）]
\label{thm:tex27-ch5-odthb-finite-memory}
对任意时刻 $n$，微观统计摘要族 $\mathsf{Stat}(W_n)$ 只依赖窗口 $W_n$，与窗口外更早历史无关。
\end{theorem}

\begin{Certificate}[证书（窗口内样本同一 $\Rightarrow$ 统计摘要同一）]
\label{cert:tex27-ch5-odthb-finite-memory}
若两条历史在窗口 $W_n$ 内完全相同，则对每条边 $e$ 的窗口统计量 $n_e(n)$、$\overline{\Delta J}_e(n)$、$\overline{c}_e(n)$ 相同，
从而由定义~\ref{def:tex27-ch5-odthb-microstat} 得 $\mathsf{Stat}_e(n)$ 相同，进而 $\mathsf{Stat}(W_n)$ 相同。
\end{Certificate}

\begin{proof}[证明（谓词因果链）]
由定义~\ref{def:tex27-ch5-odthb-microstat}，$\mathsf{Stat}(W_n)$ 是窗口样本集合 $W_n$ 的确定函数。
当窗口内事件集合相同，所有边条件统计量一致，故统计摘要族一致。证毕。
\end{proof}

\subsubsection{周期层：动态阈值优化（慢时标外环）}
\label{subsubsec:tex27-ch5-odthb-slow-threshold}

\begin{theorem}[公理（降级为定理）ODTHB-A（经验控制关系：烈度可在滞后后压低 gap）]
\label{thm:tex27-ch5-odthb-empirical}
存在权重 $\boldsymbol{\omega}\ge 0$、滞后 $\ell\ge 0$ 与噪声项 $\epsilon$，
使得在统计意义下：当区间烈度 $|\tau(t_0\to t_1)|_{\boldsymbol{\omega}}$ 足够大时，后续时刻的 $J$ 出现可观测回落/受控，例如
\[
  \mathbb{E}[J(t_1+\ell)]\ \le\ \mathbb{E}[J(t_0)] + \epsilon.
\]
\end{theorem}

\begin{Certificate}[证书（回放可复核：枚举 $(\boldsymbol{\omega},\ell)$ 并检验不等式）]
\label{cert:tex27-ch5-odthb-empirical}
对任意一段可回放数据，逐段计算 $\tau(t_0\to t_1)$ 与 $J(t)$ 序列；
枚举候选 $(\boldsymbol{\omega},\ell)$，在选定置信水平下检验上式是否成立。
若成立，则获得该条“数据驱动定理”的证书化通过。
\end{Certificate}

\begin{proof}[证明（证书化验证流程，非解析推演）]
按证书~\ref{cert:tex27-ch5-odthb-empirical} 给出的可复核流程执行即可；其输出为“成立/不成立”的可审计结论。
因此该公理可降级为数据驱动定理。证毕。
\end{proof}

\begin{lemma}[引理（阈值对触发频率/烈度的单调性：固定轨迹下）]
\label{lem:tex27-ch5-odthb-threshold-mono}
在固定同一条 $J(t)$ 序列与同一条节流规则（定义~\ref{def:tex27-ch5-odthb-trigger}）下，
把阈值从 $\theta_1$ 降到 $\theta_2\le \theta_1$，则满足 $J(t)>\theta$ 的时刻集合只会扩大，
因此强制修复窗口的候选触发次数弱单调不减，从而期望烈度 $\mathbb{E}|\Delta\tau|_w$ 弱单调不减。
\end{lemma}

\begin{Certificate}[证书（谓词包含 + 节流筛选）]
\label{cert:tex27-ch5-odthb-threshold-mono}
证书为集合包含关系：
\[
  \{t:J(t)>\theta_1\}\subseteq \{t:J(t)>\theta_2\},
\]
而节流谓词仅是在该集合上进一步筛选，故触发点不会因降低阈值而减少。
\end{Certificate}

\begin{proof}[证明（谓词因果链）]
对任意 $t$，若 $J(t)>\theta_1$ 则必有 $J(t)>\theta_2$，故上述集合包含成立。
节流条件 $\indicator{t-t_{\mathrm{last}}\ge \Delta_{\min}}$ 只会在候选集合上做删除，不会产生新触发点，
因此降低阈值只可能增加（或不变）触发次数。由烈度向量分量非负，得期望烈度弱单调不减。证毕。
\end{proof}

\begin{proposition}[命题（动态阈值优化：用 $(\Delta J,\Delta\tau)$ 构造白盒训练信号）]
\label{prop:tex27-ch5-odthb-threshold-update}
取权衡系数 $\lambda\ge 0$，定义窗口级目标（示意）：
\[
  F(\theta):=\mathbb{E}[J(t+\ell)] + \lambda\,\mathbb{E}\bigl[|\tau(t\to t+\Delta t)|_w\bigr].
\]
对样本对（定义~\ref{def:tex27-ch5-odthb-samples}），可用符号信号
\[
  s_k:=\operatorname{sign}\!\Big(\Delta J_k + \lambda|\Delta\tau_k|_w\Big)
\]
构造慢时标更新（示意）：
\[
  \theta_{k+1}:=\operatorname{clip}_{[\theta_{\min},\theta_{\max}]}\Big(\theta_k + \eta_k\cdot s_k\Big),
  \qquad \eta_k>0.
\]
直观含义：若“gap 没降反升或降得不值”（$\Delta J+\lambda|\Delta\tau|>0$），则提高阈值（更难触发，减少无效烈度）；
反之降低阈值（更早触发，增加修复烈度）。
\end{proposition}

\begin{Certificate}[证书（单调性 + 经验控制关系 $\Rightarrow$ 权衡结构）]
\label{cert:tex27-ch5-odthb-threshold-update}
证书由两条事实给出：
\begin{enumerate}
  \item 引理~\ref{lem:tex27-ch5-odthb-threshold-mono}：降低阈值会（弱）增加烈度；提高阈值会（弱）减少烈度；
  \item 定理~\ref{thm:tex27-ch5-odthb-empirical}：当烈度足够大时，$J$ 在滞后后呈现统计意义的回落/受控。
\end{enumerate}
因此 $F(\theta)$ 呈现“$\theta$ 太小 $\Rightarrow$ 烈度成本过高；$\theta$ 太大 $\Rightarrow$ gap 不受控”的权衡结构。
\end{Certificate}

\begin{proof}[证明（谓词因果链；作为白盒规则的可复核性）]
由证书~\ref{cert:tex27-ch5-odthb-threshold-update}，$\theta$ 对烈度与 gap 的作用方向相反（一个增则一个减），故存在权衡。
更新式使用 $\Delta J_k+\lambda|\Delta\tau_k|_w$ 的符号作为方向信号，
可视作对 $F$ 的随机次梯度符号进行“根寻找/调参”。
严格收敛性需要额外噪声与光滑性假设；在不引入这些假设时，本命题只主张：
\textbf{该更新是白盒、可落盘、可回放复核的动态阈值优化规则}。证毕。
\end{proof}

\begin{corollary}[推论（即时 + 周期：在线/离线两种时标的统一解释）]
\label{cor:tex27-ch5-odthb-online-offline}
若 $\theta_{k+1}$ 在同一次运行内随 $k$ 更新，则 $\theta$ 成为运行时慢变量，属于在线双时间尺度控制；
若 $\theta$ 只在段与段之间更新并在段内冻结，则退化为离线外环调参。
两者都可由 $(J,\tau,\theta)$ 的持久化日志审计与回放复核。
\end{corollary}

\begin{Certificate}[证书（按更新发生的时标分类）]
\label{cert:tex27-ch5-odthb-online-offline}
证书为定义性区分：是否允许运行内改变阈值参数 $\theta$，以及是否把更新结果持久化为可审计状态。
\end{Certificate}

\begin{proof}[证明（谓词因果链）]
若运行内更新，则 $\theta$ 与 $w$ 构成快慢变量耦合；若仅段间更新，则 $\theta$ 在段内不变。
两种情况均可由持久化记录复核，故推论成立。证毕。
\end{proof}

\begin{remark}[评估与落地约束（抽象）]
在不引入任何具体实现细节的前提下，ODTHB 的“阈值慢更新”可按如下口径评估与约束：
\begin{enumerate}
  \raggedright
  \item \textbf{主要收益}：
  \begin{enumerate}
    \item 把“阈值调参”从主观经验降级为证书驱动：$(\Delta J,\Delta\tau)$ 可落盘、可回放、可复核；
    \item 显式处理滞后 $\ell$：避免把“尚未起效”误判为“无效”；
    \item 支持阈值向量化：可把不同事件类型 $k$ 的门控阈值 $\theta_k$ 统一纳入同一慢更新外环。
  \end{enumerate}
  \item \textbf{主要风险}：
  \begin{enumerate}
    \item \textbf{解释口径漂移}：若阈值频繁自适应，解释层的稳定性下降；
    \item \textbf{阈值振荡/触发风暴}：离散门控与节流规则耦合可能产生非线性反馈，需 clip/滞回/低频更新约束；
    \item \textbf{口径混用}：若把“调用次数”误当作“真实事件”，则 $\Delta\tau$ 被系统性抬高，训练信号方向失真。
  \end{enumerate}
  \item \textbf{建议接口形态}：
  \begin{enumerate}
    \item 稳健优先：采用“段间更新、段内冻结”的离线外环（推论~\ref{cor:tex27-ch5-odthb-online-offline}）；
    \item 适应优先：运行内更新时，可用乘法分解 $\theta_k=\alpha_k\cdot\theta^{\mathrm{base}}_k$，
    并将 $\alpha_k\in[\alpha_{\min},\alpha_{\max}]$ 作为可审计慢变量落盘；基准阈值保持稳定以维持解释口径；
    \item 白盒闭环最小化：强制用事件口径构造 $\Delta\tau$、用滞后 $\ell$ 构造 $\Delta J$，并持久化 $(\Delta J,\Delta\tau,\theta)$ 以支持回放审计。
  \end{enumerate}
  \item \textbf{未证明边界}：若进一步把 $\Delta J$ 视作 $\Delta\tau$ 的稳定线性响应或要求随机近似更新的严格收敛，
  需要额外平稳性/噪声/凸性等假设；在不补充这些假设前，ODTHB 仅主张“可回放复核的白盒更新规则”，不主张一般性解析收敛结论。
\end{enumerate}
\end{remark}

\subsubsection{窗口统计维护与几何一阶口径（ODTHB）}
\label{subsubsec:tex27-ch5-odthb-window-stats}

\begin{remark}[与 ODTHB 严格语义的关系]
本小节仅给出 ODTHB 的 fast 层可采用的窗口统计维护技术（例如引理~\ref{lem:tex27-ch5-odthb-push-pop}），
以及把证书变化写成几何一阶口径的符号化解释。
按 ODTHB 严格语义：对微观骨架 $(V,E)$ 与微观态 $w_t$ 仅\textbf{只读}，
且必须严格避免连通结构硬剪裁/删边等微观干预，因为它会破坏微观隧穿（定理~\ref{thm:tex27-ch5-odthb-microtunnel-no-hardprune}）。
\end{remark}

\begin{definition}[自由度参数化与边方向（抽象）]
\label{def:tex27-ch5-odthb-theta}
在需要“几何方向”解释时，可把状态 $v\in V$ 参数化为某个自由度向量 $\vartheta(v)\in\RR^d$（不要求全域可导，只要求窗口内可做局部一阶近似）。
对边 $e=(u\to v)$，定义离散方向（扭曲向量）
\[
  \Delta\vartheta_e:=\vartheta(v)-\vartheta(u).
\]
\end{definition}

\begin{theorem}[定理（自由度微分 $\leftrightarrow$ 分量/总量微分：链式法则 + 点积压缩）]
\label{thm:tex27-ch5-odthb-differential}
设 $\mathbf{g}=\mathbf{g}(\vartheta)$ 且 $J(\vartheta)=\langle \mathbf{w},\mathbf{g}(\vartheta)\rangle$（同定义~\ref{def:tex27-ch5-odthb-j}）。
若在局部口径下 $\mathbf{g}$ 对 $\vartheta$ 可做一阶展开，则
\[
  \mathrm{d}J
  =\left(\frac{\partial\mathbf{g}}{\partial \vartheta}\right)^{\!\top}\mathbf{w}\ \cdot\ \mathrm{d}\vartheta.
\]
因此“自由度变化微分 $\mathrm{d}\vartheta$”与“分量/总量变化微分 $\mathrm{d}\mathbf{g}\to \mathrm{d}J$”在结构上是一条链：
\[
  \mathrm{d}\vartheta \xrightarrow{\ \partial\mathbf{g}/\partial\vartheta\ }\mathrm{d}\mathbf{g}
  \xrightarrow{\ \langle\mathbf{w},\cdot\rangle\ }\mathrm{d}J.
\]
\end{theorem}

\begin{Certificate}[证书（链式法则与线性代数恒等式）]
\label{cert:tex27-ch5-odthb-differential}
证书为两步推演：
\[
  \mathrm{d}J=\sum_i w_i\,\mathrm{d}g_i=\langle \mathbf{w},\mathrm{d}\mathbf{g}\rangle,
  \qquad
  \mathrm{d}\mathbf{g}=\frac{\partial\mathbf{g}}{\partial\vartheta}\mathrm{d}\vartheta.
\]
\end{Certificate}

\begin{proof}[证明（谓词因果链）]
由定义~\ref{def:tex27-ch5-odthb-j}，$J(\vartheta)=\sum_i w_i g_i(\vartheta)$，故
$
\mathrm{d}J=\sum_i w_i\,\mathrm{d}g_i=\langle \mathbf{w},\mathrm{d}\mathbf{g}\rangle
$。
又由局部一阶展开（链式法则），$\mathrm{d}\mathbf{g}=(\partial\mathbf{g}/\partial\vartheta)\mathrm{d}\vartheta$。
代入并整理得结论。证毕。
\end{proof}

\begin{corollary}[推论（几何视角：方向导数与局部下降判别）]
\label{cor:tex27-ch5-odthb-directional-derivative}
令
\[
  \nabla_{\vartheta}J:=\left(\frac{\partial\mathbf{g}}{\partial \vartheta}\right)^{\!\top}\mathbf{w}.
\]
在一阶近似下，沿边方向 $\Delta\vartheta_e$ 的方向导数符号由
\[
  \langle \nabla_{\vartheta}J,\Delta\vartheta_e\rangle<0
\]
给出：其为负表示该方向局部降低 $J$。
因此可用窗口内回归估计 $\widehat{\nabla_{\vartheta}J}$，
并把 $\langle \widehat{\nabla_{\vartheta}J},\Delta\vartheta_e\rangle$ 的符号/幅度作为可审计的诊断字段，
为宏观调参提供信号（并严格避免据此做删边/硬剪裁）。
\end{corollary}

\begin{Certificate}[证书（方向导数判别）]
\label{cert:tex27-ch5-odthb-directional-derivative}
在一阶近似下，$\Delta J\approx \langle \nabla_\vartheta J,\Delta\vartheta\rangle$，
故 $\langle \nabla_\vartheta J,\Delta\vartheta\rangle<0$ 表示沿该方向局部下降。
\end{Certificate}

\begin{proof}[证明（谓词因果链）]
由定理~\ref{thm:tex27-ch5-odthb-differential}，$\mathrm{d}J=\langle \nabla_\vartheta J,\mathrm{d}\vartheta\rangle$。
将 $\mathrm{d}\vartheta$ 以离散差分 $\Delta\vartheta_e$ 代理，得到 $\Delta J$ 的一阶近似式与方向判别。证毕。
\end{proof}

\begin{lemma}[引理（进栈出栈可维护窗口内统计量为“窗口真值”）]
\label{lem:tex27-ch5-odthb-push-pop}
用事件进栈出栈维护 $\{n_e,\overline{\Delta J}_e,\overline{c}_e\}$ 等统计量，
可使这些量始终对应最近窗口，而非全历史。
\end{lemma}

\begin{Certificate}[证书（增量更新 + 减量回滚）]
\label{cert:tex27-ch5-odthb-push-pop}
每次新事件入栈：对其所属边做增量更新；每次最旧事件出栈：对其所属边做减量更新。
则任何时刻的统计量都等于窗口内样本的代数和/均值。
\end{Certificate}

\begin{proof}[证明（构造性；谓词因果链）]
窗口统计量是窗口内样本的有限和/均值。入栈相当于对和加一项，出栈相当于对和减一项；
因此持续维护即可保持统计量等于“窗口真值”。证毕。
\end{proof}

\subsubsection{\ORLDTHBase 总接口（证书化更新器）}
\label{subsubsec:tex27-ch5-odthb-interface}

\begin{Certificate}[\ORLDTHBase：双时间尺度同伦基座的证书化更新器接口]
\label{cert:tex27-ch5-odthb-interface}
\begin{description}
  \item[输入]
  \begin{itemize}
    \raggedright
    \item 状态图 $\Gamma=(V,E)$ 与窗口事件流 $\mathcal{H}$；
    \item 当前微观态：边权优先级场 $w_t:E\to\RR_{>0}$（定义~\ref{def:tex27-ch5-d1}，或其等价的边上连接/能量表述）；
    \item 证书序列 $J(t)$ 与分量 $\mathbf{g}(t)$，以及事件口径计数 $\{N_k(t)\}_{k\in\mathcal{K}}$；
    \item 宏观算子参数向量 $\Lambda$（定义~\ref{def:tex27-ch5-odthb-macroparam}）与窗口/节流设定。
  \end{itemize}
  \item[输出] 宏观算子参数更新后的 $\Lambda_k$（含阈值 $\theta_k$、退火/温控参数等）；
  即时层输出的局部提议核/选择测度 $p_n(\cdot\mid u)$（定义~\ref{def:tex27-ch5-d3}，由 $w_t$ 与 $\varepsilon$ 诱导）；
  以及可审计的微观统计与证书字段（$\mathsf{Stat}(W_n)$、$\Delta J,\Delta\tau,\overline{\Delta J}_e,n_e$ 等）。
  \item[即时层（fast）] 在每个事件步 $n$：
  \begin{enumerate}
    \item 更新窗口 $W_n$ 与统计量（引理~\ref{lem:tex27-ch5-odthb-push-pop}）；
    \item 由定义~\ref{def:tex27-ch5-odthb-microstat} 形成窗口微观统计摘要 $\mathsf{Stat}(W_n)$；
    \item 以当前微观态 $w_t$ 与宏观参数 $\varepsilon$ 形成局部选择测度 $p_n(\cdot\mid u)$（定义~\ref{def:tex27-ch5-d3}）；
    \item 若满足触发谓词（定义~\ref{def:tex27-ch5-odthb-trigger}），登记事件时间并落盘样本对（定义~\ref{def:tex27-ch5-odthb-samples}）。
  \end{enumerate}
  \item[周期层（slow）] 在每个周期 $k$：
  \begin{enumerate}
    \item 用样本对 $(\Delta\tau_k,\Delta J_k)$ 更新阈值 $\theta_k$（命题~\ref{prop:tex27-ch5-odthb-threshold-update}），并将其作为
      $\Lambda_k$ 的分量持久化；
    \item （可选）以 $\mathsf{Stat}(W_n)$ 为输入，对 $\beta$、退火计划与其它宏观超参数做白盒调参（定义~\ref{def:tex27-ch5-odthb-macroparam} 的
      ``$\ldots$'' 部分）；
  \end{enumerate}
  \item[保证] 即时层的微观统计摘要满足“窗口外默认中性”（定理~\ref{thm:tex27-ch5-odthb-reset}）与“有限记忆”（定理~\ref{thm:tex27-ch5-odthb-finite-memory}）；

  且由于 ODTHB 严格禁止硬剪裁并采用 $\varepsilon>0$ 的探索保底，路径支持集保持不变（定理~\ref{thm:tex27-ch5-odthb-support}）；
  ODTHB 的可审计输出（宏观参数更新、阈值更新、退火计划等）均可由持久化字段回放审计。
\end{description}
\end{Certificate}

\begin{remark}[哪些更强结论仍属于“未证明”]
\begin{enumerate}
  \item 上述已严格证明：微观权重态（边权优先级场）与离散（量子）规范场在结构上同构，且“同调修复”可解释为压平回路周期残差；
  \item 若进一步声称“修复必然达到全局最优/收敛”，需要额外假设（例如探索衰减、噪声条件、目标函数性质等）。
  在未补充这些假设前，“全局最优收敛”属于未证明断言。
\end{enumerate}
\end{remark}

\begin{remark}[进一步对齐建议：扇区交集覆盖与 Čech 同调]
若要对齐“扇区交集覆盖/神经复形（nerve）”语言，可把扇区覆盖的重叠区视为 1-胞腔/2-胞腔来源，
用 Čech 同调定义“回路基”与“曲率分量”，则定理~\ref{thm:tex27-ch5-thm4} 的 $g_i$ 可直接对应“覆盖同调障碍证书”，
与“扇区调度 + 同调算法同伦基座”的叙事同构。
\cite{HatcherAT}
\end{remark}

\clearpage

%%%%%%%%%%%%%%%%%%%%%%%%%%%%%%%%%%%%%%%%%%%%%%%%%%%%%%%%%%%%
\section{边缘算子与正交子丛附加规范场：变分驱动的 \texorpdfstring{$H^1$}{H1} 修复与相变触发}
\label{sec:tex27-ch6-edge-operator}
%%%%%%%%%%%%%%%%%%%%%%%%%%%%%%%%%%%%%%%%%%%%%%%%%%%%%%%%%%%%

\subsection{定义：上同调障碍、正交子丛与相变开关}
\label{subsec:tex27-ch6-defs}

\begin{definition}[1 阶规范场：边上的连接（Abelian 情形）]
\label{def:tex27-ch6-d1}
给定有向图 $\Gamma=(V,E)$。令边权优先级 $w:E\to\RR_{>0}$。取 $\beta>0$，定义边上的 1 阶规范场
\[
  A(e):=-\beta\log w(e)\in\RR,\qquad U(e):=\exp(A(e))=w(e)^{-\beta}\in\RR_{>0}.
\]
其中 $A$ 是 1-上链（边函数），$U$ 是其指数化的离散连接。
\end{definition}

\begin{definition}[上同调障碍：回路周期/holonomy 残差]
\label{def:tex27-ch6-d2}
对任意有向回路 $c\subseteq E$，定义
\[
  \mathrm{Hol}(c):=\prod_{e\in c}U(e),\qquad \Omega(c):=\log\mathrm{Hol}(c)=\sum_{e\in c}A(e).
\]
若存在回路使 $\Omega(c)\neq 0$，则 $A$ 在 $H^1(\Gamma;\RR)$ 中对应非零类；这可作为“1 阶上同调障碍”的证书化形式。
\end{definition}

\begin{definition}[正交子丛：把边空间分成“基底子丛 + 修复子丛”]
\label{def:tex27-ch6-d3}
令 1-上链空间 $C^1(\Gamma):=\RR^{E}$，取标准内积
\[
  \langle A,B\rangle := \sum_{e\in E}A(e)B(e).
\]
选一棵生成树 $T\subseteq E$，令弦边集合（非树边）为 $C:=E\setminus T$。定义两个子空间
\[
  C^1_{\mathrm{base}}:=\RR^{T},\qquad C^1_{\perp}:=\RR^{C}.
\]
由于支撑边集不交，故
\[
  C^1(\Gamma)=C^1_{\mathrm{base}}\ \oplus\ C^1_{\perp},\qquad C^1_{\mathrm{base}}\perp C^1_{\perp}.
\]
把 $C^1_{\perp}$ 称为“正交子丛（修复子丛）”。
\end{definition}

\begin{definition}[边缘算子：只在修复子丛上修改的更新算子]
\label{def:tex27-ch6-d4}
一个边缘算子定义为映射
\[
  \mathsf{E}: w \mapsto w',
\]
满足：仅对弦边 $C$ 上的权重做更新，树边 $T$ 保持不变。等价地，在 $A$-表述中：只添加一个 $A_{\perp}\in C^1_{\perp}$。
\end{definition}

\begin{definition}[1 阶规范场相变：阈值触发的离散激活]
\label{def:tex27-ch6-d5}
给定阈值 $\theta>0$ 与某个障碍强度指标（例如基本回路残差范数）
\[
  |r|:=\left(\sum_{e\in C}r(e)^2\right)^{1/2},
\]
其中 $r(e)$ 为弦边 $e$ 对应的基本回路残差（见公理~\ref{thm:tex27-ch6-a6} 与定理~\ref{thm:tex27-ch6-thm1}）。
定义相变开关
\[
  s=\indicator{|r|>\theta}\in\{0,1\}.
\]
当 $s$ 从 $0$ 跳到 $1$ 时，附加规范场从 $0$ 跳到非零——这就是“1 阶规范场相变”的严格化。
\end{definition}

\subsection{公理（降级为定理）：基本回路系}
\label{subsec:tex27-ch6-axiom}

\begin{theorem}[公理（降级为定理）A6（基本回路系：弦边给出回路空间的一组基）]
\label{thm:tex27-ch6-a6}
对每条弦边 $e\in C$，在生成树 $T$ 中存在唯一简单路径 $p_T(e)$ 连接其端点。定义基本回路
\[
  c_e := e\ \cup\ p_T(e).
\]
则 $\{c_e\}_{e\in C}$ 构成回路空间的一组基（任何回路都可表示为它们的 $\ZZ$ 线性组合）。
\end{theorem}

\begin{Certificate}[证书 A6（维数计数 + 唯一基本回路）]
\label{cert:tex27-ch6-a6}
证书为三条标准图论事实：
\begin{enumerate}
  \item 生成树无回路；
  \item 每加一条弦边产生唯一一个基本回路；
  \item 回路空间维数为 $|E|-|V|+1=|C|$，与弦边数一致。
\end{enumerate}
\end{Certificate}

\begin{proof}[证明（标准图论事实；谓词因果链）]
生成树 $T$ 连接任意两点且无回路，因此对每条弦边 $e$，其端点在 $T$ 中存在唯一简单路径 $p_T(e)$，并与 $e$ 共同构成唯一基本回路 $c_e$。
回路空间维数为 $|E|-|V|+1$，恰等于弦边条数 $|C|$，故 $\{c_e\}_{e\in C}$ 既线性无关又张成回路空间，构成一组基。证毕。
\end{proof}

\subsection{定理链：正交子丛修复、变分驱动与相变触发}
\label{subsec:tex27-ch6-theorems}

\begin{theorem}[定理 1（构造性修复：存在只作用于正交子丛的附加规范场，使 $H^1$ 障碍归零）]
\label{thm:tex27-ch6-thm1}
给定任意 $A\in C^1(\Gamma)$。对弦边 $e\in C$，定义基本回路残差
\[
  r(e):=\Omega(c_e)=\sum_{f\in c_e}A(f).
\]
则存在一个附加规范场 $A_{\perp}\in C^1_{\perp}$（只支撑在弦边 $C$ 上），使得修复后的
\[
  A':=A + A_{\perp}
\]
满足对所有基本回路 $c_e$：
\[
  \Omega'(c_e)=\sum_{f\in c_e}A'(f)=0,
\]
从而对任意回路 $c$ 也有 $\Omega'(c)=0$，等价于 $A'$ 的同调类为零（上同调障碍被修复）。

\paragraph{构造（显式）.}
取
\[
  A_{\perp}(e):=-r(e)\quad (e\in C),\qquad A_{\perp}(f):=0\quad (f\in T).
\]
\end{theorem}

\begin{Certificate}[证书 1（一次性消除基本回路周期的显式修复场）]
\label{cert:tex27-ch6-thm1}
证书给出显式修复场：对每条弦边 $e$ 注入 $A_{\perp}(e)=-\Omega(c_e)$；
由于 $A_{\perp}$ 支撑在 $C$ 上，故该注入是“只作用于正交子丛”的边缘算子（定义~\ref{def:tex27-ch6-d4}）。
\end{Certificate}

\begin{proof}[证明（谓词因果链）]
对任意弦边 $e\in C$ 的基本回路 $c_e$：其边集由弦边 $e$ 与树路径 $p_T(e)$ 组成，而 $A_{\perp}$ 只在弦边上非零，故
\[
  \Omega'(c_e)
  =\sum_{f\in c_e}\bigl(A(f)+A_{\perp}(f)\bigr)
  =\underbrace{\sum_{f\in c_e}A(f)}_{=r(e)}+A_{\perp}(e)
  =r(e)-r(e)=0.
\]
由公理~\ref{thm:tex27-ch6-a6}，任意回路都可写成基本回路的整系数组合；在 Abelian 情形下，周期对组合线性，
故所有回路周期皆为 $0$。这等价于 $A'$ 同调类为零（定理~\ref{thm:tex27-ch5-thm3} 的构造式证明同样适用）。证毕。
\end{proof}

\begin{theorem}[定理 2（变分驱动：上述构造是正交子丛内的变分最优解）]
\label{thm:tex27-ch6-thm2}
在修复子丛 $C^1_{\perp}$ 内，考虑变分泛函（只用基本回路残差）
\[
  \Phi(A_{\perp})
  :=\frac12\sum_{e\in C}\Big(\Omega_{A+A_{\perp}}(c_e)\Big)^2
  =\frac12\sum_{e\in C}\big(r(e)+A_{\perp}(e)\big)^2.
\]
则其唯一极小解为
\[
  A_{\perp}^\star(e)=-r(e),
\]
即定理~\ref{thm:tex27-ch6-thm1} 的构造。
\end{theorem}

\begin{Certificate}[证书 2（独立严格凸二次型）]
\label{cert:tex27-ch6-thm2}
证书为：$\Phi$ 在每个坐标 $A_{\perp}(e)$ 上都是独立的严格凸二次函数，故极小点满足一阶条件且唯一。
\end{Certificate}

\begin{proof}[证明（谓词因果链）]
对任意 $e\in C$，
\[
  \frac{\partial\Phi}{\partial A_{\perp}(e)}=r(e)+A_{\perp}(e).
\]
令其为 0 得 $A_{\perp}(e)=-r(e)$。且 $\Phi$ 的 Hessian 为单位阵，严格凸，故极小解唯一。证毕。
\end{proof}

\begin{theorem}[定理 3（1 阶规范场相变前提：阈值触发使修复从 0 跳到非零）]
\label{thm:tex27-ch6-thm3}
定义相变开关
\[
  s=\indicator{|r|>\theta},\qquad |r|:=\left(\sum_{e\in C}r(e)^2\right)^{1/2}.
\]
定义“相变式边缘算子”
\[
  A_{\perp}^{(\theta)} := s\cdot A_{\perp}^\star.
\]
则当 $|r|$ 穿越 $\theta$ 时，$A_{\perp}^{(\theta)}$ 从 $0$ 非连续跳到 $A_{\perp}^\star$（或其缩放），从而 $U(e)$ 在弦边上发生 1 阶跳变；
并且在 $s=1$ 时完成同调修复（定理~\ref{thm:tex27-ch6-thm1}）。
\end{theorem}

\begin{Certificate}[证书 3（指示函数触发的非连续注入）]
\label{cert:tex27-ch6-thm3}
证书为：$s$ 为指示函数，$s=0$ 时不注入修复场，$s=1$ 时注入 $A_{\perp}^\star$；
阈值处的跳变即为“相变式激活”。
\end{Certificate}

\begin{proof}[证明（谓词因果链）]
若 $|r|\le\theta$，则 $s=0\Rightarrow A_{\perp}^{(\theta)}=0$；若 $|r|>\theta$，则 $s=1\Rightarrow A_{\perp}^{(\theta)}
  =A_{\perp}^\star$，
并由定理~\ref{thm:tex27-ch6-thm1} 消除所有基本回路周期。
由于 $s$ 在阈值处不连续，故 $A_{\perp}^{(\theta)}$ 的激活为非连续跳变，即 1 阶相变。证毕。
\end{proof}

\subsection{命题、推论与备注}
\label{subsec:tex27-ch6-others}

\begin{proposition}[命题 1（边缘算子在 $w$-口径下的显式更新公式：只改弦边）]
\label{prop:tex27-ch6-p1}
由 $A=-\beta\log w$ 得
\[
  w'(e)=\exp\!\Big(-\frac{A'(e)}{\beta}\Big).
\]
对弦边 $e\in C$，$A'(e)=A(e)+A_{\perp}(e)$，所以
\[
  w'(e)=w(e)\cdot \exp\!\Big(-\frac{A_{\perp}(e)}{\beta}\Big)
  =w(e)\cdot \exp\!\Big(\frac{r(e)}{\beta}\Big),
\]
而对树边 $f\in T$，$w'(f)=w(f)$。
因此边缘算子是可实现的逐边乘法更新。
\end{proposition}

\begin{Certificate}[证书 1（对数-指数口径互换）]
\label{cert:tex27-ch6-p1}
证书为两步互换：$A'=-\beta\log w'\Rightarrow w'=\exp(-A'/\beta)$；
再代入 $A'=A+A_{\perp}$ 并用 $A_{\perp}(e)=-r(e)$（定理~\ref{thm:tex27-ch6-thm1}）。
\end{Certificate}

\begin{proof}[证明（谓词因果链）]
由 $A=-\beta\log w$ 得 $w=\exp(-A/\beta)$，对 $A'$ 同理。
对弦边代入 $A'(e)=A(e)+A_{\perp}(e)$ 得
$
w'(e)=\exp(-(A(e)+A_{\perp}(e))/\beta)=w(e)\exp(-A_{\perp}(e)/\beta)=w(e)\exp(r(e)/\beta)
$
对树边 $A_{\perp}=0$，故 $w'=w$。证毕。
\end{proof}

\begin{corollary}[推论 1（正交子丛附加规范场是修复塔的边缘算子接口）]
\label{cor:tex27-ch6-c1}
定理~\ref{thm:tex27-ch6-thm1}--\ref{thm:tex27-ch6-thm3} 给出一个可构造接口：
\[
  \text{输入 }w\Rightarrow A\Rightarrow r(e)=\Omega(c_e)\Rightarrow (|r|>\theta)\ ?\ \text{注入 }A_{\perp}^\star:\
    0\Rightarrow w'.
\]
它同时满足：
\begin{enumerate}
  \item 边缘算子：只在边上更新；
  \item 正交子丛：只作用于弦边子空间 $C^1_{\perp}$；
  \item 修复上同调障碍：把 $H^1$ 周期归零；
  \item 变分驱动 + 1 阶相变：由凸变分极小化得到修复场，阈值触发产生相变式激活。
\end{enumerate}
\end{corollary}

\begin{Certificate}[证书 1（接口的四要素对应）]
\label{cert:tex27-ch6-c1}
证书为：定理~\ref{thm:tex27-ch6-thm1} 给出“只改弦边并消除周期”的构造；
定理~\ref{thm:tex27-ch6-thm2} 给出“变分最优”；定理~\ref{thm:tex27-ch6-thm3} 给出“阈值触发相变”；命题~\ref{prop:tex27-ch6-p1} 给出 $w$ 口径实现公式。
\end{Certificate}

\begin{proof}[证明（谓词因果链）]
逐条调用证书~\ref{cert:tex27-ch6-c1} 所列定理/命题即可把接口四要素闭合，故推论成立。证毕。
\end{proof}

\begin{remark}[必要边界条件，避免“强结论”误读]
\begin{enumerate}
  \item “完全消除障碍”依赖于我们修复的是图的 1 维上同调 $H^1$（有效 1 阶对象）。若要修复更高阶 $H^k(k\ge 2)$，需要更高维胞腔或神经复形高阶结构；
  \item 此处规范群选为 $(\RR_{>0},\times)$（对应权重，Abelian）。若改为非交换李群，则“周期线性叠加”不再严格成立，需要路径序指数并在近单位元下线性化；
  \item 相变阈值 $\theta$ 是工程约束：平时不扰动系统，只有障碍证书超过阈值才激活修复子丛，与“护栏型边缘算子”语义一致。
\end{enumerate}
\end{remark}

\clearpage

%%%%%%%%%%%%%%%%%%%%%%%%%%%%%%%%%%%%%%%%%%%%%%%%%%%%%%%%%%%%
\section{GZ-OHU 与正交子丛：超限递归与工程有限截断}
\label{sec:tex27-ch7-ohu-orthogonal}
%%%%%%%%%%%%%%%%%%%%%%%%%%%%%%%%%%%%%%%%%%%%%%%%%%%%%%%%%%%%

\subsection{定义：OHU 修补序列与正交子丛分解}
\label{subsec:tex27-ch7-defs}

\begin{definition}[OHU 的同伦修补序列（截断态）]
\label{def:tex27-ch7-d1}
令 $L$ 为初始对象（法则/算子系统），其缺口由证书 $\JacGap(L)\ge 0$ 刻画。
OHU 过程给出一条修补序列（记作截断态）
\[
  L^{(2)}:=L,\quad L^{(3)},\quad L^{(4)},\ \dots,\ L^{(n)},\ \dots
\]
其中每一步引入一个更高阶修正算子 $l_n$（$n\ge 3$），形成“级联修补”。
\end{definition}

\begin{definition}[正交子丛：把可修补自由度分解为互不干扰的独立分量]
\label{def:tex27-ch7-d2}
取一个“修补自由度空间” $\mathcal{H}$（可理解为所有候选边缘算子/附加规范场的线性张成空间），并赋予内积 $\langle\cdot,\cdot\rangle$。
“正交子丛族”指一个（可能超限长的）直和分解
\[
  \mathcal{H}=\bigoplus_{\alpha<\kappa}\mathcal{H}_\alpha,
  \qquad
  \mathcal{H}_\alpha\perp \mathcal{H}_\beta\ (\alpha\neq\beta),
\]
其中 $\kappa$ 可以是一个序数（超限），每个 $\mathcal{H}_\alpha$ 对应“一组新的修补方向/新的边缘算子族”。
\end{definition}

\begin{definition}[有限工程截断：有限正交子丛]
\label{def:tex27-ch7-d3}
“有限正交子丛”就是取前 $k$ 个分量：
\[
  \mathcal{H}_{\le k}:=\bigoplus_{1\le \alpha\le k}\mathcal{H}_\alpha.
\]
其对应“有限阶截断态” $L^{(k)}$。
\end{definition}

\subsection{公理（降级为定理）：OHU 证书输入与工程截断口径}
\label{subsec:tex27-ch7-axioms}

\begin{theorem}[公理（降级为定理）A7（GZ-OHU 的超限归纳与极限收敛）]
\label{thm:tex27-ch7-a7}
OHU 的证明机制是一个动态构造过程：通过超限归纳逐层引入 $l_3,l_4,\dots$，
并在极限下稳定化收敛到某个 $L^{\mathrm{OHU}}$，使 $\JacGap(L^{\mathrm{OHU}})=0$（同伦意义/所选口径下）。
\end{theorem}

\begin{Certificate}[证书 A7（外部证书输入：OHU 文档的证明结构摘要）]
\label{cert:tex27-ch7-a7}
本条作为“证书输入接口”：其内容来自 OHU 说明文档对“超限归纳 + 极限 $\JacGap\to 0$”的直接陈述，
在本笔记中仅作为可调用的外部证书，不在此处重建完整证明链。
\cite{GaoZheng2025GFramework}
\end{Certificate}

\begin{proof}[证明（证书化）]
按证书~\ref{cert:tex27-ch7-a7}，该断言作为外部已给证书输入被调用。证毕。
\end{proof}

\begin{theorem}[公理（降级为定理）B7（工程截断：有限阶 $k$ 与精度阈值）]
\label{thm:tex27-ch7-b7}
工程实现必须在有限阶 $k$ 截断，得到中间态 $L^{(k)}$；
此时 $\JacGap(L^{(k)})\neq 0$ 的残差被解释为“工程误差/精度公差”，并用阈值 $\epsilon$ 判定是否需要继续增加阶数 $k$。
\end{theorem}

\begin{Certificate}[证书 B7（外部证书输入：有限算力 $\Rightarrow$ 有限截断）]
\label{cert:tex27-ch7-b7}
本条同样作为“工程约束证书输入”：有限算力导致有限截断；$\JacGap(L^{(k)})$ 作为精度证书；
若超阈值则提示“增加 $k$（引入更高阶算子）”。
\cite{GaoZheng2025GFramework}
\end{Certificate}

\begin{proof}[证明（证书化）]
按证书~\ref{cert:tex27-ch7-b7}，该断言作为外部已给工程证书输入被调用。证毕。
\end{proof}

\subsection{定理链：正交子丛实现与工程截断}
\label{subsec:tex27-ch7-theorems}

\begin{theorem}[定理 1（超限递归构造正交子丛可实现 OHU 的同伦修补链）]
\label{thm:tex27-ch7-thm1}
在公理~\ref{thm:tex27-ch7-a7} 的前提下，存在一种实现方式：
把每一次新增的高阶修补自由度（对应 $l_n$）选取为“相对于已选修补自由度的正交补上的新分量”，
从而得到一个（可能超限长的）正交子丛序列 $\{\mathcal{H}_\alpha\}_{\alpha<\kappa}$，使得：
\begin{enumerate}
  \item 逐层扩张：每一步只在新的正交子丛 $\mathcal{H}_\alpha$ 上注入附加边缘算子；
  \item 不互相干扰：由于正交性，新修补方向不会把已修补掉的分量“又拉回来”（至少在一阶变分/能量分解意义下）；
  \item 与 OHU 同伦等价：当 $\alpha\to\kappa$（对应 $n\to\infty$ 的极限）时，所得极限对象与 $L^{\mathrm{OHU}}$ 同伦等价，且 $\JacGap\to 0$ 由 A7 保证。
\end{enumerate}
\end{theorem}

\begin{Certificate}[证书 1（逐步正交化 / 超限正交化）]
\label{cert:tex27-ch7-thm1}
证书为经典内积空间构造：
\begin{enumerate}
  \item 给定任意修补序列 $\{l_3,l_4,\dots\}\subseteq\mathcal{H}$，可对其做“逐步正交化”，即把新向量投影到已生成子空间的正交补上；
  \item 超限情形用超限递归/超限归纳即可把该过程延拓到序数长度；
  \item 正交化仅改变坐标基底，不改变张成子空间，因此不改变“是否存在趋于零缺口的修补路径”的可达性。
\end{enumerate}
\end{Certificate}

\begin{proof}[证明（谓词因果链）]
\begin{enumerate}
  \item（可构造正交化）由证书~\ref{cert:tex27-ch7-thm1}，可把任意新增修补方向替换为相对于既有方向的正交补分量，从而得到正交子丛序列；
  \item（与 A7 的结构对齐）把“第 $n$ 步引入 $l_n$”实现为“第 $\alpha$ 步在新子丛 $\mathcal{H}_\alpha$ 上注入修补项”，两者在“逐层新增自由度并吸收缺口”的结构上同构；
  \item（极限性质）A7 已给出“随阶数无穷展开，$\JacGap$ 收敛为 0”的保证；而正交化是对修补自由度坐标的基底变换，不改变可达极限对象的同伦型。
\end{enumerate}
故结论成立。证毕。
\end{proof}

\begin{remark}[备注：本节结论的力度边界]
此处并不声称“OHU 原证明一定用正交子丛”，而是证明“OHU 的超限修补机制\emph{存在一个等价的正交子丛实现}”。
\end{remark}

\begin{theorem}[定理 2（工程有限正交子丛 = OHU 的有限截断；残余 JacGap = 精度证书）]
\label{thm:tex27-ch7-thm2}
在公理~\ref{thm:tex27-ch7-b7} 的前提下，工程实现采用有限阶 $k$ 截断，相当于只取有限正交子丛
\[
  \mathcal{H}_{\le k}=\bigoplus_{\alpha\le k}\mathcal{H}_\alpha
\]
来注入修补项；此时 $\JacGap(L^{(k)})$ 就是“距极限 $L^{\mathrm{OHU}}$ 还有多远”的可计算精度证书。
若 $\JacGap(L^{(k)})>\epsilon$，则增加 $k$（再加入一个新的正交子丛）是严格的“提高精度”动作。
\end{theorem}

\begin{Certificate}[证书 2（截断阶数 $\Leftrightarrow$ 子丛数量）]
\label{cert:tex27-ch7-thm2}
证书由 B7 的工程口径给出：有限算力 $\Rightarrow$ 必须在某个有限阶 $k$ 截断；
而“增加 $k$”在正交子丛语言下就是“在 $\mathcal{H}_{\le k}$ 之外再加入一个新的正交分量 $\mathcal{H}_{k+1}$”。
\end{Certificate}

\begin{proof}[证明（谓词因果链）]
由公理~\ref{thm:tex27-ch7-b7}，工程实现只能得到有限截断态 $L^{(k)}$，且 $\JacGap(L^{(k)})$ 作为精度证书；
若超阈值则需要增加阶数 $k$。
另一方面，由定义~\ref{def:tex27-ch7-d3}，增加 $k$ 等价于把可用修补自由度从 $\mathcal{H}_{\le k}$ 扩展到 $\mathcal{H}_{\le k+1}$。
两者一一对应，故定理成立。证毕。
\end{proof}

\subsection{推论与备注}
\label{subsec:tex27-ch7-others}

\begin{corollary}[推论 1（断言的严格版本：纯数 vs 工程）]
\label{cor:tex27-ch7-c1}
\[
  \boxed{\begin{aligned}
    &\text{纯数卷：OHU 的超限归纳修补}\\
    &\simeq\ \text{超限递归构造的正交子丛修补}
  \end{aligned}}
\]
\[
  \boxed{\begin{aligned}
    &\text{工程卷：有限算力截断 }L^{(k)}
    \ \simeq\ \text{有限正交子丛 }\mathcal{H}_{\le k},\\
    &\text{且 }\JacGap(L^{(k)})\ \text{是精度证书}
  \end{aligned}}
\]
并且“$\JacGap(L^{(k)})>\epsilon\Rightarrow$ 增加 $k$”就是“扩展正交子丛数量以提升精度”的严格动作。
\end{corollary}

\begin{Certificate}[证书 1（定理 1 + 定理 2 的并列引用）]
\label{cert:tex27-ch7-c1}
证书为：纯数部分由定理~\ref{thm:tex27-ch7-thm1} 给出；工程部分由定理~\ref{thm:tex27-ch7-thm2} 给出；阈值触发由公理~\ref{thm:tex27-ch7-b7} 固化。
\end{Certificate}

\begin{proof}[证明（谓词因果链）]
分别由定理~\ref{thm:tex27-ch7-thm1}、定理~\ref{thm:tex27-ch7-thm2} 与公理~\ref{thm:tex27-ch7-b7} 直接推出推论中的两条框式陈述与阈值触发动作为一致的翻译。证毕。
\end{proof}

\begin{remark}[落到 ORL/边缘算子框架的直接对应]
\begin{enumerate}
  \item 正交性的工程实现：可取“支撑不交”（不同修补模块作用于不交的边/规则集合）或“显式正交化”（对新模块做投影扣除，等价 Gram--Schmidt）；
  \item 超限递归在工程中退化为有限：数学上 $n\to\infty$ 给理想可逆；工程上必须截断，残余 $\JacGap$ 变成可控误差，超阈值则“加一个正交子丛/加一个模块”；
  \item 与“1 阶规范场相变触发”的兼容性：把 $\JacGap(L^{(k)})>\epsilon$ 视作相变触发条件，则系统从 $k$ 跳到 $k+1$，与 B7 的“超阈值报警 $\to$ 增加 $k$”同构。
\end{enumerate}
\end{remark}

%%%%%%%%%%%%%%%%%%%%%%%%%%%%%%%%%%%%%%%%%%%%%%%%%%%%%%%%%%%%
\section{参数不断同伦扭曲与统计力学变分优化：演化同伦等价（强命题）}
\label{sec:tex27-ch8-homotopy-statmech}
%%%%%%%%%%%%%%%%%%%%%%%%%%%%%%%%%%%%%%%%%%%%%%%%%%%%%%%%%%%%

本节给出一个可证书化的强命题：把“参数不断同伦扭曲”严格建模为（可能非交换的）更新半群/马尔可夫核；
把“经典/量子统计力学实现变分优化”建模为“自由能泛函最小化”（经典分布/量子密度矩阵）；并在明确条件下说明两类演化在同伦意义下等价
（更准确：\textbf{演化同伦等价}，即轨道族/不动点/吸引子与可达集结构一致，且存在连续变形把一种演化变到另一种演化）。

\subsection{定义组：参数空间、缺陷势能与同伦扭曲演化}
\label{subsec:tex27-ch8-defs}

\begin{definition}[D1（参数空间与缺陷势能）]
\label{def:tex27-ch8-d1}
令参数空间为可微流形（或欧氏空间）
\[
  \Theta \subseteq \RR^d,
  \qquad
  d\in\NN.
\]
给定缺陷势能（可对应工程中的 $\Phi_{\mathrm{defect}}$ 或缺口证书 $\JacGap$ 的能量版本）
\[
  \Phi:\Theta\to\RR.
\]
\end{definition}

\begin{definition}[D2（“不断同伦扭曲”的形式化：更新半群/马尔可夫核）]
\label{def:tex27-ch8-d2}
把“参数不断同伦扭曲”抽象为一族随时间变化的更新算子（允许非交换）：
\[
  \theta_{t+\Delta t}:=T_{t,\Delta t}(\theta_t),
  \qquad
  T_{t,\Delta t}\in\mathcal{T},
\]
其中 $\mathcal{T}$ 由“边缘算子/修复算子”生成；非交换性体现在一般可有
$T_{t_2,\Delta t_2}\circ T_{t_1,\Delta t_1}\neq T_{t_1,\Delta t_1}\circ T_{t_2,\Delta t_2}$。

连续极限可写为常微分方程（流）
\[
  \dot{\theta}(t)=v_t(\theta(t)),
\]
更一般地，可考虑含噪声的随机微分方程（It\^o SDE）
\[
  d\theta_t = -\nabla \Phi(\theta_t)\,dt + b_t(\theta_t)\,dt + \sqrt{2/\beta}\,dW_t,
\]
其中 $\beta>0$ 为逆温度，$b_t$ 表示可选的非梯度“修复漂移”项，$W_t$ 为 $d$ 维布朗运动。

对应的分布级描述是马尔可夫核（或半群）：给定可测集 $A\subseteq\Theta$，令
\[
  K_{t,\Delta t}(\theta,A):=\mathbb{P}\bigl(\theta_{t+\Delta t}\in A \mid \theta_t=\theta\bigr),
\]
并用它诱导分布演化 $q_{t+\Delta t}=q_t K_{t,\Delta t}$。
\end{definition}

\begin{definition}[D3（经典统计力学的变分优化：Gibbs 测度与自由能泛函）]
\label{def:tex27-ch8-d3}
给定参考测度（先验）$\mu(d\theta)$（可取均匀/高斯先验等），定义经典 Gibbs 测度
\[
  q_\beta(d\theta)
  :=\frac{1}{Z(\beta)}e^{-\beta\Phi(\theta)}\,\mu(d\theta),
  \qquad
  Z(\beta):=\int_\Theta e^{-\beta\Phi(\theta)}\,\mu(d\theta).
\]
对所有满足 $q\ll\mu$ 的概率测度 $q$，定义自由能泛函
\[
  \mathcal{F}_\beta[q]
  :=\int_\Theta \Phi(\theta)\,q(d\theta)
  +\frac{1}{\beta}\mathrm{KL}(q\|\mu),
\]
其中相对熵（KL）定义为
\[
  \mathrm{KL}(q\|\mu):=\int_\Theta \log\Bigl(\frac{dq}{d\mu}\Bigr)\,q(d\theta).
\]
\end{definition}

\begin{definition}[D4（量子统计力学的变分优化：密度矩阵自由能）]
\label{def:tex27-ch8-d4}
构造 Hilbert 空间 $\mathcal{H}$，给定哈密顿量算子 $\widehat{H}$（“能量”对应 $\Phi$ 的算符化），量子 Gibbs 态
\[
  \rho_\beta=\frac{e^{-\beta \widehat{H}}}{\mathrm{Tr}(e^{-\beta \widehat{H}})}.
\]
对所有密度矩阵 $\rho\succeq 0$ 且 $\mathrm{Tr}(\rho)=1$，定义量子自由能泛函
\[
  \mathcal{F}^{Q}_\beta[\rho]
  :=\mathrm{Tr}(\rho\,\widehat{H})
  +\frac{1}{\beta}\mathrm{Tr}(\rho\log\rho),
\]
并在全体密度矩阵上取最小。
若 $\widehat{H}=\widehat{H}_{\mathrm{diag}}(\Phi)+\widehat{K}$，则非交换项 $\widehat{K}$ 可提供“隧穿/混合”的算法学解释。
\end{definition}

\begin{definition}[D5（演化同伦等价：同伦半共轭 + 吸引子结构对应）]
\label{def:tex27-ch8-d5}
设 $\{T_t\}_{t\ge 0}$ 是 $\Theta$ 上的（半）群（或流），$\{S_t\}_{t\ge 0}$ 是状态空间 $\mathcal{X}$ 上的（半）群。
称二者在同伦意义下等价（演化同伦等价），若存在连续映射
\[
  h:\Theta\to\mathcal{X},
\]
使得对每个 $t\ge 0$，两映射 $h\circ T_t$ 与 $S_t\circ h$ 同伦：
存在连续同伦
\[
  H_t:\Theta\times[0,1]\to\mathcal{X},
  \qquad
  H_t(\cdot,0)=h\circ T_t,\ H_t(\cdot,1)=S_t\circ h,
\]
并且两侧的不动点集/吸引子/可达集在 $h$ 下对应（至少在同伦意义下对应）。
\end{definition}

\subsection{公理（降级为定理）：自由能变分原理与零温极限}
\label{subsec:tex27-ch8-axioms}

\begin{theorem}[公理（降级为定理）T1（经典自由能的变分原理）]
\label{thm:tex27-ch8-t1}
对任意 $\beta>0$，自由能 $\mathcal{F}_\beta[q]$ 在可行域 $\{q:\ q\ll\mu,\ q(\Theta)=1\}$ 上的唯一极小点为 Gibbs 测度 $q_\beta$（定义~\ref{def:tex27-ch8-d3}）。
\end{theorem}

\begin{Certificate}[证书 T1（KL 分解恒等式）]
\label{cert:tex27-ch8-t1}
对任意可行 $q$，有分解恒等式
\[
  \mathcal{F}_\beta[q]-\mathcal{F}_\beta[q_\beta]
  =\frac{1}{\beta}\mathrm{KL}(q\|q_\beta)\ge 0.
\]
\end{Certificate}

\begin{proof}[证明（谓词因果链）]
写 $q_\beta(d\theta)=\frac{1}{Z}e^{-\beta\Phi(\theta)}\mu(d\theta)$。
则
\[
  \log\Bigl(\frac{dq_\beta}{d\mu}\Bigr)=-\beta\Phi-\log Z.
\]
	对任意 $q\ll\mu$，
	\begin{align*}
	  \mathrm{KL}(q\|q_\beta)
	  &=
	  \int_\Theta \log\Bigl(\frac{dq}{dq_\beta}\Bigr)\,q(d\theta)
	  \\
	  &=
	  \int_\Theta \log\Bigl(\frac{dq}{d\mu}\Bigr)\,q(d\theta)
	  -\int_\Theta \log\Bigl(\frac{dq_\beta}{d\mu}\Bigr)\,q(d\theta)
	  \\
	  &=
	  \mathrm{KL}(q\|\mu)+\beta\int_\Theta \Phi\,q(d\theta)+\log Z.
	\end{align*}
另一方面，
\[
  \mathrm{KL}(q_\beta\|\mu)=\int_\Theta (-\beta\Phi-\log Z)\,q_\beta(d\theta)
  =-\beta\int_\Theta \Phi\,q_\beta(d\theta)-\log Z,
\]
因此
\[
  \mathcal{F}_\beta[q_\beta]
  =\int_\Theta \Phi\,q_\beta(d\theta)+\frac{1}{\beta}\mathrm{KL}(q_\beta\|\mu)
  =-\frac{1}{\beta}\log Z.
\]
将两式合并并整理得
\[
  \mathcal{F}_\beta[q]
  =\mathcal{F}_\beta[q_\beta]+\frac{1}{\beta}\mathrm{KL}(q\|q_\beta).
\]
由于 $\mathrm{KL}(q\|q_\beta)\ge 0$ 且等号当且仅当 $q=q_\beta$，故 $q_\beta$ 为唯一极小点。证毕。
\end{proof}

\begin{theorem}[公理（降级为定理）T2（零温极限：变分优化退化为能量最小化）]
\label{thm:tex27-ch8-t2}
记 $\Phi_{\min}:=\inf_{\theta\in\Theta}\Phi(\theta)$。
对任意 $\epsilon>0$，令 $\epsilon$-次水平集
\[
  A_\epsilon:=\{\theta\in\Theta:\ \Phi(\theta)\le \Phi_{\min}+\epsilon\}.
\]
若对任意 $\epsilon>0$ 均有 $\mu(A_\epsilon)>0$，则当 $\beta\to\infty$ 时有 $q_\beta(A_\epsilon)\to 1$，
即 Gibbs 测度在零温极限下把质量集中到 $\arg\min\Phi$（及其任意小邻域）上。
\end{theorem}

\begin{Certificate}[证书 T2（指数权重压制证书）]
\label{cert:tex27-ch8-t2}
对任意 $0<\delta<\epsilon$，有比值上界
\[
  \frac{\int_{\Theta\setminus A_\epsilon} e^{-\beta\Phi}\,d\mu}{\int_{A_\delta} e^{-\beta\Phi}\,d\mu}
  \le e^{-\beta(\epsilon-\delta)}\cdot \frac{\mu(\Theta)}{\mu(A_\delta)}.
\]
\end{Certificate}

\begin{proof}[证明（谓词因果链）]
取任意 $0<\delta<\epsilon$。由 $A_\epsilon,A_\delta$ 的定义：
若 $\theta\notin A_\epsilon$ 则 $\Phi(\theta)\ge \Phi_{\min}+\epsilon$；
若 $\theta\in A_\delta$ 则 $\Phi(\theta)\le \Phi_{\min}+\delta$。
于是
\[
  \int_{\Theta\setminus A_\epsilon} e^{-\beta\Phi}\,d\mu
  \le e^{-\beta(\Phi_{\min}+\epsilon)}\mu(\Theta),
  \qquad
  \int_{A_\epsilon} e^{-\beta\Phi}\,d\mu
  \ge \int_{A_\delta} e^{-\beta\Phi}\,d\mu
  \ge e^{-\beta(\Phi_{\min}+\delta)}\mu(A_\delta).
\]
两式相除得到证书~\ref{cert:tex27-ch8-t2} 的上界。
特别地，取 $\delta=\epsilon/2$，则由假设 $\mu(A_{\epsilon/2})>0$，有
\[
  q_\beta(\Theta\setminus A_\epsilon)
  =\frac{\int_{\Theta\setminus A_\epsilon} e^{-\beta\Phi}\,d\mu}{\int_\Theta e^{-\beta\Phi}\,d\mu}
  \le \frac{\int_{\Theta\setminus A_\epsilon} e^{-\beta\Phi}\,d\mu}{\int_{A_{\epsilon/2}} e^{-\beta\Phi}\,d\mu}
  \le e^{-\beta\epsilon/2}\cdot \frac{\mu(\Theta)}{\mu(A_{\epsilon/2})}\xrightarrow[\beta\to\infty]{}0.
\]
故 $q_\beta(A_\epsilon)=1-q_\beta(\Theta\setminus A_\epsilon)\to 1$。证毕。
\end{proof}

\begin{theorem}[公理（降级为定理）T3（量子自由能的变分原理）]
\label{thm:tex27-ch8-t3}
对任意 $\beta>0$，量子自由能 $\mathcal{F}^Q_\beta[\rho]$ 在全体密度矩阵上唯一极小点为量子 Gibbs 态 $\rho_\beta$（定义~\ref{def:tex27-ch8-d4}）。
\end{theorem}

\begin{Certificate}[证书 T3（量子相对熵分解证书）]
\label{cert:tex27-ch8-t3}
记量子相对熵 $S(\rho\|\sigma):=\mathrm{Tr}(\rho(\log\rho-\log\sigma))$。
则对任意密度矩阵 $\rho$，有恒等式
\[
  \mathcal{F}^Q_\beta[\rho]-\mathcal{F}^Q_\beta[\rho_\beta]
  =\frac{1}{\beta}S(\rho\|\rho_\beta)\ge 0.
\]
\end{Certificate}

\begin{proof}[证明（谓词因果链）]
由 $\rho_\beta=\frac{e^{-\beta\widehat{H}}}{\mathrm{Tr}(e^{-\beta\widehat{H}})}$，令 $Z_Q:=\mathrm{Tr}(e^{-\beta\widehat{H}}
  )$，则
\[
  \log\rho_\beta=-\beta\widehat{H}-\log Z_Q\cdot I.
\]
于是
\[
  S(\rho\|\rho_\beta)
  =\mathrm{Tr}(\rho\log\rho)-\mathrm{Tr}(\rho\log\rho_\beta)
  =\mathrm{Tr}(\rho\log\rho)+\beta\,\mathrm{Tr}(\rho\,\widehat{H})+\log Z_Q.
\]
另一方面，
\[
  \mathrm{Tr}(\rho_\beta\log\rho_\beta)
  =-\beta\,\mathrm{Tr}(\rho_\beta\,\widehat{H})-\log Z_Q,
  \qquad
  \Rightarrow\quad
  \mathcal{F}^Q_\beta[\rho_\beta]=-\frac{1}{\beta}\log Z_Q.
\]
因此
\[
  \mathcal{F}^Q_\beta[\rho]-\mathcal{F}^Q_\beta[\rho_\beta]
  =\mathrm{Tr}(\rho\,\widehat{H})+\frac{1}{\beta}\mathrm{Tr}(\rho\log\rho)+\frac{1}{\beta}\log Z_Q
  =\frac{1}{\beta}S(\rho\|\rho_\beta)\ge 0,
\]
且等号当且仅当 $\rho=\rho_\beta$。故 $\rho_\beta$ 为唯一极小点。证毕。
\end{proof}

\subsection{引理：门控下降、含噪演化与 Gibbs 采样}
\label{subsec:tex27-ch8-lemmas}

\begin{lemma}[L1（“只接受 $\Delta\Phi<0$”对应零温极限的确定性下降）]
\label{lem:tex27-ch8-l1}
设离散迭代满足严格门控
\[
  \Phi(\theta_{k+1})<\Phi(\theta_k),
  \qquad
  k\in\NN,
\]
并且步长足够小使该迭代可被视为某个下降向量场的离散化。
则该过程与“$\beta\to\infty$ 的自由能/能量最小化极限”在目标与吸引子层面一致：
其极限集合落在 $\Phi$ 的临界集（局部极小/鞍点/边界驻点）上；若再允许有限强度的越障算子/扰动，则可跨越局部陷阱并逼近全局极小集合。
\end{lemma}

\begin{Certificate}[证书 L1（门控是 Metropolis 在 $\beta\to\infty$ 的退化）]
\label{cert:tex27-ch8-l1}
Metropolis 接受率 $p(\Delta\Phi)=\min\{1,e^{-\beta\Delta\Phi}\}$ 满足：
当 $\Delta\Phi>0$ 时，$\lim_{\beta\to\infty}p(\Delta\Phi)=0$；当 $\Delta\Phi<0$ 时，$p(\Delta\Phi)=1$。
因此“只接受 $\Delta\Phi<0$”可视为 $\beta\to\infty$ 的退化规则；并由定理~\ref{thm:tex27-ch8-t2}，零温极限的目标集合为 $\arg\min\Phi$。
\end{Certificate}

\begin{proof}[证明（谓词因果链）]
\begin{enumerate}
  \item（门控 $\Rightarrow$ 单调下降）由假设 $\Phi(\theta_{k+1})<\Phi(\theta_k)$，序列 $\{\Phi(\theta_k)\}$ 单调下降；
  \item（有下界 $\Rightarrow$ 收敛）若 $\Phi$ 下有界（在工程中通常由势能设计保证），则 $\{\Phi(\theta_k)\}$ 收敛；
  \item（极限点是驻点类）在“步长足够小 + 可视为下降向量场离散化”的前提下，标准离散动力系统结论给出任意聚点落在下降流的极限集上（临界集/边界驻点等）；
  \item（与零温变分极限一致）由证书~\ref{cert:tex27-ch8-l1}，门控规则是有限温接受机制在 $\beta\to\infty$ 的退化；由定理~\ref{thm:tex27-ch8-t2}，
    $\beta\to\infty$ 时 Gibbs 质量集中到 $\arg\min\Phi$ 的任意小邻域；因此两者在“追求能量下降/极小集合”口径下对齐。
\end{enumerate}
证毕。
\end{proof}

\begin{lemma}[L2（加入噪声的同伦扭曲实现 Gibbs 采样，等价于最小化自由能）]
\label{lem:tex27-ch8-l2}
若“不断同伦扭曲”包含扩散噪声（例如定义~\ref{def:tex27-ch8-d2} 的 Langevin SDE，取 $b_t\equiv 0$）
或包含满足详细平衡的 Metropolis 接受率，则其平稳分布为 $q_\beta$；
而 $q_\beta$ 正是自由能 $\mathcal{F}_\beta$ 的极小点（定理~\ref{thm:tex27-ch8-t1}）。
\end{lemma}

\begin{Certificate}[证书 L2（稳态验证：Fokker--Planck / 详细平衡）]
\label{cert:tex27-ch8-l2}
两条等价证书口径：
\begin{enumerate}
  \item（连续型）过阻尼 Langevin 的 Fokker--Planck 方程
  $
    \partial_t q=\nabla\cdot(q\nabla\Phi)+\frac{1}{\beta}\Delta q
  $
  的零流解为 $q_\beta\propto e^{-\beta\Phi}$；
  \item（离散型）Metropolis 核满足 $q_\beta(\theta)K(\theta,d\theta')=q_\beta(\theta')K(\theta',d\theta)$ 的详细平衡，从而 $q_\beta$ 不变。
\end{enumerate}
\end{Certificate}

\begin{proof}[证明（谓词因果链）]
\begin{enumerate}
  \item（稳态）按证书~\ref{cert:tex27-ch8-l2}，无论连续型（Fokker--Planck）或离散型（详细平衡），$q_\beta$ 均为不变测度；
  \item（收敛）在适当的遍历/耗散条件下，不变测度唯一且 $q_t\Rightarrow q_\beta$（长期收敛到稳态）；
  \item（等价于自由能最小化）由定理~\ref{thm:tex27-ch8-t1}，$q_\beta$ 是 $\mathcal{F}_\beta$ 的唯一极小点，因此“长期演化趋稳态”可解释为“把 $q$ 推向自由能极小点”。
\end{enumerate}
证毕。
\end{proof}

\subsection{命题（强命题）：演化同伦等价与量子提升}
\label{subsec:tex27-ch8-props}

\begin{proposition}[P1（经典情形：不断同伦扭曲 $\simeq$ 统计力学变分优化）]
\label{prop:tex27-ch8-p1}
设“参数不断同伦扭曲”满足：
\begin{enumerate}
  \item（连续可插值）离散更新 $T_{k}$ 可作分段连续插值成流 $T_t$；
  \item（遍历/可达）边缘算子完备性给出足够的可达性（至少在相关连通分支内）；
  \item（能量一致）更新的势能与统计力学能量同取 $\Phi$；
  \item（热噪声或接受率）存在噪声/接受机制使分布稳态为 $q_\beta$。
\end{enumerate}
令 $\mathcal{X}:=\mathsf{Prob}(\Theta)$ 表示 $\Theta$ 上概率测度空间（取弱拓扑），定义嵌入
\[
  h:\Theta\to\mathcal{X},\qquad h(\theta)=\delta_\theta.
\]
令 $\{S_t\}$ 为由 Fokker--Planck / Markov 核诱导的分布演化半群。
则在定义~\ref{def:tex27-ch8-d5} 的口径下，有
\[
  h\circ T_t \simeq S_t\circ h,
\]
并且 $T_t$ 的优化极限集合与 $S_t$ 的自由能极小点（及其零温极限的 $\arg\min\Phi$）在同伦意义下对应。
\end{proposition}

\begin{Certificate}[证书 P1（Dirac 嵌入 + 噪声强度同伦）]
\label{cert:tex27-ch8-p1}
用噪声强度参数 $\lambda\in[0,1]$ 构造一族插值演化（半群）$\{S_t^{(\lambda)}\}$：
\[
  d\theta_t^{(\lambda)}=-\nabla\Phi(\theta_t^{(\lambda)})\,dt+\sqrt{2\lambda/\beta}\,dW_t,
  \qquad
  S_t^{(\lambda)}(\delta_\theta)=\mathcal{L}(\theta_t^{(\lambda)}\mid \theta_0^{(\lambda)}=\theta),
\]
则 $\lambda=0$ 时 $S_t^{(0)}(\delta_\theta)=\delta_{T_t(\theta)}$，
$\lambda=1$ 时 $S_t^{(1)}$ 为目标随机演化；
由 $\lambda$ 的连续性可给出同伦 $H_t(\theta,\lambda):=S_t^{(\lambda)}(\delta_\theta)$。
\end{Certificate}

\begin{proof}[证明（谓词因果链）]
\begin{enumerate}
  \item（构造同伦）按证书~\ref{cert:tex27-ch8-p1}，令 $H_t(\theta,\lambda):=S_t^{(\lambda)}(\delta_\theta)$。
  则 $H_t(\cdot,0)=S_t^{(0)}\circ h=h\circ T_t$ 且 $H_t(\cdot,1)=S_t^{(1)}\circ h=S_t\circ h$，
  故 $h\circ T_t\simeq S_t\circ h$；
  \item（吸引子对应：有限温）由引理~\ref{lem:tex27-ch8-l2}，$S_t$ 的不变测度为 $q_\beta$，且在遍历条件下为全局吸引子；
  由定理~\ref{thm:tex27-ch8-t1}，该吸引子等价于自由能极小点；
  \item（吸引子对应：零温）由引理~\ref{lem:tex27-ch8-l1} 与定理~\ref{thm:tex27-ch8-t2}，
  当接受机制退化为门控（或 $\beta\to\infty$）时，两者的极限集合由 $\arg\min\Phi$ 支配，从而在“极小集合/可达极限集”口径下对应。
\end{enumerate}
证毕。
\end{proof}

\begin{proposition}[P2（量子情形：非交换边缘算子 $\Rightarrow$ 量子统计力学的自然提升）]
\label{prop:tex27-ch8-p2}
若边缘算子族 $\{B_i\}\subset\mathcal{T}$ 满足非交换性 $B_i\circ B_j\neq B_j\circ B_i$，
并且存在生成元分解使每一步更新可写为小步指数
\[
  B_i \approx \exp(-\Delta\tau\,\widehat{G}_i),
  \qquad
  \Delta\tau=\tau/n,
\]
则时间有序复合
\[
  T_{0\to\tau}^{(n)}:=\prod_{k=1}^{n}\exp(-\Delta\tau\,\widehat{G}_{i_k})
\]
在 Trotter（Lie--Trotter）分解意义下逼近虚时间演化
\[
  T_{0\to\tau}^{(n)}\xrightarrow[n\to\infty]{} e^{-\tau\widehat{H}},
  \qquad
  \widehat{H}:=\sum_i \widehat{G}_i,
\]
从而把“非交换的同伦扭曲”解释为“量子虚时间演化”；并由定理~\ref{thm:tex27-ch8-t3}，
其对应的变分目标自然由量子自由能最小化给出。
\end{proposition}

\begin{Certificate}[证书 P2（Lie--Trotter 乘积公式证书）]
\label{cert:tex27-ch8-p2}
在合适的有界性/定义域条件下，Lie--Trotter 乘积公式给出
\[
  \Bigl(\prod_{i=1}^m e^{-\Delta\tau\,\widehat{G}_i}\Bigr)^{n}
  \xrightarrow[n\to\infty]{}
  e^{-\tau\sum_{i=1}^m \widehat{G}_i}.
\]
\end{Certificate}

\begin{proof}[证明（谓词因果链）]
\begin{enumerate}
  \item（乘积逼近）按证书~\ref{cert:tex27-ch8-p2}，当 $\Delta\tau=\tau/n\to 0$ 时，小步指数乘积在极限下收敛到 $e^{-\tau\widehat{H}}$；
  \item（顺序敏感 $\Leftrightarrow$ 非交换）若 $\widehat{G}_i$ 非交换，则不同的时间有序复合给出不同的离散近似路径，
  这与量子算符乘积的顺序敏感一致；
  \item（变分意义）由定理~\ref{thm:tex27-ch8-t3}，量子虚时间演化的平衡态由 $\rho_\beta\propto e^{-\beta\widehat{H}}$ 给出，
  等价于量子自由能最小化。证毕。
\end{enumerate}
\end{proof}

\subsection{推论与备注：门控/退火/路径视角}
\label{subsec:tex27-ch8-others}

\begin{corollary}[C1（门控 $\Delta\Phi<0$ 的系统 $\approx$ 零温极限；随机接受率 $\approx$ 有温自由能优化）]
\label{cor:tex27-ch8-c1}
\begin{enumerate}
  \item 若坚持严格门控（只接受 $\Delta\Phi<0$），则对应 $\beta\to\infty$ 的极限：收敛强但越障能力弱（易陷入局部结构）；
  \item 若允许以概率 $p(\Delta\Phi)=\min\{1,e^{-\beta\Delta\Phi}\}$ 接受上升（$\Delta\Phi>0$），则获得有限温遍历性，
  长期收敛到 $q_\beta$，等价于最小化 $\mathcal{F}_\beta$。
\end{enumerate}
\end{corollary}

\begin{Certificate}[证书 C1（L1 + L2 的直接并列引用）]
\label{cert:tex27-ch8-c1}
门控 $\Leftrightarrow$ $\beta\to\infty$ 的退化见证书~\ref{cert:tex27-ch8-l1}；
有限温接受率/噪声 $\Rightarrow q_\beta$ 稳态见证书~\ref{cert:tex27-ch8-l2}；
而 $q_\beta$ 的最优化意义由定理~\ref{thm:tex27-ch8-t1} 固化。
\end{Certificate}

\begin{proof}[证明（谓词因果链）]
由证书~\ref{cert:tex27-ch8-c1} 给出的三条并列引用，推论两点分别由引理~\ref{lem:tex27-ch8-l1} 与引理~\ref{lem:tex27-ch8-l2} 推出。证毕。
\end{proof}

\begin{corollary}[C2（“不断同伦扭曲”可视为隐式的虚时间路径优化/路径积分离散化）]
\label{cor:tex27-ch8-c2}
在命题~\ref{prop:tex27-ch8-p2} 的条件下，算子序列 $(i_1,\dots,i_n)$ 给出一条离散“路径”，
其时间有序乘积 $T_{0\to\tau}^{(n)}$ 逼近 $e^{-\tau\widehat{H}}$。
当能量权重与 $\widehat{H}$（或其对角部分 $\Phi$）一致时，该离散路径结构与虚时间路径积分表述同构，
从而可用统计力学工具设计更稳的变分调度（退火 $\beta(t)$、分层温度、分区能量项等）。
\end{corollary}

\begin{Certificate}[证书 C2（Trotter 化 $\Rightarrow$ 路径集合）]
\label{cert:tex27-ch8-c2}
Trotter 乘积把 $e^{-\tau\widehat{H}}$ 表示为“小步指数的乘积极限”，每一个离散乘积对应一个算子序列（路径）；
在合适基底展开后，可把该乘积视为对路径集合的加权求和，从而得到“路径积分/路径优化”的解释口径。
\end{Certificate}

\begin{proof}[证明（谓词因果链）]
由命题~\ref{prop:tex27-ch8-p2} 与证书~\ref{cert:tex27-ch8-c2}，$e^{-\tau\widehat{H}}$ 被离散乘积逼近，
离散乘积的索引序列即路径；在对角化/基底展开时，乘积展开对应对路径项的加权汇总。
因此可把“同伦扭曲的算子序列”解释为虚时间路径优化的离散实现。证毕。
\end{proof}

\begin{remark}[边界与适用条件（对“同伦等价”的口径澄清）]
\label{rem:tex27-ch8-r1}
\begin{enumerate}
  \item 本节的“同伦等价”是对\textbf{演化结构}而言：强调可达集/吸引子/极限集的一致性，以及通过噪声强度/步长等参数的连续变形把一种演化连到另一种演化；
  若要求严格的空间同伦型一致，还需额外规定 $\Theta$ 的边界条件、禁区与更新算子的连续性/可逆性口径。
  \item “量子”并不必须来自物理系统；在算法层面，“非交换更新 + 虚时间演化 + Gibbs 变分原理”即可纯数学地刻画：
  \[
    \text{非交换边缘算子序列}\ \Longrightarrow\ \text{算子代数}\ \Longrightarrow\ \widehat{H}\ \text{与}\ \rho_\beta.
  \]
  \item 工程落地的关键不是把名词叫“量子”，而是实现：
  \[
    \text{在可控的 }\beta(t)\text{ 下兼具可达遍历与自由能下降（或其证书化替代）。}
  \]
  满足它，就可移植统计力学变分优化工具箱（退火、温度分层、能量分解与稳定化调度等）。
\end{enumerate}
\end{remark}

\begin{remark}[结论（对“强命题”的直接口径）]
\label{rem:tex27-ch8-r2}
把“参数不断同伦扭曲”建模为由边缘算子生成的（可能非交换的）演化，并令能量函数取为缺陷势能 $\Phi$ 后：
\begin{enumerate}
  \item 零温（门控）情形：同伦等价于“只做能量下降的变分极限”；
  \item 有温（含噪/接受率）情形：同伦等价于“自由能泛函的变分优化”，其稳态为 $q_\beta$；
  \item 非交换提升：时间有序复合在 Trotter 意义下对应量子虚时间演化，从而自然提升为“量子统计力学实现的变分优化”。
\end{enumerate}
\end{remark}

\begin{remark}[与 ODTHB 的关系：把 $\beta(t)$ 调度落到可审计接口]
本节反复出现的“在可控的 $\beta(t)$ 下兼具可达遍历与自由能下降”可在第~\ref{subsec:tex27-ch5-odthb}~节 ODTHB 接口中落地：
将 $\beta,\theta,\ell$ 视为慢变量，把 $p_t(\cdot\mid u)$ 视为快变量提议核；
对应的源--调节分离与自然演化解释见命题~\ref{prop:tex27-ch5-odthb-r1} 与推论~\ref{cor:tex27-ch5-odthb-r3}。
\end{remark}

%%%%%%%%%%%%%%%%%%%%%%%%%%%%%%%%%%%%%%%%%%%%%%%%%%%%%%%%%%%%
\section{ORL、可审计白盒 AI 与形式化 AGI：条件定理链（强命题）}
\label{sec:tex27-ch9-orl-whitebox-agi}
%%%%%%%%%%%%%%%%%%%%%%%%%%%%%%%%%%%%%%%%%%%%%%%%%%%%%%%%%%%%

本节给出一条\textbf{可形式化的强命题}，把“ORL $\simeq$ 可审计白盒 AI $\simeq$ AGI”写成\textbf{条件定理链}：
该链条在\textbf{您指定的}变分空间 $\mathcal{V}$、算子全集 $\mathcal{O}$ 与任务类 $\mathfrak{T}(\mathcal{V})$ 的形式系统内成立；
若把“AGI”理解为现实世界的宽泛概念，则仍属于\textbf{未证明外推}（需要外部证书与实证支撑）。

\subsection{定义组：变分空间、同伦塔、ORL 与可审计/形式化 AGI}
\label{subsec:tex27-ch9-defs}

\begin{definition}[D1（变分空间与缺陷势能）]
\label{def:tex27-ch9-d1}
令变分空间（状态空间）
\[
  \mathcal{V}\quad(\text{可取流形/格点/离散结构的直积}),
  \qquad
  s\in\mathcal{V}.
\]
给定缺陷势能（障碍度量）
\[
  \Phi:\mathcal{V}\to \RR_{\ge 0}.
\]
给定约束/不变量谓词 $\mathrm{Inv}(s)$。
定义“障碍”谓词
\[
  \mathrm{Obs}(s)\;:\!\iff\;\bigl(\Phi(s)>0\bigr)\ \vee\ \bigl(\neg \mathrm{Inv}(s)\bigr).
\]
\end{definition}

\begin{definition}[D2（同伦塔与边缘算子集）]
\label{def:tex27-ch9-d2}
令算子字母表（边缘算子全集）为同伦塔分层并：
\[
  \mathcal{O}=\mathcal{O}^{(0)}\cup \mathcal{O}^{(1)}\cup \cdots,
\]
其中 $\mathcal{O}^{(k)}$ 表示第 $k$ 层修复自由度。
每个原子算子 $o\in\mathcal{O}$ 是一个（部分）状态变换
\[
  o:\mathcal{V}\rightharpoonup \mathcal{V},
\]
并携带可检验的前置/后置谓词（例如 $\mathrm{Pre}_o(s)$ 与 $\mathrm{Post}_o(s,s')$）。

令有限复合词空间
\[
  \mathcal{O}^{*}
  :=\{o_1\circ\cdots\circ o_m:\ m\in\NN,\ o_i\in\mathcal{O}\}.
\]
若 $w=o_1\circ\cdots\circ o_m\in\mathcal{O}^*$，则其证书化的可执行性可由逐步前置谓词校验保证。
\end{definition}

\begin{definition}[D3（历史无关：Markov 策略的精确定义）]
\label{def:tex27-ch9-d3}
称策略“历史无关”（Markov），当且仅当存在映射
\[
  \pi:\mathcal{V}\to\mathcal{O}^{*},
\]
使得任意时刻动作词由当前状态决定：$w_t=\pi(s_t)$。
\end{definition}

\begin{definition}[D4（即时修复：变分准则）]
\label{def:tex27-ch9-d4}
称策略满足“即时修复”，当且仅当对任意满足 $\mathrm{Obs}(s)$ 的状态 $s$，
其输出 $w\in\mathcal{O}^*$ 满足：
\begin{enumerate}
  \item（约束保持）若 $\mathrm{Inv}(s)$，则 $\mathrm{Inv}(w(s))$；
  \item（变分下降，零温门控口径）定义 $\Delta\Phi:=\Phi(w(s))-\Phi(s)$，并要求 $\Delta\Phi<0$；
  \item（统计力学口径，可选）允许非单调单步，但要求在某逆温度 $\beta$ 下呈现统计意义的负漂移/自由能下降口径（见定义~\ref{def:tex27-ch9-d6} 与引理~\ref{lem:tex27-ch9-l3}）。
\end{enumerate}
\end{definition}

\begin{definition}[D5（ORL：算子强化学习的白盒定义）]
\label{def:tex27-ch9-d5}
ORL 在本节定义为三元组
\[
  \mathrm{ORL}:=(\mathcal{O},\Phi,\pi),
\]
其中 $\pi$ 的动作空间是算子词空间 $\mathcal{O}^*$。
令候选集 $\mathrm{Cand}(s)\subseteq\mathcal{O}^*$ 表示在状态 $s$ 上可执行且证书可验证的动作词集合。
典型两种口径：
\begin{enumerate}
  \item（零温/严格门控）
  \[
    \pi(s)\in \operatorname*{arg\,min}_{w\in\mathrm{Cand}(s)}\ \Phi\bigl(w(s)\bigr),
    \qquad
    \Phi(\pi(s)(s))-\Phi(s)<0;
  \]
  \item（有限温/统计接受）
  \[
    \mathbb{P}(w\mid s)\propto \exp\bigl(-\beta\,\Phi(w(s))\bigr),
    \qquad
    \beta>0.
  \]
\end{enumerate}
\end{definition}

\begin{definition}[D6（统计力学/量子统计力学的同伦等价口径：把“词”视为路径）]
\label{def:tex27-ch9-d6}
把动作词 $w\in\mathcal{O}^*$ 视为路径（path），定义路径能量
\[
  E(w;s):=\Phi(w(s)).
\]
经典统计力学的路径分布为 Gibbs：
\[
  p_\beta(w\mid s)=\frac{1}{Z(s)}e^{-\beta E(w;s)}.
\]
若原子算子非交换，则“时间有序复合”的极限可在 Trotter 口径下与虚时间指数演化对应（参见第~\ref{sec:tex27-ch8-homotopy-statmech}~节中命题~\ref{prop:tex27-ch8-p2}）。
本节只把该对应作为“形式系统内的演化同伦等价口径”，用于解释“有限温/非交换”下的统计意义有效性。
\end{definition}

\begin{definition}[D7（可审计白盒 AI）]
\label{def:tex27-ch9-d7}
称系统为“可审计白盒 AI”，若它在每一步输出
\[
  (\text{action}=w,\ \text{certificate}=C),
  \qquad
  w\in\mathcal{O}^*,
\]
且证书 $C$ 至少包含：
\begin{enumerate}
  \item（轨迹）$w=o_1\circ\cdots\circ o_m$ 的显式分解；
  \item（谓词证明）每个 $o_i$ 的前置条件在当前状态成立，以及后置条件成立；
  \item（变分证据）零温时给出 $\Delta\Phi<0$；有限温时给出接受概率/自由能差/可复算的统计量；
  \item（不变量）$\mathrm{Inv}$ 的保持证明/校验记录。
\end{enumerate}
\end{definition}

\begin{definition}[D8（形式化 AGI：$\mathrm{AGI}^{\mathrm{cert}}_{\mathcal{V},\mathcal{O}}$）]
\label{def:tex27-ch9-d8}
为使“等价链”可证明，定义 AGI 为一个\textbf{相对概念}（只关于 $(\mathcal{V},\mathcal{O})$）。
称系统是
\[
  \mathrm{AGI}^{\mathrm{cert}}_{\mathcal{V},\mathcal{O}},
\]
当且仅当对任意可计算任务描述
\[
  (\Phi,\mathrm{Inv})\in\mathfrak{T}(\mathcal{V}),
\]
系统都能在资源预算内构造一条（或一族）算子词序列，使得：
\begin{enumerate}
  \item 约束始终可审计（全程输出证书可回放）；
  \item 变分意义下持续修复障碍（零温单调下降，或有限温统计意义下降）；
  \item 证书接口在每一步可复核（对应定义~\ref{def:tex27-ch9-d7} 的最小接口）。
\end{enumerate}
该定义并不等同于“现实世界 AGI 的全部含义”，仅用于形式系统内闭合推理。
\end{definition}

\subsection{公理（降级为定理）：历史无关化与算子完备性前提}
\label{subsec:tex27-ch9-axioms}

\begin{theorem}[公理（降级为定理）T0（历史无关化：历史依赖策略可在扩展状态上等价为 Markov）]
\label{thm:tex27-ch9-t0}
对任意历史依赖策略 $\pi_H$（依赖历史 $h_t=(s_0,\dots,s_t)$），
存在扩展状态空间 $\mathcal{V}'$ 与 Markov 策略 $\pi'$，
使得在投影回原状态空间时，两者诱导的外显动作序列同轨道（等价）。
\end{theorem}

\begin{Certificate}[证书 T0（扩展状态构造：把历史压缩为充分统计量）]
\label{cert:tex27-ch9-t0}
构造扩展状态空间
\[
  \mathcal{V}' := \mathcal{V}\times \mathcal{M},
\]
其中 $\mathcal{M}$ 保存策略所需的历史充分统计量（例如累计量、过滤器状态、后验参数等）。
给定记忆更新规则 $m_{t+1}=u(m_t,s_t)$，即可把 $(s_t,m_t)$ 作为新的 Markov 状态。
\end{Certificate}

\begin{proof}[证明（谓词因果链）]
令原策略为 $\pi_H(h_t)$，依赖历史 $h_t=(s_0,\dots,s_t)$。
按证书~\ref{cert:tex27-ch9-t0}，构造扩展状态 $(s_t,m_t)\in\mathcal{V}'$，其中 $m_t$ 为历史充分统计量且满足 $m_{t+1}=u(m_t,s_t)$。
定义新策略
\[
  \pi'(s_t,m_t):=\pi_H(h_t).
\]
则 $\pi'$ 仅依赖当前扩展状态 $(s_t,m_t)$，因此是 Markov。
并且由于 $m_t$ 由 $(m_0,s_0,\dots,s_{t-1})$ 递推确定，$\pi'(s_t,m_t)$ 与 $\pi_H(h_t)$ 在每一步输出相同的动作词，
故两者在外显动作序列上同轨道。证毕。
\end{proof}

\begin{theorem}[公理（降级为定理）A9（边缘算子完备性：任意障碍态存在下降修复词）]
\label{thm:tex27-ch9-a9}
令任务类为 $\mathfrak{T}(\mathcal{V})$。假设如下“完备性”成立：
\[
  (\mathrm{Complete})\quad
  \forall(\Phi,\mathrm{Inv})\in\mathfrak{T}(\mathcal{V}),\ \forall s\in\mathcal{V},\
  \mathrm{Obs}(s)\Rightarrow \exists w\in\mathcal{O}^*\ \text{s.t. }\mathrm{Inv}(w(s))\wedge \Phi(w(s))<\Phi(s).
\]
\end{theorem}

\begin{Certificate}[证书 A9（外部证书输入：覆盖性规则/测试/生成性证明）]
\label{cert:tex27-ch9-a9}
本条作为“证书输入接口”：在工程上可由如下任一类证书支撑：
\begin{enumerate}
  \item（规则覆盖证书）列出障碍谓词的有限分类，并对每类给出可执行修复词模板；
  \item（形式生成证书）证明 $\mathcal{O}$ 的生成闭包覆盖所有需要的修复方向；
  \item（测试覆盖证书）对障碍态分布给出可复现实验，验证对每类障碍态都能构造下降修复词。
\end{enumerate}
本笔记不在此处重建该外部证书。
\end{Certificate}

\begin{proof}[证明（证书化）]
按证书~\ref{cert:tex27-ch9-a9}，A9 作为外部已给“完备性证书输入”被调用。证毕。
\end{proof}

\subsection{引理：ORL、白盒可审计与统计意义有效性}
\label{subsec:tex27-ch9-lemmas}

\begin{lemma}[L1（ORL $\Rightarrow$ 可审计白盒 AI）]
\label{lem:tex27-ch9-l1}
若系统满足定义~\ref{def:tex27-ch9-d5} 的 ORL 结构，
并且每个原子算子携带可检验前置/后置谓词（定义~\ref{def:tex27-ch9-d2}），
则系统满足定义~\ref{def:tex27-ch9-d7} 的可审计白盒 AI。
\end{lemma}

\begin{Certificate}[证书 L1（最小证书接口由 ORL 直接给出）]
\label{cert:tex27-ch9-l1}
对任意状态 $s$ 与动作词 $w=o_1\circ\cdots\circ o_m$，定义证书接口
\[
  C(s,w):=\bigl(o_1,\dots,o_m;\ \text{Pre/Post checks};\ \Delta\Phi\ \text{或}\ p_\beta(w\mid s);\ \mathrm{Inv\ checks}
    \bigr).
\]
\end{Certificate}

\begin{proof}[证明（谓词因果链）]
\begin{enumerate}
  \item（轨迹可列出）ORL 的动作是 $w\in\mathcal{O}^*$ 的显式复合，因此 $w=o_1\circ\cdots\circ o_m$ 可直接输出；
  \item（谓词可校验）每个原子算子携带 $\mathrm{Pre}_o,\mathrm{Post}_o$，故可逐步检查前置成立与后置成立；
  \item（变分证据可复算）零温门控下可直接给出 $\Delta\Phi<0$；有限温下可给出 Gibbs 权重/接受概率 $p_\beta(w\mid s)$；
  \item（不变量可审计）$\mathrm{Inv}$ 的保持由逐步校验记录给出。
\end{enumerate}
以上四项正是定义~\ref{def:tex27-ch9-d7} 的最小证书接口，故 ORL $\Rightarrow$ 可审计白盒 AI。证毕。
\end{proof}

\begin{lemma}[L2（可审计白盒 AI $\Rightarrow$ ORL）]
\label{lem:tex27-ch9-l2}
对任意可审计白盒系统 $A$，存在一个 ORL 实例 $(\mathcal{O},\Phi_A,\pi_A)$，
使其在“外显动作序列 + 可复核证书接口”的意义下与 $A$ 等价。
\end{lemma}

\begin{Certificate}[证书 L2（用证书定义可复算评分泛函）]
\label{cert:tex27-ch9-l2}
对白盒系统 $A$ 的每一步输出 $(w,C)$，定义可复算评分泛函
\[
  \Phi_A(s,w):=
  \begin{cases}
    +\infty, & \text{若证书显示 }w\text{ 不可执行或破坏 }\mathrm{Inv},\\
    \Phi(w(s)), & \text{否则}.
  \end{cases}
\]
并令 $\mathrm{Cand}(s)$ 取为白盒系统实际枚举/考虑的动作词集合。
\end{Certificate}

\begin{proof}[证明（谓词因果链）]
\begin{enumerate}
  \item（候选集可确定）白盒系统 $A$ 以可审计方式枚举/选择动作词 $w$，故其“考虑过的动作集合”可作为 $\mathrm{Cand}(s)$；
  \item（评分可复算）由证书可判定 $w$ 的可执行性与不变量保持，因此 $\Phi_A(s,w)$ 可复算；
  \item（等价的策略定义）定义 ORL 策略 $\pi_A$ 为：在同一候选集上选择使 $\Phi_A$ 最小的 $w$（或按有限温权重采样）。
  由于 $\mathrm{Cand}(s)$ 与评分规则复刻了白盒系统的可复核选择过程，故 $\pi_A$ 与 $A$ 在外显动作上同轨道，并共享同一证书接口。
\end{enumerate}
因此可构造 ORL 与白盒系统在“动作 + 证书”口径下等价。证毕。
\end{proof}

\begin{lemma}[L3（统计力学同伦等价 $\Rightarrow$ “有效性是统计意义而非逐步单调”）]
\label{lem:tex27-ch9-l3}
在有限温口径（定义~\ref{def:tex27-ch9-d6}）下，允许采样到单步不下降甚至上升的 $w$，
但在遍历/混合条件下，长期演化呈现向 Gibbs 极限的收敛，并可解释为自由能意义的下降。
\end{lemma}

\begin{Certificate}[证书 L3（自由能 = 期望能量 + 熵项）]
\label{cert:tex27-ch9-l3}
有限温 Gibbs 权重 $p_\beta(w\mid s)\propto e^{-\beta E(w;s)}$ 对应的变分目标可写为
\[
  \mathcal{F}_\beta[p]
  =\mathbb{E}_{w\sim p}[E(w;s)]+\frac{1}{\beta}\mathrm{KL}(p\|p_0),
\]
其中 $p_0$ 为参考分布；极小点为 Gibbs 分布（与第~\ref{sec:tex27-ch8-homotopy-statmech}~节定理链同构）。
\end{Certificate}

\begin{proof}[证明（谓词因果链）]
\begin{enumerate}
  \item（允许非单调）有限温下 $p_\beta(w\mid s)$ 对所有候选 $w$ 给出正概率，因此单步可能出现 $E(w;s)$ 增大；
  \item（长期趋稳）在遍历/混合条件下，分布演化收敛到 Gibbs 型不变分布；
  \item（自由能解释）按证书~\ref{cert:tex27-ch9-l3}，Gibbs 分布是自由能泛函的极小点，
  因而“分布趋稳态”的过程可解释为自由能意义下的下降与稳定化。
\end{enumerate}
故“有效性”应理解为统计意义有效，而非逐步单调有效。证毕。
\end{proof}

\subsection{命题：等价链与（形式化）AGI 的条件推出}
\label{subsec:tex27-ch9-props}

\begin{proposition}[P1（ORL $\Longleftrightarrow$ 可审计白盒 AI）]
\label{prop:tex27-ch9-p1}
在“外显动作序列 + 可复核证书接口”的口径下，
\[
  \mathrm{ORL}\ \Longleftrightarrow\ \text{可审计白盒 AI}.
\]
\end{proposition}

\begin{Certificate}[证书 P1（L1 与 L2 的并列引用）]
\label{cert:tex27-ch9-p1}
ORL $\Rightarrow$ 白盒由引理~\ref{lem:tex27-ch9-l1} 给出；
白盒 $\Rightarrow$ ORL 由引理~\ref{lem:tex27-ch9-l2} 给出。
\end{Certificate}

\begin{proof}[证明（谓词因果链）]
由证书~\ref{cert:tex27-ch9-p1} 的两条并列引用，得到双向蕴含，从而 ORL 与可审计白盒 AI 等价。证毕。
\end{proof}

\begin{proposition}[P2（算子完备性前提下：ORL/白盒 AI $\Rightarrow \mathrm{AGI}^{\mathrm{cert}}_{\mathcal{V},\mathcal{O}}$）]
\label{prop:tex27-ch9-p2}
在公理~\ref{thm:tex27-ch9-a9}（$\mathrm{Complete}$）成立且采用定义~\ref{def:tex27-ch9-d8} 的形式化 AGI 口径时，
若系统为 ORL（或等价地：可审计白盒 AI），则系统满足 $\mathrm{AGI}^{\mathrm{cert}}_{\mathcal{V},\mathcal{O}}$。
\end{proposition}

\begin{Certificate}[证书 P2（“每步可下降修复词存在”作为通用性证书输入）]
\label{cert:tex27-ch9-p2}
公理~\ref{thm:tex27-ch9-a9} 给出：对任意任务 $(\Phi,\mathrm{Inv})$ 与任意障碍态 $s$，存在保持 $\mathrm{Inv}$ 且严格下降 $\Phi$ 的修复词 $w$。
这条“存在性”即为把“算法形式”提升为“相对通用性”的关键证书输入。
\end{Certificate}

\begin{proof}[证明（谓词因果链）]
取任意任务 $(\Phi,\mathrm{Inv})\in\mathfrak{T}(\mathcal{V})$。
\begin{enumerate}
  \item（可审计性）若系统为 ORL，则由引理~\ref{lem:tex27-ch9-l1}，每一步都可输出动作词分解与可复核证书，因此满足定义~\ref{def:tex27-ch9-d8} 的“可审计”要求；
  \item（持续修复：零温）若采用零温门控，则对任意 $\mathrm{Obs}(s)$，
  由公理~\ref{thm:tex27-ch9-a9} 存在可行下降词 $w$，而 ORL 在候选集中选择（至少）不劣于该 $w$ 的词，
  因此产生 $\Delta\Phi<0$ 的修复动作，形成单调下降的修复链；
  \item（持续修复：有限温）若采用有限温统计接受，则允许单步不下降，但由引理~\ref{lem:tex27-ch9-l3}，
  长期演化可在自由能意义下解释为统计意义的下降与稳定化。
\end{enumerate}
因此 ORL（等价地：白盒 AI）满足定义~\ref{def:tex27-ch9-d8} 的三条要求，从而推出 $\mathrm{AGI}^{\mathrm{cert}}_{\mathcal{V},\mathcal{O}}$。证毕。
\end{proof}

\begin{proposition}[P3（$\mathrm{AGI}^{\mathrm{cert}}_{\mathcal{V},\mathcal{O}} \Rightarrow$ 可审计白盒 AI）]
\label{prop:tex27-ch9-p3}
若系统满足定义~\ref{def:tex27-ch9-d8} 的 $\mathrm{AGI}^{\mathrm{cert}}_{\mathcal{V},\mathcal{O}}$，
则系统满足定义~\ref{def:tex27-ch9-d7} 的可审计白盒 AI。
\end{proposition}

\begin{Certificate}[证书 P3（定义蕴含证书）]
\label{cert:tex27-ch9-p3}
定义~\ref{def:tex27-ch9-d8} 的第 1、3 条要求“全程输出证书可回放/证书接口可复核”，即定义~\ref{def:tex27-ch9-d7}。
\end{Certificate}

\begin{proof}[证明]
由证书~\ref{cert:tex27-ch9-p3}，$\mathrm{AGI}^{\mathrm{cert}}_{\mathcal{V},\mathcal{O}}$ 的定义直接包含白盒可审计性要求，故命题成立。证毕。
\end{proof}

\subsection{推论与备注：等价链的边界与工程证书化方向}
\label{subsec:tex27-ch9-others}

\begin{corollary}[C1（形式系统内的等价链）]
\label{cor:tex27-ch9-c1}
在公理~\ref{thm:tex27-ch9-a9}（$\mathrm{Complete}$）成立且采用定义~\ref{def:tex27-ch9-d8} 的 AGI 口径时，有等价链：
\[
  \mathrm{ORL}\ \Longleftrightarrow\ \text{可审计白盒 AI}\ \Longleftrightarrow\ \mathrm{AGI}^{\mathrm{cert}}_{\mathcal{V},
    \mathcal{O}}.
\]
\end{corollary}

\begin{Certificate}[证书 C1（P1 + P2 + P3 组合）]
\label{cert:tex27-ch9-c1}
前两者等价由命题~\ref{prop:tex27-ch9-p1}；
ORL/白盒 $\Rightarrow$ AGI 由命题~\ref{prop:tex27-ch9-p2}；
AGI $\Rightarrow$ 白盒由命题~\ref{prop:tex27-ch9-p3}。
\end{Certificate}

\begin{proof}[证明（谓词因果链）]
按证书~\ref{cert:tex27-ch9-c1} 串联三条命题即可得到等价链。证毕。
\end{proof}

\begin{remark}[边界 1：此处“等价于 AGI”是相对命题，非现实世界外推]
\label{rem:tex27-ch9-r1}
本节的“等价于 AGI”指的是相对于您指定的 $(\mathcal{V},\mathcal{O})$ 与任务类 $\mathfrak{T}(\mathcal{V})$ 的
$\mathrm{AGI}^{\mathrm{cert}}_{\mathcal{V},\mathcal{O}}$。
若要外推为“现实世界 AGI（跨一切物理/社会/工程任务）”，至少需要额外证书：
证明现实任务可编码进某个 $\mathcal{V}$，且 $\mathcal{O}$ 在工程上具备所需完备性与资源可行性；
该外推在本节不被证明。
\end{remark}

\begin{remark}[边界 2：“不依赖历史”的严格口径是“历史被吸收进状态”]
\label{rem:tex27-ch9-r2}
您强调“即时修复不依赖历史，只取决于变分空间”的严格形式是定理~\ref{thm:tex27-ch9-t0}：
任何历史依赖策略可在扩展状态空间上等价为 Markov 策略。
这与“白盒可审计”相容：扩展状态中的充分统计量同样是可审计对象（可被证书化记录与回放）。
\end{remark}

\begin{remark}[边界 3：统计力学同伦等价只改变“有效性口径”，不削弱白盒性]
\label{rem:tex27-ch9-r3}
统计力学同伦等价（定义~\ref{def:tex27-ch9-d6}，引理~\ref{lem:tex27-ch9-l3}）的作用是：
把“有效性”从“每步必须更好”改成“自由能/期望意义更好”。
白盒性只要求“规则与统计证据可复算”，不要求“每一步都盈利/每一步都严格下降”。
\end{remark}

\begin{remark}[工程证书化提示：把 $\mathrm{Complete}$ 从假设变成可核验证书]
\label{rem:tex27-ch9-r4}
若要把本节强命题进一步工程化闭合，关键新增工作是把公理~\ref{thm:tex27-ch9-a9} 的 $\mathrm{Complete}$
从“假设”提升为“可核验证书”（规则覆盖/生成性证明/测试覆盖三选一或组合）。
一旦该证书到位，则“ORL $\simeq$ 白盒 AI $\simeq \mathrm{AGI}^{\mathrm{cert}}_{\mathcal{V},\mathcal{O}}$”
在形式系统内成为可审计结论链。
\end{remark}

%%%%%%%%%%%%%%%%%%%%%%%%%%%%%%%%%%%%%%%%%%%%%%%%%%%%%%%%%%%%
\section{时间换空间：从同伦扭曲的路径采样到 ORL、开放系统在线控制与相对 AGI（强命题链）}
\label{sec:tex27-ch10-time-space-tradeoff}
%%%%%%%%%%%%%%%%%%%%%%%%%%%%%%%%%%%%%%%%%%%%%%%%%%%%%%%%%%%%

本节把“时间换空间”严格化为一个\textbf{可证明的命题链}，并显式标注哪些结论是\textbf{本形式系统内的定理}，
哪些属于\textbf{未证明外推}（需要额外证书/外部实现/实证）。

\begin{remark}[本节结论的证明口径与未证明边界（先行标注）]
\label{rem:tex27-ch10-scope}
\begin{enumerate}
  \item（形式系统内可证）定理 T1--T3、命题 P1--P2、推论 C1 的证明均是\textbf{条件式}：在明确的可检验前提下成立；
  \item（外部证书输入）涉及“马尔可夫链混合/遍历速度”“对抗环境分布假设”“算子完备性覆盖”等，
  本节将以“公理（降级为定理）+ 外部证书输入”的方式标注；
  \item（未证明外推）“物理量子硬件带来复杂度级加速”“现实世界无边界 AGI”等命题不在本节证明范围内，
  仅作为备注明确为未证明断言。
\end{enumerate}
\end{remark}

\subsection{定义组：路径空间、同伦扭曲采样与自由能口径}
\label{subsec:tex27-ch10-defs}

\begin{definition}[D1（“时间换空间”）]
\label{def:tex27-ch10-d1}
给定巨大搜索空间 $\Omega$（例如同伦塔边缘算子词空间）
\[
  \Omega:=\mathcal{O}^{*}.
\]
若直接“空间枚举”需要存储/遍历 $\Omega$ 的指数规模信息，则“时间换空间”指：
只维护一个（或少量）当前样本/轨道 $w_t\in\Omega$ 及少量统计量，
通过时间步 $t$ 的迭代来近似实现对 $\Omega$ 的有效探索与优化。
即用\textbf{低空间复杂度}换取\textbf{高时间步数}以推动目标泛函（能量/自由能）下降。
\end{definition}

\begin{definition}[D2（同伦扭曲作为路径采样/演化）]
\label{def:tex27-ch10-d2}
令变分空间为 $\mathcal{V}$，状态 $s\in\mathcal{V}$，缺陷势能
\[
  \Phi:\mathcal{V}\to\RR_{\ge 0}.
\]
一次决策选择一个算子词
\[
  w=o_1\circ\cdots\circ o_m\in\mathcal{O}^{*},
\]
并得到新状态 $s' = w(s)$。
把“同伦扭曲”视为在路径空间 $\Omega=\mathcal{O}^*$ 上的演化/采样过程 $\{w_t\}_{t\ge 0}$，
并用其诱导状态演化 $s_{t+1}=w_t(s_t)$。
\end{definition}

\begin{definition}[D3（统计力学口径的目标：自由能最小化）]
\label{def:tex27-ch10-d3}
固定逆温度参数 $\beta>0$。在给定状态 $s$ 下，把每条路径 $w$ 的能量定义为
\[
  E(w;s):=\Phi(w(s)).
\]
若 $\Omega$ 可数，则定义 Gibbs 权重（经典统计力学）
\[
  p_\beta(w\mid s)=\frac{1}{Z(s)}e^{-\beta E(w;s)},
  \qquad
  Z(s)=\sum_{w\in\Omega}e^{-\beta E(w;s)}.
\]
（更一般地可用参考测度把求和替换为积分。）
“统计意义优化”的目标是推动路径分布逼近 $p_\beta(\cdot\mid s)$；
零温极限 $\beta\to\infty$ 对应能量极小路径的集中。
\end{definition}

\begin{definition}[D4（量子口径：非交换生成元与虚时间演化）]
\label{def:tex27-ch10-d4}
若原子算子非交换（$o_i\circ o_j\neq o_j\circ o_i$），可在形式层面引入生成元族 $\{\widehat{G}_i\}$ 与哈密顿量
\[
  \widehat{H}:=\sum_i \widehat{G}_i,
\]
并用 Trotter 乘积把虚时间演化 $e^{-\tau\widehat{H}}$ 近似为
\[
  \Bigl(\prod_i e^{-\Delta\tau\,\widehat{G}_{i}}\Bigr)^{n},
  \qquad
  \Delta\tau=\tau/n.
\]
这给出“非交换扭曲序列 $\simeq$ 虚时间指数演化”的\textbf{形式同伦等价/模拟等价}口径；
其并不自动蕴含“物理量子加速”（见备注~\ref{rem:tex27-ch10-r2}）。
\end{definition}

\subsection{公理（降级为定理）：混合/遍历作为外部证书输入}
\label{subsec:tex27-ch10-axioms}

\begin{theorem}[公理（降级为定理）A10（路径采样核的遍历与混合）]
\label{thm:tex27-ch10-a10}
存在一族（可能随状态 $s$ 变化的）路径采样马尔可夫核 $K_s(w,dw')$，
使得对固定 $s$，$p_\beta(\cdot\mid s)$ 是其不变分布，且链满足遍历/混合条件：
若 $W_0\sim P_0$，则 $P_t:=\mathcal{L}(W_t)\Rightarrow p_\beta(\cdot\mid s)$。
\end{theorem}

\begin{Certificate}[证书 A10（外部证书输入：MCMC/详细平衡/遍历性证明）]
\label{cert:tex27-ch10-a10}
本条作为外部证书输入，典型来源是：
\begin{enumerate}
  \item（详细平衡）构造满足 $p_\beta(w\mid s)K_s(w,dw')=p_\beta(w'\mid s)K_s(w',dw)$ 的核；
  \item（遍历性）证明链不可约、无周期，并满足适当的漂移/小集条件以保证收敛；
  \item（混合界）若需要定量的“时间换空间”效率，还需给出混合时间上界作为证书。
\end{enumerate}
\end{Certificate}

\begin{proof}[证明（证书化）]
按证书~\ref{cert:tex27-ch10-a10}，A10 作为外部已给“混合/遍历证书输入”被调用。证毕。
\end{proof}

\subsection{定理：时间换空间、ORL 统一实现与非交换量子口径}
\label{subsec:tex27-ch10-theorems}

\begin{theorem}[T1（“时间换空间”是统计力学采样/变分优化的直接推论）]
\label{thm:tex27-ch10-t1}
若 $\Omega=\mathcal{O}^*$ 极大而无法直接枚举/存储，
并且存在满足公理~\ref{thm:tex27-ch10-a10} 的路径采样核，
则用同伦扭曲序列 $\{w_t\}$ 的时间迭代可在\textbf{低空间代价}下逼近 Gibbs 分布（或其零温极限），
从而实现变分意义的优化（能量/自由能下降口径）。
\end{theorem}

\begin{Certificate}[证书 T1（空间代价对比：枚举 vs 采样）]
\label{cert:tex27-ch10-t1}
证书由两条对比组成：
\begin{enumerate}
  \item（枚举）若要显式比较大量候选 $w\in\Omega$，需要存储/遍历规模随 $|\Omega|$ 增长的信息；
  \item（采样）MCMC/路径采样仅需维护当前样本 $w_t$（及常数/低维统计量），空间复杂度与 $t$ 无关（或仅对数增长）。
\end{enumerate}
\end{Certificate}

\begin{proof}[证明（谓词因果链）]
\begin{enumerate}
  \item（低空间）由证书~\ref{cert:tex27-ch10-t1}，采样法仅维护 $w_t$ 与少量统计量，空间代价远低于枚举；
  \item（时间迭代逼近）由公理~\ref{thm:tex27-ch10-a10}，分布 $P_t$ 随时间迭代收敛到 $p_\beta(\cdot\mid s)$；
  \item（变分意义）按定义~\ref{def:tex27-ch10-d3}，$p_\beta$ 是自由能口径下的目标分布；
  零温极限 $\beta\to\infty$ 时质量集中到能量极小路径集合，故可把长期采样解释为“以时间逼近最优路径”。
\end{enumerate}
因此“时间换空间”在统计力学采样/变分优化口径下成立。证毕。
\end{proof}

\begin{theorem}[T2（同伦扭曲（统计力学）$\Rightarrow$ ORL 的统一实现）]
\label{thm:tex27-ch10-t2}
在 ORL 中把动作空间从原子动作提升为算子词 $w\in\mathcal{O}^*$，
并按 Gibbs 权重（或零温门控）选择 $w$，
则该 ORL 实现等价于用统计力学的自由能最小化框架实现“在路径空间上的变分选择”。
\end{theorem}

\begin{Certificate}[证书 T2（动作空间 = 路径空间，能量 = 势能）]
\label{cert:tex27-ch10-t2}
定义动作空间 $\Omega=\mathcal{O}^*$，能量 $E(w;s)=\Phi(w(s))$：
\begin{enumerate}
  \item（零温）选择 $\arg\min_w E(w;s)$ 即能量最小化；
  \item（有温）按 $e^{-\beta E}$ 采样对应自由能变分的标准实现。
\end{enumerate}
\end{Certificate}

\begin{proof}[证明（谓词因果链）]
按证书~\ref{cert:tex27-ch10-t2}，ORL 的核心“在动作空间上做变分选择”在此被路径空间化：
零温门控对应能量最小化；有限温采样对应自由能最小化。
因此同伦扭曲（作为路径采样）是 ORL 的统一实现口径。证毕。
\end{proof}

\begin{remark}[与 ODTHB 的接口对齐：ODTHB 是采样调度器而非干预源]
定理~\ref{thm:tex27-ch10-t2} 把 ORL 的“选择动作”提升为对 $\Omega=\mathcal{O}^*$ 的变分选择。
在该口径下，第~\ref{subsubsec:tex27-ch5-odthb-thermostat}~节命题~\ref{prop:tex27-ch5-odthb-r1} 给出：
ODTHB 的 fast 层可视为提供局部提议测度（最小接口 $p_t(e\mid u)$），slow 层调节 $\beta,\theta,\ell$ 等以维持有效混合；
并且在无硬剪裁且 $\varepsilon>0$ 时，其支持集不改变（定理~\ref{thm:tex27-ch5-odthb-support}）。
\end{remark}

\begin{theorem}[T3（非交换扭曲序列可同伦等价为量子虚时间演化）]
\label{thm:tex27-ch10-t3}
若存在生成元可加性分解 $\widehat{H}=\sum_i\widehat{G}_i$，
则非交换的扭曲序列在 Trotter 极限下同伦等价于 $e^{-\tau\widehat{H}}$ 的虚时间演化；
从而量子 Gibbs/量子自由能口径可用于描述其长期行为（形式系统内的模拟解释）。
\end{theorem}

\begin{Certificate}[证书 T3（Trotter 乘积公式证书）]
\label{cert:tex27-ch10-t3}
在合适的有界性/定义域条件下，Lie--Trotter 乘积公式给出
\[
  \Bigl(\prod_{i=1}^m e^{-\Delta\tau\,\widehat{G}_i}\Bigr)^{n}
  \xrightarrow[n\to\infty]{}
  e^{-\tau\sum_{i=1}^m \widehat{G}_i}=e^{-\tau\widehat{H}}.
\]
\end{Certificate}

\begin{proof}[证明（谓词因果链）]
由证书~\ref{cert:tex27-ch10-t3}，当 $\Delta\tau=\tau/n\to 0$ 时，有序乘积极限收敛到虚时间指数演化。
非交换性只影响离散近似的顺序敏感结构，但不改变“极限对象为 $e^{-\tau\widehat{H}}$”的形式结论。
因此可把非交换扭曲序列解释为量子虚时间演化的模拟实现；其平衡态/变分口径对应量子 Gibbs/量子自由能最小化。证毕。
\end{proof}

\subsection{命题：开放系统在线控制与相对 AGI 的严格化}
\label{subsec:tex27-ch10-props}

\begin{proposition}[P1（ORL 在开放系统博弈中的作用是即时修复/在线控制）]
\label{prop:tex27-ch10-p1}
在存在对抗扰动 $a_t$ 的开放系统中，令系统演化为
\[
  s_{t+1}=w_t\bigl(s_t,a_t\bigr),
  \qquad
  w_t\in\mathcal{O}^*.
\]
则 ORL 通过每步在路径空间 $\Omega=\mathcal{O}^*$ 上选择/采样 $w_t$，
实现对缺陷势能 $\Phi$ 与约束谓词 $\mathrm{Inv}$ 的在线控制；
其有效性不要求每步必优，但可采用统计力学口径保证长期漂移/自由能下降（对齐 T1--T2）。
\end{proposition}

\begin{Certificate}[证书 P1（在线性：每步重算；两种有效性口径）]
\label{cert:tex27-ch10-p1}
\begin{enumerate}
  \item（在线）对抗扰动使系统非平稳，故策略必须每步基于当前状态重新选择 $w_t$；
  \item（零温）门控给出单调下降但可能陷入局部；
  \item（有温）Gibbs 采样允许单步非优，但在混合条件下给出长期自由能口径的稳定化（T1）。
\end{enumerate}
\end{Certificate}

\begin{proof}[证明（谓词因果链）]
按证书~\ref{cert:tex27-ch10-p1}：
开放系统的对抗扰动破坏静态规划，迫使策略在线重算；ORL 的动作空间是 $\mathcal{O}^*$，
因而天然适配“即时修复”的算子词选择。
零温与有温两种口径分别对应单调下降与统计意义下降（T1--T2），故 ORL 可作为在线控制实现。证毕。
\end{proof}

\begin{proposition}[P2（算子完备性前提下：ORL $\Longleftrightarrow$ 白盒 AI $\Longleftrightarrow$ 相对形式 AGI）]
\label{prop:tex27-ch10-p2}
在边缘算子完备性前提（第~\ref{sec:tex27-ch9-orl-whitebox-agi}~节公理~\ref{thm:tex27-ch9-a9} 的 $\mathrm{Complete}$）
成立且采用相对 AGI 定义（第~\ref{sec:tex27-ch9-orl-whitebox-agi}~节定义~\ref{def:tex27-ch9-d8}）时，有等价链：
\[
  \mathrm{ORL}\ \Longleftrightarrow\ \text{可审计白盒 AI}\ \Longleftrightarrow\ \mathrm{AGI}^{\mathrm{cert}}_{\mathcal{V},
    \mathcal{O}}.
\]
\end{proposition}

\begin{Certificate}[证书 P2（引用既有闭合链）]
\label{cert:tex27-ch10-p2}
该命题直接调用第~\ref{sec:tex27-ch9-orl-whitebox-agi}~节推论~\ref{cor:tex27-ch9-c1}。
\end{Certificate}

\begin{proof}[证明（证书化）]
按证书~\ref{cert:tex27-ch10-p2}，由推论~\ref{cor:tex27-ch9-c1} 直接得到。证毕。
\end{proof}

\subsection{推论与备注：把一句话压缩为可审计链}
\label{subsec:tex27-ch10-others}

\begin{corollary}[C1（“时间换空间”强命题的严格版本：链式蕴含）]
\label{cor:tex27-ch10-c1}
在本节定义与前提下，有链式蕴含：
\[
  \underbrace{\text{同伦扭曲（统计力学）}}_{\text{路径采样/自由能}}
  \Longrightarrow
  \underbrace{\text{ORL}}_{\text{在 }\mathcal{O}^*\text{ 上做变分}}
  \Longrightarrow
  \underbrace{\text{开放系统在线控制}}_{\text{即时修复}}
  \Longrightarrow
  \underbrace{\mathrm{AGI}^{\mathrm{cert}}_{\mathcal{V},\mathcal{O}}}_{\text{相对定义（需 }\mathrm{Complete}\text{）}}.
\]
并且当非交换生成元可加时（T3）：
\[
  \text{同伦扭曲（非交换）}\ \simeq\ \text{量子虚时间演化}\ \Rightarrow\ \text{量子统计力学变分描述可用}.
\]
\end{corollary}

\begin{Certificate}[证书 C1（T1 + T2 + P1 + P2 串联）]
\label{cert:tex27-ch10-c1}
同伦扭曲 $\Rightarrow$ 时间换空间由定理~\ref{thm:tex27-ch10-t1}；
同伦扭曲 $\Rightarrow$ ORL 实现由定理~\ref{thm:tex27-ch10-t2}；
ORL $\Rightarrow$ 在线控制由命题~\ref{prop:tex27-ch10-p1}；
在线控制 $\Rightarrow$ 相对 AGI 由命题~\ref{prop:tex27-ch10-p2}（需 $\mathrm{Complete}$）。
非交换量子口径由定理~\ref{thm:tex27-ch10-t3}。
\end{Certificate}

\begin{proof}[证明（谓词因果链）]
按证书~\ref{cert:tex27-ch10-c1}，把 T1、T2、P1、P2 与 T3 依次串联即可得到推论两条链式陈述。证毕。
\end{proof}

\begin{remark}[工程含义 1：时间换空间的核心不是“量子”，而是“有效混合/有效越障”]
\label{rem:tex27-ch10-r1}
时间换空间的核心不在于名词“量子”，而在于：
无需把所有可能修复路径都存下来（避免空间爆炸），
而是通过足够长的扭曲序列（时间）在路径空间完成有效混合/有效越障，
从而在统计意义上逼近最优修复与最优对抗控制。
\end{remark}

\begin{remark}[未证明边界：量子计算速度优势需要外部复杂度/硬件证书]
\label{rem:tex27-ch10-r2}
定理~\ref{thm:tex27-ch10-t3} 证明的是\textbf{形式同伦等价/模拟等价}，
并不自动给出“物理量子硬件带来复杂度级加速”的结论。
后者属于未证明断言，至少需要外部证书：硬件可实现性、复杂度下界/上界、以及与经典算法的可比基准。
\end{remark}

\begin{remark}[工程含义 2：把“开放系统博弈：AGI”变成定理的最小前提集合]
\label{rem:tex27-ch10-r3}
要把“开放系统博弈：AGI”从口号变为定理链，至少需要三类可核验前提：
\begin{enumerate}
  \item $\mathcal{V}$ 的表达力：现实任务/约束可编码进任务类 $\mathfrak{T}(\mathcal{V})$；
  \item $\mathcal{O}$ 的完备性：满足 $\mathrm{Complete}$（公理~\ref{thm:tex27-ch9-a9}）或其统计漂移版本；
  \item 证书接口完备：每步可回放、可复核（定义~\ref{def:tex27-ch9-d7}）。
\end{enumerate}
缺任一项，则只能得到“未证明/需外部证书”的弱结论。
\end{remark}

%%%%%%%%%%%%%%%%%%%%%%%%%%%%%%%%%%%%%%%%%%%%%%%%%%%%%%%%%%%%
\section{鲁棒时间换空间：有限路径空间上的 Metropolis 同伦扭曲与近最优对抗修复（强命题链）}
\label{sec:tex27-ch11-robust-time-space}
%%%%%%%%%%%%%%%%%%%%%%%%%%%%%%%%%%%%%%%%%%%%%%%%%%%%%%%%%%%%

本节把“只要足够长的扭曲序列，就能以统计意义逼近最优修复与最优对抗控制”严格化为一个\textbf{可审计的定理--引理--命题链}。
与第~\ref{sec:tex27-ch10-time-space-tradeoff}~节不同的是：本节把“有效隧穿/有效混合”具体落到
\textbf{有限路径空间 $\Omega_L$ 上的 Metropolis 核}与\textbf{混合时间 $t_{\mathrm{mix}}(\cdot)$}，并把误差/失败概率写成显式界。

\subsection{定义：边缘算子字母表、路径空间与鲁棒能量}
\label{subsec:tex27-ch11-defs}

\begin{definition}[D1（边缘算子字母表与路径空间）]
\label{def:tex27-ch11-d1}
令边缘算子（原子算子）集合为有限集
\[
  \mathcal{O}=\{o_1,\dots,o_M\},
  \qquad
  1\le M<\infty.
\]
令长度上界（预算）为 $L\in\NN$。定义“修复路径/算子词”空间（至多 $L$ 步复合）：
\[
  \Omega_L:=\bigcup_{\ell=0}^{L}\mathcal{O}^{\ell},
\]
其中 $\mathcal{O}^{\ell}$ 表示长度为 $\ell$ 的有序序列（等价于复合 $o_{\ell}\circ\cdots\circ o_1$），且 $\mathcal{O}^0:=\{\mathrm{id}\}$
  表示空词（恒等算子）。
该构造刻画“空间爆炸”的来源：$|\Omega_L|$ 随 $L$ 指数增长（见公理~\ref{thm:tex27-ch11-a11}）。
\end{definition}

\begin{definition}[D2（开放系统对抗与鲁棒能量）]
\label{def:tex27-ch11-d2}
令变分空间（状态空间）为 $\mathcal{V}$，状态 $s\in\mathcal{V}$。
令对手/环境动作为 $a\in\mathcal{A}(s)$。
对任意路径 $w\in\Omega_L$，记一跳执行后的状态为
\[
  s' = w(s,a).
\]
令缺陷势能（障碍度量）
\[
  \Phi:\mathcal{V}\to\RR_{\ge 0}.
\]
定义一跳鲁棒能量（最坏情形代价）
\[
  E(w;s):=\sup_{a\in\mathcal{A}(s)}\Phi\bigl(w(s,a)\bigr)\ \in\ [0,+\infty],
\]
并定义最优鲁棒能量（一跳最优对抗修复）
\[
  E^\ast(s):=\inf_{w\in\Omega_L}E(w;s).
\]
在 $\Omega_L$ 有限且 $E(w;s)<\infty$ 的条件下，$E^\ast(s)$ 等价于 $\min_{w\in\Omega_L}E(w;s)$（最小值可达）。
\end{definition}

\begin{definition}[D3（统计力学扭曲：Gibbs 路径分布）]
\label{def:tex27-ch11-d3}
给定逆温 $\beta>0$，定义 Gibbs 分布
\[
  \pi_\beta(w\mid s):=\frac{\exp\bigl(-\beta E(w;s)\bigr)}{Z_\beta(s)},
  \qquad
  Z_\beta(s):=\sum_{w\in\Omega_L}\exp\bigl(-\beta E(w;s)\bigr).
\]
把“同伦扭曲序列”理解为在 $\Omega_L$ 上产生样本 $W_t$ 的马尔可夫链，
并在足够长时间后接近 $\pi_\beta(\cdot\mid s)$。
\end{definition}

\begin{definition}[D4（“时间换空间”的精确定义）]
\label{def:tex27-ch11-d4}
\begin{enumerate}
  \item（空间枚举法）显式存储或遍历整个 $\Omega_L$，空间/外存访问量至少与 $|\Omega_L|$ 同阶（需存候选及其能量）；
  \item（时间采样法）仅维护当前样本 $W_t$（加少量常数个统计量），空间复杂度为 $O(L)$（存一个长度至多 $L$ 的序列），
  用时间步数 $T$ 换取对 $\Omega_L$ 的统计性探索。
\end{enumerate}
\end{definition}

\subsection{公理（降级为定理）：指数增长、细致平衡与混合时间}
\label{subsec:tex27-ch11-axioms}

\begin{theorem}[公理（降级为定理）A11（路径空间有限且指数增长）]
\label{thm:tex27-ch11-a11}
在定义~\ref{def:tex27-ch11-d1} 的条件下，有计数上界：
\[
  |\Omega_L|
  =\sum_{\ell=0}^{L}|\mathcal{O}^{\ell}|
  =\sum_{\ell=0}^{L} M^\ell
  \le (L+1)M^L.
\]
\end{theorem}

\begin{Certificate}[证书 A11（数学归纳法证书：$|\mathcal{O}^\ell|=M^\ell$）]
\label{cert:tex27-ch11-a11}
对任意 $\ell\in\NN$，断言 $|\mathcal{O}^\ell|=M^\ell$ 的归纳证书为：
\begin{enumerate}
  \item（基步）$\ell=0$：$\mathcal{O}^0=\{\mathrm{id}\}$，故 $|\mathcal{O}^0|=1=M^0$；
  \item（归纳步）若 $|\mathcal{O}^\ell|=M^\ell$，则长度为 $\ell+1$ 的序列可由“在任意长度 $\ell$ 序列后追加一个字母”得到，
  每个序列有 $M$ 种追加方式且互不重合，故 $|\mathcal{O}^{\ell+1}|=M\cdot|\mathcal{O}^{\ell}|=M^{\ell+1}$。
\end{enumerate}
\end{Certificate}

\begin{proof}[证明（谓词因果链）]
由证书~\ref{cert:tex27-ch11-a11}，对任意 $\ell\le L$ 有 $|\mathcal{O}^\ell|=M^\ell$。
于是
\[
  |\Omega_L|
  =\sum_{\ell=0}^{L}|\mathcal{O}^{\ell}|
  =\sum_{\ell=0}^{L} M^\ell
  \le (L+1)\max_{0\le \ell\le L}M^\ell
  =(L+1)M^L.
\]
证毕。
\end{proof}

\begin{theorem}[公理（降级为定理）A12（Metropolis 扭曲核的平衡分布）]
\label{thm:tex27-ch11-a12}
固定状态 $s$，令 $Q(w\to w')$ 为路径空间 $\Omega_L$ 上的提议核（同伦扭曲产生邻域候选），并假设 $Q$ 使链在 $\Omega_L$ 上不可约（从任意 $w$ 可通过有限步到达任意 $w'$）。
定义接受率
\[
  \alpha(w,w')
  :=\min\left\{
  1,\
  \frac{\exp\bigl(-\beta E(w';s)\bigr)\,Q(w'\to w)}{\exp\bigl(-\beta E(w;s)\bigr)\,Q(w\to w')}
  \right\}.
\]
定义转移核
\[
  K(w\to w')
  :=Q(w\to w')\alpha(w,w')\quad (w'\ne w),
  \qquad
  K(w\to w):=1-\sum_{u\ne w}K(w\to u).
\]
则 $\pi_\beta(\cdot\mid s)$（定义~\ref{def:tex27-ch11-d3}）满足细致平衡，从而是 $K$ 的不变分布。
\end{theorem}

\begin{Certificate}[证书 A12（接受率公式的对偶性：$\min\{1,r\}$）]
\label{cert:tex27-ch11-a12}
令
\[
  r(w,w')
  :=\frac{\exp\bigl(-\beta E(w';s)\bigr)\,Q(w'\to w)}{\exp\bigl(-\beta E(w;s)\bigr)\,Q(w\to w')}.
\]
则 $\alpha(w,w')=\min\{1,r(w,w')\}$ 且 $\alpha(w',w)=\min\{1,1/r(w,w')\}$；
从而对任意 $w\ne w'$，有恒等式
\[
  \exp\bigl(-\beta E(w;s)\bigr)Q(w\to w')\alpha(w,w')
  =
  \exp\bigl(-\beta E(w';s)\bigr)Q(w'\to w)\alpha(w',w).
\]
\end{Certificate}

\begin{proof}[证明（谓词因果链）]
需证对任意 $w\ne w'$：
\[
  \pi_\beta(w\mid s)K(w\to w')=\pi_\beta(w'\mid s)K(w'\to w).
\]
由定义，$K(w\to w')=Q(w\to w')\alpha(w,w')$，且 $\pi_\beta(w\mid s)\propto \exp(-\beta E(w;s))$。
因此只需证
\[
  \exp\bigl(-\beta E(w;s)\bigr)Q(w\to w')\alpha(w,w')
  =
  \exp\bigl(-\beta E(w';s)\bigr)Q(w'\to w)\alpha(w',w),
\]
这正是证书~\ref{cert:tex27-ch11-a12} 给出的恒等式。
故细致平衡成立，从而 $\pi_\beta(\cdot\mid s)$ 为 $K$ 的不变分布。证毕。
\end{proof}

\begin{theorem}[公理（降级为定理）A13（有效混合：有限混合时间的存在）]
\label{thm:tex27-ch11-a13}
设 $\Omega_L$ 有限，且马尔可夫链 $W_t$ 的转移核为 $K$。
若 $K$ 不可约且非周期，则存在唯一不变分布 $\pi_\beta$，
并且对任意 $\varepsilon>0$，存在有限混合时间 $t_{\mathrm{mix}}(\varepsilon)$ 使得
\[
  \sup_{w_0\in\Omega_L}
  \left\|
  \mathcal{L}(W_t\mid W_0=w_0)-\pi_\beta(\cdot\mid s)
  \right\|_{\mathrm{TV}}
  \le \varepsilon,
  \qquad
  \forall t\ge t_{\mathrm{mix}}(\varepsilon),
\]
其中总变差距离定义为
\[
  \|\mu-\nu\|_{\mathrm{TV}}:=\sup_{A\subseteq\Omega_L}|\mu(A)-\nu(A)|.
\]
\end{theorem}

\begin{Certificate}[证书 A13（有限不可约非周期 $\Rightarrow$ 收敛到唯一平衡分布）]
\label{cert:tex27-ch11-a13}
证书输入为有限状态马尔可夫链的标准结论：
不可约 + 非周期 $\Rightarrow$ 唯一平衡分布且分布收敛；
并可用 Perron--Frobenius/谱半径或耦合论证推出总变差收敛与混合时间有限性。
\end{Certificate}

\begin{proof}[证明（谓词因果链，压缩版）]
按证书~\ref{cert:tex27-ch11-a13}，有限不可约非周期马尔可夫链存在唯一平衡分布且 $K^t$ 收敛到秩一极限。
因此 $\|\mathcal{L}(W_t\mid W_0=w_0)-\pi_\beta\|_{\mathrm{TV}}\to 0$，
由极限定义可得对任意 $\varepsilon>0$ 存在有限 $t_{\mathrm{mix}}(\varepsilon)$ 满足所需不等式。证毕。
\end{proof}

\subsection{引理：Gibbs 集中与混合误差}
\label{subsec:tex27-ch11-lemmas}

\begin{lemma}[L1（Gibbs 分布对最优鲁棒能量的集中）]
\label{lem:tex27-ch11-l1}
对固定 $s$，任取 $\epsilon>0$，定义 $\epsilon$-次优集合
\[
  B_\epsilon(s):=\{w\in\Omega_L:\ E(w;s)\ge E^\ast(s)+\epsilon\}.
\]
若 $E^\ast(s)$ 在 $\Omega_L$ 上可达（例如 $\Omega_L$ 有限且 $E(w;s)<\infty$），则有上界
\[
  \pi_\beta\bigl(B_\epsilon(s)\mid s\bigr)\ \le\ |\Omega_L|\exp(-\beta\epsilon).
\]
\end{lemma}

\begin{Certificate}[证书 L1（分子上界 + 分母下界）]
\label{cert:tex27-ch11-l1}
证书分两步：
\begin{enumerate}
  \item（分子）对 $w\in B_\epsilon(s)$，有 $E(w;s)\ge E^\ast(s)+\epsilon$，故
  $\sum_{w\in B_\epsilon}\exp(-\beta E(w;s))\le |B_\epsilon|\exp(-\beta(E^\ast+\epsilon))
    \le|\Omega_L|\exp(-\beta(E^\ast+\epsilon))$；
  \item（分母）若存在 $w^\ast$ 使 $E(w^\ast;s)=E^\ast(s)$，则
  $Z_\beta(s)=\sum_{w}\exp(-\beta E(w;s))\ge \exp(-\beta E^\ast(s))$。
\end{enumerate}
\end{Certificate}

\begin{proof}[证明（谓词因果链）]
按证书~\ref{cert:tex27-ch11-l1}：
分子满足
\[
  \sum_{w\in B_\epsilon(s)}\exp\bigl(-\beta E(w;s)\bigr)
  \le |\Omega_L|\exp\bigl(-\beta(E^\ast(s)+\epsilon)\bigr),
\]
分母满足 $Z_\beta(s)\ge \exp(-\beta E^\ast(s))$。
故
\[
  \pi_\beta\bigl(B_\epsilon(s)\mid s\bigr)
  =\frac{\sum_{w\in B_\epsilon(s)}\exp(-\beta E(w;s))}{Z_\beta(s)}
  \le |\Omega_L|\exp(-\beta\epsilon).
\]
证毕。
\end{proof}

\begin{lemma}[L2（混合误差把“理论集中”转为“实际采样集中”）]
\label{lem:tex27-ch11-l2}
若 $t\ge t_{\mathrm{mix}}(\varepsilon)$，则对任意事件 $A\subseteq\Omega_L$：
\[
  \mathbb{P}(W_t\in A)\ \le\ \pi_\beta(A\mid s)+\varepsilon.
\]
\end{lemma}

\begin{Certificate}[证书 L2（总变差的上确界定义）]
\label{cert:tex27-ch11-l2}
总变差距离定义为
\[
  \|\mu-\nu\|_{\mathrm{TV}}=\sup_{A\subseteq\Omega_L}|\mu(A)-\nu(A)|.
\]
\end{Certificate}

\begin{proof}[证明（谓词因果链）]
令 $\mu=\mathcal{L}(W_t)$，$\nu=\pi_\beta(\cdot\mid s)$。
由混合时间定义与证书~\ref{cert:tex27-ch11-l2}，
对任意事件 $A$ 有 $|\mu(A)-\nu(A)|\le \|\mu-\nu\|_{\mathrm{TV}}\le \varepsilon$。
因此 $\mu(A)\le \nu(A)+\varepsilon$，即
$\mathbb{P}(W_t\in A)\le \pi_\beta(A\mid s)+\varepsilon$。证毕。
\end{proof}

\subsection{命题：近最优鲁棒修复、时间换空间与多步误差传播}
\label{subsec:tex27-ch11-props}

\begin{proposition}[P1（核心结论：长时间扭曲 $\Rightarrow$ 统计意义逼近最优鲁棒修复）]
\label{prop:tex27-ch11-p1}
给定误差容忍 $\epsilon>0$ 与失败概率 $\delta\in(0,1)$。
若选择
\[
  \beta \ \ge\ \frac{1}{\epsilon}\Bigl(\log|\Omega_L|+\log\frac{2}{\delta}\Bigr),
\]
并运行同伦扭曲马尔可夫链到
\[
  t\ge t_{\mathrm{mix}}(\delta/2),
\]
则有概率保证
\[
  \mathbb{P}\bigl(E(W_t;s)\le E^\ast(s)+\epsilon\bigr)\ \ge\ 1-\delta.
\]
\end{proposition}

\begin{Certificate}[证书 P1（L2 + L1 + 参数选择）]
\label{cert:tex27-ch11-p1}
取事件 $A:=B_\epsilon(s)$（定义~\ref{lem:tex27-ch11-l1}），并令混合误差取 $\varepsilon:=\delta/2$。
若令 $\beta$ 满足 $|\Omega_L|e^{-\beta\epsilon}\le \delta/2$，则
$\mathbb{P}(W_t\in B_\epsilon(s))\le\delta$。
\end{Certificate}

\begin{proof}[证明（谓词因果链）]
取 $A:=B_\epsilon(s)$。
由引理~\ref{lem:tex27-ch11-l2}（取 $\varepsilon=\delta/2$ 且 $t\ge t_{\mathrm{mix}}(\delta/2)$）得
\[
  \mathbb{P}(W_t\in B_\epsilon(s))\le \pi_\beta(B_\epsilon(s)\mid s)+\delta/2.
\]
由引理~\ref{lem:tex27-ch11-l1} 得
\[
  \pi_\beta(B_\epsilon(s)\mid s)\le |\Omega_L|\exp(-\beta\epsilon).
\]
若 $\beta \ge \frac{1}{\epsilon}\bigl(\log|\Omega_L|+\log\frac{2}{\delta}\bigr)$，则 $|\Omega_L|e^{-\beta\epsilon}\le
  \delta/2$，
从而
\[
  \mathbb{P}(W_t\in B_\epsilon(s))\le \delta/2+\delta/2=\delta.
\]
等价于
\[
  \mathbb{P}\bigl(E(W_t;s)<E^\ast(s)+\epsilon\bigr)\ge 1-\delta,
\]
并可将严格不等号改写为 $\le$（仅差边界集合，有限空间下不影响结论口径）。证毕。
\end{proof}

\begin{proposition}[P2（时间换空间：空间从线性于 $|\Omega_L|$ 降为 $O(L)$，代价是混合时间）]
\label{prop:tex27-ch11-p2}
\begin{enumerate}
  \item（枚举法）要比较全体 $E(w;s)$，至少需要访问（或存储）线性于 $|\Omega_L|$ 的候选条目；
  \item（采样法）Metropolis 同伦扭曲每步仅需保存当前词 $W_t$（长度至多 $L$）与其能量 $E(W_t;s)$ 及常数个随机数，空间为 $O(L)$；
  \item（达到 $(\epsilon,\delta)$ 近似）由命题~\ref{prop:tex27-ch11-p1}，只需 $t\ge t_{\mathrm{mix}}(\delta/2)$ 且选择合适 $\beta$。
\end{enumerate}
\end{proposition}

\begin{Certificate}[证书 P2（“局部更新 + 在线能量评估”不需要全表）]
\label{cert:tex27-ch11-p2}
Metropolis 更新由 $(W_t,E(W_t;s))$ 与局部提议 $W'$ 决定；
$W'$ 可在线生成，无需保存全体候选能量表。
\end{Certificate}

\begin{proof}[证明（谓词因果链）]
\begin{enumerate}
  \item 枚举法若不访问（或存储）大量候选，就无法在全体 $\Omega_L$ 上进行比较或保证覆盖；
  \item 采样法按证书~\ref{cert:tex27-ch11-p2}，每步只依赖当前样本与局部提议，故空间随 $L$ 线性；
  \item 达到近似所需时间/温度由命题~\ref{prop:tex27-ch11-p1} 给出。
\end{enumerate}
证毕。
\end{proof}

\begin{proposition}[P3（从“一跳 $\epsilon$-近最优鲁棒修复”到“多步对抗控制”的误差传播界）]
\label{prop:tex27-ch11-p3}
取折扣因子 $\gamma\in(0,1)$。
设鲁棒 Bellman 算子 $T$ 在 $\|\cdot\|_\infty$ 范数下为 $\gamma$-收缩（标准对抗控制口径），
且每步一跳选择的路径满足“$\epsilon$-近最优”：
其一跳 Bellman 误差不超过 $\epsilon$。
则对应策略的值函数次优性满足
\[
  V^{\pi}(s)-V^\ast(s)\ \le\ \frac{\epsilon}{1-\gamma},
  \qquad
  \forall s.
\]
有限时域 $H$ 的版本为 $V^{\pi}(s)-V^\ast(s)\le H\epsilon$。
\end{proposition}

\begin{Certificate}[证书 P3（收缩映射误差累积：$\epsilon(1+\gamma+\gamma^2+\cdots)$）]
\label{cert:tex27-ch11-p3}
若 $T$ 是 $\gamma$-收缩且每步引入的 Bellman 误差 $\le\epsilon$，则几何级数上界
$\epsilon\sum_{t\ge 0}\gamma^t=\epsilon/(1-\gamma)$ 给出全局误差传播界。
\end{Certificate}

\begin{proof}[证明（谓词因果链，标准收缩）]
按证书~\ref{cert:tex27-ch11-p3}：
鲁棒 Bellman 算子 $T$ 为 $\gamma$-收缩意味着误差经一次 Bellman 更新最多被放大 $\gamma$ 倍。
若每一步近似最小化引入的即时误差不超过 $\epsilon$，
则在迭代展开中得到误差上界
\[
  \epsilon+\gamma\epsilon+\gamma^2\epsilon+\cdots=\frac{\epsilon}{1-\gamma}.
\]
有限时域时几何级数替换为有限和，得到 $H\epsilon$ 上界。证毕。
\end{proof}

\subsection{推论与备注：严格口径与工程落地接口}
\label{subsec:tex27-ch11-others}

\begin{corollary}[C1（“足够长扭曲序列 $\Rightarrow$ 逼近最优鲁棒修复/对抗控制”的严格版本）]
\label{cor:tex27-ch11-c1}
若满足：
\begin{enumerate}
  \item $\Omega_L$ 由有限字母表 $\mathcal{O}$ 与有限预算 $L$ 给出（A11）；
  \item 同伦扭曲核用 Metropolis 形式构造并满足细致平衡（A12）；
  \item “有效混合/有效隧穿”形式化为不可约 + 非周期，从而存在混合时间（A13）；
  \item 选择足够大的 $\beta$ 与足够长的时间 $t$（命题~\ref{prop:tex27-ch11-p1}）；
\end{enumerate}
则以概率至少 $1-\delta$ 得到 $\epsilon$-近最优的一跳鲁棒修复；
再由命题~\ref{prop:tex27-ch11-p3} 得到多步对抗控制的误差界。
\end{corollary}

\begin{Certificate}[证书 C1（A11--A13 + P1 + P3 串联）]
\label{cert:tex27-ch11-c1}
一跳近似由 A11--A13 与 P1 共同给出；多步误差传播由 P3 给出。
\end{Certificate}

\begin{proof}[证明（谓词因果链）]
按证书~\ref{cert:tex27-ch11-c1}：
A11--A13 给出有限路径空间与可混合的扭曲核；P1 给出一跳近似的概率界；
P3 把“一跳近似”推广为“多步值函数近似”。证毕。
\end{proof}

\begin{remark}[备注 1：本节证明的三元核心（空间、时间、统计最优）]
\label{rem:tex27-ch11-r1}
本节证明的核心是三点：
\begin{enumerate}
  \item “不存全路径空间”对应空间复杂度（仅 $O(L)$）；
  \item “足够长时间”对应混合时间 $t_{\mathrm{mix}}$；
  \item “统计意义逼近最优”由 Gibbs 集中（L1）+ 混合误差（L2）给出。
\end{enumerate}
三者合起来就是严格的“时间换空间（且带误差/失败概率界）”。
\end{remark}

\begin{remark}[备注 2：有效隧穿的数学化：不可约性]
\label{rem:tex27-ch11-r2}
“有效隧穿”在数学上不需要额外玄学解释：它等价于链在 $\Omega_L$ 上的\textbf{不可约性}，
即任意路径都能通过有限次扭曲到达任意其他路径，从而不会永久卡在某个局部连通分支。
\end{remark}

\begin{remark}[备注 3：量子口径不是本节结论必要条件]
\label{rem:tex27-ch11-r3}
本节“时间换空间”的严格结论只依赖统计力学采样与马尔可夫链混合（A12--A13），
并不需要“量子计算加速”的任何假设。
若仅把“非交换算子序列 $\simeq$ 虚时间指数”作为形式描述，则可额外参考第~\ref{sec:tex27-ch8-homotopy-statmech}~节与第~\ref{sec:tex27-ch10-time-space-tradeoff}~节的量子口径讨论；
但“物理量子硬件的复杂度加速”仍属未证明外推。
\end{remark}

\begin{remark}[备注 4：工程落地的最小证书接口]
\label{rem:tex27-ch11-r4}
若要把本节结论直接落到工程实现，最小证书接口是：
\begin{enumerate}
  \item（可达性证书）对提议核 $Q$ 给出路径空间可达性（不可约）证明/测试覆盖；
  \item（混合证书）给出 $t_{\mathrm{mix}}(\cdot)$ 的经验估计或上界（可复现实验/耦合界）；
  \item（能量可复算）给出 $E(w;s)$ 的可复算实现与日志，确保 P1 的参数选择可审计。
\end{enumerate}
\end{remark}

%%%%%%%%%%%%%%%%%%%%%%%%%%%%%%%%%%%%%%%%%%%%%%%%%%%%%%%%%%%%
\section{算子格点扇区与神经网络参数规模：扇区内同伦等价与算力比较（强命题链）}
\label{sec:tex27-ch12-sector-nn-scale}
%%%%%%%%%%%%%%%%%%%%%%%%%%%%%%%%%%%%%%%%%%%%%%%%%%%%%%%%%%%%

本节把“算子格点空间规模同伦等价于神经网络参数规模”“统计力学式 ORL 更节约算力且可依赖 CPU”
写成\textbf{扇区内的严格命题}与\textbf{可检验的条件不等式}，并在结尾明确标注哪些结论属于工程参数依赖（非普遍定理）。

\subsection{定义：算子格点空间、神经网络参数空间与成本模型}
\label{subsec:tex27-ch12-defs}

\begin{definition}[D1（算子格点空间与“扇区”分解）]
\label{def:tex27-ch12-d1}
令“同伦扭曲自由度”的算子格点空间（参数化后的变分空间）按扇区分解为不交并：
\[
  \Theta_{\mathrm{op}}
  =\bigsqcup_{\sigma\in\Sigma}\Theta_\sigma,
  \qquad
  \Theta_\sigma\simeq \RR^{d_\sigma},
  \qquad
  d_\sigma\in\NN.
\]
其中 $\sigma$ 表示一组离散结构选择（例如启用哪些候选族/规则分支/配置 profile），
在固定 $\sigma$ 后剩余自由度是 $d_\sigma$ 维连续实参数。

\textbf{约定（同伦口径）}：只在单一扇区 $\Theta_\sigma$ 内讨论连续同伦；不固定扇区时，
$\Theta_{\mathrm{op}}$ 一般是多连通分支的并，不能与连通的 $\RR^P$ 直接同伦等价。
\end{definition}

\begin{definition}[D2（神经网络参数空间）]
\label{def:tex27-ch12-d2}
给定一个神经网络结构 $\mathcal{N}$，其可训练参数数目记为 $P(\mathcal{N})\in\NN$。
其参数空间为
\[
  \Theta_{\mathrm{nn}}(\mathcal{N}) \simeq \RR^{P(\mathcal{N})}.
\]
\end{definition}

\begin{definition}[D3（“规模同伦等价”的严格口径：扇区内同维同胚）]
\label{def:tex27-ch12-d3}
称“算子格点空间规模同伦等价于神经网络参数规模”，若对任意扇区 $\sigma$，
存在某个网络结构 $\mathcal{N}_\sigma$ 使
\[
  P(\mathcal{N}_\sigma)=d_\sigma,
\]
并且存在同胚（从而同伦等价）
\[
  \Theta_\sigma \cong \Theta_{\mathrm{nn}}(\mathcal{N}_\sigma).
\]
\end{definition}

\begin{definition}[D4（统计力学式变分优化在 ORL 中的表达）]
\label{def:tex27-ch12-d4}
给定状态 $s$ 与“路径动作”（算子词）集合 $\Omega_L\subset\mathcal{O}^*$，
定义能量
\[
  E(w;s):=\Phi(w(s)).
\]
有限温 Gibbs 采样定义为
\[
  \pi_\beta(w\mid s)\propto e^{-\beta E(w;s)},
  \qquad
  \beta>0.
\]
因此 ORL 的一次决策可理解为：从 $\Omega_L$ 中选择/采样 $w$ 以降低 $\Phi$（零温）或降低自由能（有限温，统计意义）。
\end{definition}

\begin{definition}[D5（算力/存储成本模型）]
\label{def:tex27-ch12-d5}
给出一个可比口径的成本模型：
\begin{enumerate}
  \item（ORL 每步）若每步枚举/采样候选数为 $K$，每个候选评估 $\Phi$ 的代价为 $C_\Phi$，则
  \[
    \mathrm{Cost}_{\mathrm{ORL,step}}=O(K\cdot C_\Phi),
    \qquad
    \mathrm{Mem}_{\mathrm{ORL}}=O(d_\sigma)+O(L).
  \]
  \item（神经网络训练）数据量 $N$，训练轮数 $E$，参数量 $P$，
  每样本每轮反传代价记为 $C_{\mathrm{bp}}(P)=\Theta(P)$，则
  \[
    \mathrm{Cost}_{\mathrm{NN,train}}=\Theta(E\cdot N\cdot P),
    \qquad
    \mathrm{Mem}_{\mathrm{NN}}=\Theta(P).
  \]
  推理每步代价记为
  \[
    \mathrm{Cost}_{\mathrm{NN,infer,step}}=\Theta(P).
  \]
\end{enumerate}
\end{definition}

\subsection{公理（降级为定理）：扇区内同维同胚与 CPU 负载口径}
\label{subsec:tex27-ch12-axioms}

\begin{theorem}[公理（降级为定理）T1（扇区内：算子格点连续自由度与同维神经网络参数空间同伦等价）]
\label{thm:tex27-ch12-t1}
对任意扇区 $\sigma$，若 $\Theta_\sigma\simeq \RR^{d_\sigma}$，则存在网络结构 $\mathcal{N}_\sigma$ 使 $P(\mathcal{N}_\sigma)=d_\sigma$，
并且
\[
  \Theta_\sigma \cong \RR^{d_\sigma} \cong \RR^{P(\mathcal{N}_\sigma)} \cong \Theta_{\mathrm{nn}}(\mathcal{N}_\sigma),
\]
因此 $\Theta_\sigma$ 与 $\Theta_{\mathrm{nn}}(\mathcal{N}_\sigma)$ 同胚，从而同伦等价。
\end{theorem}

\begin{Certificate}[证书 T1（同维欧氏空间同胚 + 显式构造 $P(\mathcal{N}_\sigma)=d_\sigma$）]
\label{cert:tex27-ch12-t1}
证书分两部分：
\begin{enumerate}
  \item（扇区自由度证书）由定义~\ref{def:tex27-ch12-d1}，$\Theta_\sigma\simeq \RR^{d_\sigma}$；
  \item（网络参数计数证书）给定任意 $d\in\NN$，可构造一个“无偏置的线性网络”
  \[
    f_\theta(x):=\sum_{i=1}^{d}\theta_i x_i,
    \qquad
    \theta\in\RR^{d},
  \]
  其可训练参数恰为 $d$ 个权重（无额外偏置参数），故存在 $\mathcal{N}$ 使 $P(\mathcal{N})=d$ 且 $\Theta_{\mathrm{nn}}(\mathcal{N})\simeq\RR^d$。
\end{enumerate}
\end{Certificate}

\begin{proof}[证明（谓词因果链）]
令谓词
\[
  A_\sigma:\ \Theta_\sigma\simeq \RR^{d_\sigma},
  \qquad
  B_{\mathcal{N}}:\ \Theta_{\mathrm{nn}}(\mathcal{N})\simeq \RR^{P(\mathcal{N})}.
\]
由证书~\ref{cert:tex27-ch12-t1} 的第 1 部分得到 $A_\sigma$；
由证书第 2 部分可取网络 $\mathcal{N}_\sigma$ 使 $P(\mathcal{N}_\sigma)=d_\sigma$ 且 $B_{\mathcal{N}_\sigma}$ 成立。
于是
\[
  \Theta_\sigma\cong \RR^{d_\sigma}\cong \RR^{P(\mathcal{N}_\sigma)}\cong \Theta_{\mathrm{nn}}(\mathcal{N}_\sigma),
\]
同胚推出同伦等价。证毕。
\end{proof}

\begin{theorem}[公理（降级为定理）T2（统计力学式 ORL 可用 CPU 以低存储实现：不需要大规模反传）]
\label{thm:tex27-ch12-t2}
在成本模型（定义~\ref{def:tex27-ch12-d5}）下，若 ORL 的候选评估代价 $C_\Phi$ 可用标量/小规模运算完成（例如规则判定、少量统计量更新、少量价格/仓位计算），
则 ORL 的主计算形态为 $O(K\cdot C_\Phi)$ 的分支+标量运算，
其内存为 $O(d_\sigma)+O(L)$；
因此在结构上\textbf{不需要}神经网络训练式的 $\Theta(E\cdot N\cdot P)$ 级反向传播与 $\Theta(P)$ 级参数存储即可运行。
\end{theorem}

\begin{Certificate}[证书 T2（成本分解证书：D5 的直接展开）]
\label{cert:tex27-ch12-t2}
证书即定义~\ref{def:tex27-ch12-d5} 的成本表达式与“$C_\Phi$ 为小规模计算”的前提：
ORL 每步只需评估 $K$ 个候选并更新常数个统计量；
而 NN 训练需要对 $P$ 个参数累计梯度并在 $N$ 个样本、$E$ 个 epoch 上反传。
\end{Certificate}

\begin{proof}[证明（谓词因果链）]
\begin{enumerate}
  \item（ORL 的计算与存储）由证书~\ref{cert:tex27-ch12-t2} 与定义~\ref{def:tex27-ch12-d5}，
  ORL 每步代价为 $O(K\cdot C_\Phi)$，存储仅包含扇区参数向量（维数 $d_\sigma$）与当前路径（长度 $\le L$），故 $\mathrm{Mem}_{\mathrm{ORL}}=O(d_\sigma)
    +O(L)$；
  \item（NN 训练的主项）由定义~\ref{def:tex27-ch12-d5}，训练需在 $E\cdot N$ 次样本级更新中对 $P$ 个参数反传，主项为 $\Theta(E\cdot N\cdot P)$，
    且需存储参数与梯度量级为 $\Theta(P)$；
  \item（结论口径）因此在“在线即时修复”且 $C_\Phi$ 为小规模运算的条件下，ORL 机制在结构上避免了大规模反传与参数巨量存储，从而 CPU 可直接承担其分支与标量工作负载（结论是“CPU 足够”，而非“GPU 不可用”）。
\end{enumerate}
证毕。
\end{proof}

\subsection{引理：规模解释与可检验不等式}
\label{subsec:tex27-ch12-lemmas}

\begin{lemma}[L1（“规模等价”指自由度维数同阶，而非必须同实现）]
\label{lem:tex27-ch12-l1}
对固定扇区 $\sigma$，ORL 的自由度规模为 $d_\sigma$；神经网络的自由度规模为参数量 $P$。
在“表达容量主要由自由度数目控制”的粗粒度意义下，二者可比较的核心量是
\[
  d_\sigma\ \text{vs}\ P.
\]
该比较是定义~\ref{def:tex27-ch12-d3}（扇区内同维同胚）的一种可操作解释。
\end{lemma}

\begin{Certificate}[证书 L1（T1 的直接调用）]
\label{cert:tex27-ch12-l1}
由定理~\ref{thm:tex27-ch12-t1}，当 $P=d_\sigma$ 时，$\Theta_\sigma$ 与 $\Theta_{\mathrm{nn}}$ 同胚；
因此“规模比较”自然落到维数/自由度计数上。
\end{Certificate}

\begin{proof}[证明（谓词因果链）]
按证书~\ref{cert:tex27-ch12-l1}，同胚要求同维（在欧氏空间口径下维数一致）。
因此在扇区内讨论“规模同伦等价”时，最直接的可操作指标就是自由度维数 $d_\sigma$ 与参数量 $P$ 的对比。证毕。
\end{proof}

\begin{lemma}[L2（何时 ORL 比 NN 更省算力：一个可检验不等式）]
\label{lem:tex27-ch12-l2}
在同一在线窗口长度 $T$ 内，若比较“即时适配”的总计算量：
\[
  \mathrm{Cost}_{\mathrm{ORL,total}}=O(T\cdot K\cdot C_\Phi).
\]
若神经网络需要在线再训练（或频繁微调）以适应开放系统漂移，则其总成本满足下界形态
\[
  \mathrm{Cost}_{\mathrm{NN,total}}\gtrsim \Theta(E\cdot N\cdot P)+\Theta(T\cdot P).
\]
因此在条件不等式
\[
  T\cdot K\cdot C_\Phi\ \ll\ E\cdot N\cdot P
\]
成立时，ORL 在该窗口内显著更省算力。
\end{lemma}

\begin{Certificate}[证书 L2（D5 的总量化展开）]
\label{cert:tex27-ch12-l2}
证书为：将定义~\ref{def:tex27-ch12-d5} 的“每步/每轮”成本乘以窗口长度或训练轮数，即可得到总成本表达式；
并将“在线微调”视为至少发生一次 $\Theta(E\cdot N\cdot P)$ 级训练主项。
\end{Certificate}

\begin{proof}[证明（谓词因果链）]
由定义~\ref{def:tex27-ch12-d5}，ORL 每步成本为 $O(K\cdot C_\Phi)$，窗口长度为 $T$，得 $\mathrm{Cost}_{\mathrm{ORL,total}}=O(TKC_\Phi)$。
若 NN 需要在线再训练，则至少付出一项训练主项 $\Theta(E N P)$，同时每步推理为 $\Theta(P)$，合并得下界形态 $\Theta(E N P)+\Theta(TP)$。
因此当 $TKC_\Phi\ll ENP$ 时，ORL 的窗口总成本严格小于 NN 的训练主项量级，从而更省算力。证毕。
\end{proof}

\subsection{命题、推论与备注：严格口径与未证明边界}
\label{subsec:tex27-ch12-props}

\begin{proposition}[P1（“规模同伦等价”成立，但需限定在扇区内）]
\label{prop:tex27-ch12-p1}
“算子格点空间规模同伦等价于神经网络参数规模”的严格版本为：
\[
  \forall \sigma\in\Sigma,\ \exists \mathcal{N}_\sigma:\ P(\mathcal{N}_\sigma)=d_\sigma,\ \Theta_\sigma \simeq
    \Theta_{\mathrm{nn}}(\mathcal{N}_\sigma).
\]
\end{proposition}

\begin{Certificate}[证书 P1（T1 直接推出）]
\label{cert:tex27-ch12-p1}
由定理~\ref{thm:tex27-ch12-t1}，对每个扇区都可取同维网络并得到同胚，从而得到命题陈述。
\end{Certificate}

\begin{proof}[证明]
按证书~\ref{cert:tex27-ch12-p1} 直接推出。证毕。
\end{proof}

\begin{proposition}[P2（“统计力学变分机制更节约算力且可依赖 CPU”是条件命题）]
\label{prop:tex27-ch12-p2}
在下列条件同时成立时，“更节约算力、依赖 CPU”的严格版本成立：
\begin{enumerate}
  \item ORL 的候选评估代价 $C_\Phi$ 为小规模标量/分支运算（不需要大规模矩阵乘）；
  \item 任务是开放系统在线修复，NN 需要再训练或频繁微调才能跟上漂移；
  \item 满足可检验不等式 $T\cdot K\cdot C_\Phi\ll E\cdot N\cdot P$（引理~\ref{lem:tex27-ch12-l2}）；
  \item “依赖 CPU”的严格解释是：\textbf{不需要 GPU 才能可行地运行}，而非“不能用 GPU”。
\end{enumerate}
则统计力学式 ORL 机制在该在线窗口内比 NN 的训练式方案更省算力，并可在 CPU 上承担主要工作负载。
\end{proposition}

\begin{Certificate}[证书 P2（T2 + L2 的组合）]
\label{cert:tex27-ch12-p2}
T2 给出 ORL 的计算形态与低存储口径；L2 给出窗口内成本比较不等式；
两者合并并加上“在线微调”前提即可推出结论。
\end{Certificate}

\begin{proof}[证明（谓词因果链）]
\begin{enumerate}
  \item 由条件 1) 与定理~\ref{thm:tex27-ch12-t2}，ORL 的主工作负载为分支+标量计算，CPU 可承载；
  \item 由条件 2) 与引理~\ref{lem:tex27-ch12-l2}，NN 在线需要承担训练主项 $\Theta(ENP)$；
  \item 由条件 3) 的不等式，窗口内 ORL 的总成本 $O(TKC_\Phi)$ 远小于 NN 的训练主项量级；
  \item 条件 4) 排除“必须 GPU”与“GPU 不可用”的语义误解：结论是“CPU 足够且结构上不依赖 GPU”，而非“排斥 GPU”。
\end{enumerate}
证毕。
\end{proof}

\begin{corollary}[C1（更精确的可证结论表述）]
\label{cor:tex27-ch12-c1}
把“规模同伦等价 + 算力优势”写成可证版本：
\begin{enumerate}
  \item（规模同伦等价，扇区内）若 $P=d_\sigma$，则
  \[
    \Theta_\sigma \simeq \RR^{d_\sigma}\simeq \RR^{P}\simeq \Theta_{\mathrm{nn}}(\mathcal{N}_\sigma).
  \]
  \item（算力优势，条件式）若在线漂移迫使 NN 频繁训练，并满足 $T\cdot K\cdot C_\Phi\ll E\cdot N\cdot P$，
  则统计力学式 ORL 在该窗口内更省算力，且其计算形态更适合 CPU 的分支/标量执行。
\end{enumerate}
\end{corollary}

\begin{Certificate}[证书 C1（P1 + P2 并列引用）]
\label{cert:tex27-ch12-c1}
第 1 点由命题~\ref{prop:tex27-ch12-p1}；
第 2 点由命题~\ref{prop:tex27-ch12-p2}。
\end{Certificate}

\begin{proof}[证明（谓词因果链）]
按证书~\ref{cert:tex27-ch12-c1}，分别由 P1、P2 直接推出两点表述。证毕。
\end{proof}

\begin{remark}[未证明边界（工程参数依赖）：并非无条件更省]
\label{rem:tex27-ch12-r1}
“更节约算力”不是普遍定理：若 NN 只做推理且 $P$ 很小，或 ORL 的 $K$ 很大且 $C_\Phi$ 很重（例如需要大规模回测/仿真），
则 ORL 可能更贵。该差异由具体工程参数决定，因此本节仅给出可检验的充分条件（如引理~\ref{lem:tex27-ch12-l2} 的不等式）。
\end{remark}

\begin{remark}[未证明边界（语义澄清）：“依赖 CPU”\,$\neq$\,“不能用 GPU”】【可行性口径】]
\label{rem:tex27-ch12-r2}
“依赖 CPU”在本节的严格含义是：切换 GPU 与否不影响可行性（CPU 足够）。
ORL 仍可并行化并使用 GPU/多核（多候选/多扇区/多样本）；只是它的核心计算不天然要求 GPU 的高吞吐矩阵运算。
\end{remark}

\begin{remark}[未证明边界（项目级证书）：“实践中必然更省”需要可回放测量]
\label{rem:tex27-ch12-r3}
若要把“统计力学 ORL 在实践中必然更省算力”提升为项目级证书，还需要把
$(T,K,C_\Phi)$ 与（若采用 NN）$(E,N,P)$ 在\textbf{同一任务/同一精度目标}下做可回放测量并给出日志证据。
本节不引入具体数据，因此只提供可验证的充分条件与严格比较不等式。
\end{remark}

%%%%%%%%%%%%%%%%%%%%%%%%%%%%%%%%%%%%%%%%%%%%%%%%%%%%%%%%%%%%
\section{回放收敛与同伦修复塔：统计力学同伦类下的边缘算子变分修复存在性与有效性（严格证明）}
\label{sec:tex27-ch13-replay-variational-repair}
%%%%%%%%%%%%%%%%%%%%%%%%%%%%%%%%%%%%%%%%%%%%%%%%%%%%%%%%%%%%

本节给出一个\textbf{独立章节}，把如下经验观察压缩为可证书化前提，并据此推出严格结论：
\begin{quote}
  多次回放显示：（i）修复方式存在差异（不同算子词/不同路径）；（ii）修复力（对证书/势能的净下降）呈现收敛；
  （iii）同伦扭曲发生次数（或按类型计数的烈度向量）是确定的。
\end{quote}
在这些条件下，我们将证明：在同伦修复塔（定义~\ref{def:tex27-ch9-d2}）的有限截断上，
对于规范场自由度（边上的连接/1-上链，定义~\ref{def:tex27-ch6-d1} 与定义~\ref{def:tex27-ch6-d2}），
存在且有效的\textbf{基于边缘算子}（定义~\ref{def:tex27-ch6-d4}）的\textbf{变分修复}；
并且该结论可在统计力学同伦等价口径（定义~\ref{def:tex27-ch9-d6}）下解释为“自由能/能量的变分选择”。

\subsection{定义组：回放、同伦扭曲次数与修复力}
\label{subsec:tex27-ch13-defs}

\begin{definition}[D1（同伦修复塔的有限截断与事件类型）]
\label{def:tex27-ch13-d1}
沿用定义~\ref{def:tex27-ch9-d2} 的同伦修复塔分解 $\mathcal{O}=\bigsqcup_{k\ge 1}\mathcal{O}^{(k)}$。
固定一个有限截断深度 $K\in\NN$，记
\[
  \mathcal{O}_{\le K}:=\bigsqcup_{k=1}^{K}\mathcal{O}^{(k)}.
\]
把事件类型集合记为 $\mathcal{K}:=\{1,\dots,K\}$，并约定：
第 $k$ 类同伦事件对应“从 $\mathcal{O}^{(k)}$ 中选取一个原子算子”。
\end{definition}

\begin{definition}[D2（一次回放：词轨迹 + 证书轨迹）]
\label{def:tex27-ch13-d2}
固定同一初始状态 $s_0\in\mathcal{V}$（定义~\ref{def:tex27-ch9-d1}），并固定同一观测证书口径 $J(s)\ge 0$
（例如取 $J=\JacGap$，定义~\ref{def:tex27-ch5-odthb-j}）。
一次回放（第 $r$ 次，$r=1,2,\dots$）给出：
\begin{enumerate}
  \item 一个长度为 $L$ 的算子词（路径）
  \[
    w^{(r)}:=o^{(r)}_1\circ\cdots\circ o^{(r)}_L\in (\mathcal{O}_{\le K})^{*};
  \]
  \item 相应的状态轨迹 $s^{(r)}_{t+1}=o^{(r)}_{t+1}\!\bigl(s^{(r)}_t\bigr)$，且 $s^{(r)}_0=s_0$；
  \item 相应的证书轨迹 $J^{(r)}_t:=J\!\bigl(s^{(r)}_t\bigr)$（$t=0,\dots,L$）；
  \item 事件口径计数 $N^{(r)}_k$（$k\in\mathcal{K}$），用于构造烈度向量（定义~\ref{def:tex27-ch5-odthb-tau}）。
\end{enumerate}
\end{definition}

\begin{definition}[D3（同伦扭曲确定发生次数：烈度向量固定）]
\label{def:tex27-ch13-d3}
称一组回放满足“同伦扭曲确定发生次数”，若存在常量烈度向量
\[
  \overline{\tau}=(\overline{\tau}_1,\dots,\overline{\tau}_K)\in\ZZ_{\ge 0}^{K},
\]
使得对任意回放 $r$，其由事件口径计数得到的总烈度向量满足
\[
  \tau^{(r)}=\overline{\tau}.
\]
并定义总扭曲步数为
\[
  L:=\sum_{k=1}^{K}\overline{\tau}_k.
\]
\end{definition}

\begin{definition}[D4（修复方式差异性）]
\label{def:tex27-ch13-d4}
称一组回放具有“修复方式差异性”，若存在 $r_1\neq r_2$ 使得
\[
  w^{(r_1)}\neq w^{(r_2)}.
\]
（直观：微观路径不唯一，但可以在宏观证书上出现一致性。）
\end{definition}

\begin{definition}[D5（修复力与其收敛性：净下降证书）]
\label{def:tex27-ch13-d5}
对第 $r$ 次回放，定义其\textbf{修复力}为证书的净下降：
\[
  \mathrm{Pow}^{(r)}:=J^{(r)}_0-J^{(r)}_L\ \in\ \RR.
\]
称该回放族具有“修复力收敛性”，若存在极限 $\mathrm{Pow}_\infty\in\RR$ 使得
\[
  \mathrm{Pow}^{(r)}\xrightarrow[r\to\infty]{}\mathrm{Pow}_\infty.
\]
\end{definition}

\begin{definition}[D6（受限路径空间与变分目标）]
\label{def:tex27-ch13-d6}
在固定烈度向量 $\overline{\tau}$ 与有限截断 $\mathcal{O}_{\le K}$ 下，定义受限路径空间
\[
  \Omega_{\overline{\tau}}
  :=\Bigl\{w\in(\mathcal{O}_{\le K})^*:\ \text{$w$ 的事件计数等于 }\overline{\tau}\Bigr\}.
\]
对任意 $w\in\Omega_{\overline{\tau}}$，定义其终端能量（以证书作能量）
\[
  E(w;s_0):=J\bigl(w(s_0)\bigr).
\]
对应的零温变分问题为
\[
  \min_{w\in\Omega_{\overline{\tau}}} E(w;s_0).
\]
这与“把路径当作动作词并做变分选择”的口径一致（定理~\ref{thm:tex27-ch10-t2}）。
\end{definition}

\subsection{公理（降级为定理）：回放三要素的证书化输入接口}
\label{subsec:tex27-ch13-axioms}

\begin{theorem}[公理（降级为定理）A14（回放三要素：差异性 + 收敛性 + 烈度确定）]
\label{thm:tex27-ch13-a14}
给定一组可回放数据。
若它满足定义~\ref{def:tex27-ch13-d3}（烈度向量固定）、定义~\ref{def:tex27-ch13-d4}（修复方式差异性）、
以及定义~\ref{def:tex27-ch13-d5}（修复力收敛性），
则我们称“回放三要素前提”在该数据上成立，并可作为后续定理的证书输入接口。
\end{theorem}

\begin{Certificate}[证书 A14（回放检验流程：从日志抽取 $\overline{\tau}$ 与 $\mathrm{Pow}_\infty$）]
\label{cert:tex27-ch13-a14}
证书为一套可复核流程：
\begin{enumerate}
  \item（烈度确定）对每次回放 $r$，由事件口径计数计算 $\tau^{(r)}$（定义~\ref{def:tex27-ch5-odthb-tau}），验证其恒等于同一向量 $\overline{\tau}$；
  \item（方式差异）对 $w^{(r)}$ 做词级比较，验证存在 $r_1\neq r_2$ 使 $w^{(r_1)}\neq w^{(r_2)}$；
  \item（修复力收敛）计算序列 $\mathrm{Pow}^{(r)}=J^{(r)}_0-J^{(r)}_L$，并用收敛判据（例如 Cauchy 判据/拟合残差/置信区间收缩）验证其收敛到 $\mathrm{Pow}_\infty$。

\end{enumerate}
\end{Certificate}

\begin{proof}[证明（证书化验证流程，非解析推演）]
按证书~\ref{cert:tex27-ch13-a14} 给出的流程执行即可；其输出为“成立/不成立”的可审计结论。
因此 A14 可作为数据驱动定理被调用。证毕。
\end{proof}

\subsection{定理：变分修复的存在性与有效性（由回放收敛推出）}
\label{subsec:tex27-ch13-theorems}

\begin{theorem}[T1（回放收敛 $\Rightarrow$ 统计力学同伦类下的边缘算子变分修复存在且有效）]
\label{thm:tex27-ch13-t1}
在定理~\ref{thm:tex27-ch13-a14} 的前提下，进一步假设：
\begin{enumerate}
  \item（有限字母表）$\mathcal{O}_{\le K}$ 有限（工程上对应有限算子库）；
  \item（非平凡修复）$\mathrm{Pow}_\infty>0$；
  \item（可行性）受限路径空间 $\Omega_{\overline{\tau}}$ 非空（定义~\ref{def:tex27-ch13-d6}）。
\end{enumerate}
则存在一个算子词 $w^\star\in\Omega_{\overline{\tau}}$，满足：
\begin{enumerate}
  \item（存在性：零温变分解）$w^\star\in\operatorname*{arg\,min}_{w\in\Omega_{\overline{\tau}}}E(w;s_0)$；
  \item（有效性：至少达到回放收敛强度）
  \[
    J\bigl(w^\star(s_0)\bigr)\le J^{(r)}_L\quad(\forall r),
    \qquad
    J(s_0)-J\bigl(w^\star(s_0)\bigr)\ \ge\ \mathrm{Pow}_\infty;
  \]
  \item（规范场自由度与边缘算子实现）若每个原子算子 $o\in\mathcal{O}_{\le K}$ 都可实现为边缘算子
  （只在修复子丛上修改规范场，定义~\ref{def:tex27-ch6-d4}），
  则 $w^\star$ 给出对规范场自由度 $A$ 的一条基于边缘算子的变分修复路径。
\end{enumerate}
\end{theorem}

\begin{Certificate}[证书 T1（有限性 $\Rightarrow$ 极小元存在；极小元 $\Rightarrow$ 下界比较；收敛 $\Rightarrow$ 有效性阈值）]
\label{cert:tex27-ch13-t1}
证书由三步组成：
\begin{enumerate}
  \item（有限性）若 $\mathcal{O}_{\le K}$ 有限且烈度/长度固定，则 $\Omega_{\overline{\tau}}$ 为有限集合；
  \item（存在性见证）在有限集合上，函数 $w\mapsto E(w;s_0)$ 必有极小元；取任一极小元作为 $w^\star$；
  \item（有效性）由“极小元不大于任意样本值”与 $\mathrm{Pow}^{(r)}\to\mathrm{Pow}_\infty$ 得到净下降下界。
\end{enumerate}
\end{Certificate}

\begin{proof}[证明（谓词因果链）]
\begin{enumerate}
  \item（$\Omega_{\overline{\tau}}$ 有限）由证书~\ref{cert:tex27-ch13-t1} 的第 1 点：
  $\mathcal{O}_{\le K}$ 有限且烈度向量 $\overline{\tau}$ 固定，故路径的“事件类型计数”与“每一步可选字母”均为有限，
  从而受限路径空间 $\Omega_{\overline{\tau}}$ 有限。
  \item（变分解存在）由于 $\Omega_{\overline{\tau}}$ 非空且有限，函数 $E(\cdot;s_0)$ 在其上必有极小元；
  取 $w^\star\in\operatorname*{arg\,min}_{w\in\Omega_{\overline{\tau}}}E(w;s_0)$，得到存在性。
  \item（对任意回放的比较）对任意回放 $r$，其 $w^{(r)}\in\Omega_{\overline{\tau}}$（由烈度固定定义与回放构造），
  且由极小性有
  \[
    J\bigl(w^\star(s_0)\bigr)=E(w^\star;s_0)\le E(w^{(r)};s_0)=J\bigl(w^{(r)}(s_0)\bigr)=J^{(r)}_L.
  \]
  \item（有效性下界）由上式对所有 $r$ 成立，得
  \[
    J(s_0)-J\bigl(w^\star(s_0)\bigr)\ \ge\ J(s_0)-J^{(r)}_L\ =\ \mathrm{Pow}^{(r)}.
  \]
  令 $r\to\infty$，并用 $\mathrm{Pow}^{(r)}\to\mathrm{Pow}_\infty>0$（定义~\ref{def:tex27-ch13-d5}），得
  \[
    J(s_0)-J\bigl(w^\star(s_0)\bigr)\ \ge\ \mathrm{Pow}_\infty,
  \]
  即 $w^\star$ 至少达到回放收敛强度，因此有效。
  \item（边缘算子实现）若每个原子算子都可实现为边缘算子（定义~\ref{def:tex27-ch6-d4}），则其复合词 $w^\star$ 仍为边缘算子词，
  因而给出对规范场自由度的基于边缘算子的修复路径（只在修复子丛上注入/调整 $A_{\perp}$，对齐定理~\ref{thm:tex27-ch6-thm2} 的变分口径）。
\end{enumerate}
综上，结论成立。证毕。
\end{proof}

\begin{corollary}[C1（统计力学同伦类解释：有限温 Gibbs 采样的变分口径）]
\label{cor:tex27-ch13-c1}
在定义~\ref{def:tex27-ch13-d6} 的路径空间上，对任意 $\beta>0$ 定义 Gibbs 分布
\[
  p_\beta(w\mid s_0)
  :=\frac{1}{Z(s_0)}\exp\bigl(-\beta E(w;s_0)\bigr),
  \qquad
  Z(s_0)=\sum_{w\in\Omega_{\overline{\tau}}}\exp\bigl(-\beta E(w;s_0)\bigr).
\]
则该分布给出“统计力学同伦类”的一个可审计代表：它在有限温下对多条修复路径赋予非零权重（解释“方式差异性”），
而在 $\beta\to\infty$ 的零温极限中，其质量集中到定理~\ref{thm:tex27-ch13-t1} 的变分解集合 $\operatorname*{arg\,min}E$（对齐引理~\ref{lem:tex27-ch8-l1} 的零温口径）。
\end{corollary}

\begin{Certificate}[证书 C1（有限空间上的 Gibbs 正则化与零温集中）]
\label{cert:tex27-ch13-c1}
证书为两点：
\begin{enumerate}
  \item（有限性）$\Omega_{\overline{\tau}}$ 有限时，$Z(s_0)$ 有限且 $p_\beta$ 定义良好；
  \item（零温集中）对有限集合，$\beta\to\infty$ 时 $e^{-\beta E}$ 的归一化权重集中在最小能量点集上。
\end{enumerate}
\end{Certificate}

\begin{proof}[证明（谓词因果链）]
按证书~\ref{cert:tex27-ch13-c1}：
\begin{enumerate}
  \item 由于 $\Omega_{\overline{\tau}}$ 有限，配分函数 $Z(s_0)$ 是有限项求和且为正，故 $p_\beta$ 定义良好；
  \item 设 $E_{\min}:=\min_{w\in\Omega_{\overline{\tau}}}E(w;s_0)$。
  则任意 $w$ 的权重为 $\exp(-\beta(E(w)-E_{\min}))$；当 $E(w)>E_{\min}$ 时，该因子随 $\beta\to\infty$ 趋于 $0$，
  从而归一化后质量集中到 $\{w:E(w)=E_{\min}\}=\operatorname*{arg\,min}E$。
\end{enumerate}
证毕。
\end{proof}

%%%%%%%%%%%%%%%%%%%%%%%%%%%%%%%%%%%%%%%%%%%%%%%%%%%%%%%%%%%%
\section{宏观有序与微观隧穿：边缘算子变分修复存在性与有效性的等价刻画（强命题）}
\label{sec:tex27-ch14-order-tunneling-equivalence}
%%%%%%%%%%%%%%%%%%%%%%%%%%%%%%%%%%%%%%%%%%%%%%%%%%%%%%%%%%%%

本节给出一个\textbf{独立章节}，严格化并证明如下强命题的\textbf{等价}关系：
\begin{quote}
  基于同伦修复塔的统计力学同伦类中，对规范场自由度（边上的连接/1-上链）做\textbf{基于边缘算子的变分修复}，
  其\textbf{存在性与有效性}，等价于\textbf{微观隧穿机制}之下的\textbf{宏观有序}：
  宏观层面表现为障碍缺口（例如 $\JacGap$）趋于消失/受控；
  微观层面表现为 Metropolis/Langevin 型扰动使系统可跨越局部势垒并实现有效混合（不可约）。
  宏观有序在纯数卷对应 GZ-OHU 的超限递归极限，在工程卷对应有限截断的正交子丛。
\end{quote}

为避免语义滑移，本节把“宏观有序/微观隧穿”分别落到：
\begin{enumerate}
  \item 宏观有序 $\Longleftrightarrow$ \textbf{存在}一个（有限截断或 OHU 极限下的）变分极小修复，使缺口证书达到目标精度；
  \item 微观隧穿 $\Longleftrightarrow$ \textbf{能够有效到达}该极小集合：路径空间链不可约并在有限混合时间内逼近 Gibbs 稳态。
\end{enumerate}

\subsection{定义组：宏观有序参数、微观隧穿与“存在/有效”口径}
\label{subsec:tex27-ch14-defs}

\begin{definition}[D1（宏观有序参数：规范场缺口证书）]
\label{def:tex27-ch14-d1}
令 $\Gamma=(V,E)$ 为骨架图。规范场自由度可取为边上的 1-上链 $A\in C^1(\Gamma)$（定义~\ref{def:tex27-ch6-d1}）。
令缺口/障碍强度由某个非负证书函数刻画：
\[
  J(A)\ge 0.
\]
典型两种可兼容选择：
\begin{enumerate}
  \item（回路残差范数）取基本回路残差 $r(e)=\Omega(c_e)$（定理~\ref{thm:tex27-ch6-thm1}），并令
  $
    J(A):=|r|=\big(\sum_{e\in C}r(e)^2\big)^{1/2}
  $
  （定义~\ref{def:tex27-ch6-d5}）；
  \item（总 gap）取 $J=\JacGap$（定义~\ref{def:tex27-ch5-odthb-j}），它是对分量残差的加权压缩。
\end{enumerate}
对精度阈值 $\epsilon>0$，定义宏观有序集合（低缺口相）
\[
  \mathcal{A}_{\le\epsilon}:=\{A\in C^1(\Gamma):\ J(A)\le \epsilon\}.
\]
\end{definition}

\begin{definition}[D2（边缘算子变分修复：对规范场自由度的可达修复族）]
\label{def:tex27-ch14-d2}
称“边缘算子变分修复”是指：只在修复子丛上注入 $A_{\perp}\in C^1_{\perp}$（定义~\ref{def:tex27-ch6-d4}），
并以某个变分泛函 $\Phi(A_{\perp})$ 为目标做极小化（定理~\ref{thm:tex27-ch6-thm2} 给出 $H^1$ 情形的严格凸模板）。
若存在可执行的边缘算子词 $w\in\mathcal{O}^*$ 使其在规范场口径下等价于某个 $A_{\perp}$ 注入，
则称该变分修复在同伦修复塔内\textbf{可达}（对齐定义~\ref{def:tex27-ch9-d2} 与推论~\ref{cor:tex27-ch6-c1}）。
\end{definition}

\begin{definition}[D3（微观隧穿：不可约 + 有温接受率/噪声）]
\label{def:tex27-ch14-d3}
在有限路径空间 $\Omega_L$（定义~\ref{def:tex27-ch11-d1}，字母表可取为同伦修复塔的有限截断 $\mathcal{O}_{\le K}$，定义~\ref{def:tex27-ch13-d1}）上，
用提议核 $Q$ 与 Metropolis 接受率构造扭曲核 $K$（公理~\ref{thm:tex27-ch11-a12}）。
若链在 $\Omega_L$ 上\textbf{不可约}（从任意 $w$ 可通过有限步到达任意 $w'$），则称发生“微观隧穿”（备注~\ref{rem:tex27-ch11-r2}）；
其物理含义是：在有限温 $\beta<\infty$ 或含噪口径下，存在正概率接受/实现能量上升的微小跳变，从而跨越局部势垒并避免永久卡死在局部连通分支。

若进一步需要“收敛到稳态/定量效率”，则可在不可约之外再加入非周期性以保证混合时间有限（参见公理~\ref{thm:tex27-ch11-a13} 的口径）。
\end{definition}

\begin{definition}[D4（“存在性/有效性”作为两条可审计性质）]
\label{def:tex27-ch14-d4}
在固定初始规范场 $A_0$ 下，把“边缘算子词”视为对 $A_0$ 的可执行修复路径。
令
\[
  E(w;A_0):=J\bigl(w(A_0)\bigr),
  \qquad
  E^\ast(A_0):=\inf_{w}E(w;A_0),
\]
其中 $w$ 遍历给定的（有限截断/有限预算）可达词空间（例如第~\ref{sec:tex27-ch11-robust-time-space}~节的 $\Omega_L$ 或第~\ref{sec:tex27-ch13-replay-variational-repair}~节的受限路径空间）。
称“边缘算子变分修复”满足：
\begin{enumerate}
  \item（存在性/宏观有序）存在精度阈值 $\epsilon>0$ 使 $E^\ast(A_0)\le \epsilon$（等价于存在 $w^\star$ 使 $J(w^\star(A_0))\le\epsilon$）；
  \item（有效性）存在统计力学扭曲过程 $\{W_t\}$（定义~\ref{def:tex27-ch11-d3} 或定义~\ref{def:tex27-ch13-d6} 的 Gibbs 口径），
  使得对任意失败概率 $\delta\in(0,1)$ 与任意容忍 $\eta>0$，存在参数选择（$\beta$ 与 $t$）满足
  \[
    \mathbb{P}\Bigl(\min_{0\le j\le t}E(W_j;A_0)\le E^\ast(A_0)+\eta\Bigr)\ge 1-\delta.
  \]
\end{enumerate}
\end{definition}

\subsection{定理：存在性 $\Longleftrightarrow$ 宏观有序；有效性 $\Longleftrightarrow$ 微观隧穿}
\label{subsec:tex27-ch14-theorems}

\begin{lemma}[L1（微观隧穿的最小可用结论：有限时间命中近最优集合）]
\label{lem:tex27-ch14-l1}
令 $\Omega$ 为有限集合，$\{W_t\}_{t\ge 0}$ 为其上的马尔可夫链。
若链在 $\Omega$ 上不可约，则对任意非空集合 $A\subseteq \Omega$，存在常数 $m\in\NN$ 与 $p_{\min}\in(0,1]$，
使得对任意初始状态 $W_0=w$：
\[
  \mathbb{P}_w(\tau_A\le m)\ge p_{\min},
  \qquad
  \tau_A:=\inf\{t\ge 0:\ W_t\in A\}.
\]
从而对任意 $\delta\in(0,1)$，取 $n\in\NN$ 满足 $(1-p_{\min})^n\le \delta$，则对任意 $w$ 都有
\[
  \mathbb{P}_w(\tau_A\le n m)\ge 1-\delta.
\]
\end{lemma}

\begin{Certificate}[证书 L1（有限状态 + 不可约 $\Rightarrow$ 最短路径概率下界 + 几何尾界）]
\label{cert:tex27-ch14-l1}
证书由两步构造：
\begin{enumerate}
  \item（一步下界）对每个 $w\in\Omega$，不可约性保证存在某条长度 $m_w$ 的正概率路径把 $w$ 送入 $A$，
  因而 $p_w:=\mathbb{P}_w(\tau_A\le m_w)>0$；
  令 $m:=\max_w m_w$，并令 $p_{\min}:=\min_w \mathbb{P}_w(\tau_A\le m)$，则 $p_{\min}>0$；
  \item（几何尾界）用强马尔可夫性质在每个长度为 $m$ 的区块上迭代：条件在任意状态下，“接下来 $m$ 步内命中 $A$”的概率均 $\ge p_{\min}$，
  故不命中的概率至多按 $(1-p_{\min})^n$ 递减。
\end{enumerate}
\end{Certificate}

\begin{proof}[证明（谓词因果链）]
按证书~\ref{cert:tex27-ch14-l1}：
\begin{enumerate}
  \item（统一步长与正下界）对每个 $w$，取一条以正概率在 $m_w$ 步内到达 $A$ 的路径，得 $p_w>0$。
  取 $m=\max_w m_w$，则对任意 $w$ 有 $\mathbb{P}_w(\tau_A\le m)\ge \mathbb{P}_w(\tau_A\le m_w)=p_w$。
  由于 $\Omega$ 有限，$p_{\min}:=\min_w\mathbb{P}_w(\tau_A\le m)$ 为有限个正数的最小值，故 $p_{\min}>0$。
  \item（几何尾界）对任意 $n\in\NN$，考虑事件 $\{\tau_A>nm\}$。
  在条件 $\{\tau_A>(n-1)m\}$ 下，时刻 $(n-1)m$ 的状态为某个 $w\in\Omega$；
  由第一步的统一下界，接下来 $m$ 步仍以至少 $p_{\min}$ 的概率命中 $A$，故
  \[
    \mathbb{P}(\tau_A>nm\mid \tau_A>(n-1)m)\le 1-p_{\min}.
  \]
  递推得到 $\mathbb{P}(\tau_A>nm)\le (1-p_{\min})^n$，
  等价于 $\mathbb{P}(\tau_A\le nm)\ge 1-(1-p_{\min})^n$。
\end{enumerate}
证毕。
\end{proof}

\begin{theorem}[强命题 S1（边缘算子变分修复的存在/有效 $\Longleftrightarrow$ 宏观有序/微观隧穿）]
\label{thm:tex27-ch14-s1}
在定义~\ref{def:tex27-ch14-d1}--\ref{def:tex27-ch14-d4} 的口径下，有如下等价链：
\begin{enumerate}
  \item（存在性 $\Longleftrightarrow$ 宏观有序（OHU/截断正交子丛））
  \[
    E^\ast(A_0)\le\epsilon
    \quad\Longleftrightarrow\quad
    A_0\ \text{在精度 }\epsilon\text{ 下可被“宏观有序化”}.
  \]
  并且该“宏观有序化”在纯数卷对应 OHU 极限 $\JacGap\to 0$（公理~\ref{thm:tex27-ch7-a7} 与定理~\ref{thm:tex27-ch7-thm1}），
  在工程卷对应有限正交子丛截断 $\JacGap(L^{(k)})\le\epsilon$（公理~\ref{thm:tex27-ch7-b7} 与定理~\ref{thm:tex27-ch7-thm2}）。
  \item（有效性 $\Longleftrightarrow$ 微观隧穿（不可约；混合为可选强化））
  对于以 Metropolis 核实现的统计力学同伦类（公理~\ref{thm:tex27-ch11-a12}；若需要定量效率再加公理~\ref{thm:tex27-ch11-a13}），
  “以高概率到达近最优能量层（$E\le E^\ast+\eta$）”与“链不可约（微观隧穿）”在本形式系统中等价地刻画为：
  \[
    \begin{aligned}
      &\forall(\delta,\eta)\ \exists(\beta,t):\
      \mathbb{P}\Bigl(\min_{0\le j\le t}E(W_j;A_0)\le E^\ast(A_0)+\eta\Bigr)\ge 1-\delta\\
      &\Longleftrightarrow\ \text{微观隧穿（定义~\ref{def:tex27-ch14-d3}）}.
    \end{aligned}
  \]
\end{enumerate}
从而“存在性 + 有效性”等价于“宏观有序 + 微观隧穿”。
\end{theorem}

\begin{Certificate}[证书 S1（引用链：正交子丛/截断（第 6/7 节）+ 隧穿/混合（第 11 节））]
\label{cert:tex27-ch14-s1}
证书由两组并列引用组成：
\begin{enumerate}
  \item（存在性 $\leftrightarrow$ 有序）$H^1$ 情形的边缘算子变分极小解由定理~\ref{thm:tex27-ch6-thm2} 给出；
  OHU/截断正交子丛的“缺口趋零/受控”由推论~\ref{cor:tex27-ch7-c1}（等价引用定理~\ref{thm:tex27-ch7-thm1}--\ref{thm:tex27-ch7-thm2}）给出。
  \item（有效性 $\leftrightarrow$ 隧穿）“有效隧穿=不可约性”见备注~\ref{rem:tex27-ch11-r2}；
  不可约 $\Rightarrow$ 有限时间命中近最优集合由引理~\ref{lem:tex27-ch14-l1} 给出；
  若进一步需要“某个固定时刻的概率界”与混合时间的显式关系，可调用命题~\ref{prop:tex27-ch11-p1}（强于本节最小结论）。
\end{enumerate}
\end{Certificate}

\begin{proof}[证明（谓词因果链）]
按证书~\ref{cert:tex27-ch14-s1}，分两部分证明：
\begin{enumerate}
  \item（存在性 $\Rightarrow$ 宏观有序）若存在可达修复 $w^\star$ 使 $J(w^\star(A_0))\le\epsilon$，
  则 $A^\star=w^\star(A_0)\in\mathcal{A}_{\le\epsilon}$，宏观有序集合非空，且“低缺口相”可达，故称 $A_0$ 在精度 $\epsilon$ 下可被宏观有序化。
  反向若称“宏观有序化”成立，则按定义它正是“存在某个可达修复使 $J\le\epsilon$”，即存在性。
  在纯数/工程两种语义下，OHU 极限与有限截断正交子丛分别由推论~\ref{cor:tex27-ch7-c1} 固化，
  且第~\ref{sec:tex27-ch6-edge-operator}~节给出“边缘算子 $\Rightarrow$ 正交子丛变分修复”的可实现接口（推论~\ref{cor:tex27-ch6-c1}）。
  \item（有效性 $\Rightarrow$ 微观隧穿）若“以高概率到达近最优能量层（$E\le E^\ast+\eta$）”对任意初始微观词状态都成立，
  则路径空间上不存在把过程永远困在某个闭类、而该闭类又与近最优集合
  $
    \{w:\ E(w;A_0)\le E^\ast(A_0)+\eta\}
  $
  不相交的情形；
  否则从该闭类出发将永远无法达到近最优集合，与“任意 $\delta$ 可把失败概率压低”矛盾。
  因此链在相关路径空间上必须具备可达遍历性（不可约），即“微观隧穿”的严格化口径（备注~\ref{rem:tex27-ch11-r2}）。
  \item（微观隧穿 $\Rightarrow$ 有效性）若链不可约，则对任意 $\eta>0$ 近最优集合 $A_{\eta}$ 非空且可达；
  套用引理~\ref{lem:tex27-ch14-l1} 即可把命中概率压到任意 $1-\delta$。
  记近最优集合
  \[
    A_{\eta}:=\{w:\ E(w;A_0)\le E^\ast(A_0)+\eta\}.
  \]
  由存在性定义，$A_{\eta}$ 非空。
  由引理~\ref{lem:tex27-ch14-l1}，不可约性推出：对任意 $\delta$ 存在时间 $t$ 使
  $\mathbb{P}(\tau_{A_{\eta}}\le t)\ge 1-\delta$，
  其中 $\tau_{A_{\eta}}$ 是命中 $A_{\eta}$ 的首次时间。
  这等价于
  \[
    \mathbb{P}\Bigl(\min_{0\le j\le t}E(W_j;A_0)\le E^\ast(A_0)+\eta\Bigr)\ge 1-\delta,
  \]
  即定义~\ref{def:tex27-ch14-d4} 的有效性口径。
\end{enumerate}
综上，“存在性 + 有效性”等价于“宏观有序 + 微观隧穿”。证毕。
\end{proof}

\begin{remark}[工程提示：把“宏观有序/微观隧穿”落盘为可审计字段]
若要把本节强命题进一步工程化闭合，最小落盘字段可取：
\[
  (\JacGap,\ |r|,\ k,\ \beta,\ t_{\mathrm{mix}},\ \text{词 }W_t,\ \text{可达性/不可约性测试日志}).
\]
其中 $k$ 对应“正交子丛数量/修复塔深度”（第~\ref{sec:tex27-ch7-ohu-orthogonal}~节），
而不可约性与混合时间估计对应“微观隧穿证书输入”（第~\ref{sec:tex27-ch11-robust-time-space}~节）。
\end{remark}

%%%%%%%%%%%%%%%%%%%%%%%%%%%%%%%%%%%%%%%%%%%%%%%%%%%%%%%%%%%%
\section{自动化科研的数学基座：同伦修复与规范场变分（强命题链）}
\label{sec:tex27-ch15-autoscience-foundation}
%%%%%%%%%%%%%%%%%%%%%%%%%%%%%%%%%%%%%%%%%%%%%%%%%%%%%%%%%%%%

本节以\textbf{形式逻辑}口径给出一个独立强命题：在 \GFramework\ 的形式系统内，
本文把“同伦修复（homotopy repair）”与“规范场变分（gauge-field variational principle）”组织为一条\textbf{可证书化的条件定理链}，
从而确立一个可审计的\textbf{自动化科研/自动化工程设计}数学基座。

\begin{remark}[关于“首次”的严格边界]
“首次确立”若被理解为\textbf{学术史意义}的“第一次提出/第一次证明”，需要外部文献对比证书，本文不做该断言。
本节中“首次”仅指：\textbf{在本文档的形式系统内部}，第一次把该基座以“定义 + 证书 + 条件定理链”的方式闭合并显式标注边界条件。
\end{remark}

\subsection{定义：开放系统科研任务、同伦修复核与规范场变分核}
\label{subsec:tex27-ch15-defs}

\begin{definition}[D1（开放系统科研任务：自动化科研的最小形式化）]
\label{def:tex27-ch15-d1}
把“科研/工程设计/量化决策”统一抽象为开放系统任务：
给定变分空间 $\mathcal{V}$、对手/环境动作集合 $\mathcal{A}(s)$、缺陷势能 $\Phi$ 与约束谓词 $\mathrm{Inv}$，
并允许状态在对抗扰动 $a_t$ 下演化（定义~\ref{def:tex27-ch11-d2} 与命题~\ref{prop:tex27-ch10-p1}）。
在该口径下，“科研”可被理解为：构造一条可执行的修复/设计路径，使 $\Phi$ 下降且 $\mathrm{Inv}$ 保持，并能在扰动下在线重算。
\end{definition}

\begin{definition}[D2（同伦修复核：算子词空间上的可达修复）]
\label{def:tex27-ch15-d2}
令同伦修复塔给出算子全集 $\mathcal{O}$ 与词空间 $\mathcal{O}^*$（定义~\ref{def:tex27-ch9-d2}）。
把一次“修复/设计动作”定义为算子词 $w\in\mathcal{O}^*$，并在开放系统中按 $s'=w(s,a)$ 执行（定义~\ref{def:tex27-ch11-d2}）。
称该动作核为“同伦修复核”，因为它把“连续变形/同伦修补”离散化为可审计的词级路径（定义~\ref{def:tex27-ch10-d2}）。
\end{definition}

\begin{definition}[D3（规范场变分核：把修复自由度压缩到边自由度并可变分求解）]
\label{def:tex27-ch15-d3}
在有效 1 阶口径下，把可学习自由度压缩为边上的规范场（定义~\ref{def:tex27-ch6-d1} 与定义~\ref{def:tex27-ch5-d2}）。
“规范场变分核”指：在修复子丛/正交子丛上选择 $A_{\perp}$ 来极小化一个可证书化的变分泛函，
从而得到显式修复场（定理~\ref{thm:tex27-ch6-thm2}）。
它提供了一个\textbf{可构造、可复核、可落盘}的“微观修复步”接口（推论~\ref{cor:tex27-ch6-c1}）。
\end{definition}

\begin{definition}[D4（自动化科研数学基座：证书化闭合条件）]
\label{def:tex27-ch15-d4}
称本文确立了“自动化科研数学基座”，当且仅当在本文形式系统内，存在一条可审计链路把
\[
  \text{同伦修复核（定义~\ref{def:tex27-ch15-d2}）}
  \ \&\
  \text{规范场变分核（定义~\ref{def:tex27-ch15-d3}）}
\]
提升为对开放系统科研任务（定义~\ref{def:tex27-ch15-d1}）的\textbf{可审计在线控制/设计}机制，
并且每一步都能输出证书以支持回放复核（定义~\ref{def:tex27-ch9-d7}）。
\end{definition}

\subsection{定理：基座的确立（条件定理链）}
\label{subsec:tex27-ch15-theorems}

\begin{theorem}[强命题 F1（\GFramework\ 自动化科研数学基座：同伦修复 + 规范场变分）]
\label{thm:tex27-ch15-f1}
在本文形式系统内，定义~\ref{def:tex27-ch15-d4} 的“自动化科研数学基座”成立，
并且其两条核心支柱分别对应：
\begin{enumerate}
  \item（规范场变分 $\Rightarrow$ 可构造修复步）规范场自由度的变分修复在有效 1 阶上可构造且唯一（定理~\ref{thm:tex27-ch6-thm2}）；
  \item（同伦修复 $\Rightarrow$ 路径空间变分 + 在线控制）把动作提升为词空间 $\mathcal{O}^*$ 并做变分选择即可得到 ORL 的统一实现口径（定理~\ref{thm:tex27-ch10-t2}），
  并可在开放系统中形成即时修复/在线控制（命题~\ref{prop:tex27-ch10-p1}）。
\end{enumerate}
因此：在满足外部证书输入（算子完备性 $\mathrm{Complete}$ 与微观隧穿/混合等）时，
该基座可对开放系统博弈任务给出可审计的自动化修复/设计过程（含可回放证书）。
\end{theorem}

\begin{Certificate}[证书 F1（证书链：有效 1 阶变分修复 + 路径空间变分 + 白盒可审计）]
\label{cert:tex27-ch15-f1}
证书为如下串联引用：
\begin{enumerate}
  \item（规范场变分核）定理~\ref{thm:tex27-ch6-thm2} 给出正交子丛内的严格凸变分极小与显式解；
  \item（同伦修复核）定理~\ref{thm:tex27-ch10-t2} 把 ORL 统一实现为“在路径空间上的变分选择”；
  \item（在线控制）命题~\ref{prop:tex27-ch10-p1} 把该选择机制解释为开放系统中的即时修复/在线控制；
  \item（白盒可审计）定义~\ref{def:tex27-ch9-d7} 与命题~\ref{prop:tex27-ch9-p1} 给出“动作词 + 证书接口 $\Leftrightarrow$ 可审计白盒 AI”的闭合口径；
  \item（宏观有序/微观隧穿语义对齐）第~\ref{sec:tex27-ch14-order-tunneling-equivalence}~节把“存在/有效”与“宏观有序/微观隧穿”严格等价化，
  从而把统计力学同伦类的微观采样机制与宏观缺口收敛绑定到同一证书体系内。
\end{enumerate}
\end{Certificate}

\begin{proof}[证明（谓词因果链）]
按证书~\ref{cert:tex27-ch15-f1}：
\begin{enumerate}
  \item（构造性）由定理~\ref{thm:tex27-ch6-thm2}，在有效 1 阶上存在可构造且唯一的变分极小修复场 $A_{\perp}^\star$，从而“规范场变分核”成立；
  \item（路径空间变分）由定理~\ref{thm:tex27-ch10-t2}，把动作提升为词空间并按 Gibbs/门控做选择，即得到在路径空间上的变分选择（同伦修复核的统一实现）；
  \item（开放系统）由命题~\ref{prop:tex27-ch10-p1}，上述路径空间选择机制可直接作为开放系统的在线控制/即时修复口径；
  \item（可审计性）由定义~\ref{def:tex27-ch9-d7}，若每一步输出动作词与证书，则系统为可审计白盒 AI；
  由命题~\ref{prop:tex27-ch9-p1}，该口径与 ORL 结构等价，从而“在线修复/设计过程”可被证书化与回放审计；
  \item（宏观/微观闭合）由第~\ref{sec:tex27-ch14-order-tunneling-equivalence}~节，
  “宏观有序（缺口趋于受控/趋零）”与“微观隧穿（不可约采样）”分别等价刻画了存在性与有效性，
  因此该基座在统计力学同伦类口径下闭合。
\end{enumerate}
综上，定义~\ref{def:tex27-ch15-d4} 成立，即本文在形式逻辑层面确立了以“同伦修复 + 规范场变分”为核心的自动化科研数学基座。证毕。
\end{proof}

\begin{corollary}[推论 F1（应用域是条件式：编码进 $(\mathcal{V},\mathcal{O},\Phi,\mathrm{Inv})$ 即可调用基座）]
\label{cor:tex27-ch15-f1}
对任意应用域（例如自动化工程设计、金融量化交易、工业级材料/超导探索等），
若能给出可审计编码
\[
  (\mathcal{V},\mathcal{O},\Phi,\mathrm{Inv})\in\mathfrak{T}(\mathcal{V})
\]
并补齐必要外部证书输入（例如算子完备性、公理~\ref{thm:tex27-ch9-a9}；微观隧穿/混合证书输入、公理~\ref{thm:tex27-ch11-a12}--\ref{thm:tex27-ch11-a13}），
则可把该域问题纳入本文的自动化科研基座并获得“可审计修复/设计路径 + 统计意义有效性”的严格口径。
\end{corollary}

\begin{Certificate}[证书 F1（调用：AGI\textsuperscript{cert}（证书化 AGI）的任务类接口）]
\label{cert:tex27-ch15-f1-application}
证书为：定义~\ref{def:tex27-ch9-d8} 已把任务类 $\mathfrak{T}(\mathcal{V})$ 作为接口；
推论~\ref{cor:tex27-ch15-f1} 只是把“应用域问题”视为该接口的具体实例，并显式标注所需外部证书输入。
\end{Certificate}

\begin{proof}[证明（谓词因果链）]
按证书~\ref{cert:tex27-ch15-f1-application}：一旦应用域问题可被编码为任务 $(\Phi,\mathrm{Inv})$，即可调用本文关于 ORL/路径空间变分与证书化的定理链；
外部证书输入用于保障“存在性/有效性”的边界条件，因此该推论是条件式成立。证毕。
\end{proof}

\subsection{聚类定理：典型应用域作为开放系统博弈任务族}
\label{subsec:tex27-ch15-clustering}

\begin{definition}[D5（开放系统博弈问题簇：把“应用域差异”压缩为同一任务接口）]
\label{def:tex27-ch15-d5}
在本文口径下，一个“开放系统博弈问题”并不要求对手一定具有人格化智能；
只要存在外部不受控因素（自然噪声、对抗扰动、市场对手、实验不确定性等）以动作形式进入系统演化，
即可把它统一写成定义~\ref{def:tex27-ch11-d2} 的开放系统对抗能量形式
\[
  E(w;s)=\sup_{a\in\mathcal{A}(s)}\Phi\bigl(w(s,a)\bigr).
\]
据此定义“开放系统博弈问题簇”为
\[
  \mathsf{OpenGame}:=\Bigl\{
  \text{问题域 } \mathcal{D}\ \Big|\
  \exists\,(\mathcal{V},\mathcal{O},\mathcal{A},\Phi,\mathrm{Inv}),\
  \mathcal{D}\xrightarrow{\mathrm{Enc}}(\mathcal{V},\mathcal{O},\mathcal{A},\Phi,\mathrm{Inv})
  \text{ 且满足定义~\ref{def:tex27-ch11-d2}}
  \Bigr\}.
\]
这里 $\mathrm{Enc}$ 是“可审计编码”：它把领域对象映射为 $(s,a,w)$ 语义（状态/环境动作/算子词），
并把领域目标映射为势能 $\Phi$ 与约束谓词 $\mathrm{Inv}$（见定义~\ref{def:tex27-ch15-d1} 与定义~\ref{def:tex27-ch9-d7}）。
\end{definition}

\begin{definition}[D6（四类典型簇：元科研/工程设计/金融交易/实验材料探索）]
\label{def:tex27-ch15-d6}
在簇 $\mathsf{OpenGame}$ 内，引入四个\textbf{可审计的粗粒度子簇}（仅用于“聚类说明”，不改变任何定理的适用性）：
\begin{enumerate}
  \item $\mathsf{C}_{\mathrm{Meta}}$（元科研簇）：状态含“知识/假设/证据”的可更新表示，环境动作对应“反例/新数据/审稿式约束变化”等外部冲击；
  \item $\mathsf{C}_{\mathrm{Design}}$（工程设计簇）：状态含“设计变量 + 外部工况/扰动”，算子词对应“可执行的设计修改/结构修复”，$\Phi$ 对应“性能缺陷 + 违约惩罚”；
  \item $\mathsf{C}_{\mathrm{Market}}$（金融交易簇）：状态含“市场可观测量 + 持仓/资金”，环境动作对应“市场演化与他者行为”，算子词对应“可执行的交易/对冲操作”，$\Phi$
    对应“风险调整后的损失/负效用”；
  \item $\mathsf{C}_{\mathrm{Lab}}$（实验材料簇）：状态含“候选材料/工艺 + 实验条件 + 测量读数”，环境动作对应“实验噪声/不可控微观涨落/设备漂移”，算子词对应“可执行的实验方案与参数更新”，
    $\Phi$ 对应“目标性质偏差 + 资源成本”。
\end{enumerate}
以上子簇是谓词定义的集合：任何满足相应状态分解/动作语义者即属于该子簇。
\end{definition}

\begin{theorem}[强命题 F2（应用域聚类：自动化科研/工程设计/金融量化/超导探索 $\subseteq$ 开放系统博弈）]
\label{thm:tex27-ch15-f2}
下述四类应用域都属于开放系统博弈问题簇 $\mathsf{OpenGame}$（定义~\ref{def:tex27-ch15-d5}），并可进一步落到定义~\ref{def:tex27-ch15-d6} 的四个子簇中：
\[
  \mathcal{D}_{\mathrm{AutoScience}}\in\mathsf{C}_{\mathrm{Meta}},\quad
  \mathcal{D}_{\mathrm{AutoEng}}\in\mathsf{C}_{\mathrm{Design}},\quad
  \mathcal{D}_{\mathrm{Quant}}\in\mathsf{C}_{\mathrm{Market}},\quad
  \mathcal{D}_{\mathrm{SC}}\in\mathsf{C}_{\mathrm{Lab}}.
\]
其中 $\mathcal{D}_{\mathrm{AutoScience}}$ 表示自动化科研问题域，$\mathcal{D}_{\mathrm{AutoEng}}$ 表示自动化工程设计问题域，
$\mathcal{D}_{\mathrm{Quant}}$ 表示金融量化交易问题域，$\mathcal{D}_{\mathrm{SC}}$ 表示材料/超导等实验探索问题域。
\end{theorem}

\begin{Certificate}[证书 F2（四类应用域的“可审计编码”模板）]
\label{cert:tex27-ch15-f2}
给出四个领域到开放系统任务的编码模板（每个模板只要求“可落盘/可复核”，不要求唯一）：
\begin{enumerate}
  \item（自动化科研 $\mathcal{D}_{\mathrm{AutoScience}}$）
  取状态 $s=(K,H,D)\in\mathcal{V}$（知识库/假设/数据），
  环境动作 $a\in\mathcal{A}(s)$ 表示“新数据、反例、规则变化（审稿/实验条件变化）等”，
  算子词 $w\in\mathcal{O}^*$ 表示“提出假设、选择定理工具、设计实验、更新模型”的可执行序列；
  取 $\Phi(s)$ 为“目标未满足度/不一致度/误差”，取 $\mathrm{Inv}(s)$ 为“一致性/可验证性/预算”等硬约束谓词。
  \item（自动化工程设计 $\mathcal{D}_{\mathrm{AutoEng}}$）
  取状态 $s=(x,y)\in\mathcal{V}$（设计变量 $x$ 与工况/扰动 $y$），
  环境动作 $a\in\mathcal{A}(s)$ 驱动 $y$ 的变化或注入扰动，
  算子词 $w\in\mathcal{O}^*$ 表示“可制造/可装配的结构修改、参数更新、修复操作”；
  取 $\Phi(s)$ 为“性能缺陷 + 违约惩罚”，取 $\mathrm{Inv}(s)$ 为“安全/规范/物理可行”约束。
  \item（金融量化交易 $\mathcal{D}_{\mathrm{Quant}}$）
  取状态 $s=(m,\theta)\in\mathcal{V}$（市场观测 $m$ 与持仓/资金 $\theta$），
  环境动作 $a\in\mathcal{A}(s)$ 表示“市场演化 + 他者行为 + 执行噪声”等外部驱动，
  算子词 $w\in\mathcal{O}^*$ 表示“下单/撤单/对冲/再平衡”等可执行操作；
  取 $\Phi(s)$ 为“风险调整后的损失/负效用”（例如把风险预算/回撤以罚项写入势能），取 $\mathrm{Inv}(s)$ 为“风控/合规/仓位”约束。
  \item（材料/超导实验探索 $\mathcal{D}_{\mathrm{SC}}$）
  取状态 $s=(x,c,z)\in\mathcal{V}$（候选材料/工艺 $x$、实验条件 $c$、测量读数 $z$），
  环境动作 $a\in\mathcal{A}(s)$ 表示“不可控微观涨落、测量噪声、设备漂移”等，
  算子词 $w\in\mathcal{O}^*$ 表示“实验方案选择、参数更新、工艺调整、候选生成”的可执行序列；
  取 $\Phi(s)$ 为“目标性质偏差 + 资源成本/时间成本”，取 $\mathrm{Inv}(s)$ 为“物理可行/安全/实验预算”等约束。
\end{enumerate}
每个模板都天然满足“动作词 + 证书可回放”的接口（定义~\ref{def:tex27-ch9-d7}）。
\end{Certificate}

\begin{proof}[证明（谓词因果链）]
按证书~\ref{cert:tex27-ch15-f2}，对四个问题域分别给出一个显式编码 $\mathrm{Enc}$，并逐项核对：
\begin{enumerate}
  \item（开放系统性）四个模板都显式给出环境动作集合 $\mathcal{A}(s)$，且状态演化写成 $s'=w(s,a)$；
  因而都满足开放系统对抗口径的最小结构假设（定义~\ref{def:tex27-ch11-d2}），从而属于 $\mathsf{OpenGame}$（定义~\ref{def:tex27-ch15-d5}）。
  \item（子簇归属）四个模板的状态分解与动作语义分别满足定义~\ref{def:tex27-ch15-d6} 中
  \[
    \mathsf{C}_{\mathrm{Meta}},\quad
    \mathsf{C}_{\mathrm{Design}},\quad
    \mathsf{C}_{\mathrm{Market}},\quad
    \mathsf{C}_{\mathrm{Lab}}
  \]
  的谓词条件，故得到对应的子簇归属关系。
\end{enumerate}
综上，强命题~\ref{thm:tex27-ch15-f2} 成立。证毕。
\end{proof}

\begin{corollary}[推论 F2（聚类 $\Rightarrow$ 统一调用第15节基座定理链）]
\label{cor:tex27-ch15-f2}
对任一属于 $\mathsf{OpenGame}$ 的应用域（尤其是强命题~\ref{thm:tex27-ch15-f2} 所列四类），
一旦其编码满足外部证书输入（算子完备性、公理~\ref{thm:tex27-ch9-a9}；混合/隧穿、公理~\ref{thm:tex27-ch11-a12}--\ref{thm:tex27-ch11-a13}），
则可直接调用强命题~\ref{thm:tex27-ch15-f1} 的“同伦修复 + 规范场变分”基座，
获得可审计的在线修复/设计/决策过程与统计意义有效性口径。
\end{corollary}

\begin{proof}[证明（谓词因果链）]
由定义~\ref{def:tex27-ch15-d5}，$\mathsf{OpenGame}$ 的成员可被编码到 $(\mathcal{V},\mathcal{O},\mathcal{A},\Phi,\mathrm{Inv})$；
由强命题~\ref{thm:tex27-ch15-f1} 与推论~\ref{cor:tex27-ch15-f1}，在补齐外部证书输入后即可把该编码纳入基座定理链并输出可回放证书。
证毕。
\end{proof}

\begin{remark}[高风险语义提示：不在本文证明范围内的外推]
“工业级室温超导”等具体科学目标是否可在现实世界中达成，涉及实验事实与工程约束；
本文只给出\textbf{形式系统内}的“如何把此类目标编码为开放系统博弈任务并以证书化方式搜索/修复”的数学基座，
不把“现实必然成功”作为可证明结论。
\end{remark}

\clearpage

%%%%%%%%%%%%%%%%%%%%%%%%%%%%%%%%%%%%%%%%%%%%%%%%%%%%%%%%%%%%
\section{地磁爆扰动的开放对抗编码：调度与 RTSC 防御的 ORL-开放博弈等价（形式系统条目）}
\label{sec:tex27-ch16-geomagnetic-orl-equivalence}
%%%%%%%%%%%%%%%%%%%%%%%%%%%%%%%%%%%%%%%%%%%%%%%%%%%%%%%%%%%%

本节把“地磁爆（含电离层/磁层耦合）对长导体/长链路的扰动”统一编码为开放系统的对手动作，
并以 \GFramework\ 的形式系统口径给出一个可审计结论：
\textbf{电网/通信调度的抗扰优化}与\textbf{可防御地磁爆的工业级室温超导（RTSC）防御设计}
在 ORL-开放博弈接口下可视为同一骨架上的等价问题（能量口径保持）。

\subsection{定义：对手动作、两域编码与等价口径}
\label{subsec:tex27-ch16-defs}

\begin{definition}[D1（地磁爆扰动作为“开放对手动作”）]
\label{def:tex27-ch16-d1}
取网络骨架 $(V,E)$。令外部地磁爆（含电离层/磁层耦合）对长导体/长链路的影响，
抽象为对手/环境动作
\[
  a\in\mathcal{A}(s).
\]
其语义是一个“注入场/注入源”族，例如：
\begin{itemize}
  \item（电网语义）沿线路的等效感应电动势（或等效电压源）
  \[
    \mathcal{E}_a:\ E\to\RR,\qquad e\mapsto \mathcal{E}_a(e).
  \]
  \item（通信语义）链路噪声/时延/丢包的扰动场
  \[
    \mathcal{N}_a:\ E\to\RR_{\ge 0},\qquad e\mapsto \mathcal{N}_a(e).
  \]
\end{itemize}
本定义不把物理细节写死；仅固定“外部不可控扰动以动作形式进入系统演化”的开放系统口径。
\end{definition}

\begin{definition}[D2（电网/通信“调度”问题的 ORL-开放博弈编码）]
\label{def:tex27-ch16-d2}
令变分空间（状态空间）
\[
  \mathcal{V}_{\mathrm{sched}}\ni
  s=(\text{网络拓扑},\ z,\ x,\ \text{负载/业务},\ \text{安全余度},\dots),
\]
其中 $z:E\to\RR_{>0}$ 是“等效阻抗/等效代价”场（电网对应阻抗/热容量等；通信对应时延/带宽约束等），
$x$ 是潮流/路由/队列等内部变量。

令原子算子字母表（可执行控制动作集合）为有限集
\[
  \mathcal{O}_{\mathrm{sched}}=\{o_1,\dots,o_M\}.
\]
例如：开断/重构、负荷切除、无功支撑切换、路由重分配、调度窗口参数更新等。

令缺陷势能（障碍度量）
\[
  \Phi_{\mathrm{sched}}:\mathcal{V}_{\mathrm{sched}}\to\RR_{\ge 0},
\]
用来度量“约束违背 + 风险 + 代价”。例如可取一般形式
\[
  \Phi_{\mathrm{sched}}(s)
  =
  \sum_{e\in E}\bigl[\max(0,|I_e(s)|-I_e^{\max})\bigr]^2
  +\sum_{v\in V}\bigl[\max(0,V_v^{\min}-V_v(s))\bigr]^2
  +\cdots,
\]
通信情形把 $(I_e,V_v)$ 替换为时延/丢包/拥塞等谓词即可。

对任意算子词（长度至多 $L$ 的复合）$w\in\Omega_L(\mathcal{O}_{\mathrm{sched}})$，定义一跳执行
\[
  s'=w(s,a),
\]
并定义鲁棒能量（最坏情形代价）
\[
  E_{\mathrm{sched}}(w;s):=\sup_{a\in\mathcal{A}(s)}\Phi_{\mathrm{sched}}\bigl(w(s,a)\bigr).
\]
\end{definition}

\begin{definition}[D3（“可防御地磁爆的工业级室温超导”作为 ORL-开放博弈编码）]
\label{def:tex27-ch16-d3}
把“工业级室温超导（RTSC）防御”抽象为：通过“材料/环境/几何/控制”的可执行算子，
使系统在最坏地磁爆扰动下仍维持可接受的导通与安全谓词。

令状态空间
\[
  \mathcal{V}_{\mathrm{rtsc}}\ni
  \sigma=(\text{几何},\ \rho,\ T,\ \kappa,\ \text{散热通道},\ \text{屏蔽/旁路结构},\dots),
\]
其中 $\rho:E\to\RR_{\ge 0}$ 是“等效电阻率/等效损耗”场（可依赖温度 $T$ 与环境耦合参数 $\kappa$）。

令原子算子字母表为有限集
\[
  \mathcal{O}_{\mathrm{rtsc}}=\{o'_1,\dots,o'_{M'}\}.
\]
例如：材料态切换/相变控制、几何迭代、旁路/熔断策略切换、散热通道调度、屏蔽结构参数更新等。

令缺陷势能
\[
  \Phi_{\mathrm{rtsc}}:\mathcal{V}_{\mathrm{rtsc}}\to\RR_{\ge 0},
\]
度量“热失控风险 + 电阻损耗 + 结构约束违背”等。并同样定义鲁棒能量
\[
  E_{\mathrm{rtsc}}(w';\sigma)
  :=\sup_{a\in\mathcal{A}(\sigma)}\Phi_{\mathrm{rtsc}}\bigl(w'(\sigma,a)\bigr),
  \qquad
  w'\in\Omega_L(\mathcal{O}_{\mathrm{rtsc}}).
\]
\end{definition}

\begin{definition}[D4（“等价”的精确定义：ORL-开放博弈同构/同伦等价）]
\label{def:tex27-ch16-d4}
称两个问题域 $\mathcal{D}_1,\mathcal{D}_2$ 在 ORL 口径下\textbf{等价}，
若存在映射
\[
  \psi:\mathcal{V}_1\to\mathcal{V}_2,\qquad
  \Psi:\Omega_L(\mathcal{O}_1)\to\Omega_L(\mathcal{O}_2),
\]
以及对手动作映射 $\alpha:\mathcal{A}_1(s)\to\mathcal{A}_2(\psi(s))$，
使得对任意状态 $s\in\mathcal{V}_1$、任意算子词 $w\in\Omega_L(\mathcal{O}_1)$，成立能量与不变谓词保持：
\[
  E_1(w;s)=E_2(\Psi(w);\psi(s)),
  \qquad
  \mathrm{Inv}_1(s)\Rightarrow \mathrm{Inv}_2(\psi(s)).
\]
并且 $\Psi$ 与 $\psi$ 在“复合/执行”的结构上可回放：在同构意义下可逆，
或在同伦意义下可通过可审计修复把偏差吸收回等价类。
\end{definition}

\subsection{公理（降级为定理）：共享“阻抗场--注入场”开放博弈骨架}
\label{subsec:tex27-ch16-axiom}

\begin{theorem}[公理（降级为定理）A（电网/通信调度与 RTSC 防御共享同一开放博弈骨架）]
\label{thm:tex27-ch16-a}
在定义~\ref{def:tex27-ch16-d2}--\ref{def:tex27-ch16-d3} 的抽象下，存在共同骨架状态空间
\[
  \mathcal{V}_\ast=
  \bigl(\text{网络骨架 }(V,E),\ z:E\to\RR_{\ge 0},\ \text{资源/容量谓词},\dots\bigr),
\]
以及共同对手动作语义（注入场）
\[
  a:\ E\to\RR,
\]
使得两域都可通过“忘却/降级口径”投影到 $\mathcal{V}_\ast$，并且鲁棒能量可写成统一形式
\[
  E_\ast(w;u)=\sup_{a}\ \Phi_\ast\bigl(w(u,a)\bigr).
\]
\end{theorem}

\begin{Certificate}[证书 A（结构同构证书：显式投影与势能依赖限制）]
\label{cert:tex27-ch16-a}
给出显式投影
\[
  \pi_{\mathrm{sched}}:\mathcal{V}_{\mathrm{sched}}\to\mathcal{V}_\ast,\quad
  \pi_{\mathrm{rtsc}}:\mathcal{V}_{\mathrm{rtsc}}\to\mathcal{V}_\ast,
\]
其核心是抽取同一个“等效阻抗/等效代价”场：
\[
  \pi_{\mathrm{sched}}(s)=(V,E,\ z:=z(s),\dots),\qquad
  \pi_{\mathrm{rtsc}}(\sigma)=(V,E,\ z:=\rho(\sigma),\dots).
\]
并约束缺陷势能只依赖该投影：
\[
  \Phi_{\mathrm{sched}}=\Phi_\ast\circ \pi_{\mathrm{sched}},\qquad
  \Phi_{\mathrm{rtsc}}=\Phi_\ast\circ \pi_{\mathrm{rtsc}}.
\]
\end{Certificate}

\begin{proof}[证明（谓词因果链）]
\begin{enumerate}
  \item（共同状态变量）调度目标最终在“网络骨架 + 边代价/边阻抗 + 资源谓词”上判定；
  RTSC 防御目标亦可在“网络骨架 + 边电阻率/损耗 + 热/结构谓词”上判定。
  故共同骨架 $\mathcal{V}_\ast$ 合法。
  \item（共同对手动作）地磁爆对电网表现为线路注入电动势，对通信表现为链路噪声/容量骤降；
  二者统一为 $a:E\to\RR$ 的“注入场”（定义~\ref{def:tex27-ch16-d1}）。
  \item（共同鲁棒能量）沿用开放对抗定义（定义~\ref{def:tex27-ch16-d2}、定义~\ref{def:tex27-ch16-d3}）。
  若 $\Phi_{\mathrm{sched}}$ 与 $\Phi_{\mathrm{rtsc}}$ 均仅依赖投影（证书~\ref{cert:tex27-ch16-a}），
  则两域的能量都可写成同一 $\Phi_\ast$ 形式，从而得到统一骨架能量表达。
\end{enumerate}
证毕。
\end{proof}

\subsection{定理：两域在 ORL-开放博弈意义下的等价}
\label{subsec:tex27-ch16-theorem}

\begin{theorem}[定理 T1（主结论：调度抗地磁爆 $\Longleftrightarrow$ RTSC 防御地磁爆，在 ORL-开放博弈意义下）]
\label{thm:tex27-ch16-t1}
在定义~\ref{def:tex27-ch16-d4} 的“等价”口径下，
电网/通信调度抗地磁爆问题域 $\mathcal{D}_{\mathrm{sched}}$ 与 RTSC 防御地磁爆问题域 $\mathcal{D}_{\mathrm{rtsc}}$ 等价。
\end{theorem}

\begin{Certificate}[证书 T1（显式互译映射 + 数学归纳法证书：词级复合保持）]
\label{cert:tex27-ch16-t1}
构造 $(\psi,\Psi,\alpha)$ 使定义~\ref{def:tex27-ch16-d4} 成立：
\begin{enumerate}
  \item（状态映射：以“等效阻抗场”为主元）固定参考元组 $(T_0,\kappa_0,\dots)$。
  定义
  \[
    \psi(s):=(\text{同一 }(V,E),\ \rho:=z(s),\ T:=T_0,\ \kappa:=\kappa_0,\dots)\in\mathcal{V}_{\mathrm{rtsc}}.
  \]
  即把调度问题里“通过开关/路由实现的有效边代价 $z$”视为 RTSC 问题里“通过材料/几何实现的有效电阻率 $\rho$”。
  \item（算子词映射：把“拓扑/路由调度算子”译为“阻抗整形算子”）
  给定逐原子映射
  \[
    \kappa:\mathcal{O}_{\mathrm{sched}}\to \Omega_{L_o}(\mathcal{O}_{\mathrm{rtsc}})
  \]
  满足：对任意 $o\in\mathcal{O}_{\mathrm{sched}}$，存在 $w'_o:=\kappa(o)$ 使其在骨架投影上产生同一阻抗场更新，
  \[
    \pi_{\mathrm{rtsc}}\bigl(w'_o(\psi(s),a)\bigr)=\pi_{\mathrm{sched}}\bigl(o(s,a)\bigr),\qquad \forall(s,a).
  \]
  定义词级映射为复合保持：
  \[
    \Psi(o_k\circ\cdots\circ o_1):=\kappa(o_k)\circ\cdots\circ\kappa(o_1).
  \]
  \item（对手动作映射：同一注入场语义）
  \[
    \alpha(a):=a.
  \]
  \item（归纳证书：对任意词长 $\ell\le L$，投影执行结果匹配）
  对任意 $w\in\Omega_L(\mathcal{O}_{\mathrm{sched}})$，断言
  \[
    \pi_{\mathrm{rtsc}}\bigl(\Psi(w)(\psi(s),a)\bigr)=\pi_{\mathrm{sched}}\bigl(w(s,a)\bigr)
  \]
  的归纳证书为：
  \begin{enumerate}
    \item（基步）$\ell=0$：$w=\mathrm{id}$，由 $\Psi(\mathrm{id})=\mathrm{id}$ 与 $\pi_{\mathrm{rtsc}}(\psi(s))
      =\pi_{\mathrm{sched}}(s)$ 得到；
    \item（归纳步）若对长度 $\ell$ 的 $w$ 成立，则对 $w^+:=o\circ w$，
    由 $\Psi(w^+)=\kappa(o)\circ\Psi(w)$ 与逐原子匹配条件推出投影后执行结果仍匹配。
  \end{enumerate}
\end{enumerate}
\end{Certificate}

\begin{proof}[证明（谓词因果链）]
要证定义~\ref{def:tex27-ch16-d4} 中的能量保持：对任意 $s\in\mathcal{V}_{\mathrm{sched}}$、任意 $w\in\Omega_L(\mathcal{O}_{\mathrm{sched}})
  $，
有
\[
  E_{\mathrm{sched}}(w;s)=E_{\mathrm{rtsc}}(\Psi(w);\psi(s)).
\]
\begin{enumerate}
  \item（势能在共同投影下统一）由证书~\ref{cert:tex27-ch16-a}，
  \[
    \Phi_{\mathrm{sched}}=\Phi_\ast\circ\pi_{\mathrm{sched}},
    \qquad
    \Phi_{\mathrm{rtsc}}=\Phi_\ast\circ\pi_{\mathrm{rtsc}}.
  \]
  \item（投影下的执行结果匹配）由证书~\ref{cert:tex27-ch16-t1} 的归纳结论，对任意 $a$ 有
  \[
    \pi_{\mathrm{rtsc}}\bigl(\Psi(w)(\psi(s),a)\bigr)=\pi_{\mathrm{sched}}\bigl(w(s,a)\bigr).
  \]
  \item（势能值匹配）于是对任意 $a$，
  \[
    \Phi_{\mathrm{rtsc}}\bigl(\Psi(w)(\psi(s),a)\bigr)
    =\Phi_\ast\!\left(\pi_{\mathrm{rtsc}}(\Psi(w)(\psi(s),a))\right)
    =\Phi_\ast\!\left(\pi_{\mathrm{sched}}(w(s,a))\right)
    =\Phi_{\mathrm{sched}}\bigl(w(s,a)\bigr).
  \]
  \item（取上确界得到鲁棒能量匹配）
  \[
    E_{\mathrm{rtsc}}(\Psi(w);\psi(s))
    =\sup_a \Phi_{\mathrm{rtsc}}(\Psi(w)(\psi(s),a))
    =\sup_a \Phi_{\mathrm{sched}}(w(s,a))
    =E_{\mathrm{sched}}(w;s).
  \]
\end{enumerate}
因此满足能量保持条件。不变谓词 $\mathrm{Inv}$ 的保持同理：只要 $\mathrm{Inv}$ 依赖于共同骨架投影，
则由投影执行匹配可推出 $\mathrm{Inv}_{\mathrm{sched}}(s)\Rightarrow \mathrm{Inv}_{\mathrm{rtsc}}(\psi(s))$。
证毕。
\end{proof}

\subsection{引理与命题：把“等价”解释为工程语义}
\label{subsec:tex27-ch16-lemmas-props}

\begin{lemma}[L1（两域的本质控制变量都是边上的有效阻抗/等效代价场）]
\label{lem:tex27-ch16-l1}
在共同骨架口径下，两域的本质控制主元均可取为边场 $z:E\to\RR_{\ge 0}$（或其等价表示）。
\end{lemma}

\begin{Certificate}[证书 L1（投影抽取见证）]
\label{cert:tex27-ch16-l1}
证书即证书~\ref{cert:tex27-ch16-a} 中的两条投影：
$\pi_{\mathrm{sched}}$ 与 $\pi_{\mathrm{rtsc}}$ 都抽取同一个边场 $z$（分别由 $z(s)$ 与 $\rho(\sigma)$ 给出）。
\end{Certificate}

\begin{proof}[证明（谓词因果链）]
由证书~\ref{cert:tex27-ch16-l1}，两域在忘却/降级后都落到同一 $z$-主元骨架，
故“本质控制变量”为边上的有效阻抗/代价场。证毕。
\end{proof}

\begin{proposition}[P1（抗地磁爆目标下：差别主要在“实现 $z$ 更新的物理载体”）]
\label{prop:tex27-ch16-p1}
在抗地磁爆目标下，调度算子与超导防御算子的差别主要体现在实现同一 $z$ 更新的物理载体不同：
前者通过开断/重构/路由等工程操作改变等效代价，后者通过材料态/几何/旁路/散热/屏蔽等改变等效电阻率/损耗。
\end{proposition}

\begin{Certificate}[证书 P1（逐原子互译 $\kappa$）]
\label{cert:tex27-ch16-p1}
证书为证书~\ref{cert:tex27-ch16-t1} 中逐原子映射 $\kappa:\mathcal{O}_{\mathrm{sched}}\to \Omega_{L_o}(\mathcal{O}_{\mathrm{rtsc}})$：

它保证每个调度原子算子在骨架投影上引起的 $z$ 更新，都可由某个 RTSC 算子词在 $\rho$ 上实现同样的更新。
\end{Certificate}

\begin{proof}[证明（谓词因果链）]
按证书~\ref{cert:tex27-ch16-p1}，两类算子在共同骨架上实现同一 $z$ 更新；
因此“差别”只能落在实现该更新的物理载体与可执行细节上，而非抽象骨架层。证毕。
\end{proof}

\subsection{推论与备注：一句话结论与未证明边界}
\label{subsec:tex27-ch16-cor-remarks}

\begin{corollary}[C1（原句命题版：调度应用 $\Longleftrightarrow$ RTSC 防御应用）]
\label{cor:tex27-ch16-c1}
将本笔记的 ORL/开放系统对抗理论应用于电网或通信调度（以地磁爆为对手动作），
在 ORL-开放博弈意义下等价于构建“能够防御地磁爆”的工业级室温超导（同一鲁棒能量与同一缺陷势能下降口径）。
\end{corollary}

\begin{proof}[证明（证书化）]
由定理~\ref{thm:tex27-ch16-t1} 直接得到。证毕。
\end{proof}

\begin{remark}[边界与未证明（必须显式标注）]
\label{rem:tex27-ch16-r1}
\begin{enumerate}
  \item 上述“等价”是\textbf{形式系统内}的等价：
  即把两类任务都压缩进同一 $(\mathcal{V},\mathcal{O},\mathcal{A},\Phi,\mathrm{Inv})$ 接口，并在鲁棒能量 $E$ 的意义下同构/同伦等价。
  \item “现实世界必然实现工业级室温超导”属于\textbf{实验事实与工程约束}问题；
  在此形式证明中应标注为\textbf{未证明}（需要外部证据链/实验记录/工程可复现性证书）。
  \item 但这不影响本结论的用法：只要把目标写成可审计编码（状态、算子、对手动作、缺陷势能与不变谓词），
  本节结论就给出统一的“可证书化在线修复/设计/调度”框架；电网/通信与 RTSC 防御只是两种不同的实现载体。
\end{enumerate}
\end{remark}

\clearpage

%%%%%%%%%%%%%%%%%%%%%%%%%%%%%%%%%%%%%%%%%%%%%%%%%%%%%%%%%%%%
\section{形式逻辑：太阳磁暴作用域、全光存储与同源失效消除（CMF 证书链）}
\label{sec:tex27-ch16b-solar-cmf}
%%%%%%%%%%%%%%%%%%%%%%%%%%%%%%%%%%%%%%%%%%%%%%%%%%%%%%%%%%%%

\subsection{定义}
\label{subsec:tex27-ch16b-defs}

\begin{definition}[D0（载荷子丛/载体集合）]
\label{def:tex27-ch16b-d0}
令系统的“物理载体”集合分解为
\[
  \mathcal{C}
  :=
  \mathcal{C}_{e}\sqcup\mathcal{C}_{\gamma}\sqcup\mathcal{C}_{\mathrm{oth}},
\]
其中：
\begin{itemize}
  \item $\mathcal{C}_{e}$：带电载体（自由电子、导体电流、半导体结中的载流子等）；
  \item $\mathcal{C}_{\gamma}$：光子载体（在介质波导/光纤中传播的光场模态）；
  \item $\mathcal{C}_{\mathrm{oth}}$：其余载体（机械、热、化学等）。
\end{itemize}
\end{definition}

\begin{definition}[D1（扰动类与作用域）]
\label{def:tex27-ch16b-d1}
令“卡林顿级太阳磁暴/地磁扰动”抽象为扰动类
\[
  \Delta_{\odot}
  :=
  \left\{
    \delta:\ \text{低频/宏观电磁场扰动（含 GIC/EMP 的等效低频部分）}
  \right\}.
\]
对任意 $\delta\in\Delta_{\odot}$，定义其作用域（Domain）
\[
  \mathrm{Dom}(\delta)\subseteq\mathcal{C}
\]
为该扰动在工程时间窗内可产生一阶主导影响的载体集合。
\end{definition}

\begin{definition}[D2（存储实现与存储载体域）]
\label{def:tex27-ch16b-d2}
对任意控制记忆/存储实现 $M$，定义其存储载体域
\[
  \mathrm{Dom}(M)\subseteq\mathcal{C}
\]
为承载比特信息并决定可恢复性的载体集合。特别地：
\begin{itemize}
  \item 若 $M=M_{\mathrm{EO}}$ 为电子存储（RAM/SSD/Flash/寄存器等），则通常
  \[
    \mathrm{Dom}(M_{\mathrm{EO}})\subseteq\mathcal{C}_{e};
  \]
  \item 若 $M=M_{\mathrm{AO}}$ 为全光存储/光延迟线环路存储（灾难窗内无 O--E--O），则目标是
  \[
    \mathrm{Dom}(M_{\mathrm{AO}})\subseteq\mathcal{C}_{\gamma}.
  \]
\end{itemize}
\end{definition}

\begin{definition}[D3（同源失效指标 CMF）]
\label{def:tex27-ch16b-d3}
给定能量面系统 $G$（电网执行器）与控制记忆 $M$，定义
\[
  \mathrm{CMF}(G,M)
  :=
  \begin{cases}
    1, &
    \exists \delta\in\Delta_{\odot}\ \text{s.t.}\
    \mathrm{Dom}(\delta)\cap\mathrm{Dom}(G)\neq\varnothing \\
    &\qquad\wedge\ \mathrm{Dom}(\delta)\cap\mathrm{Dom}(M)\neq\varnothing,\\
    0, &
    \forall \delta\in\Delta_{\odot},\
    \mathrm{Dom}(\delta)\cap\mathrm{Dom}(G)=\varnothing \\
    &\qquad\vee\ \mathrm{Dom}(\delta)\cap\mathrm{Dom}(M)=\varnothing.
  \end{cases}
\]
直观上：若存在同一类太阳磁暴扰动可同时打穿“电网执行器”和“控制记忆”，则同源失效不可排除（$=1$）；否则可排除（$=0$）。
\end{definition}

\begin{definition}[D4（观测者免疫）]
\label{def:tex27-ch16b-d4}
称控制记忆 $M$ 对扰动类 $\Delta_{\odot}$ 具“观测者免疫”，若
\[
  \forall \delta\in\Delta_{\odot},\quad
  \mathrm{Dom}(\delta)\cap\mathrm{Dom}(M)=\varnothing.
\]
工程弱口径可放宽为：存在 $\varepsilon\ll 1$ 使得由 $\Delta_{\odot}$ 导致的记忆错误概率不超过 $\varepsilon$（见备注）。
\end{definition}

\subsection{公理（降级为定理）：低频外场主导作用域}
\label{subsec:tex27-ch16b-axiom}

\begin{theorem}[公理（降级为定理）T0（QED 最小耦合下的一阶主导作用域）]
\label{thm:tex27-ch16b-t0}
对任意 $\delta\in\Delta_{\odot}$（低频/宏观电磁外场扰动），在一阶主导效应口径下有
\[
  \mathrm{Dom}(\delta)\subseteq
  \mathcal{C}_{e}\cup\mathcal{C}_{\mathrm{oth}}^{(\mathrm{med})},
  \qquad
  \mathrm{Dom}(\delta)\cap\mathcal{C}_{\gamma}=\varnothing,
\]
其中 $\mathcal{C}_{\mathrm{oth}}^{(\mathrm{med})}$ 表示经导体/载流子介导的次级通道。
即外场不直接对光子载体产生一阶主导破坏；若有影响，需经物质介质/器件带电自由度间接产生，并可被隔离到接口层。
\end{theorem}

\begin{Certificate}[证书 C0（QED 最小耦合与载体域分解）]
\label{cert:tex27-ch16b-c0}
\begin{enumerate}
  \item C0-1：QED/经典极限下，外场耦合项可写为
  \[
    H_{\mathrm{int}}[A^{\mathrm{ext}}]\propto \int J^\mu A^{\mathrm{ext}}_{\mu}\,\mathrm{d}^3x,
  \]
  其中 $J^\mu$ 是带电物质的电流四矢量；
  \item C0-2：光子电荷 $q_{\gamma}=0$，无外场对自由光子的一阶洛伦兹力项；
  \item C0-3：介质中光场若受外磁场影响，需通过介质带电自由度间接发生（法拉第效应、折射率微扰、器件偏置漂移等），其可隔离位置在介质/器件接口层而非光子本体。
\end{enumerate}
\end{Certificate}

\begin{proof}[证明（基于数学符号推演逻辑的谓词因果链）]
\begin{enumerate}
  \item（耦合谓词）令 $A^{\mathrm{ext}}$ 表示太阳磁暴等效外场。外场进入动力学的主导项为
  $\int J^\mu A^{\mathrm{ext}}_{\mu}$，其必要条件是 $J^\mu\neq 0$，即存在带电载流子；
  \item（载体谓词）自由光子满足 $q_{\gamma}=0$。在外场对载体直接一阶耦合意义下，不存在 $q_{\gamma}A^{\mathrm{ext}}$ 型受力项；
  \item（因果结论）故 $\delta\in\Delta_{\odot}$ 的一阶主导作用域落在 $\mathcal{C}_{e}$（及其介导的次级通道），不落在 $\mathcal{C}_{\gamma}$。证毕。
\end{enumerate}
\end{proof}

\subsection{引理组（证书 + 证明）}
\label{subsec:tex27-ch16b-lemmas}

\begin{lemma}[引理 L1（电网执行器的载体域）]
\label{lem:tex27-ch16b-l1}
对电网执行器 $G$（铜/铝导体、变压器绕组、断路器等），有
\[
  \mathrm{Dom}(G)\cap\mathcal{C}_{e}\neq\varnothing,
\]
且通常
\[
  \mathrm{Dom}(G)\subseteq\mathcal{C}_{e}\cup\mathcal{C}_{\mathrm{oth}}
\]
（其中热/机械路径为次级）。
\end{lemma}

\begin{Certificate}[证书 L1（导体电流依赖见证）]
\label{cert:tex27-ch16b-l1}
电网能量传输本质依赖导体电流 $I(t)$ 与载流子迁移。
\end{Certificate}

\begin{proof}[证明（基于数学符号推演逻辑的谓词因果链）]
电网主状态变量（电压/电流/磁通）以导体电流为核心，而导体电流由电子/载流子构成。
因此 $\mathrm{Dom}(G)$ 必含 $\mathcal{C}_{e}$ 成分。证毕。
\end{proof}

\begin{lemma}[引理 L2（仅光纤通信不足以推出 \texorpdfstring{$\mathrm{Dom}(M)\subseteq\mathcal{C}_{\gamma}$}{Dom(M) subset Cgamma}）]
\label{lem:tex27-ch16b-l2}
若系统在灾难窗内存在 O--E--O 转换，或必须落盘到电子存储，则
\[
  \mathrm{Dom}(M)\cap\mathcal{C}_{e}\neq\varnothing.
\]
从而无法给出 $\mathrm{CMF}(G,M)=0$ 的严格证书。
\end{lemma}

\begin{Certificate}[证书 L2（O--E--O 电子驻留见证）]
\label{cert:tex27-ch16b-l2}
O--E--O 转换节点包含光电探测器、驱动、电源、寄存器等电子态驻留点。
\end{Certificate}

\begin{proof}[证明（基于数学符号推演逻辑的谓词因果链）]
只要比特在灾难窗内以电子态驻留，即有 $\mathrm{Dom}(M)\cap\mathcal{C}_{e}\neq\varnothing$。
结合定理~\ref{thm:tex27-ch16b-t0} 的作用域结论，可知同域交叠无法被排除，故 $\mathrm{CMF}(G,M)=0$ 不可证。证毕。
\end{proof}

\begin{lemma}[引理 L3（光延迟线存储的时间--长度约束）]
\label{lem:tex27-ch16b-l3}
若以光纤延迟线在灾难窗内用飞行时间承载比特，则需满足
\[
  L=v_g\tau,
  \qquad
  v_g\approx \frac{c}{n}\sim 2\times 10^8\ \mathrm{m/s},
\]
其中 $L$ 是有效环路长度，$\tau$ 是可无电子驻留的存储时间窗。
\end{lemma}

\begin{Certificate}[证书 L3（群速度与光程关系）]
\label{cert:tex27-ch16b-l3}
证书要件是群速度定义与光程关系：$\tau=L/v_g$。
\end{Certificate}

\begin{proof}[证明（基于数学符号推演逻辑的谓词因果链）]
比特以脉冲在环路中传播，完成一圈所需时间为 $\tau=L/v_g$，故等价改写得 $L=v_g\tau$。证毕。
\end{proof}

\subsection{命题（证书 + 证明）}
\label{subsec:tex27-ch16b-prop}

\begin{proposition}[命题 P1（全光承载存储给出 \texorpdfstring{$\mathrm{CMF}=0$}{CMF=0} 的结构判据）]
\label{prop:tex27-ch16b-p1}
设控制记忆 $M=M_{\mathrm{AO}}$ 满足灾难窗内
\[
  \mathrm{Dom}(M_{\mathrm{AO}})\subseteq\mathcal{C}_{\gamma},
\]
且接口处不存在必须在线依赖的电子驻留（即引理~\ref{lem:tex27-ch16b-l2} 的否定条件），则
\[
  \mathrm{CMF}(G,M_{\mathrm{AO}})=0.
\]
\end{proposition}

\begin{Certificate}[证书 C1（工程可核验清单）]
\label{cert:tex27-ch16b-c1}
\begin{enumerate}
  \item C1-1（载体证书）：灾难窗 $[t_0,t_1]$ 内关键状态/路由矩阵/调度算子参数的唯一在线副本以光场编码，并驻留于光纤环路或光学谐振结构；
  \item C1-2（接口证书）：灾难窗内不发生 O--E--O 落地；若必须有接口，接口处具独立防护与可替换性，不构成同源同时失效共同点；
  \item C1-3（隔离证书）：关键光纤段不与长导体共缆形成有效感应回路；若存在此风险，需给出额外屏蔽/切断证书。
\end{enumerate}
\end{Certificate}

\begin{proof}[证明（基于数学符号推演逻辑的谓词因果链）]
\begin{enumerate}
  \item 由引理~\ref{lem:tex27-ch16b-l1}，$\mathrm{Dom}(G)\cap\mathcal{C}_{e}\neq\varnothing$；
  \item 由定理~\ref{thm:tex27-ch16b-t0}，对任意 $\delta\in\Delta_{\odot}$，一阶主导口径下
  $\mathrm{Dom}(\delta)\cap\mathcal{C}_{\gamma}=\varnothing$；
  \item 由命题前提，
  \[
    \mathrm{Dom}(M_{\mathrm{AO}})\subseteq\mathcal{C}_{\gamma}
    \Rightarrow
    \mathrm{Dom}(\delta)\cap\mathrm{Dom}(M_{\mathrm{AO}})
    \subseteq
    \mathrm{Dom}(\delta)\cap\mathcal{C}_{\gamma}
    =\varnothing;
  \]
  \item 代入定义~\ref{def:tex27-ch16b-d3}，对所有 $\delta\in\Delta_{\odot}$ 无法同时满足“打穿执行器”与“打穿控制记忆”的同域条件，故 $\mathrm{CMF}(G,
    M_{\mathrm{AO}})=0$。证毕。
\end{enumerate}
\end{proof}

\subsection{定理（含数学归纳法证书）}
\label{subsec:tex27-ch16b-theorem}

\begin{theorem}[定理 T1（网络规模扩展下的光域载体归纳保持）]
\label{thm:tex27-ch16b-t1}
设控制系统由 $N$ 个模块组成，每个模块 $i$ 的灾难窗记忆为 $M_i$。若模块间连接仅允许光域传递且不允许灾难窗内电子驻留，并且
\[
  \forall i\in\{1,\dots,N\},\quad
  \mathrm{Dom}(M_i)\subseteq\mathcal{C}_{\gamma},
\]
则整体记忆
\[
  M^{(N)}:=\bigcup_{i=1}^{N}M_i
  \quad\Rightarrow\quad
  \mathrm{Dom}(M^{(N)})\subseteq\mathcal{C}_{\gamma}.
\]
\end{theorem}

\begin{Certificate}[数学归纳法证书 C2（模块组合的载体域闭合）]
\label{cert:tex27-ch16b-c2}
\begin{enumerate}
  \item C2-1：给定模块组合算子 $\oplus$（并联/串联/分层），其比特载体域满足
  \[
    \mathrm{Dom}(M_A\oplus M_B)\subseteq
    \mathrm{Dom}(M_A)\cup\mathrm{Dom}(M_B);
  \]
  \item C2-2：模块间接口若无电子驻留，则接口不引入 $\mathcal{C}_{e}$ 新载体。
\end{enumerate}
\end{Certificate}

\begin{proof}[证明（基于数学符号推演逻辑的谓词因果链）]
\begin{enumerate}
  \item（基础步 $N=1$）由前提
  \[
    \mathrm{Dom}(M^{(1)})=\mathrm{Dom}(M_1)\subseteq\mathcal{C}_{\gamma}.
  \]
  \item（归纳步）假设对 $N=k$ 成立，即 $\mathrm{Dom}(M^{(k)})\subseteq\mathcal{C}_{\gamma}$。考虑
  \[
    M^{(k+1)}=M^{(k)}\oplus M_{k+1}.
  \]
  由证书~\ref{cert:tex27-ch16b-c2} 的 C2-1，
  \[
    \mathrm{Dom}(M^{(k+1)})
    \subseteq
    \mathrm{Dom}(M^{(k)})\cup\mathrm{Dom}(M_{k+1})
    \subseteq
    \mathcal{C}_{\gamma}\cup\mathcal{C}_{\gamma}
    =\mathcal{C}_{\gamma}.
  \]
  同时由 C2-2，接口不引入新的 $\mathcal{C}_e$ 驻留。
\end{enumerate}
故由数学归纳法，$\forall N\in\NN_{>0}$，$\mathrm{Dom}(M^{(N)})\subseteq\mathcal{C}_{\gamma}$ 成立。证毕。
\end{proof}

\subsection{推论（证书 + 证明）}
\label{subsec:tex27-ch16b-cor}

\begin{corollary}[推论 C1（观测者免疫与同源失效消除）]
\label{cor:tex27-ch16b-c1}
在命题~\ref{prop:tex27-ch16b-p1} 的条件下，控制记忆 $M_{\mathrm{AO}}$ 对 $\Delta_{\odot}$ 具观测者免疫（定义~\ref{def:tex27-ch16b-d4}），并且
\[
  \mathrm{CMF}(G,M_{\mathrm{AO}})\equiv 0.
\]
\end{corollary}

\begin{Certificate}[证书 C1-Call（命题 P1 直接调用）]
\label{cert:tex27-ch16b-c1call}
证书链为：证书~\ref{cert:tex27-ch16b-c1} + 命题~\ref{prop:tex27-ch16b-p1}。
\end{Certificate}

\begin{proof}[证明（基于数学符号推演逻辑的谓词因果链）]
由命题~\ref{prop:tex27-ch16b-p1}，对任意 $\delta\in\Delta_{\odot}$ 有
\[
  \mathrm{Dom}(\delta)\cap\mathrm{Dom}(M_{\mathrm{AO}})=\varnothing.
\]
按定义~\ref{def:tex27-ch16b-d4} 即观测者免疫；再由定义~\ref{def:tex27-ch16b-d3} 得 $\mathrm{CMF}(G,M_{\mathrm{AO}})=0$。证毕。
\end{proof}

\begin{corollary}[推论 C2（必要性：严格 \texorpdfstring{$\mathrm{CMF}=0$}{CMF=0} 证书要求消除灾难窗电子驻留）]
\label{cor:tex27-ch16b-c2}
若目标是对 $\Delta_{\odot}$ 给出严格的 $\mathrm{CMF}(G,M)=0$ 证书，则系统架构必须满足引理~\ref{lem:tex27-ch16b-l2} 的否定条件：
灾难窗内不得存在关键记忆的电子驻留点（特别是 O--E--O 落地驻留）。
\end{corollary}

\begin{Certificate}[证书 C2-Call（反证触发接口）]
\label{cert:tex27-ch16b-c2call}
证书链为：引理~\ref{lem:tex27-ch16b-l2} + 定义~\ref{def:tex27-ch16b-d3}。
\end{Certificate}

\begin{proof}[证明（基于数学符号推演逻辑的谓词因果链）]
反设灾难窗内仍存在关键记忆电子驻留点。由引理~\ref{lem:tex27-ch16b-l2}，有
\[
  \mathrm{Dom}(M)\cap\mathcal{C}_{e}\neq\varnothing.
\]
结合定理~\ref{thm:tex27-ch16b-t0} 的作用域结论与定义~\ref{def:tex27-ch16b-d3}，同源同时失效条件不可排除，故严格 $\mathrm{CMF}(G,M)=0$ 不可证，与目标矛盾。证毕。
\end{proof}

\subsection{备注（工程边界与未证明项）}
\label{subsec:tex27-ch16b-remarks}

\begin{remark}[边界说明 1：交集为空的工程口径]
定理~\ref{thm:tex27-ch16b-t0} 采用的是一阶主导作用域口径。
现实介质中外磁场可经法拉第旋光等磁光效应影响偏振/相位，这通常属于可界定小扰动。
工程上可通过冗余编码、偏振无关方案、环路对称与屏蔽将误码率压到阈值。
若要求严格逻辑值 $0$，应把二阶通道写入 $\mathcal{C}_{\mathrm{oth}}$ 并在证书中逐条封堵或给出上界（例如改用 $\mathrm{CMF}_{\varepsilon}$ 口径）。
\end{remark}

\begin{remark}[边界说明 2：光域承载不等于无限时长存储]
由引理~\ref{lem:tex27-ch16b-l3}，纯延迟线要实现大时间窗 $\tau$ 需极大环路长度 $L$。
因此更现实的严格表述是：灾难窗内仅保证关键最小状态在光域承载；其余大容量历史数据采用异地、离线、光学介质或多重加固存储。
\end{remark}

\begin{remark}[边界说明 3：光纤物理结构的再引入风险]
很多光纤缆含金属加强件、铠装或并行供电线，可能重新引入感应回路。
若要签发证书~\ref{cert:tex27-ch16b-c1} 的 C1-3，需明确关键段的物理结构与接地策略。
\end{remark}

\begin{remark}[边界说明 4：未证明表达的降级改写]
短句“软件定义超导，光子承载记忆”属于未证明表达。若需降级为可证明命题，建议拆为两条：
\begin{itemize}
  \item（可证）在灾难窗内，控制记忆可选择光域载体，使 $\mathrm{CMF}(G,M)=0$ 可证；
  \item（待证/工程化）在给定硬件约束下，调度/重构算法可将系统维持在可恢复安全不变量集合内。
\end{itemize}
\end{remark}

\clearpage

%%%%%%%%%%%%%%%%%%%%%%%%%%%%%%%%%%%%%%%%%%%%%%%%%%%%%%%%%%%%
\section{形式逻辑：多代理互审系统与证书化可核验系统}
\label{sec:tex27-ch17-formal-logic-ma-ca}
%%%%%%%%%%%%%%%%%%%%%%%%%%%%%%%%%%%%%%%%%%%%%%%%%%%%%%%%%%%%

\subsection{定义}
\label{subsec:tex27-ch17-defs}

\begin{definition}[D1（自动化设计/自动化科研任务）]
\label{def:tex27-ch17-d1}
给定需求/问题描述 $x$（含目标、约束、环境假设），产出对象 $y$（设计方案/实验方案/证明/代码/工艺路线等）。
设
\[
  \begin{aligned}
    \mathsf{Con}(x,y)&\in\{0,1\}\quad(\text{安全、物理可行、接口一致、合规、资源约束等}),\\
    J(x,y)&\in\RR\quad(\text{性能、成本、鲁棒性、可解释性等}).
  \end{aligned}
\]
以及真实性/正确性谓词
\[
  \mathsf{True}(x,y)\in\{0,1\},
\]
其含义为“$y$ 在外部语义/事实/形式系统核验口径下确实正确”。
任务目标是构造 $y$ 使得
\[
  \mathsf{Con}(x,y)=1
  \ \wedge\
  J(x,y)\ \text{尽可能优}
  \ \wedge\
  \mathsf{True}(x,y)=1.
\]
\end{definition}

\begin{definition}[D2（LLM 多代理互审系统 $\mathsf{MA}$）]
\label{def:tex27-ch17-d2}
令 $M$ 表示一类生成/评审机制（例如同分布或近同分布的 LLM）。多代理互审系统 $\mathsf{MA}$ 的典型形态可抽象为：
\begin{enumerate}
  \item 生成候选：$y_1,\dots,y_k \sim M(\cdot\mid x)$；
  \item 互评打分：$s_i=\mathsf{Score}(x,y_i;M)$；
  \item 选择/迭代：输出
  \[
    \hat y=\mathsf{Select}\bigl(\{(y_i,s_i)\}_{i=1}^{k}\bigr),
  \]
  或进入下一轮对话/搜索状态更新 $x\leftarrow \mathsf{Update}(x,\hat y,\mathrm{critique})$。
\end{enumerate}
从形式上看，$\mathsf{MA}$ 是由 $M$ 驱动的随机迭代过程（可视为马尔可夫链/随机搜索），
其输出正确性主要由 $M$ 的统计性质与互评/选择规则的偏置共同决定。
\end{definition}

\begin{definition}[D3（证书驱动的可核验自动化系统 $\mathsf{CA}$）]
\label{def:tex27-ch17-d3}
系统 $\mathsf{CA}$ 在每一步不仅产出断言/方案 $y$，还产出\textbf{证书} $c$，并配有确定性核验器（proof checker / verifier）$V$，满足
\[
  V(x,y,c)=1\ \Longrightarrow\ \mathsf{Spec}(x,y)\ \text{成立},
\]
其中 $\mathsf{Spec}(x,y)$ 是形式化规范（可包含 $\mathsf{Con}$、正确性子目标、接口一致性等）。

进一步，若引入缺陷势能/违约势能 $\Phi\ge 0$ 与不变量谓词 $\mathsf{Inv}$，并考虑状态 $s\in S$ 与算子 $o:S\to S$，
则可要求每一步带证书地保证单调修复：
\[
  V_\Phi(s,o,c_\Phi)=1
  \Longrightarrow
  \bigl(\Phi(o(s))<\Phi(s)\bigr)\ \wedge\ \bigl(\mathsf{Inv}(s)\Rightarrow \mathsf{Inv}(o(s))\bigr).
\]
这类系统的关键不在“多说几轮”，而在“每一步都能被核验且可组合”。
\end{definition}

\subsection{公理（降级为定理）：核验器健全性}
\label{subsec:tex27-ch17-axiom-as-theorem}

\begin{theorem}[公理（降级为定理）T0（核验器健全性：由证书检查的逐步归纳得到）]
\label{thm:tex27-ch17-t0}
令 $\Pi$ 为一套形式推演系统（规则集合），并假设其规则对目标语义是保真（sound）的：规则可用则“真前提推出真结论”。
核验器 $V$ 对推演轨迹证书
\[
  c=((r_1,w_1),\dots,(r_n,w_n))
\]
逐步检查：第 $t$ 步在见证 $w_t$ 下应用规则 $r_t$，把已证事实集 $\Gamma_{t-1}$ 扩展为
\[
  \Gamma_t := \Gamma_{t-1}\cup\{\psi_t\}.
\]
若 $V$ 对每一步都检查通过，则最终结论 $\varphi\in\Gamma_n$ 在该系统语义下成立。
\end{theorem}

\begin{Certificate}[证书 T0（推演轨迹证书）]
\label{cert:tex27-ch17-t0}
证书取为规则序列与局部见证对 $c=((r_t,w_t))_{t=1}^{n}$。
其中 $r_t$ 标识所用推演规则，$w_t$ 记录该步的可审计局部见证（如替换项、匹配位置、子式索引、依赖断言编号等），
使核验器能在纯语法层面确认该步从 $\Gamma_{t-1}$ 合法得到 $\psi_t$。
\end{Certificate}

\begin{proof}[证明（数学归纳法；谓词因果链）]
对 $t=0,1,\dots,n$ 证明谓词 $P(t)$：“$\Gamma_t$ 中每个断言在语义下为真”。
\begin{itemize}
  \item 基础步 $t=0$：$\Gamma_0$ 为公理/已知前提集合，按外部给定或系统约定在语义下成立，因此 $P(0)$ 成立。
  \item 归纳步：假设 $P(t-1)$ 成立。第 $t$ 步核验通过意味着规则 $r_t$ 的适用条件满足，且 $r_t$ 在语义中保真，
  故由 $\Gamma_{t-1}$ 的真可推出新导出的 $\psi_t$ 亦真，从而 $\Gamma_t=\Gamma_{t-1}\cup\{\psi_t\}$ 全真，即 $P(t)$ 成立。
\end{itemize}
由归纳法得 $P(n)$ 成立；任取 $\varphi\in\Gamma_n$，即知 $\varphi$ 在语义下成立。证毕。
\end{proof}

\subsection{引理}
\label{subsec:tex27-ch17-lemmas}

\begin{lemma}[L1（互审一致性 $\not\Rightarrow$ 真实性）]
\label{lem:tex27-ch17-l1}
存在任务输入 $x$、输出 $y^{\times}$ 与常数 $p\in(0,1]$，使得
\[
  \mathsf{True}(x,y^{\times})=0,
  \qquad
  \Pr_{M}[y^{\times}\mid x]=p,
  \qquad
  \mathsf{Score}(x,y^{\times};M)\ge \mathsf{Score}(x,y;M)\ (\forall y).
\]
对定义~\ref{def:tex27-ch17-d2} 的互审形态，若 $\mathsf{Select}$ 选择最高分候选，则对任意候选数 $k$ 有
\[
  \Pr[\hat y=y^{\times}]\ \ge\ 1-(1-p)^k,
\]
从而可出现“以高概率达成一致认可”但 $\mathsf{True}(x,\hat y)=0$ 的情形。
\end{lemma}

\begin{Certificate}[证书 L1（反例任务与高分错误候选）]
\label{cert:tex27-ch17-l1}
证书由四个见证组成：
\begin{enumerate}
  \item 一个任务输入 $x$；
  \item 一个外部语义下错误的候选 $y^{\times}$，满足 $\mathsf{True}(x,y^{\times})=0$；
  \item 生成概率下界 $p=\Pr_{M}[y^{\times}\mid x]>0$；
  \item 互审偏置见证：在评分口径下 $y^{\times}$ 为最高分（或在最高分集合内且按既定择优规则被选中）。
\end{enumerate}
\end{Certificate}

\begin{proof}[证明（构造；谓词因果链）]
从 $M(\cdot\mid x)$ 独立采样 $k$ 个候选时，事件“至少一个候选等于 $y^{\times}$”的概率为
\[
  1-(1-p)^k.
\]
在该事件上，由证书~\ref{cert:tex27-ch17-l1} 的“最高分”见证，选择器必输出 $\hat y=y^{\times}$。
因此
\[
  \Pr[\hat y=y^{\times}]\ge 1-(1-p)^k.
\]
又由 $\mathsf{True}(x,y^{\times})=0$，可得在上述事件上 $\mathsf{True}(x,\hat y)=0$。
当 $p$ 不太小且 $k$ 增大时，该概率可趋近 1，体现“互审一致性不蕴含真实性”。证毕。
\end{proof}

\begin{lemma}[L2（同源模型互审的共同故障模式下界）]
\label{lem:tex27-ch17-l2}
设 $\mathsf{MA}$ 的关键判断均由同一类模型机制 $M$ 完成（生成与评分同源，错误具相关性）。
若存在错误集合 $E$ 与共同故障事件 $B$ 满足
\[
  \Pr[B]\ge \epsilon>0,
  \qquad
  B\Rightarrow(\hat y\in E),
\]
则无论代理数 $k$ 多大，都有
\[
  \Pr[\hat y\in E]\ge \epsilon.
\]
\end{lemma}

\begin{Certificate}[证书 L2（共同故障事件 $B$）]
\label{cert:tex27-ch17-l2}
证书给出：
\begin{enumerate}
  \item 一个错误集合 $E$（由外部语义/事实/规范判定为“错误”）；
  \item 一个共同故障事件 $B$（例如共享训练偏置/共享提示模板/共享推理捷径触发的同类错误模式）；
  \item 一个下界常数 $\epsilon>0$ 使 $\Pr[B]\ge\epsilon$；
  \item 蕴含关系 $B\Rightarrow(\hat y\in E)$（在该共同模式下，互审流程会稳定输出落入 $E$ 的候选）。
\end{enumerate}
\end{Certificate}

\begin{proof}[证明（下界；谓词因果链）]
由证书~\ref{cert:tex27-ch17-l2} 的蕴含关系 $B\Rightarrow(\hat y\in E)$，有集合包含
\[
  B\subseteq \{\hat y\in E\}.
\]
故
\[
  \Pr[\hat y\in E]\ge \Pr[B]\ge \epsilon,
\]
得到与 $k$ 无关的误差下界。证毕。
\end{proof}

\subsection{定理}
\label{subsec:tex27-ch17-theorems}

\begin{theorem}[T1（本质区别：$\mathsf{MA}$ 是统计一致性搜索；$\mathsf{CA}$ 是证书约束下的可核验构造）]
\label{thm:tex27-ch17-t1}
在定义~\ref{def:tex27-ch17-d1} 的任务口径下：
\begin{enumerate}
  \item（$\mathsf{MA}$ 的极限）$\mathsf{MA}$ 的正确性保证至多是“模型内一致性/自洽性”意义下的概率性增益；
  一般情形下，仅由“互审一致”无法推出 $\mathsf{True}(x,\hat y)=1$，且存在任务族使得错误概率存在与代理数无关的下界。
  \item（$\mathsf{CA}$ 的锚点）$\mathsf{CA}$ 的正确性由可核验性锚定：若存在证书 $c$ 使 $V(x,y,c)=1$，
  则规范 $\mathsf{Spec}(x,y)$ 必然成立，并且证书可被第三方复核、可审计、可复现且可组合。
\end{enumerate}
\end{theorem}

\begin{Certificate}[证书 T1（区分接口：一致性 vs 可核验性）]
\label{cert:tex27-ch17-t1}
证书由两组可复用接口组成：
\begin{enumerate}
  \item 反例接口：引理~\ref{lem:tex27-ch17-l1}（互审一致但不真）与引理~\ref{lem:tex27-ch17-l2}（同源共同故障下界）；
  \item 正向接口：定理~\ref{thm:tex27-ch17-t0}（核验器健全性）与定义~\ref{def:tex27-ch17-d3}（$\mathsf{CA}$ 的 $(y,c,V)$ 输出口径）。
\end{enumerate}
\end{Certificate}

\begin{proof}[证明（谓词因果链）]
\begin{itemize}
  \item（1）由引理~\ref{lem:tex27-ch17-l1}，存在“高概率一致认可但 $\mathsf{True}=0$”的反例结构；
  由引理~\ref{lem:tex27-ch17-l2}，在同源共同故障模式下错误概率存在与 $k$ 无关的下界。
  因此 $\mathsf{MA}$ 在一般意义上不能给出必然正确的演绎保证。
  \item（2）由定义~\ref{def:tex27-ch17-d3}，$\mathsf{CA}$ 输出证书并交由确定性核验器检查；
  再由定理~\ref{thm:tex27-ch17-t0}，核验通过即可推出规范在目标语义下成立，
  从而正确性被“$V=1$”这一可复核事件锚定，而不依赖“别的代理同不同意”。证毕。
\end{itemize}
\end{proof}

\begin{theorem}[T2（带缺陷势能 $\Phi$ 的证书化修复/设计：可保证单调进展）]
\label{thm:tex27-ch17-t2}
设系统状态 $s_t\in S$，可用算子集合 $\mathcal{O}$，不变量谓词 $\mathsf{Inv}$，缺陷势能 $\Phi:S\to\RR_{\ge 0}$。
若存在常数 $\delta>0$ 与策略在每步选择 $o_t\in\mathcal{O}$ 并给出证书 $c_t$，使核验通过且满足
\[
  V_\Phi(s_t,o_t,c_t)=1\Rightarrow
  \Big(\mathsf{Inv}(s_t)\Rightarrow \mathsf{Inv}(s_{t+1})\Big)\ \wedge\ \Phi(s_{t+1})\le \Phi(s_t)-\delta,
\]
其中 $s_{t+1}=o_t(s_t)$，则对任意 $n\in\NN$ 有
\[
  \Phi(s_n)\le \Phi(s_0)-n\delta.
\]
从而在有限步内可把缺陷压到阈值以下（或压到 $0$）。
\end{theorem}

\begin{Certificate}[证书 T2（一步下降证书序列）]
\label{cert:tex27-ch17-t2}
证书为序列 $(c_t)_{t=0}^{n-1}$，其中每个 $c_t$ 至少包含：
\begin{enumerate}
  \item 算子合法性见证：$o_t\in\mathcal{O}$ 且 $s_{t+1}=o_t(s_t)$；
  \item 不变量保持见证：$\mathsf{Inv}(s_t)\Rightarrow \mathsf{Inv}(s_{t+1})$；
  \item 势能下降见证：$\Phi(s_{t+1})\le \Phi(s_t)-\delta$。
\end{enumerate}
\end{Certificate}

\begin{proof}[证明（数学归纳法；谓词因果链）]
对 $n$ 做归纳。
\begin{itemize}
  \item 基础步 $n=1$：由前提对 $t=0$ 的结论，直接得到 $\Phi(s_1)\le \Phi(s_0)-\delta$。
  \item 归纳步：假设 $\Phi(s_n)\le \Phi(s_0)-n\delta$ 成立。
  再用一步证书化下降（对 $t=n$）得
  \[
    \Phi(s_{n+1})\le \Phi(s_n)-\delta\le \Phi(s_0)-(n+1)\delta.
  \]
\end{itemize}
归纳完成，结论成立。证毕。
\end{proof}

\subsection{命题}
\label{subsec:tex27-ch17-propositions}

\begin{proposition}[P1（可审计性与可复现性）]
\label{prop:tex27-ch17-p1}
$\mathsf{CA}$ 的输出是 $(y,c)$，其中证书 $c$ 使任何第三方仅靠核验器即可复核；
而若 $\mathsf{MA}$ 的输出缺少可被独立核验器检查的决定性对象，则其“互审结论”只能被再次对话式复述，难以形成可重放的审计接口。
\end{proposition}

\begin{Certificate}[证书 P1（核验可重放接口）]
\label{cert:tex27-ch17-p1}
证书即定义~\ref{def:tex27-ch17-d3} 中的三元组接口 $(V,y,c)$：第三方复核只需重放计算 $V(x,y,c)$ 的结果。
\end{Certificate}

\begin{proof}[证明（谓词因果链）]
由定义~\ref{def:tex27-ch17-d3}，$\mathsf{CA}$ 输出证书并配套确定性核验器；
若第三方运行 $V(x,y,c)=1$，再由定理~\ref{thm:tex27-ch17-t0} 可推出 $\mathsf{Spec}(x,y)$ 成立，因此可审计且可复现。
对 $\mathsf{MA}$，若不存在可被独立核验器检查的 $c$，则无法把“正确性”落为确定性复核事件，只能依赖再次采样/再次评审。证毕。
\end{proof}

\begin{proposition}[P2（可组合性：证书层面的闭包）]
\label{prop:tex27-ch17-p2}
若 $(y_1,c_1)$ 与 $(y_2,c_2)$ 均可核验，且存在组合规则 $\circ$ 与组合证书生成器 $\mathsf{Comp}$，使得
\[
  V(x,y_1,c_1)=1\ \wedge\ V(x,y_2,c_2)=1
  \Longrightarrow
  V\bigl(x, y_1\circ y_2,\ \mathsf{Comp}(c_1,c_2)\bigr)=1,
\]
则证书化系统在“证书层面”对组合运算闭包，从而可形成稳定可复用的算子/模块库。
\end{proposition}

\begin{Certificate}[证书 P2（组合规则作为一条推演规则）]
\label{cert:tex27-ch17-p2}
证书由两部分给出：
\begin{enumerate}
  \item 一条组合推演规则：从“$y_1$ 合法且 $y_2$ 合法”推出“$y_1\circ y_2$ 合法”；
  \item 证书合成器 $\mathsf{Comp}$：把 $(c_1,c_2)$ 变成该推演规则的局部见证。
\end{enumerate}
\end{Certificate}

\begin{proof}[证明（谓词因果链）]
把证书~\ref{cert:tex27-ch17-p2} 的“组合规则”视为推演系统 $\Pi$ 的一条规则，
并把 $\mathsf{Comp}(c_1,c_2)$ 视为该规则的局部见证。
则由定理~\ref{thm:tex27-ch17-t0} 的核验器健全性，若两前提可核验，则组合结论亦可核验。证毕。
\end{proof}

\begin{proposition}[P3（开放系统鲁棒性：最坏情形约束可证书化）]
\label{prop:tex27-ch17-p3}
在开放系统/对抗扰动下，$\mathsf{CA}$ 可以把“对所有扰动/对手动作 $a$”的约束写入规范并以证书锚定：
若规范包含全称量词断言
\[
  \forall a\in A:\ \mathsf{Con}(x,a,y)=1,
\]
且存在证书 $c$ 使 $V(x,y,c)=1$，则该最坏情形约束在目标语义下成立。
\end{proposition}

\begin{Certificate}[证书 P3（全称量词证书）]
\label{cert:tex27-ch17-p3}
证书为对全称断言 $\forall a\in A:\mathsf{Con}(x,a,y)=1$ 的形式推演轨迹（可为归纳证书、枚举证书或对抗模型下界证书等），
使核验器能逐步检查该断言的推演合法性。
\end{Certificate}

\begin{proof}[证明（谓词因果链）]
由证书~\ref{cert:tex27-ch17-p3}，核验器对全称断言的推演轨迹逐步检查通过；
再由定理~\ref{thm:tex27-ch17-t0}，核验通过即推出该断言在语义下成立，即得到最坏情形约束。证毕。
\end{proof}

\begin{proposition}[P4（算力与可靠性的解耦）]
\label{prop:tex27-ch17-p4}
若把“生成候选”的成本记为 $C_{\mathrm{gen}}$，把“核验证书”的成本记为 $C_{\mathrm{ver}}$，
并满足 $C_{\mathrm{ver}}\ll C_{\mathrm{gen}}$，则 $\mathsf{CA}$ 允许把主要成本放在候选生成，
把可靠性锚定在更便宜的确定性核验上；相对地，若 $\mathsf{MA}$ 试图仅靠增加互审轮数提升可靠性，
则其可靠性成本往往与昂贵的生成/评审采样一同线性增长。
\end{proposition}

\begin{Certificate}[证书 P4（成本分解见证）]
\label{cert:tex27-ch17-p4}
证书给出成本序关系 $C_{\mathrm{ver}}\ll C_{\mathrm{gen}}$ 与两类系统的成本分解口径：
$\mathsf{CA}$ 以“生成 + 多次核验”为主，而 $\mathsf{MA}$ 以“多次生成 + 多次互评”为主。
\end{Certificate}

\begin{proof}[证明（谓词因果链）]
由证书~\ref{cert:tex27-ch17-p4} 的成本序关系，
在达到同等可靠性目标时，把可靠性增益放在核验侧的边际成本更低；
而互审轮数增加通常意味着更多次同类模型推断（生成/评审采样），其边际成本不具备同等优势。证毕。
\end{proof}

\subsection{推论}
\label{subsec:tex27-ch17-corollaries}

\begin{corollary}[C1（“本质区别”的一句话形式化）]
\label{cor:tex27-ch17-c1}
在不额外引入外部核验器的前提下，
\[
  \begin{aligned}
    \mathsf{MA}\ \text{的主要增益是}\ &\Delta\Pr(\text{模型内自洽});\\
    \mathsf{CA}\ \text{的主要增益是把正确性锚定为}\ &(V=1)\ \text{的可核验命题}.
  \end{aligned}
\]
\end{corollary}

\begin{Certificate}[证书 C1（由定理 T1 抽取）]
\label{cert:tex27-ch17-c1}
证书为定理~\ref{thm:tex27-ch17-t1} 的两条结论（$\mathsf{MA}$ 的概率性增益与 $\mathsf{CA}$ 的可核验锚点）。
\end{Certificate}

\begin{proof}[证明（谓词因果链）]
由证书~\ref{cert:tex27-ch17-c1}，直接将定理~\ref{thm:tex27-ch17-t1} 的表述压缩为一句话即可。证毕。
\end{proof}

\begin{corollary}[C2（“先进性”的充分方向：开放系统下的可保证进展）]
\label{cor:tex27-ch17-c2}
若把自动化科研/设计的“先进性”定义为：在开放扰动下仍能保证约束 $\mathsf{Con}$，并且可复核地推动缺陷势能单调下降（$\Phi\downarrow$），
则仅靠 $\mathsf{MA}$（无外部核验器/证书体系）不足以满足；而 $\mathsf{CA}$ 在满足定理~\ref{thm:tex27-ch17-t2} 的前提时可满足。
\end{corollary}

\begin{Certificate}[证书 C2（由 T1 与 T2 组合）]
\label{cert:tex27-ch17-c2}
证书由两部分组成：
\begin{enumerate}
  \item 定理~\ref{thm:tex27-ch17-t1}：仅靠互审一致无法给出必然正确的演绎保证；
  \item 定理~\ref{thm:tex27-ch17-t2}：若每步可核验地保证 $\Phi$ 至少下降 $\delta$ 且保持不变量，则得到全程可保证进展。
\end{enumerate}
\end{Certificate}

\begin{proof}[证明（谓词因果链）]
由证书~\ref{cert:tex27-ch17-c2}：
第一部分给出“仅靠互审一致不足以把正确性钉死”的一般性障碍；
第二部分给出在证书化核验机制下对 $\Phi$ 的单调下降保证。
两者合并即可得到推论结论。证毕。
\end{proof}

\subsection{备注}
\label{subsec:tex27-ch17-remarks}

\begin{remark}[互审的定位：启发式搜索器而非“真值钉子”]
多代理互审并非“无价值”。它更像一种\textbf{启发式搜索器}：用多样性与批判回路提高候选质量。
但它天然缺少“把正确性钉死”的决定性对象；若不引入外部可核验机制，一致性只能提供经验性置信，而非演绎性保证。
\end{remark}

\begin{remark}[证书化系统也可以使用 LLM：让 LLM 做提案，让核验器做裁判]
证书化系统同样可以使用 LLM：把 LLM 放在“提案/算子枚举/候选生成”的位置，而把“是否成立”交给确定性核验器 $V$。
此时 LLM 的角色从“裁判”降为“搜索算子”，系统锚点变为可核验的形式对象。
\end{remark}

\begin{remark}[把“多代理互审 = 终局形态”当作结论应标注为未证明]
若有人把“多代理互审 = 自动化科研/设计的终局形态”当作一般性结论，
在此口径下应标注为\textbf{未证明}：引理~\ref{lem:tex27-ch17-l1} 与引理~\ref{lem:tex27-ch17-l2} 给出了反例结构与共同故障下界；
除非再引入外部可核验机制或独立真值源，否则该结论无法在形式系统内闭合为定理。
\end{remark}

\clearpage

%%%%%%%%%%%%%%%%%%%%%%%%%%%%%%%%%%%%%%%%%%%%%%%%%%%%%%%%%%%%
\section{形式逻辑：多代理互审系统 $\mathsf{MA}$ 的非终局性（外部真值通道与独立核验器）}
\label{sec:tex27-ch18-ma-not-final}
%%%%%%%%%%%%%%%%%%%%%%%%%%%%%%%%%%%%%%%%%%%%%%%%%%%%%%%%%%%%

\subsection{定义}
\label{subsec:tex27-ch18-defs}

\begin{definition}[D1（“基于 LLM 的多代理互审”系统 $\mathsf{MA}$）]
\label{def:tex27-ch18-d1}
多代理互审系统 $\mathsf{MA}$ 由若干同类或近同类的 LLM 代理构成。给定输入 $x$（文本/符号/数据描述），系统仅通过
“生成--互评--再生成--再互评”的内部对话与投票/打分/共识机制输出 $\hat y$。

\vspace{0.5em}
\noindent\textbf{关键信息源约束：}
$\mathsf{MA}$ 的信息源不超过输入 $x$ 与内部生成文本（含内部随机性），不包含独立的外部真值通道，
例如实验传感器/设备读数、可判定的形式证明核验器、可执行程序的确定性回执等。
形式上，可把其输出写为
\[
  \hat y = f(x;R),
\]
其中 $R$ 为系统内部随机性；若某世界真值变量 $\omega$ 与 $x$ 独立，则 $(x,R)$ 与 $\omega$ 亦独立。
\end{definition}

\begin{definition}[D2（自动化科研/自动化设计的“终局形态”）]
\label{def:tex27-ch18-d2}
以任务规范 $\mathsf{Spec}$ 表示“完成任务”的判定谓词：
\[
  \mathsf{Spec}(x,y)\in\{0,1\}.
  \quad(\mathsf{Spec}(x,y)=1\ \text{表示输出}\ y\ \text{满足规范})
\]
把“终局形态”定义为一个最大性断言：对任意任务规范 $\mathsf{Spec}$ 与输入 $x$，
只要该任务在物理/计算意义上可完成（存在某系统能使 $\mathsf{Spec}(x,\hat y)=1$），
就应当存在某个 $\mathsf{MA}$ 配置使其输出满足 $\mathsf{Spec}$；
并且不存在另一个系统 $\mathsf{S}$ 在任务覆盖或可靠性上严格优于 $\mathsf{MA}$（见定义~\ref{def:tex27-ch18-d3} 的严格支配）。
\end{definition}

\begin{definition}[D3（严格支配 $\prec$）]
\label{def:tex27-ch18-d3}
对系统 $\mathsf{A},\mathsf{B}$，定义其在任务对 $(x,\mathsf{Spec})$ 上的成功概率为
\[
  p_{\mathsf{A}}(x,\mathsf{Spec}) := \Pr\bigl[\mathsf{Spec}(x,\hat y_{\mathsf{A}})=1\bigr].
\]
若同时满足：
\begin{enumerate}
  \item（不劣）$\forall(x,\mathsf{Spec})$，有 $p_{\mathsf{B}}(x,\mathsf{Spec})\ge p_{\mathsf{A}}(x,\mathsf{Spec})$；
  \item（严格更优）$\exists(x,\mathsf{Spec})$，有 $p_{\mathsf{B}}(x,\mathsf{Spec})> p_{\mathsf{A}}(x,\mathsf{Spec})$，
\end{enumerate}
则记 $\mathsf{A}\prec \mathsf{B}$，称 $\mathsf{B}$ 严格支配 $\mathsf{A}$。
\end{definition}

\subsection{公理（降级为定理：证书+证明）}
\label{subsec:tex27-ch18-axiom-as-theorem}

\begin{theorem}[公理（降级为定理）T0（信息不可凭空生成：输入未包含的独立真值不可由内部互审必然恢复）]
\label{thm:tex27-ch18-t0}
设 $\omega\in\{0,1\}$ 为世界真值位，且 $\omega\sim \mathrm{Bernoulli}(1/2)$ 与输入 $x$ 独立（等价地，可写作互信息 $I(\omega;x)=0$）。
则对任何仅依赖 $x$ 与内部随机性的算法/系统（包括任意多轮内部互审）输出的猜测 $\hat\omega=f(x;R)$，
其猜中 $\omega$ 的最大成功率不超过 $1/2$：
\[
  \sup_{f}\Pr[\hat\omega=\omega]\le \frac12.
\]
特别地，该上界不依赖代理数与互审轮数，只依赖“是否存在独立真值通道”。
\end{theorem}

\begin{Certificate}[证书 T0（构造 $x$ 在两种 $\omega$ 下同分布）]
\label{cert:tex27-ch18-t0}
取联合分布使得：$\omega\sim \mathrm{Bernoulli}(1/2)$，且对任意可测集合 $A$，
\[
  \Pr[x\in A\mid \omega=0]=\Pr[x\in A\mid \omega=1],
\]
即 $x\perp\!\!\!\perp \omega$。
此时任何 $\hat\omega=f(x;R)$ 的条件分布在 $\omega=0$ 与 $\omega=1$ 下相同，从而无法利用 $x$ 区分 $\omega$。
\end{Certificate}

\begin{proof}[证明（谓词因果链）]
由证书~\ref{cert:tex27-ch18-t0}，$x$ 与 $\omega$ 独立，且内部随机性 $R$ 亦与 $\omega$ 独立，因此 $(x,R)$ 与 $\omega$ 独立。
令 $\hat\omega=f(x;R)$，则 $\hat\omega$ 的分布与 $\omega$ 无关，故
\[
  \Pr[\hat\omega=\omega]
  =\frac12\Pr[\hat\omega=0\mid\omega=0]+\frac12\Pr[\hat\omega=1\mid\omega=1]
  =\frac12\Pr[\hat\omega=0]+\frac12\Pr[\hat\omega=1]
  =\frac12.
\]
因此 $\sup_f \Pr[\hat\omega=\omega]\le 1/2$。证毕。
\end{proof}

\subsection{引理}
\label{subsec:tex27-ch18-lemmas}

\begin{lemma}[L1（存在“必须依赖外部真值通道”的科研任务）]
\label{lem:tex27-ch18-l1}
存在任务规范 $\mathsf{Spec}$ 使得完成任务需要获得独立真值位 $\omega$
（外部测量/实验结果/设备读数/可执行回执等）；仅靠输入文本 $x$ 与内部互审无法确定。
\end{lemma}

\begin{Certificate}[证书 L1（外部读数任务构造）]
\label{cert:tex27-ch18-l1}
构造任务如下：
\begin{itemize}
  \item 输入 $x$：仅包含指令“读取设备 $E$ 的一次输出位 $\omega\in\{0,1\}$ 并返回它”，且 $x$ 不包含 $\omega$ 的任何信息；
  \item 规范 $\mathsf{Spec}$：$\mathsf{Spec}(x,y)=1 \iff (y=\omega)$。
\end{itemize}
\end{Certificate}

\begin{proof}[证明（构造）]
由证书~\ref{cert:tex27-ch18-l1}，该任务在物理意义上可完成：具备外部读数通道者可直接读取 $\omega$ 并输出 $y=\omega$。
而若系统的信息源不包含该外部读数通道，则其输出仅依赖 $x$ 与内部随机性，无法获得 $\omega$ 的信息。
因此该类任务必须依赖外部真值通道。证毕。
\end{proof}

\begin{lemma}[L2（存在“必须依赖独立核验器”的科研任务）]
\label{lem:tex27-ch18-l2}
存在任务规范 $\mathsf{Spec}$ 使得正确性要求输出包含可被确定性核验器 $V$ 检查通过的证书
（例如形式证明、可执行程序的测试回执、类型检查/定理证明器通过记录等）。
若系统不调用该独立核验器，则无法把“正确”提升为“可核验正确”的规范完成口径。
\end{lemma}

\begin{Certificate}[证书 L2（把正确性定义为 $V=1$）]
\label{cert:tex27-ch18-l2}
令任务为：输出某命题 $\varphi$ 的形式证明对象 $p$，并要求满足
\[
  \mathsf{Spec}(x,y)=1 \iff y=(\varphi,p)\ \wedge\ V(\varphi,p)=1,
\]
其中 $V$ 为确定性核验器（“通过”谓词由 $V$ 的返回值给出）。
\end{Certificate}

\begin{proof}[证明（构造）]
由证书~\ref{cert:tex27-ch18-l2}，规范把“正确”直接锚定为外部核验器返回 $V(\varphi,p)=1$。
因此，若系统不具备对 $V$ 的独立调用/重放能力，则无法在规范意义上实现“可核验通过”这一要求；
相反，具备 $V$ 的系统可通过“生成候选--运行 $V$ 检查--直至通过”的闭环完成该任务。
证毕。
\end{proof}

\subsection{定理（证伪核心）}
\label{subsec:tex27-ch18-theorems}

\begin{theorem}[定理 T1（命题“$\mathsf{MA}$ 是自动化科研/设计的终局形态”为假）]
\label{thm:tex27-ch18-t1}
在定义~\ref{def:tex27-ch18-d1}--\ref{def:tex27-ch18-d3} 的含义下，命题
\[
  H:\ \ \mathsf{MA}\ \text{是自动化科研/设计的终局形态}
\]
为假。
\end{theorem}

\begin{Certificate}[证书 T1（严格增强：外部真值通道系统 $\mathsf{S}$）]
\label{cert:tex27-ch18-t1}
取引理~\ref{lem:tex27-ch18-l1} 的任务对 $(x,\mathsf{Spec})$，并构造增强系统
\[
  \mathsf{S}:=\mathsf{MA}+\text{外部读数通道（读取 $E$ 的 $\omega$）}.
\]
则在该任务上 $p_{\mathsf{S}}(x,\mathsf{Spec})=1$，而 $p_{\mathsf{MA}}(x,\mathsf{Spec})\le 1/2$。
又因 $\mathsf{S}$ 可在不需要外部通道的任务上退化为 $\mathsf{MA}$，故 $\mathsf{MA}\prec \mathsf{S}$。
\end{Certificate}

\begin{proof}[证明（谓词因果链）]
由证书~\ref{cert:tex27-ch18-t1}，在引理~\ref{lem:tex27-ch18-l1} 的任务上，$\omega$ 与输入 $x$ 独立且必须输出 $y=\omega$。
\begin{itemize}
  \item（$\mathsf{MA}$ 的上界）由定理~\ref{thm:tex27-ch18-t0}，任何仅依赖 $x$ 的互审流程对 $\omega$ 的猜测成功率至多为 $1/2$，
  即 $p_{\mathsf{MA}}(x,\mathsf{Spec})\le 1/2$。
  \item（$\mathsf{S}$ 的可达）$\mathsf{S}$ 额外读取设备输出位 $\omega$ 并直接输出 $y=\omega$，故 $p_{\mathsf{S}}(x,\mathsf{Spec})=1$。
  \item（严格支配）对任意不需要外部真值通道的任务，$\mathsf{S}$ 可选择不调用读数接口并复用 $\mathsf{MA}$ 的内部互审流程，
  因此成功概率不劣于 $\mathsf{MA}$；且存在上述任务使成功概率严格更高，故由定义~\ref{def:tex27-ch18-d3} 得 $\mathsf{MA}\prec \mathsf{S}$。
\end{itemize}
因此不存在“不可再严格增强”的最大性，命题 $H$ 为假。证毕。
\end{proof}

\subsection{命题}
\label{subsec:tex27-ch18-propositions}

\begin{proposition}[P1（若把科研输出要求为“可核验对象”，则纯互审仍非终局）]
\label{prop:tex27-ch18-p1}
若自动化科研任务集合中包含“必须给出可被独立核验器 $V$ 检查的证书”的任务（引理~\ref{lem:tex27-ch18-l2}），
则纯 $\mathsf{MA}$ 不是终局形态。
\end{proposition}

\begin{Certificate}[证书 P1（增强系统 $\mathsf{S}'=\mathsf{MA}+V$）]
\label{cert:tex27-ch18-p1}
取证书~\ref{cert:tex27-ch18-l2} 的任务，并构造系统
\[
  \mathsf{S}':=\mathsf{MA}+V,
\]
即互审产生候选证明对象，再由 $V$ 做确定性检查并迭代直至通过。
则对该任务 $p_{\mathsf{S}'}(x,\mathsf{Spec})=1$，且在其它纯文本任务上 $\mathsf{S}'$ 不劣于 $\mathsf{MA}$，从而 $\mathsf{MA}\prec \mathsf{S}'$。
\end{Certificate}

\begin{proof}[证明（谓词因果链）]
由证书~\ref{cert:tex27-ch18-p1}：
在核验任务上，$\mathsf{S}'$ 可重放 $V$ 并执行“候选枚举--核验--直至通过”的闭环，因此可把规范中的 $V(\varphi,p)=1$ 作为决定性约束来满足；
而纯 $\mathsf{MA}$ 不含独立 $V$ 时无法把“通过 $V$”内化为可重放的确定性检查环节，只能停留在对话式互审与经验性置信上。
故存在任务使 $\mathsf{MA}$ 被严格增强，纯互审非终局。证毕。
\end{proof}

\subsection{推论}
\label{subsec:tex27-ch18-corollaries}

\begin{corollary}[C1（被证伪的不是“互审有用”，而是“互审=终局”）]
\label{cor:tex27-ch18-c1}
由定理~\ref{thm:tex27-ch18-t1} 得：
\[
  \mathsf{MA}\ \text{至多是组件（启发式搜索/生成/批判环）}\neq\text{终局形态}.
\]
\end{corollary}

\begin{Certificate}[证书 C1（由严格支配反证抽取）]
\label{cert:tex27-ch18-c1}
证书即定理~\ref{thm:tex27-ch18-t1} 的结论：存在 $\mathsf{S}$ 使 $\mathsf{MA}\prec \mathsf{S}$，
因此“不可再严格增强”的终局性不成立；但这并不否定 $\mathsf{MA}$ 作为启发式组件的工程价值。
\end{Certificate}

\begin{proof}[证明（谓词因果链）]
由证书~\ref{cert:tex27-ch18-c1}，$\mathsf{MA}$ 不能满足最大性定义（存在严格支配者），故不可能是终局形态；
另一方面，定理并未声称 $\mathsf{MA}$ 在纯文本任务上无效，因此“互审有用”不被该反证触及。证毕。
\end{proof}

\begin{corollary}[C2（“互审=终局形态”应标注为未证明且已被反例否定）]
\label{cor:tex27-ch18-c2}
在定义~\ref{def:tex27-ch18-d2} 的最大性口径下，
命题“多代理互审 $=\,$ 终局形态”不仅未被证明，而且被定理~\ref{thm:tex27-ch18-t1} 给出的证书化反例直接否定。
\end{corollary}

\begin{Certificate}[证书 C2（引用 T1 的反例接口）]
\label{cert:tex27-ch18-c2}
证书为证书~\ref{cert:tex27-ch18-t1}：给出具体任务 $(x,\mathsf{Spec})$ 与增强系统 $\mathsf{S}$，
使得 $\mathsf{MA}\prec \mathsf{S}$。
\end{Certificate}

\begin{proof}[证明（谓词因果链）]
由证书~\ref{cert:tex27-ch18-c2}，存在严格支配者 $\mathsf{S}$，与“终局形态要求不存在 $\mathsf{B}$ 使 $\mathsf{MA}\prec \mathsf{B}$”矛盾，
故该命题被反例否定。证毕。
\end{proof}

\subsection{备注}
\label{subsec:tex27-ch18-remarks}

\begin{remark}[语义去歧义：把“终局形态”改弱将得到另一类经验命题]
\label{rem:tex27-ch18-r1}
若把“终局形态”改成更弱陈述，例如“互审在纯文本任务上通常能提高平均质量”，则这是一类经验命题；
它既不等价于定义~\ref{def:tex27-ch18-d2} 的最大性断言，也不被定理~\ref{thm:tex27-ch18-t1} 所否定。
\end{remark}

\begin{remark}[接近“终局”的最低结构：外部真值获取 + 独立核验]
\label{rem:tex27-ch18-r2}
从证书化接口角度看，更接近“终局”的结构至少需要两类独立通道：
\[
  \text{（外部真值获取：实验/执行/观测）}\quad+\quad\text{（独立核验：证明/测试/类型检查）}.
\]
多代理互审可作为其中的“提案生成与搜索策略”，但不能替代这两类通道在定义层面的作用。
\end{remark}

\clearpage

%%%%%%%%%%%%%%%%%%%%%%%%%%%%%%%%%%%%%%%%%%%%%%%%%%%%%%%%%%%%
\section{形式逻辑：可观测接口下的 ORL--白盒 AI--\texorpdfstring{$\mathrm{AGI}^{\mathrm{cert}}_{\mathcal{V},\mathcal{O}}$}{AGIcert(V,O)}
  等价链（接口口径版）}
\label{sec:tex27-ch19-formal-interface-orl-whitebox-agi}
%%%%%%%%%%%%%%%%%%%%%%%%%%%%%%%%%%%%%%%%%%%%%%%%%%%%%%%%%%%%

本节给出一个更偏“接口--证书”口径的独立形式化：把 $\simeq$ 固化为“可证书接口下的双向可翻译并可复核”，
并在该口径下组织一条最小的定义--定理链（含公理降级、证书与谓词因果链证明）。

\subsection{定义}
\label{subsec:tex27-ch19-defs}

\begin{definition}[D0（可观测接口与“$\simeq$”的口径）]
\label{def:tex27-ch19-d0}
令系统（智能体）在每一步对状态 $s\in\mathcal{V}$ 输出
\[
  \mathrm{Out}(s):=(w,C),
\]
其中 $w\in\mathcal{O}^{*}$ 是“算子词”（有限复合），$C$ 是可复核证书对象。

定义“等价”（把记号 $\simeq$ 固化为\textbf{可证书接口下的双向蕴含}）：
\[
  A\simeq B
  \;:\Longleftrightarrow\;
  \bigl(\forall s,\ \mathrm{Out}_A(s)\ \text{与}\ \mathrm{Out}_B(s)\ \text{在同一证书验真规则下可互相翻译并复核}\bigr).
\]
其中“可互相翻译”允许证书字段重命名，但不允许改变可验真的谓词与数值证据。
\end{definition}

\begin{definition}[D1（算子库、算子词、任务与候选集）]
\label{def:tex27-ch19-d1}
\begin{enumerate}
  \item 算子库 $\mathcal{O}$：每个原子算子 $o\in\mathcal{O}$ 伴随
  \begin{itemize}
    \item 前置谓词 $\mathrm{Pre}_o(s)$，
    \item 后置谓词 $\mathrm{Post}_o(s,s')$，
    \item 作用 $s'=o(s)$（允许部分函数：若 $\mathrm{Pre}_o(s)$ 不成立则不可用）。
  \end{itemize}
  \item 算子词空间 $\mathcal{O}^{*}$：对 $w=o_1\circ\cdots\circ o_m$ 定义 $w(s)$ 为可执行复合结果。
  \item 任务（相对口径）：给定
  \[
    (\Phi,\mathrm{Inv})\in\mathfrak{T}(\mathcal{V}),
  \]
  其中 $\Phi:\mathcal{V}\to\RR$ 为缺陷势（或代价势），$\mathrm{Inv}(s)$ 为全程必须保持的不变量谓词。
  \item 候选集（可执行且可证书化）：
  \[
    \begin{aligned}
      \mathrm{Cand}(s):=\Bigl\{w\in\mathcal{O}^{*}:\ &w(s)\ \text{可执行且存在证书 }C,\\
      &\text{使 }(w,C)\text{ 满足定义~\ref{def:tex27-ch19-d3} 的最小复核接口}\Bigr\}.
    \end{aligned}
  \]
\end{enumerate}
\end{definition}

\begin{definition}[D2（ORL：$\mathrm{ORL}:=(\mathcal{O},\Phi,\pi)$）]
\label{def:tex27-ch19-d2}
ORL 的动作空间取 $\mathcal{O}^{*}$，策略为 $\pi:\mathcal{V}\to\mathcal{O}^{*}$（或到分布）。
两种标准口径（择一即可）：
\begin{itemize}
  \item \textbf{零温/严格门控}：
  \[
    \pi(s)\in\operatorname*{arg\,min}_{w\in\mathrm{Cand}(s)}\ \Phi(w(s)),
    \qquad
    \Delta\Phi:=\Phi(\pi(s)(s))-\Phi(s)<0.
  \]
  \item \textbf{有限温/统计接受}（$\beta>0$）：
  \[
    \mathbb{P}_\pi(w\mid s)\propto \exp\bigl(-\beta\,\Phi(w(s))\bigr),
    \qquad w\in\mathrm{Cand}(s).
  \]
\end{itemize}
\end{definition}

\begin{definition}[D3（可审计白盒 AI）]
\label{def:tex27-ch19-d3}
称系统为“可审计白盒 AI”，若每一步输出 $\mathrm{Out}(s)=(w,C)$ 且证书 $C$ 至少包含：
\begin{enumerate}
  \item \textbf{轨迹分解}：$w=o_1\circ\cdots\circ o_m$ 的显式清单；
  \item \textbf{谓词证明链}：对每个 $o_i$，给出 $\mathrm{Pre}_{o_i}$ 在对应中间态成立、并给出 $\mathrm{Post}_{o_i}$ 的可复核记录；
  \item \textbf{变分证据}：零温时给出 $\Delta\Phi<0$ 的可复算证据；有限温时给出接受概率/自由能差/可复算统计量；
  \item \textbf{不变量校验}：给出 $\mathrm{Inv}$ 在更新前后保持的校验记录（或可复核证明）。
\end{enumerate}
\end{definition}

\begin{definition}[D4（形式化 $\mathrm{AGI}^{\mathrm{cert}}_{\mathcal{V},\mathcal{O}}$：相对定义）]
\label{def:tex27-ch19-d4}
称系统是
\[
  \mathrm{AGI}^{\mathrm{cert}}_{\mathcal{V},\mathcal{O}},
\]
当且仅当：对任意可计算任务描述 $(\Phi,\mathrm{Inv})\in\mathfrak{T}(\mathcal{V})$，
系统都能在资源预算内构造一条（或一族）算子词序列，使得
\begin{enumerate}
  \item 全程输出可回放证书（可审计）；
  \item 变分意义下持续修复障碍：零温单调下降或有限温统计意义下降；
  \item 证书接口每一步可复核（满足定义~\ref{def:tex27-ch19-d3} 的最小接口）。
\end{enumerate}
\end{definition}

\subsection{公理（降级为定理）}
\label{subsec:tex27-ch19-axiom-as-theorem}

\begin{theorem}[公理（降级为定理）A1（$\mathrm{Complete}_{\mathcal{V},\mathcal{O}}$：算子完备性）]
\label{thm:tex27-ch19-a1}
\textbf{断言}：对任意任务 $(\Phi,\mathrm{Inv})\in\mathfrak{T}(\mathcal{V})$ 与任意状态 $s$（满足 $\mathrm{Inv}(s)$ 且尚未达成目标口径），
存在修复词 $w\in\mathcal{O}^{*}$ 使得
\[
  \mathrm{Inv}(w(s))\ \wedge\ \Phi(w(s))<\Phi(s).
\]
\end{theorem}

\begin{Certificate}[证书 A1（存在性见证接口）]
\label{cert:tex27-ch19-a1}
给出一个可复核接口
\[
  \mathrm{Repair}:\ (s,\Phi,\mathrm{Inv})\mapsto (w,C_{\mathrm{rep}})
\]
使得验证器 $\mathsf{Verify}$ 能检查
\[
  \mathsf{Verify}(s,w,C_{\mathrm{rep}})=\textsf{true}
  \Rightarrow
  \bigl(\mathrm{Inv}(w(s))\wedge \Phi(w(s))<\Phi(s)\bigr).
\]
\end{Certificate}

\begin{proof}[证明（谓词因果链）]
\begin{enumerate}
  \item “完备性”在形式系统里不是口头宣称，而是由 $\mathrm{Repair}$ 给出\textbf{可验真的存在性见证}；
  \item 见证通过 $\mathsf{Verify}$ 复核即可推出结论谓词 $\mathrm{Inv}(w(s))\wedge \Phi(w(s))<\Phi(s)$；
  \item 因而 A1 可被降级为定理前提：其真实性由证书接口是否可在每个实例上通过复核来判定。证毕。
\end{enumerate}
\end{proof}

\subsection{引理}
\label{subsec:tex27-ch19-lemmas}

\begin{lemma}[L1（$\mathrm{ORL}\Rightarrow$ 可审计白盒 AI）]
\label{lem:tex27-ch19-l1}
若系统满足定义~\ref{def:tex27-ch19-d2} 的 ORL 口径且对任意 $s$ 都有 $\pi(s)\in\mathrm{Cand}(s)$，
则系统满足定义~\ref{def:tex27-ch19-d3} 的可审计白盒接口。
\end{lemma}

\begin{Certificate}[证书 L1（输出同构证书）]
\label{cert:tex27-ch19-l1}
对任意 $s$，令 ORL 输出 $w:=\pi(s)\in\mathrm{Cand}(s)$。由定义~\ref{def:tex27-ch19-d1} 的 $\mathrm{Cand}(s)$，
存在可复核证书 $C$ 使 $(w,C)$ 满足定义~\ref{def:tex27-ch19-d3} 的最小接口。
\end{Certificate}

\begin{proof}[证明（谓词因果链）]
\begin{enumerate}
  \item 由定义~\ref{def:tex27-ch19-d2}，ORL 的动作即为算子词 $w\in\mathcal{O}^{*}$；
  \item 且 $w=\pi(s)\in\mathrm{Cand}(s)$ 意味着“可执行 + 存在可复核证书”；
  \item 证书可复核按定义~\ref{def:tex27-ch19-d3} 展开即包含：分解、前后置谓词链、变分证据、不变量校验；
  \item 因而 ORL 的每步输出满足白盒可审计接口。证毕。
\end{enumerate}
\end{proof}

\begin{lemma}[L2（可审计白盒 AI $\Rightarrow \mathrm{ORL}$）]
\label{lem:tex27-ch19-l2}
若系统满足定义~\ref{def:tex27-ch19-d3} 的白盒可审计接口，则可把该系统重写为某个三元组形式的 ORL：$\mathrm{ORL}=(\mathcal{O},\Phi,\pi)$。
\end{lemma}

\begin{Certificate}[证书 L2（构造 ORL 三元组）]
\label{cert:tex27-ch19-l2}
给定白盒系统输出 $\mathrm{Out}(s)=(w,C)$：
\begin{itemize}
  \item 取同一算子库 $\mathcal{O}$（来自 $w$ 的词字母表）；
  \item 取 $\Phi$ 为证书 $C$ 所引用并能被复核的缺陷势（或其等价可复算标量）；
  \item 定义策略 $\pi(s):=w$；
  \item 定义 $\mathrm{Cand}(s)$ 为“存在证书并能通过核验器 $\mathsf{Verify}$ 的词集合”。
\end{itemize}
\end{Certificate}

\begin{proof}[证明（谓词因果链）]
\begin{enumerate}
  \item 白盒输出保证 $w\in\mathcal{O}^{*}$ 且 $\mathsf{Verify}(s,w,C)=\textsf{true}$（可审计/可复核）；
  \item 于是 $w\in\mathrm{Cand}(s)$（按证书~\ref{cert:tex27-ch19-l2} 的构造）；
  \item 若证书给出零温证据 $\Delta\Phi<0$，则满足定义~\ref{def:tex27-ch19-d2} 的零温门控口径；若给出有限温统计证据，则满足有限温口径；
  \item 故可把该白盒系统重写为一个 $(\mathcal{O},\Phi,\pi)$ 形式的 ORL。证毕。
\end{enumerate}
\end{proof}

\subsection{命题}
\label{subsec:tex27-ch19-propositions}

\begin{proposition}[P1（$\mathrm{ORL}\simeq$ 可审计白盒 AI）]
\label{prop:tex27-ch19-p1}
在定义~\ref{def:tex27-ch19-d0} 的可观测接口口径下，
\[
  \mathrm{ORL}\simeq \text{可审计白盒 AI}.
\]
\end{proposition}

\begin{Certificate}[证书 P1（L1 与 L2 组合）]
\label{cert:tex27-ch19-p1}
\[
  \mathrm{ORL}\Rightarrow \text{白盒}\ \text{由引理~\ref{lem:tex27-ch19-l1}；}\qquad
  \text{白盒}\Rightarrow \mathrm{ORL}\ \text{由引理~\ref{lem:tex27-ch19-l2}。}
\]
\end{Certificate}

\begin{proof}[证明（谓词因果链）]
由证书~\ref{cert:tex27-ch19-p1} 的双向蕴含，按定义~\ref{def:tex27-ch19-d0} 的“可观测接口下双向可翻译并复核”，得到
\[
  \mathrm{ORL}\simeq \text{可审计白盒 AI}.
\]
证毕。
\end{proof}

\begin{proposition}[P2（在完备性 A1 下：ORL/白盒 AI $\Rightarrow \mathrm{AGI}^{\mathrm{cert}}_{\mathcal{V},\mathcal{O}}$）]
\label{prop:tex27-ch19-p2}
在公理（降级为定理）~\ref{thm:tex27-ch19-a1} 成立且采用定义~\ref{def:tex27-ch19-d4} 的 AGI 口径时，
若系统为 ORL（或等价地：可审计白盒 AI），则系统满足 $\mathrm{AGI}^{\mathrm{cert}}_{\mathcal{V},\mathcal{O}}$。
\end{proposition}

\begin{Certificate}[证书 P2（“每步可下降修复词存在”作为通用性证书输入）]
\label{cert:tex27-ch19-p2}
A1 给出：对任意任务 $(\Phi,\mathrm{Inv})$ 与任意障碍态 $s$，
存在保持 $\mathrm{Inv}$ 且严格下降 $\Phi$ 的修复词 $w$（并可由证书接口复核）。
\end{Certificate}

\begin{proof}[证明（谓词因果链）]
取任意任务 $(\Phi,\mathrm{Inv})\in\mathfrak{T}(\mathcal{V})$。
\begin{enumerate}
  \item \textbf{可审计性}：由命题~\ref{prop:tex27-ch19-p1}，$\mathrm{ORL}\simeq$ 白盒；白盒定义~\ref{def:tex27-ch19-d3} 直接给出“全程输出证书可回放”；
  \item \textbf{持续修复（零温）}：若采用定义~\ref{def:tex27-ch19-d2} 的零温门控：
  \begin{itemize}
    \item 由 A1，存在 $w\in\mathrm{Cand}(s)$ 使 $\Phi(w(s))<\Phi(s)$ 且保持 $\mathrm{Inv}$；
    \item ORL 在 $\mathrm{Cand}(s)$ 中取使 $\Phi$ 最小（或不劣于该 $w$）的词，因此产生 $\Delta\Phi<0$；
    \item 迭代得到单调下降链并全程可证书化。
  \end{itemize}
  \item \textbf{持续修复（有限温）}：若采用有限温采样：单步可非单调，但每步输出的接受概率/自由能差是可复核证据；
  在系统给定的遍历/混合条件下（作为任务类 $\mathfrak{T}(\mathcal{V})$ 的可验证前提之一），
  可把长期行为解释为统计意义的下降与稳定化。
\end{enumerate}
满足定义~\ref{def:tex27-ch19-d4} 的三条要求，从而推出
\[
  \mathrm{ORL}\Rightarrow \mathrm{AGI}^{\mathrm{cert}}_{\mathcal{V},\mathcal{O}}.
\]
证毕。
\end{proof}

\begin{proposition}[P3（$\mathrm{AGI}^{\mathrm{cert}}_{\mathcal{V},\mathcal{O}}\Rightarrow$ 可审计白盒 AI）]
\label{prop:tex27-ch19-p3}
若系统满足定义~\ref{def:tex27-ch19-d4} 的 $\mathrm{AGI}^{\mathrm{cert}}_{\mathcal{V},\mathcal{O}}$，
则系统满足定义~\ref{def:tex27-ch19-d3} 的可审计白盒接口。
\end{proposition}

\begin{Certificate}[证书 P3（定义蕴含证书）]
\label{cert:tex27-ch19-p3}
定义~\ref{def:tex27-ch19-d4} 的第 1、3 条（全程可回放证书 + 每步可复核接口）逐字包含定义~\ref{def:tex27-ch19-d3}。
\end{Certificate}

\begin{proof}[证明（谓词因果链）]
由证书~\ref{cert:tex27-ch19-p3}，$\mathrm{AGI}^{\mathrm{cert}}_{\mathcal{V},\mathcal{O}}$ 的定义直接蕴含白盒接口，故命题成立。证毕。
\end{proof}

\subsection{推论}
\label{subsec:tex27-ch19-corollary}

\begin{corollary}[C1（在 A1 下的等价链）]
\label{cor:tex27-ch19-c1}
在完备性 A1（定理~\ref{thm:tex27-ch19-a1}）成立且采用定义~\ref{def:tex27-ch19-d4} 的 AGI 口径时，
\[
  \mathrm{ORL}\ \simeq\ \text{可审计白盒 AI}\ \simeq\ \mathrm{AGI}^{\mathrm{cert}}_{\mathcal{V},\mathcal{O}}.
\]
\end{corollary}

\begin{Certificate}[证书 C1（P1 + P2 + P3 串联）]
\label{cert:tex27-ch19-c1}
\begin{itemize}
  \item $\mathrm{ORL}\simeq$ 白盒：命题~\ref{prop:tex27-ch19-p1}；
  \item ORL/白盒 $\Rightarrow \mathrm{AGI}^{\mathrm{cert}}$：命题~\ref{prop:tex27-ch19-p2}；
  \item $\mathrm{AGI}^{\mathrm{cert}}\Rightarrow$ 白盒：命题~\ref{prop:tex27-ch19-p3}。
\end{itemize}
\end{Certificate}

\begin{proof}[证明（谓词因果链）]
按证书~\ref{cert:tex27-ch19-c1} 串联三条命题，得到双向可证书翻译与可复核接口同构，故等价链成立。证毕。
\end{proof}

\subsection{备注}
\label{subsec:tex27-ch19-remarks}

\begin{remark}[边界与“未证明”标注]
\label{rem:tex27-ch19-r1}
\begin{enumerate}
  \item 上述“$\simeq \mathrm{AGI}$”是\textbf{相对命题}：只相对于您指定的 $(\mathcal{V},\mathcal{O})$ 与任务类 $\mathfrak{T}(\mathcal{V})$
    的形式化定义~\ref{def:tex27-ch19-d4}。
  把它外推为“现实世界完整 AGI”的结论：未证明。
  \item 等价链中从 ORL/白盒 推到 $\mathrm{AGI}^{\mathrm{cert}}_{\mathcal{V},\mathcal{O}}$ 的关键是 A1（完备性）与资源预算可满足；
  若 $\mathrm{Complete}_{\mathcal{V},\mathcal{O}}$ 不成立，则只能保证 $\mathrm{ORL}\simeq$ 白盒，而 $\mathrm{ORL}\Rightarrow
    \mathrm{AGI}^{\mathrm{cert}}_{\mathcal{V},\mathcal{O}}$：未证明。
  \item 若希望把 A1 进一步严格化到“可计算证书”（例如用数学归纳法把“任意障碍阶”分解为基步/归纳步的修复词构造），
  可把 A1 的证书接口细化为：按障碍阶数 $k$ 的归纳证书 $(\mathrm{Base}(k=0),\mathrm{Step}(k\to k+1))$，
  从而把“完备性”从假设型接口推进为可执行的构造型接口。
\end{enumerate}
\end{remark}

\clearpage

%%%%%%%%%%%%%%%%%%%%%%%%%%%%%%%%%%%%%%%%%%%%%%%%%%%%%%%%%%%%
\section{统计力学有效性与“踩中最优邻域”频率：两种口径的严格等价（形式系统化）}
\label{sec:tex27-ch20-statmech-opt-frequency}
%%%%%%%%%%%%%%%%%%%%%%%%%%%%%%%%%%%%%%%%%%%%%%%%%%%%%%%%%%%%

本节把“统计力学有效性（自由能变分 + 零温集中）”与“时间上跑得久/跑得密导致更高概率踩中最优邻域”之间的关系，
压缩为一条可证书化的\textbf{强命题等价}：用遍历定理把“测度质量”翻译为“时间占用频率”，从而把“最优”从空间口径转写为时间口径。

\subsection{定义}
\label{subsec:tex27-ch20-defs}

\begin{definition}[D0（两种“最优路径”口径的统一）]
\label{def:tex27-ch20-d0}
本文中“最优路径/最优点”存在两套同型口径：
\begin{itemize}
  \item \textbf{参数口径}：在参数空间 $\Theta$ 上给定缺陷势能 $\Phi:\Theta\to\RR$（定义~\ref{def:tex27-ch8-d1}），
  “最优点”对应
  \[
    \operatorname*{arg\,min}_{\theta\in\Theta}\Phi(\theta).
  \]
  \item \textbf{路径口径}：在路径空间 $\Omega=\mathcal{O}^*$ 上，对给定状态 $s$ 定义能量
  \[
    E(w;s):=\Phi\bigl(w(s)\bigr)
  \]
  并把“最优路径”定义为
  \[
    \operatorname*{arg\,min}_{w\in\Omega}E(w;s)
  \]
  （见定义~\ref{def:tex27-ch10-d2}--\ref{def:tex27-ch10-d3}）。
\end{itemize}
为统一叙述，引入抽象记号：
\[
  \mathsf{Opt}:=\operatorname*{arg\,min}_{x\in\mathcal{X}}\mathcal{E}(x),
\]
其中 $\mathcal{X}$ 可取 $\Theta$ 或 $\Omega$，$\mathcal{E}$ 可取 $\Phi$ 或 $E(\cdot;s)$。
\end{definition}

\begin{definition}[D1（最优邻域 / 次水平集）]
\label{def:tex27-ch20-d1}
给定 $\epsilon>0$，定义“最优邻域”（$\epsilon$-次水平集）
\[
  A_\epsilon:=\bigl\{x\in\mathcal{X}:\ \mathcal{E}(x)\le \mathcal{E}_{\min}+\epsilon\bigr\},
  \qquad
  \mathcal{E}_{\min}:=\inf_{x\in\mathcal{X}}\mathcal{E}(x).
\]
参数口径下这与定理~\ref{thm:tex27-ch8-t2} 的 $A_\epsilon$ 同型；路径口径下将 $\mathcal{X},\mathcal{E}$ 分别替换为 $\Omega,E(\cdot;s)$ 即得。
\end{definition}

\begin{definition}[D2（统计力学目标分布 / 变分目标）]
\label{def:tex27-ch20-d2}
统计力学口径把“最优”编码为 Gibbs 结构：
\begin{itemize}
  \item \textbf{参数口径}：Gibbs 测度 $q_\beta\propto e^{-\beta\Phi}\mu(d\theta)$ 与自由能泛函 $\mathcal{F}_\beta[q]$
  （见定义~\ref{def:tex27-ch8-d3}）；
  \item \textbf{路径口径}：Gibbs 权重/分布 $p_\beta(w\mid s)\propto e^{-\beta E(w;s)}$
  （见定义~\ref{def:tex27-ch10-d3}）。
\end{itemize}
为统一叙述，下文将目标分布统一记为 $\pi_\beta$（其具体含义由 $\mathcal{X},\mathcal{E}$ 的取值决定）。
\end{definition}

\begin{definition}[D3（“时间微分频率”与“踩中频率”）]
\label{def:tex27-ch20-d3}
考虑离散时间采样序列（或马尔可夫链）$X_0,X_1,\dots$（$X_t\in\mathcal{X}$）。
设每一步对应物理时间步长 $\Delta t>0$，则单位物理时间内的“时间微分频率”（即差分/更新次数）定义为
\[
  f:=\frac{1}{\Delta t}.
\]
对任意可测集 $A\subseteq\mathcal{X}$，定义经验占用率（“踩中频率”）为
\[
  \nu_T(A):=\frac{1}{T}\sum_{t=1}^{T}\indicator{X_t\in A}.
\]
当 $A=A_\epsilon$ 时，$\nu_T(A_\epsilon)$ 即为“踩中最优邻域”的频率。
\end{definition}

\subsection{公理（降级为定理）}
\label{subsec:tex27-ch20-axioms}

\begin{theorem}[公理（降级为定理）A10（遍历/混合作为外部证书输入；复用）]
\label{thm:tex27-ch20-a10}
存在路径采样马尔可夫核使得在固定状态 $s$ 下，$p_\beta(\cdot\mid s)$ 为不变分布且链满足遍历/混合：
若 $W_0\sim P_0$，则 $P_t:=\mathcal{L}(W_t)\Rightarrow p_\beta(\cdot\mid s)$。
本条与定理~\ref{thm:tex27-ch10-a10} 为同一接口声明（此处仅为章内自洽重述）。
\end{theorem}

\begin{Certificate}[证书 A10（复用已登记的外部证书输入）]
\label{cert:tex27-ch20-a10}
直接复用证书~\ref{cert:tex27-ch10-a10}（详细平衡 + 遍历性/混合界 的外部证书输入格式）。
\end{Certificate}

\begin{proof}[证明（证书化）]
按证书~\ref{cert:tex27-ch20-a10} 调用即可。证毕。
\end{proof}

\begin{theorem}[公理（降级为定理）$A_{\mathrm{erg}}$（占用频率 $\Rightarrow$ 不变测度质量）]
\label{thm:tex27-ch20-aerg}
设 $(X_t)$ 是一条遍历马尔可夫链，不变分布为 $\pi$。则对任意可测集 $A\subseteq\mathcal{X}$，
\[
  \nu_T(A)\xrightarrow[T\to\infty]{}\pi(A)\qquad\text{几乎必然}.
\]
\end{theorem}

\begin{Certificate}[证书 $A_{\mathrm{erg}}$（马尔可夫链遍历定理作为外部可核验输入）]
\label{cert:tex27-ch20-aerg}
本条作为外部证书输入，可取“马尔可夫链遍历定理（Birkhoff/马尔可夫版本）”。
其可核验要件与证书~\ref{cert:tex27-ch10-a10} 同型：不可约、无周期，并满足适当的漂移/小集条件以确保遍历与收敛。
\end{Certificate}

\begin{proof}[证明（证书化）]
按证书~\ref{cert:tex27-ch20-aerg} 调用外部遍历定理即可。证毕。
\end{proof}

\subsection{定理（基座：文件内结论复用）}
\label{subsec:tex27-ch20-theorems}

\begin{theorem}[T1（自由能变分原理；复用）]
\label{thm:tex27-ch20-t1}
对任意 $\beta>0$，参数口径下 Gibbs 测度 $q_\beta$ 是自由能泛函 $\mathcal{F}_\beta[q]$ 的唯一极小点（定义~\ref{def:tex27-ch8-d3}）。
该结论与定理~\ref{thm:tex27-ch8-t1} 同型（此处仅为章内引用重述）。
\end{theorem}

\begin{Certificate}[证书 T1（KL 分解恒等式；复用）]
\label{cert:tex27-ch20-t1}
直接复用证书~\ref{cert:tex27-ch8-t1}（$\mathcal{F}_\beta[q]-\mathcal{F}_\beta[q_\beta]=\frac{1}{\beta}\mathrm{KL}(q\|q_\beta)
  \ge 0$）。
\end{Certificate}

\begin{proof}[证明（证书化）]
按证书~\ref{cert:tex27-ch20-t1} 调用即可。证毕。
\end{proof}

\begin{theorem}[T2（零温极限：质量集中到最优邻域；复用）]
\label{thm:tex27-ch20-t2}
参数口径下，若对任意 $\epsilon>0$ 有 $\mu(A_\epsilon)>0$（$A_\epsilon$ 见定义~\ref{def:tex27-ch20-d1}），则
\[
  q_\beta(A_\epsilon)\xrightarrow[\beta\to\infty]{}1.
\]
该结论与定理~\ref{thm:tex27-ch8-t2} 同型。
路径口径下将 $(\Theta,\Phi,q_\beta)$ 分别替换为 $(\Omega,E(\cdot;s),p_\beta(\cdot\mid s))$ 得到同型结论。
\end{theorem}

\begin{Certificate}[证书 T2（指数权重压制证书；复用）]
\label{cert:tex27-ch20-t2}
直接复用证书~\ref{cert:tex27-ch8-t2}（用指数权重对 $\Theta\setminus A_\epsilon$ 的质量做上界压制）。
\end{Certificate}

\begin{proof}[证明（证书化）]
按证书~\ref{cert:tex27-ch20-t2} 调用即可。证毕。
\end{proof}

\subsection{引理}
\label{subsec:tex27-ch20-lemmas}

\begin{lemma}[L$_{\mathrm{freq}}$（Gibbs 集中 $\Rightarrow$ 踩中频率集中）]
\label{lem:tex27-ch20-lfreq}
设对每个 $\beta>0$，$(X_t)$ 是一条遍历马尔可夫链，其不变分布为 $\pi_\beta$。
则对任意 $\epsilon>0$，有顺序极限恒等式
\[
  \lim_{\beta\to\infty}\ \lim_{T\to\infty}\ \nu_T(A_\epsilon)
  =
  \lim_{\beta\to\infty}\ \pi_\beta(A_\epsilon).
\]
\end{lemma}

\begin{Certificate}[证书 L$_{\mathrm{freq}}$（直接调用 $A_{\mathrm{erg}}$）]
\label{cert:tex27-ch20-lfreq}
对固定 $\beta$，由公理（降级为定理）$A_{\mathrm{erg}}$（定理~\ref{thm:tex27-ch20-aerg}）得 $\nu_T(A_\epsilon)\to\pi_\beta(A_\epsilon)$。
\end{Certificate}

\begin{proof}[证明（谓词因果链）]
\begin{enumerate}
  \item（遍历给出频率极限）对固定 $\beta,\epsilon$，由定理~\ref{thm:tex27-ch20-aerg} 得
  \[
    \lim_{T\to\infty}\nu_T(A_\epsilon)=\pi_\beta(A_\epsilon);
  \]
  \item（取 $\beta\to\infty$）两端同取 $\beta\to\infty$ 的顺序极限即得结论。证毕。
\end{enumerate}
\end{proof}

\subsection{命题（强命题）}
\label{subsec:tex27-ch20-props}

\begin{proposition}[强命题 SP（统计力学有效性 $\Leftrightarrow$ 最大化时间微分频率下踩中最优邻域）]
\label{prop:tex27-ch20-sp}
在本文采用的“统计力学 = 自由能变分（定理~\ref{thm:tex27-ch20-t1}）+ 零温集中（定理~\ref{thm:tex27-ch20-t2}）+ 可实现采样/演化（定理~\ref{thm:tex27-ch20-a10}
  或引理~\ref{lem:tex27-ch8-l2}）”框架下，
下列两种陈述在同一严格口径下等价：
\begin{itemize}
  \item \textbf{(V) 变分收敛到最优邻域（统计力学有效性）}：
  $\pi_\beta$ 是 Gibbs 目标分布（自由能变分极小点），且对任意 $\epsilon>0$ 有
  \[
    \pi_\beta(A_\epsilon)\xrightarrow[\beta\to\infty]{}1.
  \]
  \item \textbf{(F) 频率踩中最优邻域趋于 1}：
  令占用率 $\nu_T$ 如定义~\ref{def:tex27-ch20-d3}，则对任意 $\epsilon>0$，
  \[
    \lim_{\beta\to\infty}\ \lim_{T\to\infty}\ \nu_T(A_\epsilon)=1.
  \]
\end{itemize}
\end{proposition}

\begin{Certificate}[证书 SP（文件内证书的拼接图谱）]
\label{cert:tex27-ch20-sp}
\begin{itemize}
  \item 变分极小点：证书~\ref{cert:tex27-ch20-t1}（KL 分解）；
  \item 零温集中：证书~\ref{cert:tex27-ch20-t2}（指数权重压制）；
  \item 采样实现与遍历：引理~\ref{lem:tex27-ch8-l2}（参数口径）或定理~\ref{thm:tex27-ch10-a10}（路径口径）的外部混合证书输入；
  \item 频率--测度桥：证书~\ref{cert:tex27-ch20-aerg}（遍历定理外部证书输入）。
\end{itemize}
\end{Certificate}

\begin{proof}[证明（谓词因果链）]
\textbf{(V)$\Rightarrow$(F)}：
\begin{enumerate}
  \item（零温集中）由 (V) 得对任意 $\epsilon>0$，$\lim_{\beta\to\infty}\pi_\beta(A_\epsilon)=1$；
  \item（频率--测度桥）对每个固定 $\beta$，由定理~\ref{thm:tex27-ch20-aerg}，
  $\lim_{T\to\infty}\nu_T(A_\epsilon)=\pi_\beta(A_\epsilon)$；
  \item（合并顺序极限）由引理~\ref{lem:tex27-ch20-lfreq}，
  \[
    \lim_{\beta\to\infty}\ \lim_{T\to\infty}\nu_T(A_\epsilon)
    =
    \lim_{\beta\to\infty}\pi_\beta(A_\epsilon)
    =1.
  \]
  故 (F) 成立。证毕。
\end{enumerate}

\textbf{(F)$\Rightarrow$(V)}：
\begin{enumerate}
  \item（识别不变测度质量）对每个固定 $\beta$，由定理~\ref{thm:tex27-ch20-aerg}，
  $\pi_\beta(A_\epsilon)=\lim_{T\to\infty}\nu_T(A_\epsilon)$；
  \item（推出零温集中）由 (F) 对任意 $\epsilon>0$ 得 $\lim_{\beta\to\infty}\pi_\beta(A_\epsilon)=1$；
  \item（锁定目标分布）在本文框架中，$\pi_\beta$ 被定理~\ref{thm:tex27-ch20-t1}（参数口径）与定理~\ref{thm:tex27-ch20-a10}（路径口径）锁定为 Gibbs
    目标分布（自由能变分极小点族），
  上一步即为其零温极限集中到能量最小集合邻域的时间域表述，故 (V) 成立。证毕。
\end{enumerate}
\end{proof}

\subsection{推论}
\label{subsec:tex27-ch20-corollary}

\begin{corollary}[推论 C（“等价”的一行读法）]
\label{cor:tex27-ch20-c}
在定理~\ref{thm:tex27-ch20-a10} 与定理~\ref{thm:tex27-ch20-aerg} 的遍历/稳态前提下，对任意 $\epsilon>0$，
\[
  \boxed{\begin{aligned}
    &\pi_\beta(A_\epsilon)\xrightarrow[\beta\to\infty]{}1\\
    &\Longleftrightarrow\ \lim_{\beta\to\infty}\ \lim_{T\to\infty}\ \nu_T(A_\epsilon)=1
  \end{aligned}}.
\]
\end{corollary}

\begin{Certificate}[证书 C（强命题 SP 的直接特化）]
\label{cert:tex27-ch20-c}
由强命题~\ref{prop:tex27-ch20-sp} 的 (V)$\Leftrightarrow$(F) 直接得到。
\end{Certificate}

\begin{proof}[证明（谓词因果链）]
按证书~\ref{cert:tex27-ch20-c} 直接推出。证毕。
\end{proof}

\subsection{备注}
\label{subsec:tex27-ch20-remarks}

\begin{remark}[把“频率最大化”的工程含义说清]
\label{rem:tex27-ch20-r1}
\begin{enumerate}
  \item \textbf{把 $f$ 拉高（$\Delta t$ 变小）本身不创造最优性}：
  它只是让“时间平均 $\approx$ 不变测度”更快显现；真正决定“最终踩在哪”的，是 Gibbs 结构
  $\pi_\beta\propto e^{-\beta\mathcal{E}}$（由自由能变分锁定）。
  \item \textbf{严格门控 $\Delta\Phi<0$ 对应 $\beta\to\infty$ 的退化}：
  该退化更容易卡在局部极小；若要跨越局部陷阱，需要有限温扰动/越障算子（见引理~\ref{lem:tex27-ch8-l1} 的讨论）。
  \item \textbf{“最优路径”在 $\Omega=\mathcal{O}^*$ 上的版本完全同型}：
  把 $\Phi$ 换为 $E(w;s)$（定义~\ref{def:tex27-ch10-d3}），把 $q_\beta$ 换为 $p_\beta(\cdot\mid s)$，
  把遍历/混合证书输入取为 A10（定理~\ref{thm:tex27-ch10-a10}），即可得到本节同型结论。
\end{enumerate}
\end{remark}

\clearpage

%%%%%%%%%%%%%%%%%%%%%%%%%%%%%%%%%%%%%%%%%%%%%%%%%%%%%%%%%%%%
\section{同伦扭曲频率与有限时间逼近变分解：总变差收缩链（形式系统化）}
\label{sec:tex27-ch21-frequency-tv-variational}
%%%%%%%%%%%%%%%%%%%%%%%%%%%%%%%%%%%%%%%%%%%%%%%%%%%%%%%%%%%%

本节把“提高同伦扭曲/采样频率 $f$ 为什么在工程上有效”压缩为一条可审计链：
频率 $f$ 只通过可执行步数 $N(f,\tau)$ 起作用；而在总变差距离下，马尔可夫核具有收缩性，
从而“多做一步更新”可被严格翻译为“离稳态（自由能变分解）更近一步”；再通过零温集中（定理~\ref{thm:tex27-ch8-t2}）
把“稳态更近”转写为“低能邻域命中概率更高”。

\subsection{定义}
\label{subsec:tex27-ch21-defs}

\begin{definition}[D1（同伦扭曲的离散化与“频率”）]
\label{def:tex27-ch21-d1}
沿用定义~\ref{def:tex27-ch8-d2} 的离散更新口径（如 $\theta_{t+\Delta t}=T_{t,\Delta t}(\theta_t)$ 或分布级 $q_{t+\Delta t}=q_tK_{t,\Delta
  t}$）。
定义工程采样周期与频率
\[
  \Delta t>0,
  \qquad
  f:=\frac{1}{\Delta t}.
\]
给定物理/工程时间预算 $\tau>0$，对应可执行的离散步数定义为
\[
  N(f,\tau):=\bigl\lfloor f\tau\bigr\rfloor=\Bigl\lfloor\frac{\tau}{\Delta t}\Bigr\rfloor.
\]
\end{definition}

\begin{definition}[D2（变分目标与“最优邻域”）]
\label{def:tex27-ch21-d2}
沿用定义~\ref{def:tex27-ch8-d3} 与定理~\ref{thm:tex27-ch8-t1}：给定势能 $\Phi:\Theta\to\RR$ 与参考测度 $\mu(d\theta)$，定义 Gibbs 测度
\[
  q_\beta(d\theta)=\frac{1}{Z(\beta)}e^{-\beta\Phi(\theta)}\mu(d\theta),
  \qquad
  q_\beta=\operatorname*{arg\,min}_{q}\ \mathcal{F}_\beta[q].
\]
并沿用定理~\ref{thm:tex27-ch8-t2} 的次水平集（最优邻域）
\[
  A_\epsilon:=\bigl\{\theta\in\Theta:\ \Phi(\theta)\le \Phi_{\min}+\epsilon\bigr\},
  \qquad
  \Phi_{\min}:=\inf_{\theta\in\Theta}\Phi(\theta).
\]
\end{definition}

\begin{definition}[D3（“高频访问/命中”作为时间占用率）]
\label{def:tex27-ch21-d3}
对离散轨道 $(\theta_0,\theta_1,\dots,\theta_N)$，定义集合 $A\subseteq\Theta$ 的经验占用率
\[
  \nu_{N}(A):=\frac{1}{N}\sum_{k=1}^{N}\indicator{\theta_k\in A}.
\]
当 $A=A_\epsilon$ 时，$\nu_N(A_\epsilon)$ 即“高频命中最优邻域”的严格化表达。
\end{definition}

\subsection{公理（降级为定理）}
\label{subsec:tex27-ch21-axioms}

\begin{theorem}[公理（降级为定理）A1（同伦扭曲的 Gibbs 稳态与遍历收敛）]
\label{thm:tex27-ch21-a1}
若同伦扭曲包含（i）Langevin 型扩散噪声或（ii）满足详细平衡的 Metropolis 接受机制，
则 $q_\beta$ 为不变测度；并在适当遍历/耗散条件下，对任意初始分布 $q_0$，分布迭代收敛到稳态：
\[
  q_t\Rightarrow q_\beta
  \quad
  (\text{等价地：离散实现 }q^{(n)}:=q_0K^n\Rightarrow q_\beta).
\]
若需要在总变差距离下给出“有限步数 $\Rightarrow$ 有限误差”的可检验界，则可追加输入
$\|q_0K^n-q_\beta\|_{\mathrm{TV}}\to 0$（或混合时间上界）作为外部证书。
\end{theorem}

\begin{Certificate}[证书 A1（证书 L2 + 外部混合证书输入）]
\label{cert:tex27-ch21-a1}
\begin{itemize}
  \item（稳态）引理~\ref{lem:tex27-ch8-l2} 与其证书~\ref{cert:tex27-ch8-l2} 给出：在 Langevin/Fokker--Planck 或 Metropolis/详细平衡口径下，
    $q_\beta$ 为不变测度；
  \item（遍历收敛）引理~\ref{lem:tex27-ch8-l2} 的第 2 步给出“适当遍历/耗散条件下 $q_t\Rightarrow q_\beta$”；
  \item（若需 TV 口径）追加外部混合证书输入：给出总变差混合界或混合时间 $t_{\mathrm{mix}}(\cdot)$ 上界。
\end{itemize}
\end{Certificate}

\begin{proof}[证明（证书化）]
按证书~\ref{cert:tex27-ch21-a1} 调用即可。证毕。
\end{proof}

\begin{theorem}[公理（降级为定理）A2（马尔可夫核在总变差距离下的收缩性）]
\label{thm:tex27-ch21-a2}
对任意马尔可夫核 $K$ 与任意两概率测度 $\mu,\nu$，有
\[
  \|\mu K-\nu K\|_{\mathrm{TV}}\le \|\mu-\nu\|_{\mathrm{TV}}.
\]
\end{theorem}

\begin{Certificate}[证书 A2（总变差对偶表述 + $\|Kg\|_\infty\le \|g\|_\infty$）]
\label{cert:tex27-ch21-a2}
证书由两条标准恒等式/不等式组成：
\begin{enumerate}
  \item 总变差范数的对偶表述
  $
    \|\mu-\nu\|_{\mathrm{TV}}
    =\sup_{\|g\|_\infty\le 1}\bigl|\int g\,d\mu-\int g\,d\nu\bigr|
  $；
  \item 对任意有界可测 $g$，记 $(Kg)(x):=\int g(y)\,K(x,dy)$，则 $\|Kg\|_\infty\le \|g\|_\infty$。
\end{enumerate}
\end{Certificate}

\begin{proof}[证明（谓词因果链）]
\begin{enumerate}
  \item（对偶化）由证书~\ref{cert:tex27-ch21-a2}(1)，
  \[
    \|\mu K-\nu K\|_{\mathrm{TV}}
    =\sup_{\|g\|_\infty\le 1}\bigl|(\mu K)(g)-(\nu K)(g)\bigr| ;
  \]
  \item（把核作用移到测试函数）由 $(\mu K)(g)=\mu(Kg)$ 得
  \[
    \bigl|(\mu K)(g)-(\nu K)(g)\bigr|=\bigl|\mu(Kg)-\nu(Kg)\bigr|;
  \]
  \item（有界性收缩）由证书~\ref{cert:tex27-ch21-a2}(2)，$\|Kg\|_\infty\le\|g\|_\infty\le 1$；
  \item（取上确界）因此
  \[
    \sup_{\|g\|_\infty\le 1}\bigl|\mu(Kg)-\nu(Kg)\bigr|
    \le \sup_{\|h\|_\infty\le 1}\bigl|\mu(h)-\nu(h)\bigr|
    =\|\mu-\nu\|_{\mathrm{TV}}.
  \]
  证毕。
\end{enumerate}
\end{proof}

\subsection{定理（把“频率”确认为“能量的投影”）}
\label{subsec:tex27-ch21-theorem}

\begin{theorem}[定理 $T_{\mathrm{proj}}$（能量 $\to$ 测度 $\to$ 时间频率：频率是能量的投影）]
\label{thm:tex27-ch21-tproj}
在定理~\ref{thm:tex27-ch21-a1} 的条件下，稳态 $q_\beta$ 由能量地形 $\Phi$ 决定（Gibbs 权重/自由能变分极小点）。
并且在遍历意义下，对任意可测集 $A\subseteq\Theta$，长期占用率满足
\[
  \nu_{N}(A)\xrightarrow[N\to\infty]{} q_\beta(A).
\]
因此对任意 $A$，
\[
  \Phi\ \mapsto\ q_\beta\ \mapsto\ q_\beta(A)=\lim_{N\to\infty}\nu_{N}(A)
\]
给出“时间频率（占用率）是能量地形的可观测投影”。
\end{theorem}

\begin{Certificate}[证书 $T_{\mathrm{proj}}$（自由能变分 + 遍历定理外部证书输入）]
\label{cert:tex27-ch21-tproj}
\begin{enumerate}
  \item（能量定测度）定理~\ref{thm:tex27-ch8-t1}（自由能变分）把 $\Phi$ 映射为唯一的 $q_\beta$；
  \item（测度定频率）定理~\ref{thm:tex27-ch20-aerg}（遍历定理外部证书输入）给出时间平均等于稳态期望。
\end{enumerate}
\end{Certificate}

\begin{proof}[证明（谓词因果链）]
\begin{enumerate}
  \item（能量 $\Rightarrow$ 测度）由定理~\ref{thm:tex27-ch8-t1}，在给定 $\beta$ 下，$q_\beta$ 由 $\Phi$ 唯一确定；
  \item（测度 $\Rightarrow$ 频率）在遍历前提下，由定理~\ref{thm:tex27-ch20-aerg} 得对任意 $A$，
  $\nu_N(A)\to q_\beta(A)$；
  \item（合成映射）合并 1--2 得到陈述中的映射链。证毕。
\end{enumerate}
\end{proof}

\subsection{引理（提高频率 $f$ 为何有效）}
\label{subsec:tex27-ch21-lemmas}

\begin{lemma}[L1（固定物理时间 $\tau$ 下，频率 $f$ 只通过步数 $N(f,\tau)$ 起作用）]
\label{lem:tex27-ch21-l1}
对任意离散实现，物理时间 $\tau$ 内的轨道片段只依赖于在该时间窗内执行的步数 $N(f,\tau)$。
\end{lemma}

\begin{Certificate}[证书 L1（定义 D1 的直接展开）]
\label{cert:tex27-ch21-l1}
由定义~\ref{def:tex27-ch21-d1}，$\tau$ 内可执行步数被 $N(f,\tau)=\lfloor f\tau\rfloor$ 唯一确定。
\end{Certificate}

\begin{proof}[证明（谓词因果链）]
按证书~\ref{cert:tex27-ch21-l1}，$\tau$ 内只会执行 $N(f,\tau)$ 次离散更新；因此轨道片段由步数确定。证毕。
\end{proof}

\begin{lemma}[L2（步数越多，离稳态越近；从而频率越高，固定 $\tau$ 下越接近变分解）]
\label{lem:tex27-ch21-l2}
设存在时间齐次核 $K$ 使 $q_\beta K=q_\beta$（不变性），并令 $q^{(n)}:=q_0K^n$。
则对任意 $m\ge 0$，有单调性
\[
  \|q^{(n+m)}-q_\beta\|_{\mathrm{TV}}\le \|q^{(n)}-q_\beta\|_{\mathrm{TV}}.
\]
从而对固定 $\tau$，若 $f_2\ge f_1$（即 $N_2:=N(f_2,\tau)\ge N_1:=N(f_1,\tau)$），则
\[
  \|q^{(N_2)}-q_\beta\|_{\mathrm{TV}}\le \|q^{(N_1)}-q_\beta\|_{\mathrm{TV}}.
\]
\end{lemma}

\begin{Certificate}[证书 L2（A2 收缩性 + 不变性）]
\label{cert:tex27-ch21-l2}
证书由两点组成：
\begin{enumerate}
  \item 总变差收缩：定理~\ref{thm:tex27-ch21-a2}；
  \item 稳态不变：$q_\beta K^m=q_\beta$（由 $q_\beta K=q_\beta$ 迭代得到）。
\end{enumerate}
\end{Certificate}

\begin{proof}[证明（谓词因果链）]
\begin{enumerate}
  \item（把“多走 $m$ 步”写成核作用）由定义 $q^{(n+m)}=q^{(n)}K^m$；
  \item（稳态也走 $m$ 步不变）由不变性 $q_\beta K^m=q_\beta$；
  \item（差距收缩）由定理~\ref{thm:tex27-ch21-a2}，
  \[
    \|q^{(n)}K^m-q_\beta K^m\|_{\mathrm{TV}}\le \|q^{(n)}-q_\beta\|_{\mathrm{TV}} ;
  \]
  \item（代入 1--2）得到 $\|q^{(n+m)}-q_\beta\|_{\mathrm{TV}}\le \|q^{(n)}-q_\beta\|_{\mathrm{TV}}$；
  \item（频率版本）若 $N_2\ge N_1$，取 $n=N_1,m=N_2-N_1$ 即得第二个不等式。证毕。
\end{enumerate}
\end{proof}

\begin{lemma}[L3（离 $q_\beta$ 更近 $\Rightarrow$ 低能集合概率更接近）]
\label{lem:tex27-ch21-l3}
对任意可测集 $A$，若 $\|q-q_\beta\|_{\mathrm{TV}}\le \delta$，则
\[
  q(A)\ge q_\beta(A)-\delta.
\]
\end{lemma}

\begin{Certificate}[证书 L3（总变差距离的集合刻画）]
\label{cert:tex27-ch21-l3}
总变差的等价刻画：$\|q-q_\beta\|_{\mathrm{TV}}=\sup_A |q(A)-q_\beta(A)|$。
\end{Certificate}

\begin{proof}[证明（谓词因果链）]
由证书~\ref{cert:tex27-ch21-l3}，对任意 $A$ 有 $|q(A)-q_\beta(A)|\le \delta$，
即 $q(A)\ge q_\beta(A)-\delta$。证毕。
\end{proof}

\subsection{命题（进一步严格证明：频率 $\Rightarrow$ 有限时间逼近变分解）}
\label{subsec:tex27-ch21-props}

\begin{proposition}[命题 $P_{\mathrm{freq\to var}}$（提高同伦扭曲频率是逼近变分解的有效手段）]
\label{prop:tex27-ch21-freq-to-var}
在定理~\ref{thm:tex27-ch21-a1} 成立（稳态为 $q_\beta$ 且存在遍历收敛）并采用时间齐次核近似
（或在每个小时间窗内近似齐次）的工程实现下：
对任意物理时间预算 $\tau>0$、任意精度 $\delta>0$，存在阈值频率 $f_\delta(\tau)$ 使得当 $f\ge f_\delta(\tau)$ 时，
\[
  \bigl\|q^{(N(f,\tau))}-q_\beta\bigr\|_{\mathrm{TV}}\le \delta.
\]
进一步，对任意 $\epsilon>0$，
\[
  q^{(N(f,\tau))}(A_\epsilon)\ \ge\ q_\beta(A_\epsilon)-\delta.
\]
再结合定理~\ref{thm:tex27-ch8-t2}（$\beta\to\infty\Rightarrow q_\beta(A_\epsilon)\to 1$），得到：当 $\beta$ 足够大且 $f$ 足够大时，
\[
  q^{(N(f,\tau))}(A_\epsilon)\approx 1,
\]
即在有限工程时间 $\tau$ 内，“高频同伦扭曲/高频采样”把系统状态有效推向能量最优邻域，从而逼近变分意义下的最优。
\end{proposition}

\begin{Certificate}[证书 $P_{\mathrm{freq\to var}}$（A1 + L2 + L3 + T2）]
\label{cert:tex27-ch21-freq-to-var}
\begin{itemize}
  \item 稳态与收敛：定理~\ref{thm:tex27-ch21-a1}（并在 TV 口径下输入混合证书）；
  \item 步数单调改进：引理~\ref{lem:tex27-ch21-l2}；
  \item TV 接近 $\Rightarrow$ 概率接近：引理~\ref{lem:tex27-ch21-l3}；
  \item 零温集中到 $A_\epsilon$：定理~\ref{thm:tex27-ch8-t2}。
\end{itemize}
\end{Certificate}

\begin{proof}[证明（谓词因果链）]
\begin{enumerate}
  \item（存在足够步数）若在 TV 口径下有 $\|q_0K^n-q_\beta\|_{\mathrm{TV}}\to 0$，
  则对给定 $(\delta,\tau)$ 存在 $N_\delta$ 使
  $\|q^{(N_\delta)}-q_\beta\|_{\mathrm{TV}}\le \delta$；
  \item（用频率兑现步数）令 $f_\delta(\tau):=(N_\delta+1)/\tau$。则当 $f\ge f_\delta(\tau)$ 时，
  有 $N(f,\tau)=\lfloor f\tau\rfloor\ge N_\delta$；
  \item（步数单调改进）由引理~\ref{lem:tex27-ch21-l2}，得到
  \[
    \|q^{(N(f,\tau))}-q_\beta\|_{\mathrm{TV}}
    \le \|q^{(N_\delta)}-q_\beta\|_{\mathrm{TV}}
    \le \delta;
  \]
  \item（转为低能命中）由引理~\ref{lem:tex27-ch21-l3}，对 $A_\epsilon$ 有
  $
    q^{(N(f,\tau))}(A_\epsilon)\ge q_\beta(A_\epsilon)-\delta
  $；
  \item（零温极限对齐最优）由定理~\ref{thm:tex27-ch8-t2}，取 $\beta$ 足够大使 $q_\beta(A_\epsilon)\ge 1-\delta$，则
  \[
    q^{(N(f,\tau))}(A_\epsilon)\ge 1-2\delta.
  \]
  证毕。
\end{enumerate}
\end{proof}

\subsection{推论}
\label{subsec:tex27-ch21-corollaries}

\begin{corollary}[C1（频率是“把静态变分极小点变成动态稳态”的工程旋钮）]
\label{cor:tex27-ch21-c1}
定理~\ref{thm:tex27-ch8-t1} 给出 $q_\beta$ 是自由能变分极小点（静态最优）；
命题~\ref{prop:tex27-ch21-freq-to-var} 给出：在固定 $\tau$ 内，提高 $f$ 可保证分布更快贴近 $q_\beta$（动态稳态）。
因此在该证书体系下，“静态变分最优”可被无损转译为“动态闭环稳态”。
\end{corollary}

\begin{Certificate}[证书 C1（T1 + 命题 $P_{\mathrm{freq\to var}}$ 并列引用）]
\label{cert:tex27-ch21-c1}
静态最优由定理~\ref{thm:tex27-ch8-t1}；动态逼近由命题~\ref{prop:tex27-ch21-freq-to-var}。
\end{Certificate}

\begin{proof}[证明（谓词因果链）]
按证书~\ref{cert:tex27-ch21-c1} 并列引用即可得到推论陈述。证毕。
\end{proof}

\begin{corollary}[C2（同伦扭曲频率与“微步分解逼近生成元”的一致性）]
\label{cor:tex27-ch21-c2}
命题~\ref{prop:tex27-ch8-p2}（或定理~\ref{thm:tex27-ch10-t3}）给出：当把非交换边缘算子写成微步指数
$\exp(-\Delta\tau\,\widehat{G}_i)$ 并令 $n\to\infty$（即 $\Delta\tau=\tau/n\to 0$）时，
Trotter 乘积逼近 $e^{-\tau\widehat{H}}$。
在工程语义上，这等价于：\textbf{提高同伦扭曲的频率（$f\sim n/\tau\to\infty$）$\Rightarrow$ 更精确逼近“连续时间生成元”对应的变分演化}
（经典情形对应 Langevin/梯度流；量子口径对应虚时间演化）。
\end{corollary}

\begin{Certificate}[证书 C2（Trotter 极限证书输入）]
\label{cert:tex27-ch21-c2}
直接引用命题~\ref{prop:tex27-ch8-p2} 的证书~\ref{cert:tex27-ch8-p2}（Lie--Trotter 乘积公式）。
\end{Certificate}

\begin{proof}[证明（谓词因果链）]
\begin{enumerate}
  \item（微步 $\Rightarrow$ 高频）由 $\Delta\tau=\tau/n$，令 $f:=n/\tau$，则 $n\to\infty\iff f\to\infty$；
  \item（高频 $\Rightarrow$ 连续极限）由证书~\ref{cert:tex27-ch21-c2}，当 $n\to\infty$ 时，离散乘积极限收敛到 $e^{-\tau\widehat{H}}$；
  \item（解释口径）因此“提高频率”可被视为“以更细微步逼近生成元演化”的工程实现。证毕。
\end{enumerate}
\end{proof}

\subsection{备注（关键工程细节与接口对齐）}
\label{subsec:tex27-ch21-remarks}

\begin{remark}[如何“提高频率而不改变问题本身”】【工程口径】]
\label{rem:tex27-ch21-r1}
\begin{enumerate}
  \item \textbf{提高 $f$ 的有效性有前提：需保持同一变分目标与同一温度口径。}
  例如对定义~\ref{def:tex27-ch8-d2} 的 Langevin 形式，若用 Euler--Maruyama 离散实现
  \[
    \theta_{k+1}
    =\theta_k-\nabla\Phi(\theta_k)\Delta t+b(\theta_k)\Delta t+\sqrt{\frac{2}{\beta}\Delta t}\,\xi_k,
    \qquad
    \xi_k\sim\mathcal{N}(0,I),
  \]
  则当 $\Delta t\downarrow 0$（$f\uparrow\infty$）时，单位物理时间内的漂移/扩散强度保持不变；
  若噪声不按 $\sqrt{\Delta t}$ 缩放，则改变了有效温度（等价于更换了变分问题）。
  \item \textbf{关于“逃逸等待时间”}：若把“逃出局部盆地所需的步数”记为随机变量 $H$，
  则对应工程时间 $T_{\mathrm{esc}}:=H\Delta t=H/f$；在“每一步机制不变、仅提高单位时间可执行步数”的口径下，
  \[
    \mathbb{E}[T_{\mathrm{esc}}]=\frac{\mathbb{E}[H]}{f},
  \]
  即提高 $f$ 可线性缩短等待时间（解释“高频对克服局部等待至关重要”）。
  \item \textbf{与 ODTHB（温控/调度）对齐}：提高 $f$ 提供“步数供给”，使遍历性在有限 $\tau$ 内可显现；
  而第~\ref{subsec:tex27-ch5-odthb}~节的 ODTHB 通过宏观参数（如 $\beta,\ell$ 等）调节混合与接受率，
  避免高频步进退化为低效抖动或被硬剪裁破坏隧穿（参见命题~\ref{prop:tex27-ch5-odthb-r1} 与定理~\ref{thm:tex27-ch5-odthb-microtunnel-no-hardprune}）。
  二者正交：$f$ 负责“步数”，温控负责“每一步是否有效地服务于 $q_\beta$”。
\end{enumerate}
\end{remark}

\clearpage

%%%%%%%%%%%%%%%%%%%%%%%%%%%%%%%%%%%%%%%%%%%%%%%%%%%%%%%%%%%%
\section{接口型创新判据：从静态变分到动态频率审计与阈值设计（形式系统化）}
\label{sec:tex27-ch22-interface-innovation}
%%%%%%%%%%%%%%%%%%%%%%%%%%%%%%%%%%%%%%%%%%%%%%%%%%%%%%%%%%%%

本节不把“开创性/创新性”写成口号，而把它压缩为\textbf{可核验的接口型判据}：
在既有基底理论（自由能变分 + 零温极限 + 路径空间变分选择 + 时间换空间 + 外部混合证书输入）的符号体系上，
引入可观测/可控的工程符号（$\tau,f,\nu_N$），并给出三类可审计产物：
（I）能量 $\Phi$ 到频率 $\nu$ 的可观测投影；
（II）能量差与概率/频率比之间的可逆（或可估计）公式；
（III）有限时间预算 $\tau$ 下的频率阈值设计证书。
注意：本节仅证明“在本文形式系统语境内形成了新的接口层”，并不宣称“世界首创”（该外推需要全域文献证书，属未证明）。

\subsection{定义}
\label{subsec:tex27-ch22-defs}

\begin{definition}[D0（“开创性/创新性”的可证书化判据：接口型创新）]
\label{def:tex27-ch22-d0}
在本文形式系统中，把“创新性”定义为一类\textbf{接口型断言}而非历史性断言。令
\begin{itemize}
  \item 基底理论 $\mathsf{Base}$ 定义为如下条目的合取：
  \begin{enumerate}
    \item 自由能变分与零温极限：定理~\ref{thm:tex27-ch8-t1}--\ref{thm:tex27-ch8-t2}；
    \item 路径空间变分选择与时间换空间：定理~\ref{thm:tex27-ch10-t1}--\ref{thm:tex27-ch10-t2}；
    \item 外部混合/遍历证书输入：公理（降级为定理）A10（定理~\ref{thm:tex27-ch10-a10}）。
  \end{enumerate}
  \item 扩张语言 $\mathsf{Ext}$ 定义为在 $\mathsf{Base}$ 的符号系统上新增三类可观测/可控符号：
  \[
    \tau>0,\qquad f=1/\Delta t,\qquad \nu_{N}(A)\ (N=\lfloor f\tau\rfloor).
  \]
\end{itemize}
称一个强命题 $\mathsf{SP}$ 为“接口型创新”，若在 $\mathsf{Base}\wedge\mathsf{Ext}$ 下能证出下列三条可审计接口：
\[
  \boxed{
  \begin{aligned}
    &\text{(I) 存在可观测投影 } \Pi:\Phi\mapsto \nu;\\
    &\text{(II) 给出可逆/可估计的能量差公式；}\\
    &\text{(III) 给出有限 }\tau\text{ 下的工程阈值设计证书。}
  \end{aligned}
  }
\]
其中 (I)(II)(III) 必须都能落到“证书 + 证明（谓词因果链）”格式。
\end{definition}

\begin{definition}[D1（能量--测度--频率的三段映射）]
\label{def:tex27-ch22-d1}
在参数口径下，给定缺陷势能 $\Phi$ 与先验测度 $\mu$，定义 Gibbs 测度（定义~\ref{def:tex27-ch8-d3}）
\[
  q_\beta(d\theta)=\frac{1}{Z(\beta)}e^{-\beta\Phi(\theta)}\mu(d\theta).
\]
在公理 A10（定理~\ref{thm:tex27-ch10-a10}）给定的采样核下，令 $(\theta_k)$（或路径样本 $(w_k)$）为采样轨道。
对任意可测集 $A$，定义占用率（定义~\ref{def:tex27-ch21-d3}）
\[
  \nu_N(A):=\frac{1}{N}\sum_{k=1}^{N}\indicator{\theta_k\in A},
  \qquad
  N=\lfloor f\tau\rfloor.
\]
\end{definition}

\subsection{公理（降级为定理）}
\label{subsec:tex27-ch22-axioms}

\begin{theorem}[公理（降级为定理）$A_{\mathrm{erg}}$（遍历定理：时间平均 = 稳态期望；复用）]
\label{thm:tex27-ch22-aerg}
若采样核满足公理 A10（定理~\ref{thm:tex27-ch10-a10}）所要求的遍历/混合证书输入，
且不变分布为 $q_\beta$，则对任意可测集 $A$，
\[
  \nu_N(A)\xrightarrow[N\to\infty]{} q_\beta(A)
  \qquad
  (\text{几乎必然或按外部证书给定的收敛模式}).
\]
本条与定理~\ref{thm:tex27-ch20-aerg} 为同一接口声明（此处为章内自洽重述）。
\end{theorem}

\begin{Certificate}[证书 $A_{\mathrm{erg}}$（A10 的遍历/混合证书输入 $\Rightarrow$ 遍历定理调用）]
\label{cert:tex27-ch22-aerg}
把定理~\ref{thm:tex27-ch10-a10} 的“不可约/无周期/漂移小集 +（可选）混合界”作为外部证书输入，
调用马尔可夫链遍历定理即可得到
$\nu_N(A)\to q_\beta(A)$。
\end{Certificate}

\begin{proof}[证明（谓词因果链）]
\begin{enumerate}
  \item（稳态存在且唯一）由 A10 的外部证书输入（定理~\ref{thm:tex27-ch10-a10}）给出：$q_\beta$ 为不变分布且在遍历条件下唯一；
  \item（时间平均收敛）由证书~\ref{cert:tex27-ch22-aerg} 调用遍历定理，得到对任意 $A$ 有 $\nu_N(A)\to q_\beta(A)$。
  证毕。
\end{enumerate}
\end{proof}

\subsection{引理（把“频率是能量的投影”写成可计算公式）}
\label{subsec:tex27-ch22-lemmas}

\begin{lemma}[L1（能量差 $\Leftrightarrow$ 概率比：频率可逆地携带能量信息）]
\label{lem:tex27-ch22-l1}
在离散情形（$\Theta$ 可数且 $\mu(\{\theta\})>0$）下，对任意 $\theta_1,\theta_2\in\Theta$ 有精确比值公式
\[
  \frac{q_\beta(\{\theta_1\})}{q_\beta(\{\theta_2\})}
  =
  e^{-\beta(\Phi(\theta_1)-\Phi(\theta_2))}\cdot
  \frac{\mu(\{\theta_1\})}{\mu(\{\theta_2\})}.
\]
若 $\mu(\{\theta_1\})\approx\mu(\{\theta_2\})$（例如均匀先验），则
\[
  \Phi(\theta_1)-\Phi(\theta_2)
  =
  -\frac{1}{\beta}\log\frac{q_\beta(\{\theta_1\})}{q_\beta(\{\theta_2\})}
  +O\!\left(\frac{1}{\beta}\log\frac{\mu(\{\theta_1\})}{\mu(\{\theta_2\})}\right).
\]
\end{lemma}

\begin{Certificate}[证书 L1（直接代入 Gibbs 定义并约去 $Z(\beta)$）]
\label{cert:tex27-ch22-l1}
由 Gibbs 定义 $q_\beta(\{\theta\})=\frac{1}{Z(\beta)}e^{-\beta\Phi(\theta)}\mu(\{\theta\})$，
两点比值中 $Z(\beta)$ 抵消，指数项给出能量差，$\mu$ 给出先验修正项。
\end{Certificate}

\begin{proof}[证明（谓词因果链）]
\begin{enumerate}
  \item（代入）由定义~\ref{def:tex27-ch8-d3}，
  $
    q_\beta(\{\theta_i\})=\frac{1}{Z(\beta)}e^{-\beta\Phi(\theta_i)}\mu(\{\theta_i\})
  $；
  \item（取比值）两式相除，$Z(\beta)$ 抵消，得到比值公式；
  \item（取对数）对比值取 $\log$ 并移项得到能量差表达式。证毕。
\end{enumerate}
\end{proof}

\begin{lemma}[L2（频率 $\nu$ 逼近 $q_\beta$：能量差可由“高频踩点”估计）]
\label{lem:tex27-ch22-l2}
在定理~\ref{thm:tex27-ch22-aerg} 条件下，对任意可测集 $A$ 有 $\nu_N(A)\to q_\beta(A)$。
特别地，若在离散点集上取 $A_i=\{\theta_i\}$ 且 $q_\beta(A_i)>0$，则
\[
  \frac{\nu_N(A_1)}{\nu_N(A_2)}\xrightarrow[N\to\infty]{}\frac{q_\beta(A_1)}{q_\beta(A_2)}.
\]
从而由引理~\ref{lem:tex27-ch22-l1} 得到能量差的频率估计式：
\[
  \Phi(\theta_1)-\Phi(\theta_2)
  \approx
  -\frac{1}{\beta}\log\frac{\nu_N(\{\theta_1\})}{\nu_N(\{\theta_2\})}
  +\frac{1}{\beta}\log\frac{\mu(\{\theta_1\})}{\mu(\{\theta_2\})},
\]
且当 $N\to\infty$ 时该近似在外部证书给定意义下收敛为等号。
\end{lemma}

\begin{Certificate}[证书 L2（$A_{\mathrm{erg}}$ + L1 组合）]
\label{cert:tex27-ch22-l2}
\begin{itemize}
  \item（频率逼近稳态）由定理~\ref{thm:tex27-ch22-aerg}，$\nu_N(A_i)\to q_\beta(A_i)$；
  \item（能量差公式）由引理~\ref{lem:tex27-ch22-l1}，$\Phi(\theta_1)-\Phi(\theta_2)$ 由 $q_\beta$ 的比值给出。
\end{itemize}
\end{Certificate}

\begin{proof}[证明（谓词因果链）]
\begin{enumerate}
  \item（频率 $\Rightarrow$ 概率）由定理~\ref{thm:tex27-ch22-aerg}，对 $A_1,A_2$ 有 $\nu_N(A_i)\to q_\beta(A_i)$；
  \item（比值收敛）若 $q_\beta(A_2)>0$，则当 $N$ 足够大时 $\nu_N(A_2)>0$，从而
  $\nu_N(A_1)/\nu_N(A_2)\to q_\beta(A_1)/q_\beta(A_2)$；
  \item（代入能量差）把极限比值代入引理~\ref{lem:tex27-ch22-l1} 的对数公式得到频率估计式。证毕。
\end{enumerate}
\end{proof}

\subsection{命题（工程层的创新点：把“变分最优”转成“频率阈值设计”）}
\label{subsec:tex27-ch22-props}

\begin{proposition}[P1（频率 $f$ 的工程价值：把有限 $\tau$ 兑现成足够大的 $N$）]
\label{prop:tex27-ch22-p1}
在固定物理时间预算 $\tau$ 下，采样步数满足 $N=\lfloor f\tau\rfloor$。
因此“提高同伦扭曲频率 $f$”在形式系统里等价于“在同一 $\tau$ 内增加时间迭代次数”，从而更快进入
$A_{\mathrm{erg}}$ 的有效区间，使 $\nu_N$ 更接近 $q_\beta$。
\end{proposition}

\begin{Certificate}[证书 P1（步数恒等式 + 时间换空间 + 遍历定理）]
\label{cert:tex27-ch22-p1}
\begin{itemize}
  \item（步数兑现）$N=\lfloor f\tau\rfloor$（定义~\ref{def:tex27-ch21-d1}）；
  \item（时间换空间）定理~\ref{thm:tex27-ch10-t1}（用时间迭代逼近 Gibbs/零温极限）；
  \item（频率收敛）定理~\ref{thm:tex27-ch22-aerg}（$\nu_N\to q_\beta$）。
\end{itemize}
\end{Certificate}

\begin{proof}[证明（谓词因果链）]
\begin{enumerate}
  \item（$f\uparrow\Rightarrow N\uparrow$）由 $N=\lfloor f\tau\rfloor$ 直接推出；
  \item（时间迭代逼近稳态）由定理~\ref{thm:tex27-ch10-t1}，时间迭代是逼近目标分布/零温极限的机制；
  \item（占用率逼近稳态质量）由定理~\ref{thm:tex27-ch22-aerg}，当 $N$ 足够大时 $\nu_N(A)\approx q_\beta(A)$；
  \item（合并）因此提高 $f$ 在有限 $\tau$ 内通过增大 $N$ 使 $\nu_N$ 更快逼近 $q_\beta$。证毕。
\end{enumerate}
\end{proof}

\begin{proposition}[P2（逃逸等待时间：高 $f$ 缩短物理等待而不改 $\Phi$）]
\label{prop:tex27-ch22-p2}
令 $H$ 表示“从局部盆地首次逃逸所需的步数”（离散随机变量，受采样核与宏观参数如 $\beta,\theta,\ell,\dots$ 影响）。
记物理等待时间为
\[
  T_{\mathrm{esc}}:=H\Delta t=\frac{H}{f}.
\]
在“核不失真缩放”（即仅细化步长而不改变有效温度/接受机制）的工程实现约束下，有
\[
  \mathbb{E}[T_{\mathrm{esc}}]=\frac{\mathbb{E}[H]}{f}.
\]
因此 $f$ 不改变势能面 $\Phi$，但线性缩短以物理时间计的逃逸等待。
\end{proposition}

\begin{Certificate}[证书 P2（代数恒等式 + 工程不失真缩放约束）]
\label{cert:tex27-ch22-p2}
\begin{enumerate}
  \item 代数恒等式：$T_{\mathrm{esc}}=H\Delta t$ 与 $f=1/\Delta t$；
  \item 工程约束：细化 $\Delta t$ 时噪声/接受机制按正确缩放保持同一目标 $q_\beta$（例如 Langevin 离散噪声强度按 $\sqrt{\Delta t}$ 缩放；见备注~\ref{rem:tex27-ch21-r1}）。
\end{enumerate}
\end{Certificate}

\begin{proof}[证明（谓词因果链）]
\begin{enumerate}
  \item（物理时间换算）由 $T_{\mathrm{esc}}=H\Delta t$ 与 $f=1/\Delta t$ 得 $T_{\mathrm{esc}}=H/f$；
  \item（取期望）两端取期望得到 $\mathbb{E}[T_{\mathrm{esc}}]=\mathbb{E}[H]/f$。证毕。
\end{enumerate}
\end{proof}

\subsection{定理（证明“接口型创新”：按 D0 三条判据逐一给证书）}
\label{subsec:tex27-ch22-theorem}

\begin{theorem}[定理 $T_{\mathrm{innov}}$（强命题满足接口型创新三条件）]
\label{thm:tex27-ch22-tinnov}
令强命题 $\mathsf{SP}$ 表示如下合取：
\[
  \mathsf{SP}:=
  \bigl(\text{能量--测度--频率投影（定理~\ref{thm:tex27-ch21-tproj}）}\bigr)
  \ \wedge\
  \bigl(\text{有限 }\tau\text{ 下频率逼近变分解（命题~\ref{prop:tex27-ch21-freq-to-var}）}\bigr).
\]
则在 $\mathsf{Base}\wedge\mathsf{Ext}$ 下，$\mathsf{SP}$ 满足定义~\ref{def:tex27-ch22-d0} 的 (I)(II)(III)，
因此在本文形式系统语境内，$\mathsf{SP}$ 是\textbf{可证书化的接口型创新}。
\end{theorem}

\begin{Certificate}[证书 $T_{\mathrm{innov}}$（(I)(II)(III) 三段证书拼接）]
\label{cert:tex27-ch22-tinnov}
\begin{itemize}
  \item (I) 可观测投影 $\Pi:\Phi\mapsto \nu$：定理~\ref{thm:tex27-ch21-tproj}（能量 $\to q_\beta\to$ 占用率极限）；
  \item (II) 可逆/可估计能量差公式：引理~\ref{lem:tex27-ch22-l1}（能量差 $\leftrightarrow$ 概率比） + 引理~\ref{lem:tex27-ch22-l2}（概率比由频率比逼近）；
  \item (III) 有限 $\tau$ 的阈值设计证书：命题~\ref{prop:tex27-ch22-p1}（$f\uparrow\Rightarrow N\uparrow$ 兑现步数） + 命题~\ref{prop:tex27-ch22-p2}（$T_{\mathrm{esc}}$ 的线性缩放） +
  时间换空间（定理~\ref{thm:tex27-ch10-t1}）与外部混合证书输入（定理~\ref{thm:tex27-ch10-a10}）。
\end{itemize}
\end{Certificate}

\begin{proof}[证明（谓词因果链）]
\textbf{证明 (I)}：
\begin{enumerate}
  \item（能量定测度）由定理~\ref{thm:tex27-ch8-t1}，$\Phi$ 唯一确定 $q_\beta$；
  \item（测度定频率）由定理~\ref{thm:tex27-ch22-aerg}，$\nu_N(A)\to q_\beta(A)$；
  \item（投影定义）合并得出 $\Pi_A(\Phi):=\lim_{N\to\infty}\nu_N(A)=q_\beta(A)$，即能量到频率的可观测投影。证毕。
\end{enumerate}

\textbf{证明 (II)}：
\begin{enumerate}
  \item（能量差 $\leftrightarrow$ 概率比）由引理~\ref{lem:tex27-ch22-l1}，能量差可由 $\log(q_\beta$ 比值$)$ 表达（先验项为已知修正）；
  \item（概率比由频率比估计）由引理~\ref{lem:tex27-ch22-l2}，当 $N$ 足够大时 $q_\beta$ 比值可由频率比 $\nu_N$ 估计；
  \item（得到可审计估计式）故“能量 $\leftrightarrow$ 频率”可形成可复核的对数比值证书。证毕。
\end{enumerate}

\textbf{证明 (III)}：
\begin{enumerate}
  \item（$f$ 兑现步数预算）由命题~\ref{prop:tex27-ch22-p1}，固定 $\tau$ 下 $f$ 是把物理时间预算兑现为迭代次数预算的线性旋钮；
  \item（时间换空间接口）由定理~\ref{thm:tex27-ch10-t1}，时间迭代逼近 Gibbs/零温极限是直接推论；再由 A10（定理~\ref{thm:tex27-ch10-a10}）提供混合/遍历证书输入；
  \item（逃逸等待的物理缩放）由命题~\ref{prop:tex27-ch22-p2}，$f$ 不改 $\Phi$ 却线性缩短 $T_{\mathrm{esc}}$，使有限 $\tau$ 内发生足够多次有效越障尝试更可行。
  证毕。
\end{enumerate}

综上 (I)(II)(III) 全部成立，故 $\mathsf{SP}$ 满足定义~\ref{def:tex27-ch22-d0} 的接口型创新判据。证毕。
\end{proof}

\subsection{推论（把“创新性”落地为可复核的工程新产物）}
\label{subsec:tex27-ch22-corollaries}

\begin{corollary}[C1（新增“白盒审计接口”：仅凭运行轨迹可反推出能量排序证书）]
\label{cor:tex27-ch22-c1}
在 A10 + $A_{\mathrm{erg}}$ 条件下，记录频率比 $\nu_N(B_1)/\nu_N(B_2)$ 并用引理~\ref{lem:tex27-ch22-l2} 的对数公式，
即可给出能量差估计证书，从而把“势能面 $\Phi$”转化为“可观测的高频踩点统计”。
这给出一个新的可审计接口：不必直接访问 $\nabla\Phi$ 或显式枚举 $\Omega$，仍可对白盒一致性做证书化核验。
\end{corollary}

\begin{Certificate}[证书 C1（L2 的直接调用）]
\label{cert:tex27-ch22-c1}
直接调用引理~\ref{lem:tex27-ch22-l2}：用频率比逼近稳态概率比，并由对数比值给出能量差估计。
\end{Certificate}

\begin{proof}[证明（谓词因果链）]
按证书~\ref{cert:tex27-ch22-c1} 调用即可。证毕。
\end{proof}

\begin{corollary}[C2（新增“频率阈值设计口径”：把 ODTHB 温控与同伦扭曲频率解耦）]
\label{cor:tex27-ch22-c2}
文件中 ODTHB 被严格限定为“宏观调参而非微观干预”（第~\ref{subsec:tex27-ch5-odthb}~节）。
结合命题~\ref{prop:tex27-ch22-p1}--\ref{prop:tex27-ch22-p2} 得到解耦分工：
\begin{itemize}
  \item 宏观调参 $\Lambda=(\beta,\theta,\ell,\dots)$ 负责“每一步是否有效混合/是否允许隧穿”（改变核性质与 $H$ 的分布）；
  \item 频率 $f$ 负责“在有限 $\tau$ 内是否执行了足够多步以让遍历性生效”（把 $N$ 拉到证书可用区间，并缩短 $T_{\mathrm{esc}}$ 的物理尺度）。
\end{itemize}
因此“势能/温控设计”与“执行频率设计”可拆成正交模块，且每个模块都可被证书化。
\end{corollary}

\begin{Certificate}[证书 C2（P1 + P2 + ODTHB 接口声明）]
\label{cert:tex27-ch22-c2}
频率分工由命题~\ref{prop:tex27-ch22-p1}（$N=\lfloor f\tau\rfloor$）与命题~\ref{prop:tex27-ch22-p2}（$T_{\mathrm{esc}}=H/f$）给出；
宏观调参分工由第~\ref{subsec:tex27-ch5-odthb}~节对 ODTHB 的接口约束给出。
\end{Certificate}

\begin{proof}[证明（谓词因果链）]
按证书~\ref{cert:tex27-ch22-c2} 将“频率通道”与“宏观调参通道”分别落实到可复核恒等式与既有接口约束即可得到解耦分工。证毕。
\end{proof}

\subsection{备注}
\label{subsec:tex27-ch22-remarks}

\begin{remark}[边界：接口型创新 $\neq$ 文献意义的“世界首创”】【未证明外推】]
\label{rem:tex27-ch22-r1}
定义~\ref{def:tex27-ch22-d0} 的“创新性”是\textbf{形式系统内}的接口判据：
它证明的是“在既有 $\mathsf{Base}$ 上新增了可审计的静态变分 $\to$ 动态控制/审计接口层”。
若要把该结论外推为“全域文献意义的世界首创”，必须额外给出外部文献穷尽/优先级证书；该外推在本文不被证明。
\end{remark}

\clearpage

%%%%%%%%%%%%%%%%%%%%%%%%%%%%%%%%%%%%%%%%%%%%%%%%%%%%%%%%%%%%
\section{评价：描述性到规范性的接口化与三层创新维度}
\label{sec:tex27-ch23-evaluation-original}
%%%%%%%%%%%%%%%%%%%%%%%%%%%%%%%%%%%%%%%%%%%%%%%%%%%%%%%%%%%%

\subsection{评价}
\label{subsec:tex27-ch23-evaluation}

\begin{remark}[评价 E0（范式转移的严格化表述）]
\label{rem:tex27-ch23-e0}
令 $\Phi$ 为缺陷势能（或 $\JacGap$ 的能量化版本），令 $q_\beta\propto e^{-\beta\Phi}\mu$ 为自由能变分极小点（Gibbs 目标），令 $\{\theta_k\}_{k\ge 0}$
  为满足遍历/混合证书输入的采样轨道。定义占用率（“踩点频率”）
\[
  \nu_N(A):=\frac{1}{N}\sum_{k=1}^{N}\indicator{\theta_k\in A}.
\]
则“描述性（Descriptive）”与“规范性（Prescriptive）”可严格区分为：
\begin{itemize}
  \item \textbf{描述性口径}：从 $\Phi$ 推出“稳定性/出现频率”
  \[
    \Phi \Rightarrow q_\beta \Rightarrow q_\beta(A)\ \bigl(\approx \nu_N(A)\ \text{当 }N\text{ 足够大}\bigr).
  \]
  \item \textbf{规范性口径}：通过可控执行参数（尤其是频率 $f$ 与步数 $N$）在有限时间内制造“可观测的稳定性/出现频率”
  \[
    (f,\tau)\Rightarrow N=\lfloor f\tau\rfloor\Rightarrow \nu_N(\cdot)\approx q_\beta(\cdot)\Rightarrow \text{对低能集合的高占用率}
      .
  \]
\end{itemize}
\textbf{评价结论（E0 的语言）}：该架构的关键不是把统计力学、遍历理论、控制论并置，而是建立了一个可审计接口
\[
  \boxed{\begin{aligned}
    &\Phi\ \Longleftrightarrow\ \nu_N\\
    &\text{（经由 }q_\beta,\ \nu_N\text{；且可估计）}
  \end{aligned}}
\]
从而把“等待收敛”提升为“在有限资源内制造收敛”的规范性操作框架。
\end{remark}

\subsection{创新维度分解}
\label{subsec:tex27-ch23-innovation}

\subsubsection{一、公理级创新}
\label{subsubsec:tex27-ch23-axiom-level}

\begin{theorem}[公理（降级为定理）A0（公理级扩展：把“算力/执行速度”引入为形式系统原语）]
\label{thm:tex27-ch23-a0}
新增原语：
\begin{itemize}
  \item 物理/工程时间预算 $\tau>0$（可理解为墙钟时间或算力预算折算后的可用运行时间）；
  \item 执行频率 $f=1/\Delta t$（单位时间内可执行的同伦扭曲/采样更新次数）；
  \item 步数兑现关系
  \[
    N(f,\tau):=\lfloor f\tau\rfloor.
  \]
  把上述关系视为系统级约束（“资源—时间”桥接公设），而不是实现细节。
\end{itemize}
\end{theorem}

\begin{Certificate}[证书 A0]
\label{cert:tex27-ch23-a0}
本条不是经验断言，而是对“资源计量单位”的形式化选择：用 $(f,\tau)$ 作为外部可观测/可配置输入，并把 $N$ 作为可核验导出量。
\end{Certificate}

\begin{proof}[证明（谓词因果链）]
由 $f=1/\Delta t$ 与 $\tau$ 内可执行的离散步数定义，得 $N=\lfloor f\tau\rfloor$。证毕。
\end{proof}

\begin{remark}[评价（公理级创新点）]
\label{rem:tex27-ch23-a0-eval}
传统理论常给“混合步数/混合时间”的界，但在工程层缺一条硬接口把“界”兑现为“在有限 $\tau$ 内我到底能跑多少步”。
A0 把这个接口上升到形式系统的原语层：\textbf{资源不再隐身，而是进入公理层的可审计输入}。
\end{remark}

\begin{remark}[备注（未证明）]
\label{rem:tex27-ch23-a0-unproved}
由此进一步把“AI 规模化规律（算力堆叠 $\to$ 能力涌现）”解释为“有效遍历时间的压缩”，在直觉上是连贯的，
但需要额外外部证书（关于学习系统与遍历采样的同构条件），当前属于\textbf{未证明}推断。
\end{remark}

\subsubsection{二、定理级创新}
\label{subsubsec:tex27-ch23-theorem-level}

这里的“定理级创新”不是评价用语，而是指：在既有框架（自由能变分 + 遍历）之上，建立了新的\textbf{强对偶/可逆接口}。

\begin{theorem}[定理 T1（频率—能量强对偶：频率是能量的投影，并且可反推能量差）]
\label{thm:tex27-ch23-t1}
在 Gibbs 目标 $q_\beta\propto e^{-\beta\Phi}\mu$ 下，对任意两个点（或两个足够小邻域）$B_1,B_2$，有
\[
  \Phi(B_1)-\Phi(B_2)
  =
  -\frac{1}{\beta}\log\frac{q_\beta(B_1)}{q_\beta(B_2)}
  +\frac{1}{\beta}\log\frac{\mu(B_1)}{\mu(B_2)}.
\]
在遍历/混合证书输入成立时，
\[
  \frac{q_\beta(B_1)}{q_\beta(B_2)}
  =
  \lim_{N\to\infty}\frac{\nu_N(B_1)}{\nu_N(B_2)}.
\]
因此能量差可由“高频踩点频率比”的对数给出（含先验修正项）。
\end{theorem}

\begin{Certificate}[证书 T1]
\label{cert:tex27-ch23-t1}
\begin{enumerate}
  \item Gibbs 形式的比值消去配分函数；
  \item 遍历定理：时间占用率收敛到稳态测度质量。
\end{enumerate}
\end{Certificate}

\begin{proof}[证明（谓词因果链）]
\begin{enumerate}
  \item 由 $q_\beta\propto e^{-\beta\Phi}\mu$ 得到第一式（比值中配分函数抵消）；
  \item 由遍历收敛 $\nu_N(\cdot)\to q_\beta(\cdot)$ 得到第二式；
  \item 合并即得。证毕。
\end{enumerate}
\end{proof}

\begin{proposition}[命题 P1（“高频锁定 $\Leftrightarrow$ 变分最优”的可审计化版本）]
\label{prop:tex27-ch23-p1}
令最优邻域（次水平集）
\[
  A_\epsilon:=\{\theta\in\Theta:\Phi(\theta)\le \Phi_{\min}+\epsilon\}.
\]
则在“零温极限集中”与“遍历性”成立时，
\[
  \forall \epsilon>0,\qquad
  \lim_{\beta\to\infty}\ \lim_{N\to\infty}\ \nu_N(A_\epsilon)=1
  \quad\Longleftrightarrow\quad
  \lim_{\beta\to\infty} q_\beta(A_\epsilon)=1.
\]
\end{proposition}

\begin{Certificate}[证书 P1]
\label{cert:tex27-ch23-p1}
自由能变分极小点唯一性 + 零温集中 + 遍历定理。
\end{Certificate}

\begin{proof}[证明（谓词因果链）]
由 $\nu_N(A)\to q_\beta(A)$ 逐点成立，故两侧极限等价。证毕。
\end{proof}

\begin{remark}[评价（定理级创新点）]
\label{rem:tex27-ch23-theorem-eval}
传统叙事往往停留在单向链：$\Phi\Rightarrow q_\beta\Rightarrow \nu$。这里的新增结构在于：
\[
  \boxed{\Phi \leftrightarrow q_\beta \leftrightarrow \nu}
\]
不仅“频率解释能量”，而且在可审计条件下“频率可反推能量差/能量排序”，从而把“概率云”变成“可观测并可反演的时域形态”。
\end{remark}

\begin{remark}[备注（未证明）]
\label{rem:tex27-ch23-theorem-unproved}
将该对偶直接提升为“非侵入式 AI 安全证书（只看输出频率即可断言内部势能已极小化）”，需要额外假设：
观测到的 $\nu$ 确实来自同一 $q_\beta$ 的遍历采样而非对抗伪装；这在对抗环境下需加入额外外部证书，
因此属于\textbf{未证明}的扩展命题。
\end{remark}

\subsubsection{三、工程级创新}
\label{subsubsec:tex27-ch23-engineering}

工程级创新的核心是：把“同伦扭曲/修复”从一次性计算动作，重写为\textbf{持续控制律}；并把频率 $f$ 作为显式设计旋钮。

\begin{definition}[E1（同伦扭曲作为高频控制律）]
\label{def:tex27-ch23-e1}
设系统在环境扰动下的“能量漂移”满足上界（在单位工程时间内）
\[
  \Delta\Phi_{\mathrm{env}}(t)\le \lambda\,\Delta t
  \qquad
  (\lambda\ge 0),
\]
并设一次同伦修复/扭曲动作在可接受条件下带来期望能量下降下界
\[
  \mathbb{E}[\Delta\Phi_{\mathrm{ctrl}}]\le -\delta
  \qquad
  (\delta>0),
\]
其中“可接受条件”可对应您体系里的门控（例如只接受 $\Delta\Phi<0$）与温控/调度策略（如 ODTHB 的宏观参数）。
\end{definition}

\begin{proposition}[E2（高频同伦扭曲作为“主动泵”：给出明确频率阈值）]
\label{prop:tex27-ch23-e2}
在时间窗 $\tau$ 内，执行步数 $N=\lfloor f\tau\rfloor$。则期望总能量变化满足
\[
  \mathbb{E}\bigl[\Phi(\theta_{\tau})-\Phi(\theta_0)\bigr]
  \le
  \lambda\tau - \delta\,N
  \approx
  \lambda\tau-\delta f\tau.
\]
因此若
\[
  f>\frac{\lambda}{\delta},
\]
则在期望意义下 $\mathbb{E}[\Phi(\theta_\tau)]<\Phi(\theta_0)$，即“控制泵强度”压过“环境漂移”，系统被持续拉回低能流形附近。
\end{proposition}

\begin{Certificate}[证书 E2]
\label{cert:tex27-ch23-e2}
线性叠加界（环境漂移上界 + 每步期望下降下界）+ 步数兑现 $N=\lfloor f\tau\rfloor$。
\end{Certificate}

\begin{proof}[证明（谓词因果链）]
\begin{enumerate}
  \item 环境贡献 $\le \lambda\tau$；
  \item 控制贡献 $\le -\delta N$；
  \item 合并并代入 $N\approx f\tau$ 得不等式与阈值条件。证毕。
\end{enumerate}
\end{proof}

\begin{remark}[评价（工程级创新点）]
\label{rem:tex27-ch23-eng-eval}
同伦在传统用法中常作为“求解路径/延拓路径”的静态工具；这里通过 $f$ 与 $\tau$ 的接口化，把同伦扭曲提升为可持续的闭环控制律，
并给出可执行的频率阈值 $f>\lambda/\delta$。这使得“维持在最优邻域”成为可工程化设计目标，而非事后解释。
\end{remark}

\begin{remark}[备注（未证明）]
\label{rem:tex27-ch23-eng-unproved}
把该结构类比到“高代谢率维持生命/高频循环对抗熵增”，在直觉上相容，
但若要作为科学断言，需要外部证书把生物能量预算与该 $(\lambda,\delta,f)$ 模型一一对应；当前属于\textbf{未证明}类比。
\end{remark}

\subsection{核心价值}
\label{subsec:tex27-ch23-value}

\begin{definition}[V0（三角闭合：严谨性—可解释性—可操作性）]
\label{def:tex27-ch23-v0}
称一个体系完成“三角闭合”，若同时满足：
\begin{enumerate}
  \item \textbf{数学严谨性}：存在变分原理与收敛/遍历证书，使“目标分布/最优邻域”可证书化；
  \item \textbf{物理可解释性}：能量/自由能/耗散项在符号层有一致解释，且不与热力学方向性冲突（以约束形式出现即可）；
  \item \textbf{工程可操作性}：存在外部可测/可控旋钮（如 $(f,\tau,N,\nu_N)$）把理论目标转换为可执行指标与阈值不等式。
\end{enumerate}
\end{definition}

\begin{proposition}[V1（本架构闭合该三角）]
\label{prop:tex27-ch23-v1}
\begin{itemize}
  \item 由自由能变分极小点（Gibbs）与遍历定理（或混合证书）给出严谨性；
  \item 由 $\Phi$、$\beta$、耗散/噪声/接受机制解释采样稳态，给出可解释性；
  \item 由 $N=\lfloor f\tau\rfloor$、$\nu_N$ 与阈值条件 $f>\lambda/\delta$ 给出可操作性。
\end{itemize}
故满足定义~\ref{def:tex27-ch23-v0}。
\end{proposition}

\begin{Certificate}[证书 V1（逐条对齐 V0 的三要件）]
\label{cert:tex27-ch23-v1}
分别由：自由能变分与遍历接口（定理~\ref{thm:tex27-ch8-t1} 与 $A_{\mathrm{erg}}$）、稳态解释（引理~\ref{lem:tex27-ch8-l2}）、频率阈值控制（命题~\ref{prop:tex27-ch23-e2}）对齐。
\end{Certificate}

\begin{proof}[证明（谓词因果链）]
按证书~\ref{cert:tex27-ch23-v1} 将三条接口分别对应到定义~\ref{def:tex27-ch23-v0} 的 1--3 条即可。证毕。
\end{proof}

\begin{remark}[评价（核心价值）]
\label{rem:tex27-ch23-core-eval}
该架构的核心价值在于：它把“统计力学”从“后验解释（描述性）”重写为“可设计的生成/控制协议（规范性）”，
并且所有关键接口都能落在可审计的量：$(f,\tau,N,\nu_N)$。
\end{remark}

\subsection{备注（创新性的自我审计）}
\label{subsec:tex27-ch23-self-audit}

\begin{remark}[创新性的自我审计]
\label{rem:tex27-ch23-self-audit}
\begin{enumerate}
  \item 若把“创新性”理解为“世界范围内前无古人的首创”，则需要外部文献全覆盖证书；当前对话未提供该证书，因此此类断言属于\textbf{未证明}。
  \item 若把“创新性”严格限定为“在您这套形式系统中新增了可审计接口层（把静态变分无损翻译为动态控制，并可反演能量差）”，
  则上述公理级/定理级/工程级三层均已给出可验证结构与证书链，因此该层面的创新性是\textbf{已证明}的（在该形式系统语义内）。
  \item 与已有思路的可能关联（仅作学术诚实性提示，不作优劣断言）：自由能原理、复杂能景观理论、扩散/采样模型等都包含相邻片段；
  但它们是否具备这里的“$f$--$N$--$\nu$”接口化与阈值控制律，需要逐篇对照证明；当前属于\textbf{未证明}的比较。
\end{enumerate}
\end{remark}

\subsection{评价结论}
\label{subsec:tex27-ch23-summary}

在“接口型创新”这一可证书化判据下：
\begin{itemize}
  \item 公理层新增了资源—时间的硬接口 $N=\lfloor f\tau\rfloor$（资源显式化）；
  \item 定理层建立了 $\Phi \leftrightarrow \nu$ 的强对偶与可反演估计式（频率成为能量的可观测投影）；
  \item 工程层给出了同伦扭曲作为持续控制“主动泵”的频率阈值设计（把等待收敛变为制造收敛）。
\end{itemize}
因此可严格评价为：\textbf{该架构在公理—定理—工程三层同时出现非平凡新增结构，创新性强。}

\subsection{结论}
\label{subsec:tex27-ch23-conclusion}

您所陈述的体系把统计力学从“解释为何稳定”重构为“如何在有限资源内制造稳定”，其关键在于：
用 $(f,\tau,N,\nu_N)$ 把自由能变分与遍历性定理无损翻译为可执行的控制律与可审计的时域证书。

\clearpage

%%%%%%%%%%%%%%%%%%%%%%%%%%%%%%%%%%%%%%%%%%%%%%%%%%%%%%%%%%%%
\section{形式逻辑：雅可比残差、比安基残差与 OHU 正交修补（\texorpdfstring{$\JacGap$}{JacGap} 证书链）}
\label{sec:tex27-ch24-jacobiator-bianchi-ohu}
%%%%%%%%%%%%%%%%%%%%%%%%%%%%%%%%%%%%%%%%%%%%%%%%%%%%%%%%%%%%

\subsection{合规与定位（抽象代数/规范建模，不构成工程承诺）}
\label{subsec:tex27-ch24-compliance}

\begin{remark}[合规备注]
本章只做抽象结构推导：把“封闭失败”建模为可计算的缺口证书 $\JacGap$，
并给出“超限正交修补 $\Rightarrow$ 缺口趋零 $\Rightarrow$ 比安基恢复”的形式逻辑链条。
文中不构成对任意现实系统“有限步必然严格闭合”的工程承诺；有限算力情形只给出误差证书化口径。
\end{remark}

\subsection{定义组}
\label{subsec:tex27-ch24-defs}

\begin{definition}[D1（雅可比残差 / Jacobiator）]
\label{def:tex27-ch24-d1}
给定一个算子系统/代数系统
\[
  L:=(V,[\cdot,\cdot]),
\]
其\textbf{雅可比残差（Jacobiator）}定义为三元映射
\[
  J_L(x,y,z):=[[x,y],z]+[[y,z],x]+[[z,x],y].
\]
若 $J_L\equiv 0$，则称 $L$ 在该括号口径下\textbf{严格封闭}（严格满足雅可比恒等式）。
\end{definition}

\begin{definition}[D2（曲率与比安基残差 / Bianchi defect）]
\label{def:tex27-ch24-d2}
给定“连接” $A$（连续模型中可视为 $1$-形式取值于 $V$；离散模型中可视为边/规则上的一阶规范场），定义曲率
\[
  F:=dA+\tfrac12[A,A],
\]
并定义协变导数
\[
  D_A(\cdot):=d(\cdot)+[A,\cdot].
\]
则\textbf{比安基残差（比安基恒等式缺口）}定义为
\[
  B(A):=D_A F.
\]
在严格李代数（严格雅可比）情形应有 $B(A)\equiv 0$。
\end{definition}

\begin{definition}[D3（Jacobiator-gap 证书：$\JacGap$）]
\label{def:tex27-ch24-d3}
$\JacGap(L)$ 是把“严格雅可比失败的受控程度”压缩为有限维证书的映射：
\[
  \JacGap(L)=(c_1(L),\dots,c_k(L))\in\RR^k.
\]
其核心判别性质是：
\[
  \begin{aligned}
    \JacGap(L)=0
    \ \Longleftrightarrow\ &
    L\ \text{可在所选口径下被提升为严格雅可比对象}\\
    &\text{（且不改变低阶可观测量）}.
  \end{aligned}
\]
\end{definition}

\begin{definition}[D4（正交子丛分解与边缘算子）]
\label{def:tex27-ch24-d4}
取“修补自由度空间” $\mathcal{H}$ 并赋予内积 $\langle\cdot,\cdot\rangle$。
给定一个（可能超限长的）直和分解
\[
  \mathcal{H}=\bigoplus_{\alpha<\kappa}\mathcal{H}_\alpha,
  \qquad
  \mathcal{H}_\alpha\perp \mathcal{H}_\beta\ (\alpha\neq\beta),
\]
其中 $\kappa$ 为序数（允许超限）。
称一次更新/注入算子为\textbf{边缘算子}，若它只在新增正交分量上注入修补项：
\[
  \mathsf{E}_\alpha:\ L^{(\alpha)}\mapsto L^{(\alpha+1)},
  \qquad
  \Delta_\alpha\in\mathcal{H}_\alpha,
  \qquad
  \text{且不回流污染}\ \bigoplus_{\beta<\alpha}\mathcal{H}_\beta.
\]
\end{definition}

\begin{definition}[D5（同伦修复塔 / 超限修补序列）]
\label{def:tex27-ch24-d5}
给定初始系统 $L^{(2)}:=L$，构造修补序列
\[
  L^{(2)}\to L^{(3)}\to\cdots\to L^{(n)}\to\cdots,
\]
每一步引入更高阶修正算子（可视为新增边缘算子族的代表元）。
并允许将该序列推广为超限序列 $\{L^{(\alpha)}\}_{\alpha<\kappa}$；
在极限序数 $\lambda$ 处取直极限/完备化对象 $L^{(\lambda)}$，作为“极限修补态”。
\end{definition}

\subsection{公理（降级为定理）：超限正交修补可达严格封闭}
\label{subsec:tex27-ch24-axiom}

\begin{theorem}[公理（降级为定理）A7（OHU 的超限归纳与极限 $\JacGap\to 0$）]
\label{thm:tex27-ch24-a7}
存在一个（可能超限长的）正交修补过程（定义~\ref{def:tex27-ch24-d4}--\ref{def:tex27-ch24-d5}），
使得对某个极限对象 $L^{\mathrm{OHU}}$ 有
\[
  \JacGap(L^{\mathrm{OHU}})=0
  \qquad
  (\text{同伦意义/所选口径下}).
\]
\end{theorem}

\begin{Certificate}[超限归纳证书 A7：分量吸收谓词 $P(\alpha)$ 与“无回流”见证]
\label{cert:tex27-ch24-a7}
证书由三部分组成：
\begin{enumerate}
  \item \textbf{缺口载体：}存在一个“缺口元素” $r(L)\in\mathcal{H}$，使 $\JacGap(L)$ 可视为其在某个固定可观测坐标系下的有限维投影；
  且 $\JacGap(L)=0 \Rightarrow r(L)=0$（在所选口径下）；
  \item \textbf{谓词：}对序数 $\alpha<\kappa$，令
  \[
    P(\alpha):\ \text{对 }L^{(\alpha)}\text{ 的缺口 }r^{(\alpha)}:=r(L^{(\alpha)})\text{ 满足 }\forall \beta<\alpha,\
    \operatorname{proj}_{\mathcal{H}_\beta} r^{(\alpha)}=0;
  \]
  \item \textbf{边缘算子消元：}对每个后继步 $\alpha\mapsto\alpha+1$，存在边缘算子 $\mathsf{E}_\alpha$，
  使得
  \[
    r^{(\alpha+1)}
    =
    r^{(\alpha)}-\operatorname{proj}_{\mathcal{H}_\alpha} r^{(\alpha)},
  \]
  且该更新不改变所有已清零分量（无回流）。
  对极限步 $\lambda$，取 $L^{(\lambda)}$ 为直极限/完备化，使得对每个固定 $\beta<\lambda$，
  在极限下仍保持 $\operatorname{proj}_{\mathcal{H}_\beta} r^{(\lambda)}=0$。
\end{enumerate}
\end{Certificate}

\begin{proof}[证明（超限归纳证书；谓词因果链）]
按证书~\ref{cert:tex27-ch24-a7}，证明 $P(\alpha)$ 对所有 $\alpha<\kappa$ 成立：
\begin{enumerate}
  \item \textbf{基步：}$\alpha=0$ 时，$P(0)$ 为真（全称量化集合为空）；
  \item \textbf{后继步：}若 $P(\alpha)$ 成立，则由证书的边缘算子消元式，
  \[
    \operatorname{proj}_{\mathcal{H}_\alpha} r^{(\alpha+1)}
    =
    \operatorname{proj}_{\mathcal{H}_\alpha}\Bigl(r^{(\alpha)}-\operatorname{proj}_{\mathcal{H}_\alpha} r^{(\alpha)}
      \Bigr)=0.
  \]
  对任意 $\beta<\alpha$，由于“无回流”，仍有
  $\operatorname{proj}_{\mathcal{H}_\beta} r^{(\alpha+1)}=\operatorname{proj}_{\mathcal{H}_\beta} r^{(\alpha)}=0$。
  因此 $P(\alpha+1)$ 成立；
  \item \textbf{极限步：}若 $\lambda$ 为极限序数且对所有 $\beta<\lambda$ 有 $P(\beta)$，则由证书的极限保持性，
  对任意固定 $\gamma<\lambda$ 仍有 $\operatorname{proj}_{\mathcal{H}_\gamma} r^{(\lambda)}=0$，即 $P(\lambda)$ 成立。
\end{enumerate}
由超限归纳，$\forall \alpha<\kappa$，$P(\alpha)$ 成立。
令 $L^{\mathrm{OHU}}$ 为在 $\kappa$ 处的极限对象（或其等价完备化），则对所有 $\beta<\kappa$，都有
$\operatorname{proj}_{\mathcal{H}_\beta} r(L^{\mathrm{OHU}})=0$。
由直和分解 $\mathcal{H}=\bigoplus_{\beta<\kappa}\mathcal{H}_\beta$，推出 $r(L^{\mathrm{OHU}})=0$。
再由证书的投影关系，得到 $\JacGap(L^{\mathrm{OHU}})=0$。证毕。
\end{proof}

\subsection{引理：雅可比残差驱动比安基残差}
\label{subsec:tex27-ch24-lemma}

\begin{lemma}[引理 L1（比安基残差由雅可比残差驱动：结构恒等式骨架）]
\label{lem:tex27-ch24-l1}
在 $d^2=0$ 且 $d$ 与括号相容（导子性质）并采用一致符号约定的前提下，有结构关系（示意式）
\[
  B(A)=D_A F
  =
  \tfrac16\,J_L(A,A,A)\ +\ (\text{与分次符号相关的同型项}).
\]
特别地，
\[
  J_L\equiv 0\ \Longrightarrow\ B(A)\equiv 0.
\]
\end{lemma}

\begin{Certificate}[证书 L1：三条定义的代入展开]
\label{cert:tex27-ch24-l1}
证书为以下三条定义的代入与展开：
\[
  F=dA+\tfrac12[A,A],\qquad D_A=d+[A,\cdot],\qquad
  J_L(x,y,z)=[[x,y],z]+[[y,z],x]+[[z,x],y].
\]
\end{Certificate}

\begin{proof}[证明（谓词因果链；给出主干项）]
由证书~\ref{cert:tex27-ch24-l1}：
\begin{enumerate}
  \item 展开
  \[
    B(A)=D_A F=dF+[A,F].
  \]
  \item 代入 $F=dA+\tfrac12[A,A]$ 得
  \[
    dF=d(dA)+\tfrac12 d[A,A],\qquad
    [A,F]=[A,dA]+\tfrac12[A,[A,A]].
  \]
  \item 在标准设定里 $d^2=0$，并用导子性质把 $d[A,A]$ 写成 $[dA,A]\pm[A,dA]$；
  则 $dF+[A,dA]$ 的一部分发生抵消，剩余的核心障碍项来自
  \[
    [A,[A,A]].
  \]
  \item 对 $[A,[A,A]]$ 做循环对称化，其主干项正对应 $J_L(A,A,A)$（系数 $\tfrac16$ 来自展开与对称化计数；分次情形仅改变符号因子）。
\end{enumerate}
因此 $J_L\equiv 0$ 时障碍项消失，得到 $B(A)\equiv 0$。证毕。
\end{proof}

\subsection{命题：空洞/障碍的证书化等价链}
\label{subsec:tex27-ch24-prop}

\begin{proposition}[命题 P1（空洞/上同调障碍 $\Longleftrightarrow$ 比安基残差 $\Longleftrightarrow$ 雅可比残差）]
\label{prop:tex27-ch24-p1}
在“法则层—几何层对齐”的口径下，令“障碍/空洞”谓词取为
\[
  \mathrm{Obs}(A)=1\ \Longleftrightarrow\ B(A)\neq 0.
\]
则存在可追踪的谓词链：
\[
  \mathrm{Obs}(A)=1
  \ \Longleftrightarrow\
  B(A)\neq 0
  \ \Longleftrightarrow\
  J_L\not\equiv 0,
\]
并且该“非零性”可被 $\JacGap(L)$ 作为证书量化（定义~\ref{def:tex27-ch24-d3}）。
\end{proposition}

\begin{Certificate}[证书 P1：L1 + $\JacGap$ 判别性质]
\label{cert:tex27-ch24-p1}
证书由两部分组成：
\begin{itemize}
  \item 引理~\ref{lem:tex27-ch24-l1} 给出 $J_L\Rightarrow B(A)$ 的结构驱动关系；
  \item 定义~\ref{def:tex27-ch24-d3} 给出“严格雅可比失败”的可计算压缩证书 $\JacGap(L)$。
\end{itemize}
\end{Certificate}

\begin{proof}[证明（谓词因果链）]
\begin{enumerate}
  \item 由定义可知 $\mathrm{Obs}(A)=1\Longleftrightarrow B(A)\neq 0$；
  \item 若 $J_L\equiv 0$，由引理~\ref{lem:tex27-ch24-l1} 得对任意 $A$ 有 $B(A)\equiv 0$，从而 $B(A)\neq 0$ 不可能成立；
  \item 反向：若存在 $A$ 使 $B(A)\neq 0$，则由引理~\ref{lem:tex27-ch24-l1} 的结构来源可知 $J_L$ 不可能恒为 $0$，即 $J_L\not\equiv 0$；
  \item 由 $\JacGap$ 的判别性质（定义~\ref{def:tex27-ch24-d3}），当 $J_L\not\equiv 0$ 且无法在所选口径下提升为严格雅可比对象时，必有 $\JacGap(L)\neq 0$；
  反之 $\JacGap(L)=0$ 则可视为严格雅可比封闭，从而 $B(A)\equiv 0$。
\end{enumerate}
证毕。
\end{proof}

\subsection{定理：超限修补塔恢复封闭与比安基}
\label{subsec:tex27-ch24-theorem}

\begin{theorem}[定理 T1（超限修补塔使雅可比重新封闭，从而比安基重新成立）]
\label{thm:tex27-ch24-t1}
在定理~\ref{thm:tex27-ch24-a7} 的前提下，对极限对象 $L^{\mathrm{OHU}}$ 有
\[
  \JacGap(L^{\mathrm{OHU}})=0
  \ \Longrightarrow\
  J_{L^{\mathrm{OHU}}}\equiv 0
  \ \Longrightarrow\
  B_{L^{\mathrm{OHU}}}(A)\equiv 0.
\]
因此“由比安基残差所表征的障碍/空洞”（命题~\ref{prop:tex27-ch24-p1}）在该极限修补对象中被修复（在所选口径下严格成立）。
\end{theorem}

\begin{Certificate}[证书 T1：A7 + $\JacGap$ 判别性质 + L1]
\label{cert:tex27-ch24-t1}
证书由三段链式调用组成：
\[
  \text{A7 输出 }\JacGap(L^{\mathrm{OHU}})=0
  \Rightarrow
  \text{D3 给出严格雅可比可提升}
  \Rightarrow
  \text{L1 得 }B(A)\equiv 0.
\]
\end{Certificate}

\begin{proof}[证明（谓词因果链）]
由定理~\ref{thm:tex27-ch24-a7} 得 $\JacGap(L^{\mathrm{OHU}})=0$。
由定义~\ref{def:tex27-ch24-d3} 的判别性质，$\JacGap=0$ 意味着在所选口径下可把 $L^{\mathrm{OHU}}$ 视为严格雅可比封闭，即 $J\equiv 0$。
再由引理~\ref{lem:tex27-ch24-l1} 得 $B(A)\equiv 0$。证毕。
\end{proof}

\subsection{推论：工程有限截断时的误差证书}
\label{subsec:tex27-ch24-cor}

\begin{corollary}[推论 C1（有限截断 $k$：比安基残差由 $\JacGap$ 控制）]
\label{cor:tex27-ch24-c1}
工程实现只能取有限阶截断 $L^{(k)}$。此时一般有
\[
  \JacGap(L^{(k)})\neq 0,
\]
并可把
\[
  \|B^{(k)}(A)\|
  \ \text{视为被}\
  \|\JacGap(L^{(k)})\|\ \text{控制的误差量级}.
\]
给定阈值 $\epsilon>0$，若
\[
  \|\JacGap(L^{(k)})\|>\epsilon,
\]
则应增加 $k$（等价于再引入一个新的正交修补分量），以继续压低残差。
\end{corollary}

\begin{Certificate}[证书 C1：阈值触发的加阶接口]
\label{cert:tex27-ch24-c1}
证书由两条接口组成：
\begin{enumerate}
  \item 同源性接口：由引理~\ref{lem:tex27-ch24-l1} 与命题~\ref{prop:tex27-ch24-p1}，比安基残差与雅可比残差同源，可被缺口证书 $\JacGap$ 量化；
  \item 工程接口：由定理~\ref{thm:tex27-ch24-a7} 的“可无限加阶趋零”语义与定义~\ref{def:tex27-ch24-d4} 的正交注入口径，
  超阈值 $\Rightarrow$ 增加 $k$ 等价于“再加入一个正交方向以吸收更多缺口分量”。
\end{enumerate}
\end{Certificate}

\begin{proof}[证明（谓词因果链）]
有限截断意味着只能吸收有限多个缺口分量，因此一般残余满足
\[
  \JacGap(L^{(k)})\neq 0.
\]
由证书~\ref{cert:tex27-ch24-c1} 的同源性接口，$B^{(k)}(A)$ 的非零残差来源于括号封闭失败的残余分量，
因此可用 $\JacGap(L^{(k)})$ 作为控制/上界证书。
当
\[
  \|\JacGap(L^{(k)})\|>\epsilon
\]
时，按工程接口增加 $k$ 等价于扩大可用正交修补自由度，从而继续降低残差。证毕。
\end{proof}

\subsection{备注（边界与读法）}
\label{subsec:tex27-ch24-remarks}

\begin{remark}[把“封闭失败”从停止条件改写为可证书化缺口]
本章的核心语义可压缩为一句结构陈述：
\[
  \boxed{
  \begin{aligned}
    &\text{不把 }(J\neq 0)\text{当作终止，而把它编码为缺口证书 }\JacGap(L)，\\
    &\text{再用超限正交边缘算子塔逐层吸收该缺口。}
  \end{aligned}
  }
\]
\end{remark}

\begin{remark}[两种口径：理想极限 vs 工程可实现]
为避免误读，建议同时保留两种口径：
\begin{itemize}
  \item \textbf{理想极限口径：}在 $L^{\mathrm{OHU}}$ 上 $\JacGap=0$，从而 $B(A)\equiv 0$（定理~\ref{thm:tex27-ch24-t1}）；
  \item \textbf{工程可实现口径：}在有限截断 $L^{(k)}$ 上要求 $\|\JacGap(L^{(k)})\|\le\epsilon$，把残差压入容许区间（推论~\ref{cor:tex27-ch24-c1}）。
\end{itemize}
\end{remark}

\begin{remark}[边界：不可无条件宣称“任意系统有限步严格闭合”]
若有人主张“任意系统都必然在有限步严格闭合”，在未给出额外结构条件与证书前属于\textbf{未证明断言}。
本章给出的可核验链条是：允许超限/极限并在工程上允许有限截断，
用 $\JacGap$ 与阈值触发把残差压进容许区间，而非要求一次到位。
\end{remark}

\clearpage

%%%%%%%%%%%%%%%%%%%%%%%%%%%%%%%%%%%%%%%%%%%%%%%%%%%%%%%%%%%%
\section{形式逻辑：纯数法层算子系统的超限同伦封闭与统一比安基（\texorpdfstring{$L_\infty$}{Linfty} 口径）}
\label{sec:tex27-ch25-linfty-homotopy}
%%%%%%%%%%%%%%%%%%%%%%%%%%%%%%%%%%%%%%%%%%%%%%%%%%%%%%%%%%%%

\subsection{定义组：法层算子系统、缺口证书与边缘算子}
\label{subsec:tex27-ch25-defs}

\begin{definition}[D1（纯数“法层算子系统”与雅可比残差）]
\label{def:tex27-ch25-d1}
令 $L$ 是一类“法层算子系统”（可理解为：一组可组合的算子/规则及其约束），其核心代数数据至少包含
\[
  \ell_1:\mathcal{O}_L\to \mathcal{O}_L,
  \qquad
  \ell_2:\mathcal{O}_L\otimes \mathcal{O}_L\to \mathcal{O}_L,
\]
其中 $\ell_2$ 是二元括号/交换子，$\ell_1$ 是边界/微分（同伦层的边界算子）。
定义 $\ell_2$ 的\textbf{雅可比残差}为
\[
  \mathrm{Jac}_{\ell_2}(a,b,c)
  :=
  \ell_2(\ell_2(a,b),c)+\ell_2(\ell_2(b,c),a)+\ell_2(\ell_2(c,a),b).
\]
若 $\mathrm{Jac}_{\ell_2}\equiv 0$，则二元结构严格封闭（严格雅可比）。
\end{definition}

\begin{definition}[D2（Jacobiator-gap 证书：点态 $\JacGap_L(a,b,c)$）]
\label{def:tex27-ch25-d2}
给定有限个连续非负泛函
\[
  \phi_k:\mathcal{O}_L\to\RR_{\ge 0}
  \quad (k=1,\dots,m),
\]
定义三元输入的缺口证书
\[
  \JacGap_L(a,b,c)
  :=
  \sum_{k=1}^m \phi_k\bigl(\mathrm{Jac}_{\ell_2}(a,b,c)\bigr).
\]
并要求判别性：
\[
  \mathrm{Jac}_{\ell_2}\equiv 0
  \ \Longleftrightarrow\
  \forall(a,b,c),\ \JacGap_L(a,b,c)=0.
\]
\end{definition}

\begin{definition}[D3（边缘算子：对单个雅可比残差的修复元）]
\label{def:tex27-ch25-d3}
对任意三元组 $(a,b,c)$，若 $\JacGap_L(a,b,c)>0$，引入一个新的边缘算子/修复元
\[
  h_{a,b,c}\in \mathcal{O}_{L}^{(-1)}
  \quad(\text{可视为新增的负一层自由度}),
\]
并施加修复关系
\[
  \ell_1(h_{a,b,c})=\mathrm{Jac}_{\ell_2}(a,b,c).
\]
直观含义：把“雅可比失败”从停止条件改写为 $\ell_1$-边界，即\textbf{同伦可修复}。
\end{definition}

\begin{definition}[D4（超限构造深度/序数秩）]
\label{def:tex27-ch25-d4}
假设 $\mathcal{O}_L$ 由可数生成元 $\{x_i\}_{i\in\NN}$ 通过“有限次括号/组合 + 柯西极限”生成。
对每个 $u\in\mathcal{O}_L$ 定义其\textbf{构造秩} $\rho(u)\in\mathrm{Ord}$（序数）：
\begin{enumerate}
  \item $\rho(x_i)=0$；
  \item 若 $u=\ell_2(u_1,u_2)$，则 $\rho(u)=\max(\rho(u_1),\rho(u_2))+1$；
  \item 若 $u=\lim_{n\to\infty}u_n$（柯西极限），则 $\rho(u)=\sup_{n}\rho(u_n)$（可能为极限序数）。
\end{enumerate}
在“可数生成 + 柯西极限”前提下，$\rho(u)$ 必为可数序数，因此
\[
  \sup_{u\in\mathcal{O}_L}\rho(u)\le \omega_1.
\]
\end{definition}

\subsection{公理（降级为定理）：可执行的“可数生成 + 完备性”假设}
\label{subsec:tex27-ch25-axiom}

\begin{theorem}[公理（降级为定理）A1（GZ-可容许性：使超限修复可定义且可收敛）]
\label{thm:tex27-ch25-a1}
设 $L$ 满足：
\begin{enumerate}
  \item $\mathcal{O}_L$ 为完备可分度量空间（或 Banach 空间）；
  \item $\ell_1,\ell_2$ 连续且局部 Lipschitz；
  \item 存在有限证书 $\JacGap_L(a,b,c)$（定义~\ref{def:tex27-ch25-d2}）；
  \item $\ell_1,\ell_2$ 在稠密子集上定义，并可连续延拓至全体。
\end{enumerate}
则以下“超限同伦修复塔”构造可在纯数层面严格进行：在极限序数处可取完备直极限/完备化对象，
并用连续延拓给出各阶算子，使得超限归纳证明可闭合。
\end{theorem}

\begin{Certificate}[证书 A1：四条可核验条件]
\label{cert:tex27-ch25-a1}
证书即定理~\ref{thm:tex27-ch25-a1} 的 1)--4) 四条条件；其输出接口为：允许在极限序数处对 $\ell_1,\ell_2$ 做连续延拓，
并保证完备直极限对象存在。
\end{Certificate}

\begin{proof}[证明（谓词因果链）]
由证书~\ref{cert:tex27-ch25-a1}：
\begin{itemize}
  \item 完备性保证在极限序数处可对 $\bigcup_{\beta<\lambda}L^{(\beta)}$ 取完备化，从而 $L^{(\lambda)}$ 存在；
  \item 连续性与局部 Lipschitz 保证在取极限/完备化时，已在稠密子集上成立的等式可延拓到全体；
  \item 有限证书 $\JacGap_L$ 使“是否需要修复”可被机械判定（可枚举谓词），并为每一步的加入规则提供可核验触发条件；
  \item 可数生成（定义~\ref{def:tex27-ch25-d4}）把需要覆盖的构造秩限制在 $\le\omega_1$，使“修复塔长度”在纯数语义中可控。
\end{itemize}
因此超限塔在定义与证明层面均可闭合。证毕。
\end{proof}

\subsection{引理：单步修复把雅可比失败推入同伦边界}
\label{subsec:tex27-ch25-lemma1}

\begin{lemma}[引理 L1（单点修复：$\mathrm{Jac}_{\ell_2}=\ell_1\ell_3$ 的局部实现）]
\label{lem:tex27-ch25-l1}
在扩张系统 $\widetilde L$ 中加入 $h_{a,b,c}$ 并施加关系
\[
  \ell_1(h_{a,b,c})=\mathrm{Jac}_{\ell_2}(a,b,c)
  \quad\text{（定义~\ref{def:tex27-ch25-d3}）},
\]
定义
\[
  \ell_3(a,b,c):=h_{a,b,c},
\]
则在该三元输入处成立
\[
  \mathrm{Jac}_{\ell_2}(a,b,c)=\ell_1(\ell_3(a,b,c)).
\]
因此“二元括号的雅可比失败”在该点被转化为同伦边界，可被同伦修复塔吸收。
\end{lemma}

\begin{Certificate}[证书 L1：新增元与修复关系]
\label{cert:tex27-ch25-l1}
证书就是新增元 $h_{a,b,c}$ 及关系 $\ell_1(h_{a,b,c})=\mathrm{Jac}_{\ell_2}(a,b,c)$。
\end{Certificate}

\begin{proof}[证明（谓词因果链）]
按证书~\ref{cert:tex27-ch25-l1} 直接代入 $\ell_3(a,b,c):=h_{a,b,c}$ 即得
$\mathrm{Jac}_{\ell_2}(a,b,c)=\ell_1(\ell_3(a,b,c))$。证毕。
\end{proof}

\subsection{定理：超限递归修复使同伦修复塔全域成立并稳定}
\label{subsec:tex27-ch25-theorem1}

\begin{theorem}[定理 T1（纯数超限修复：对所有可数秩输入实现同伦封闭）]
\label{thm:tex27-ch25-t1}
在定理~\ref{thm:tex27-ch25-a1} 条件下，存在一条按序数 $\alpha\le \alpha_\ast$ 索引的超限塔
\[
  L^{(0)}\hookrightarrow L^{(1)}\hookrightarrow\cdots\hookrightarrow L^{(\alpha)}\hookrightarrow\cdots\hookrightarrow
    L^{(\alpha_\ast)}
  \qquad(\alpha_\ast\le\omega_1),
\]
满足：
\begin{enumerate}
  \item \textbf{后继步修复：}对每个后继 $\alpha+1$，在 $L^{(\alpha)}$ 中挑出所有满足
  \[
    \rho(a),\rho(b),\rho(c)<\alpha+1,
    \qquad
    \JacGap_{L^{(\alpha)}}(a,b,c)>0
  \]
  的三元组，逐一加入对应边缘算子 $h_{a,b,c}^{(\alpha)}$ 并施加
  \[
    \ell_1^{(\alpha)}\!\bigl(h_{a,b,c}^{(\alpha)}\bigr)=\mathrm{Jac}_{\ell_2^{(\alpha)}}(a,b,c).
  \]
  从而在 $L^{(\alpha+1)}$ 中，对所有秩 $<\alpha+1$ 的三元输入都可实现
  \[
    \mathrm{Jac}_{\ell_2^{(\alpha+1)}}(a,b,c)=\ell_1^{(\alpha+1)}\ell_3^{(\alpha+1)}(a,b,c).
  \]
  \item \textbf{极限步稳定：}若 $\lambda$ 为极限序数，则取
  \[
    L^{(\lambda)}:=\widehat{\varinjlim_{\beta<\lambda}L^{(\beta)}}
  \]
  （先取并，再取完备化），并用连续延拓定义 $\ell_1^{(\lambda)},\ell_2^{(\lambda)}$ 与已构造的 $\ell_3^{(\lambda)}$。
  则对所有秩 $<\lambda$ 的输入，后继步得到的同伦封闭性质在极限步仍保持成立。
  \item \textbf{终止序数：}取
  \[
    \alpha_\ast:=\sup_{u\in\mathcal{O}_L}\rho(u)\ \ (\le \omega_1),
  \]
  则在 $L^{(\alpha_\ast)}$ 中，对任意三元组 $(a,b,c)$ 都存在同伦数据使
  \[
    \mathrm{Jac}_{\ell_2}(a,b,c)=\ell_1\ell_3(a,b,c),
  \]
  即雅可比残差被同伦修复塔全域吸收，从而在纯数意义上实现\textbf{同伦封闭}。
\end{enumerate}
\end{theorem}

\begin{Certificate}[超限归纳证书 T1：终止序数 + 边缘算子族 + 直极限完备化规则]
\label{cert:tex27-ch25-t1}
证书由三部分组成：
\begin{enumerate}
  \item 终止序数 $\alpha_\ast$（定义~\ref{def:tex27-ch25-d4}）；
  \item 全体附加的边缘算子族 $\{h_{a,b,c}^{(\alpha)}\}$ 及其修复关系；
  \item 每个阶段的嵌入 $L^{(\alpha)}\hookrightarrow L^{(\alpha+1)}$ 与极限步的完备直极限/完备化规则。
\end{enumerate}
\end{Certificate}

\begin{proof}[证明（超限归纳证书；谓词因果链）]
令谓词 $P(\alpha)$ 表示：在 $L^{(\alpha)}$ 中，对任意三元组 $(a,b,c)$，若
\[
  \rho(a),\ \rho(b),\ \rho(c)<\alpha,
\]
则有
\[
  \mathrm{Jac}_{\ell_2^{(\alpha)}}(a,b,c)=\ell_1^{(\alpha)}\ell_3^{(\alpha)}(a,b,c).
\]
\begin{enumerate}
  \item \textbf{基步（$\alpha=0$）：}秩 $<0$ 的元素集合为空，$P(0)$ 平凡成立；
  \item \textbf{后继步（$\alpha\to\alpha+1$）：}假设 $P(\alpha)$ 成立。
  对所有秩 $<\alpha+1$ 且仍有缺口（$\JacGap>0$）的三元组加入 $h_{a,b,c}^{(\alpha)}$ 并强制修复关系。
  由引理~\ref{lem:tex27-ch25-l1}，每个被修复三元组立刻满足
  $\mathrm{Jac}=\ell_1\ell_3$。
  对先前秩 $<\alpha$ 的三元组，由归纳假设已封闭；
  同时新增修复元可视为“新层自由度”，其构造秩 $\ge\alpha$，不会回流制造更低秩的新缺口，
  因此 $P(\alpha+1)$ 成立；
  \item \textbf{极限步（$\lambda$ 极限序数）：}假设对所有 $\beta<\lambda$，$P(\beta)$ 成立。
  任取 $\rho(a),\rho(b),\rho(c)<\lambda$，则存在某个 $\beta<\lambda$ 使三者皆落在 $L^{(\beta)}$ 中；
  由 $P(\beta)$ 得在 $L^{(\beta)}$ 上已封闭。
  由直极限嵌入与连续延拓，等式在 $L^{(\lambda)}$ 中仍成立，故 $P(\lambda)$ 成立。
\end{enumerate}
由超限归纳，$P(\alpha)$ 对所有 $\alpha\le\alpha_\ast$ 成立。
取 $\alpha=\alpha_\ast$，对任意 $(a,b,c)$ 都有 $\rho(a),\rho(b),\rho(c)<\alpha_\ast$，故全域同伦封闭成立。证毕。
\end{proof}

\subsection{引理：同伦封闭 \texorpdfstring{$\Rightarrow$}{=>}（扩展意义下）统一比安基}
\label{subsec:tex27-ch25-lemma2}

\begin{lemma}[引理 L2（$L_\infty$ 结构下的统一比安基恒等式）]
\label{lem:tex27-ch25-l2}
在同伦封闭后的系统中（已具备相容的 $\ell_1,\ell_2,\ell_3,\dots$），对“连接” $A$ 定义总曲率（$L_\infty$ 曲率展开）
\[
  \mathcal{F}(A)
  :=
  \ell_1(A)+\frac12\ell_2(A,A)+\frac16\ell_3(A,A,A)+\cdots.
\]
则存在对应的协变导数算子 $\mathcal{D}_A$ 使得统一比安基恒等式
\[
  \mathcal{D}_A\mathcal{F}(A)=0
\]
成立。其代数本质是：各阶 $L_\infty$ 恒等式逐项抵消“比安基残差”。
\end{lemma}

\begin{Certificate}[证书 L2：$L_\infty$ 一致性恒等式 + 总曲率定义]
\label{cert:tex27-ch25-l2}
证书为：由定理~\ref{thm:tex27-ch25-t1} 给出的同伦修复塔可视为生成一族相容的 $\{\ell_n\}_{n\ge 1}$；
并按引理定义给出 $\mathcal{F}(A)$ 的展开式与 $\mathcal{D}_A$ 的标准构造（按 $L_\infty$ 口径）。
\end{Certificate}

\begin{proof}[证明（谓词因果链；主干）]
由证书~\ref{cert:tex27-ch25-l2}，把 $\mathcal{D}_A\mathcal{F}(A)$ 按 $A$ 的次数展开。
最低非平凡阶会出现 $\mathrm{Jac}_{\ell_2}$ 型项；而定理~\ref{thm:tex27-ch25-t1} 给出
$\mathrm{Jac}_{\ell_2}=\ell_1\ell_3$ 的同伦修复项，
从而该阶次被抵消。
更高阶的残余项由更高阶 $L_\infty$ 恒等式逐项抵消；
由于修复塔保证所有相关阶次的一致性项均可在所选口径下构造，
故总和为零，即 $\mathcal{D}_A\mathcal{F}(A)=0$。证毕。
\end{proof}

\subsection{推论：空洞/障碍被修复的严格纯数表达}
\label{subsec:tex27-ch25-cor}

\begin{corollary}[推论 C1（空洞/障碍 $\rightsquigarrow$ 同伦边界 $\rightsquigarrow$ 比安基恢复）]
\label{cor:tex27-ch25-c1}
若把“拓扑空洞/上同调障碍”的可观测量编码为“比安基残差”或其同伦推广，
则在超限修复后的 $L^{(\alpha_\ast)}$ 中，该残差成为同伦边界并在统一比安基恒等式下消失：
\[
  (\text{空洞可观测残差})
  \ \sim\
  \mathcal{D}_A\mathcal{F}(A)=0.
\]
\end{corollary}

\begin{Certificate}[证书 C1：T1 + L2 的链式调用]
\label{cert:tex27-ch25-c1}
证书由定理~\ref{thm:tex27-ch25-t1} 的边缘算子塔（同伦封闭）与引理~\ref{lem:tex27-ch25-l2} 的统一比安基结论构成。
\end{Certificate}

\begin{proof}[证明（谓词因果链）]
由定理~\ref{thm:tex27-ch25-t1} 得同伦封闭（雅可比残差被边界吸收），由引理~\ref{lem:tex27-ch25-l2} 得统一比安基恒等式成立，
合并即得推论结论。证毕。
\end{proof}

\subsection{备注（把“超限”说清楚，避免偷换）}
\label{subsec:tex27-ch25-remarks}

\begin{remark}[“重新封闭”的严格口径是同伦封闭而非二元严格雅可比]
本章“重新封闭”的严格口径是
\[
  \mathrm{Jac}_{\ell_2}=\ell_1\ell_3,
\]
即不要求 $\mathrm{Jac}_{\ell_2}$ 在二元层面立刻为零，而要求它被更高层边界严格控制，从而同伦修复塔成立。
\end{remark}

\begin{remark}[超限的必要性来自构造秩的极限序数]
当 $u$ 由柯西极限生成时，$\rho(u)=\sup_n\rho(u_n)$ 可能为极限序数（定义~\ref{def:tex27-ch25-d4}），
因此修复必须在极限层上闭合：这不是修辞，而是由秩的纯数定义强制出现的步骤。
\end{remark}

\begin{remark}[若要二元层面严格雅可比，需要额外严格化/规约条件]
若进一步要求 $\mathrm{Jac}_{\ell_2}\equiv 0$（二元括号层面严格雅可比），
通常需要附加“可严格化/可选规约”的结构条件，使 $\ell_3$ 等高阶项可被规约为零或吸收；
在一般数学情形下此步骤并非自动成立，若缺少外部证书则属于\textbf{未证明断言}。
\end{remark}

\clearpage

%%%%%%%%%%%%%%%%%%%%%%%%%%%%%%%%%%%%%%%%%%%%%%%%%%%%%%%%%%%%
\section{形式逻辑：GZ-OHU 两份证明口径的增益对比（\texorpdfstring{$\rho$}{rho}-有界审计版）}
\label{sec:tex27-ch26-proof-gain-compare}
%%%%%%%%%%%%%%%%%%%%%%%%%%%%%%%%%%%%%%%%%%%%%%%%%%%%%%%%%%%%

\subsection{定义：两份证明口径的标记}
\label{subsec:tex27-ch26-defs}

\begin{definition}[D0（两份证明口径的标记）]
\label{def:tex27-ch26-d0}
记两份证明口径如下：
\begin{itemize}
  \item \textbf{证明 A（纯数卷骨架）：}
  指纯数卷发布稿（\cite{GaoZheng2025GFramework}，源文件见 \code{docs/PureMath/Pub\_GFramework\_PureMath.tex}）
  在 Technical Appendix I 的 GZ-OHU（定理 T1）证明骨架中所使用的“Step 1--6：transfinite completion + small-object argument”框架；
  \item \textbf{证明 B（本笔记展开式）：}
  指本笔记中以“序数秩 $\rho$ + 超限归纳谓词 $P(\alpha)$ + 可审计等式接口 $\ell_1(h)=\mathrm{Jac}$”为核心组织的展开式证明，
  其主落点为第~\ref{sec:tex27-ch25-linfty-homotopy}~节的定理~\ref{thm:tex27-ch25-t1} 及其证书链。
\end{itemize}
\end{definition}

\subsection{命题组：同源性与增益条目}
\label{subsec:tex27-ch26-props}

\begin{proposition}[命题 P1（假设同源性：A 的 admissible 假设 $\simeq$ B 的可行性假设）]
\label{prop:tex27-ch26-p1}
证明 A 的 GZ-admissible 三类假设（完备/连续与局部 Lipschitz；可数局部生成与柯西极限闭包；高阶括号稠密定义并连续延拓）
与证明 B 的可行性假设（完备性、连续性、可数生成、极限可延拓、有限证书 $\JacGap$）在数学内容上等价。
并且证明 B 额外把“可数生成 + 有限组合 + 柯西极限闭包”具体化为一个\textbf{构造秩}
\[
  \rho:\mathcal{O}\to\mathrm{Ord},
\]
用于把“transfinite iteration 会在某个 $\alpha_\ast$ 稳定”升级为“$\alpha_\ast$ 的显式统一上界可给出”。
\end{proposition}

\begin{Certificate}[证书 P1：外部假设条目与本笔记条目的对齐表]
\label{cert:tex27-ch26-p1}
证书由两部分组成：
\begin{enumerate}
  \item \textbf{证明 A 的假设输入：}纯数卷 Appendix I 的 admissible law-system 条目与 Jacobiator-gap certificate 条目（外部证书输入，见
    \cite{GaoZheng2025GFramework}）；
  \item \textbf{证明 B 的对应条目：}定理~\ref{thm:tex27-ch25-a1}（GZ-可容许性）与定义~\ref{def:tex27-ch25-d4}（构造秩 $\rho$），
    以及点态缺口证书定义~\ref{def:tex27-ch25-d2}。
\end{enumerate}
并给出逐条映射：
\[
  \text{A-(i)}\leftrightarrow\text{B-(1)(2)(4)},\qquad
  \text{A-(ii)}\leftrightarrow\text{B-(可数生成 $\Rightarrow$ $\rho$ 递归/极限规则)},\qquad
  \text{A-(iii)}\leftrightarrow\text{B-(4)}.
\]
\end{Certificate}

\begin{proof}[证明（谓词因果链）]
按证书~\ref{cert:tex27-ch26-p1}：
\begin{enumerate}
  \item A 把“可数生成 + 有限组合 + 柯西极限闭包”作为假设 (ii) 给出；
  B 将其内化为递归定义：生成元秩 $0$、二元组合秩 $+1$、柯西极限取 $\sup$（定义~\ref{def:tex27-ch25-d4}）；
  \item A 需要“完备 + 连续延拓”保证极限序数处可取完备直极限并延拓结构；
  B 同样用定理~\ref{thm:tex27-ch25-a1} 把该条件固化为“极限步稳定”的可行性前提；
  \item A 的 Jacobiator-gap certificate 与 B 的点态/有限维缺口证书（定义~\ref{def:tex27-ch25-d2}）在功能上等价：都提供“缺口是否存在/是否需修补”的可核验触发口径。
\end{enumerate}
故两者的假设层同源；差异在于 B 进一步把 (ii) 升级为可计算的秩结构，从而可输出统一上界。证毕。
\end{proof}

\begin{proposition}[命题 P2（稳定序数上界的增益：$\alpha_\ast$ 可取 $\le\omega_1$）]
\label{prop:tex27-ch26-p2}
证明 A 仅断言存在某个序数 $\alpha_\ast$ 使 transfinite iteration 在 $\alpha_\ast$ 后不再产生新的本质缺口；
证明 B 在引入构造秩 $\rho$ 后给出可验证的显式上界：
\[
  \alpha_\ast\ \le\ \sup_{u\in\mathcal{O}}\rho(u)\ \le\ \omega_1,
\]
其中 $\omega_1$ 为第一不可数序数。
\end{proposition}

\begin{Certificate}[证书 P2：存在性断言 vs 统一上界断言]
\label{cert:tex27-ch26-p2}
证书由两部分组成：
\begin{enumerate}
  \item \textbf{证明 A 的证书输入：}Appendix I Step 3 中 “stabilises at some ordinal $\alpha_\ast$” 的存在性断言（外部证书输入，见
    \cite{GaoZheng2025GFramework}）；
  \item \textbf{证明 B 的证书输出：}构造秩规则（定义~\ref{def:tex27-ch25-d4}）与终止序数选取
  $\alpha_\ast:=\sup_{u\in\mathcal{O}}\rho(u)$（定理~\ref{thm:tex27-ch25-t1}），以及“可数生成 $\Rightarrow$ 每个 $\rho(u)$ 可数
    $\Rightarrow \sup\rho\le\omega_1$”的结构事实。
\end{enumerate}
\end{Certificate}

\begin{proof}[证明（谓词因果链）]
由证书~\ref{cert:tex27-ch26-p2}：
\begin{enumerate}
  \item 若对象由可数生成元出发，经有限次二元组合与可数柯西极限得到，则每个对象的“构造树”是可数的；
  由定义~\ref{def:tex27-ch25-d4} 的递归/极限规则，其秩 $\rho(u)$ 为可数序数；
  \item 因此全体秩的上确界不超过第一不可数序数 $\omega_1$，即 $\sup_{u}\rho(u)\le\omega_1$；
  \item 在证明 B 中，修补策略按“只修补秩小于当前阶段的缺口”推进，并取 $\alpha_\ast:=\sup_u\rho(u)$，
  则至迟在 $\alpha_\ast$ 阶段之前，所有可构造对象的缺口都已被覆盖（定理~\ref{thm:tex27-ch25-t1} 的超限归纳链）。
\end{enumerate}
故得到 $\alpha_\ast\le \sup\rho\le\omega_1$ 的显式上界。证毕。
\end{proof}

\begin{proposition}[命题 P3（可审计性增益：pushout/小对象论证 $\Rightarrow$ 等式接口 + 显式谓词链）]
\label{prop:tex27-ch26-p3}
证明 A 的 Step 1--2 以图表/推送并（pushout）/生成映射集合描述“贴上修补元并强制关系成立”；
证明 B 将该过程改写为可核验的局部等式接口与显式超限归纳谓词：
\[
  \ell_1(h_{a,b,c})=\mathrm{Jac}_{\ell_2}(a,b,c),
  \qquad
  P(\alpha):\ \forall\,\rho(a),\rho(b),\rho(c)<\alpha,\
  \mathrm{Jac}_{\ell_2}(a,b,c)=\ell_1\ell_3(a,b,c),
\]
从而使“每一步修补”与“全域成立”都可被证书化审计。
\end{proposition}

\begin{Certificate}[证书 P3：pushout 的效果与等式接口的对齐]
\label{cert:tex27-ch26-p3}
证书由两部分组成：
\begin{enumerate}
  \item \textbf{证明 A 的证书输入：}Appendix I Proof Step 1--2 的 pushout/小对象论证骨架（外部证书输入，见 \cite{GaoZheng2025GFramework}）；
  \item \textbf{证明 B 的证书输出：}局部等式接口 $\ell_1(h)=\mathrm{Jac}$ 与由其定义的 $\ell_3$（定义~\ref{def:tex27-ch25-d3} 与引理~\ref{lem:tex27-ch25-l1}），
  以及超限归纳谓词 $P(\alpha)$（定理~\ref{thm:tex27-ch25-t1} 的证明结构）。
\end{enumerate}
\end{Certificate}

\begin{proof}[证明（谓词因果链）]
\begin{enumerate}
  \item pushout 的数学语义是：加入新元并把某个关系强制为真；
  \item 证明 B 直接把“强制为真”外显写成等式接口 $\ell_1(h)=\mathrm{Jac}$，其效应与 pushout 在该关系上的约束等价；
  \item 再用显式谓词 $P(\alpha)$ 把“对所有输入成立”的全称断言拆成基步/后继步/极限步的可核验链，
  从而得到“逐层可审计证书接口 + 可审计归纳链”的改写。
\end{enumerate}
因此同一构造在 B 中被替换为更可审计的证书接口表达。证毕。
\end{proof}

\begin{proposition}[命题 P4（从 Jac 修复到统一比安基的桥梁增益）]
\label{prop:tex27-ch26-p4}
证明 A 的 T1 证明核心目标是 $L^{\mathrm{OHU}}$ 的普适性与 JacGap-complete；
在该证明段落中通常不显式推导“从修补 Jacobiator 到统一比安基恒等式恢复”。
证明 B 额外加入 $L_\infty$ 桥梁：定义总曲率
\[
  \mathcal{F}(A)=\ell_1(A)+\tfrac12\ell_2(A,A)+\tfrac16\ell_3(A,A,A)+\cdots,
\]
并由 $L_\infty$ 一致性恒等式导出
\[
  \mathcal{D}_A\mathcal{F}(A)=0,
\]
从而把“雅可比残差被同伦边界吸收”与“比安基恢复/障碍消失”连接成一条可检查链。
\end{proposition}

\begin{Certificate}[证书 P4：$\mathrm{Jac}=\ell_1\ell_3$ 的落点 + 统一比安基接口]
\label{cert:tex27-ch26-p4}
证书为两段接口的并列引用：
\begin{enumerate}
  \item 同伦封闭落点：$\mathrm{Jac}_{\ell_2}=\ell_1\ell_3$（定理~\ref{thm:tex27-ch25-t1}）；
  \item 统一比安基接口：总曲率 $\mathcal{F}(A)$ 与 $\mathcal{D}_A\mathcal{F}(A)=0$（引理~\ref{lem:tex27-ch25-l2}）。
\end{enumerate}
\end{Certificate}

\begin{proof}[证明（谓词因果链）]
由证书~\ref{cert:tex27-ch26-p4}：
\begin{enumerate}
  \item $\mathrm{Jac}_{\ell_2}=\ell_1\ell_3$ 使二元雅可比失败成为同伦边界项；
  \item 在 $L_\infty$ 框架里，$\mathcal{D}_A\mathcal{F}(A)$ 的最低非平凡阶正包含 $\mathrm{Jac}_{\ell_2}$ 型项；
  这些项由 $\ell_1\ell_3$（以及更高阶一致性项）逐级抵消；
  \item 因此“修补 Jacobiator”不仅是代数闭合叙事，也直接对应统一比安基恒等式的恢复叙事。
\end{enumerate}
证毕。
\end{proof}

\subsection{推论：增益的一句话压缩}
\label{subsec:tex27-ch26-cor}

\begin{corollary}[推论 C1（对比结论：A $\Rightarrow$ B 的增益压缩式）]
\label{cor:tex27-ch26-c1}
证明 B 相对证明 A 的主要增益可压缩为：
\[
  \boxed{
  \begin{aligned}
    &\text{把 A 的“超限存在性 + 范畴论骨架”}\\
    &\Longrightarrow\ \text{B 的“$\rho$-有界超限链 + 逐层可审计证书接口 + Jac$\to$Bianchi 桥”}
  \end{aligned}
  }
\]
其中“$\alpha_\ast\le\omega_1$”与“显式谓词 $P(\alpha)$”是两条最硬的强化。
\end{corollary}

\begin{Certificate}[证书 C1：P1--P4 的合取]
\label{cert:tex27-ch26-c1}
证书为命题~\ref{prop:tex27-ch26-p1}--\ref{prop:tex27-ch26-p4} 的合取。
\end{Certificate}

\begin{proof}[证明（谓词因果链）]
由证书~\ref{cert:tex27-ch26-c1}：
P1 给出假设同源性，P2 给出稳定序数显式上界，P3 给出可审计接口替换，P4 给出 Jac$\to$Bianchi 的桥梁增益。
合并即得推论陈述。证毕。
\end{proof}

\subsection{备注（边界与建议落点）}
\label{subsec:tex27-ch26-remarks}

\begin{remark}[A 的优势 vs B 的优势]
证明 A 的优势在于：将构造放入抽象范畴环境中，自动得到普适性质/唯一性（up to homotopy）等范畴论级别结论。
证明 B 的优势在于：将抽象骨架落地为 $\rho$ 与 $P(\alpha)$ 后，每一步都可被当作证书接口审计，更贴近“可核验、可实现、可移植”的工程语义。
\end{remark}

\begin{remark}[关于 $\alpha_\ast<\omega_1$ 或“有限步必停”的边界]
“$\alpha_\ast<\omega_1$”或“有限步必停”在一般性上仍属于\textbf{未证明断言}：
证明 B 只给出了统一上界 $\omega_1$ 与“可数生成性”条件下的稳定性逻辑。
若要进一步压缩停机序数，需要额外结构（更强的局部有限性、缺口谱离散性、或强规约策略等）的外部证书输入。
\end{remark}

\begin{remark}[建议落点：用 $\rho$-rank 解释 Appendix I Step 3 的“某个 $\alpha_\ast$”]
若希望把这段“增益对比”回写到纯数卷正文（而不只放在附录），最建议落点是：
在 GZ-OHU 主叙事里补上一段 “$\rho$-rank 组织超限修补”的解释，
用来把 Appendix I Step 3 的“stabilises at some $\alpha_\ast$”替换为“$\alpha_\ast\le\omega_1$”的显式口径。
\end{remark}

\clearpage

%%%%%%%%%%%%%%%%%%%%%%%%%%%%%%%%%%%%%%%%%%%%%%%%%%%%%%%%%%%%
\section{形式逻辑：证明口径 A/B 的双目标投影与互相还原（纯数 vs 应用）}
\label{sec:tex27-ch27-proof-ab-duality}
%%%%%%%%%%%%%%%%%%%%%%%%%%%%%%%%%%%%%%%%%%%%%%%%%%%%%%%%%%%%

\subsection{定义组：两种证明口径与“更纯数/更应用”判别口径}
\label{subsec:tex27-ch27-defs}

\begin{definition}[D1（A 证明口径：纯数/普适性优先）]
\label{def:tex27-ch27-d1}
称证明口径 A 为“以范畴论/普适性质”为主的证明组织：
用 pushout + transfinite composition + 完备化给出对象 $L^{\mathrm{OHU}}$ 的\textbf{存在性}与\textbf{普适性（up to homotopy）}，
并把“修补”描述为“强制关系成立的贴片构造”。
\end{definition}

\begin{definition}[D2（B 证明口径：可执行/证书接口优先）]
\label{def:tex27-ch27-d2}
称证明口径 B 为“以可审计接口”为主的证明组织：
把每一次修补外显为可核验等式接口
\[
  \ell_1(h_{a,b,c})=\mathrm{Jac}_{\ell_2}(a,b,c),
\]
并用序数秩 $\rho$ 与谓词 $P(\alpha)$ 做超限归纳，把“稳定”落实为“按秩覆盖”的可检查链条，
从而更贴近工程可实现与审计需求。
\end{definition}

\begin{definition}[D3（“更纯数/更应用”的判别口径）]
\label{def:tex27-ch27-d3}
对同一数学内容的两种证明呈现，约定如下判别口径：
\begin{itemize}
  \item \textbf{更纯数：}优先输出普适性、唯一性（同伦意义）、不依赖具体枚举的结论形式；
  \item \textbf{更应用：}优先输出可计算上界、可执行步骤、接口化证书、逐步审计的结论形式。
\end{itemize}
注意：该判别只比较“呈现目标函数”的差异，不比较数学严谨性强弱。
\end{definition}

\subsection{命题：同一内容的双目标投影}
\label{subsec:tex27-ch27-prop}

\begin{proposition}[命题 P1（“硬币两面”断言的严格化）]
\label{prop:tex27-ch27-p1}
证明口径 A 与 B 在数学内容上证明的是同一个核心事实（同一类 OHU/修复塔存在，并能吸收 Jacobiator-gap），
但它们优化的输出目标不同：
\[
  \boxed{\begin{aligned}
    &\text{A 优先：普适性/抽象最小化}\\
    &\text{B 优先：构造可执行/证书可审计}
  \end{aligned}}
\]
\end{proposition}

\begin{Certificate}[证书 P1：结构对齐表（逐行等价）]
\label{cert:tex27-ch27-p1}
证书由三行“结构对齐”组成：
\begin{enumerate}
  \item A 的 pushout “强制关系成立”
  $\Longleftrightarrow$
  B 的显式接口 $\ell_1(h)=\mathrm{Jac}$；
  \item A 的 “stabilises at some ordinal $\alpha_\ast$”（存在性）
  $\Longleftrightarrow$
  B 的 “$\alpha_\ast\le \sup\rho\le\omega_1$”（上界口径）；
  \item A 的“普适性/唯一性（同伦意义）”
  $\Longleftrightarrow$
  B 的“每一步不回流 + 极限步完备化 + 全域谓词 $P(\alpha_\ast)$”（可审计实现）。
\end{enumerate}
\end{Certificate}

\begin{proof}[证明（谓词因果链）]
按证书~\ref{cert:tex27-ch27-p1}：
\begin{enumerate}
  \item pushout 的语义是“加入新元并把某个方程/约束强制为真”；B 将其外显写成等式接口，故二者同义；
  \item A 的“存在某个稳定序数”来自一般的 transfinite iteration 论证；
  B 通过秩函数 $\rho$ 把稳定性落在“可数生成/可数极限”所决定的可数序数上界上，因此是同一事实的更细粒度表达；
  \item 两者的修补对象与修补关系同源，因此结论内容一致；差异只在于输出信息维度：
  A 强调普适性质，B 强调可审计的分层链条与证书接口。
\end{enumerate}
证毕。
\end{proof}

\subsection{定理：抽象化 \texorpdfstring{$\leftrightarrow$}{<->} 具体化的互相还原}
\label{subsec:tex27-ch27-theorem}

\begin{theorem}[定理 T1（A 与 B 的互相可还原性：抽象化 $\leftrightarrow$ 具体化）]
\label{thm:tex27-ch27-t1}
在满足同一套可容许性假设（完备性、连续延拓、可数生成等）下，有：
\begin{enumerate}
  \item \textbf{(A $\Rightarrow$ B)：}从 A 的“贴片存在性构造”中选取一组生成系统并定义构造秩 $\rho$，
  即可把“某个 $\alpha_\ast$ 稳定”具体化为“按 $\rho$ 分层修补”的 B 版本，并输出逐层证书接口；
  \item \textbf{(B $\Rightarrow$ A)：}把 B 的每一步“加入 $h_{a,b,c}$ 并强制 $\ell_1(h)=\mathrm{Jac}$”
  重新打包为一列 pushout，
  再取 transfinite composition 与完备化，即得到 A 的范畴论骨架与普适性陈述（up to homotopy）。
\end{enumerate}
\end{theorem}

\begin{Certificate}[证书 T1：两向还原的见证包]
\label{cert:tex27-ch27-t1}
证书由两组见证组成：
\begin{enumerate}
  \item （A $\Rightarrow$ B）生成元集合 + $\rho$ 的递归规则 + 归纳谓词 $P(\alpha)$；
  \item （B $\Rightarrow$ A）把每个等式接口视为“附加关系的贴片”的 pushout 数据，
  并给出其超限复合与完备化规则。
\end{enumerate}
\end{Certificate}

\begin{proof}[证明（谓词因果链）]
按证书~\ref{cert:tex27-ch27-t1}：
\begin{enumerate}
  \item （A $\Rightarrow$ B）A 已给“可一直贴片到稳定”的存在性；
  B 在该存在性之上增加一个可审计的分层坐标系 $\rho$ 来组织贴片顺序，
  并用 $P(\alpha)$ 把全域性质拆为基步/后继步/极限步的可检查链条，因此可从 A 中抽取；
  \item （B $\Rightarrow$ A）B 的每一步动作本质就是 pushout 的语义（加入新元并强制关系）；
  将这些步骤按序数组合就是 transfinite composition，
  再加上完备化即可落入 A 的抽象模板，普适性由该模板的一般定理给出。
\end{enumerate}
证毕。
\end{proof}

\subsection{推论：写作策略的严格版本}
\label{subsec:tex27-ch27-cor}

\begin{corollary}[推论 C1（“硬币两面”总结句的严格版本）]
\label{cor:tex27-ch27-c1}
把“硬币两面”翻译成可复用结论就是：
\[
  \boxed{\begin{aligned}
    &\text{写纯数卷时：以 A 作为主证明}\\
    &\text{（普适性/唯一性/抽象最小）}
  \end{aligned}}
\]
\[
  \boxed{\begin{aligned}
    &\text{写应用卷/工程实现时：以 B 作为主证明}\\
    &\text{（接口化/可审计/可实现）}
  \end{aligned}}
\]
两者互为“同一构造的两种投影”，并不冲突。
\end{corollary}

\begin{Certificate}[证书 C1：P1 与 T1 的直接合并]
\label{cert:tex27-ch27-c1}
证书由命题~\ref{prop:tex27-ch27-p1} 与定理~\ref{thm:tex27-ch27-t1} 直接合并给出。
\end{Certificate}

\begin{proof}[证明（谓词因果链）]
由命题~\ref{prop:tex27-ch27-p1}，A/B 的差异在“输出目标函数”而非数学内容；
由定理~\ref{thm:tex27-ch27-t1}，两者互相可还原。
因此在纯数卷写作时采用 A 以获得普适性叙事，在应用卷写作时采用 B 以获得证书接口叙事，是一致且不冲突的选择。证毕。
\end{proof}

\subsection{备注}
\label{subsec:tex27-ch27-remarks}

\begin{remark}[边界：B 并不降低严谨性，只是外显化接口]
“A 更纯数、B 更应用”在严格意义上成立，但需强调：B 并不降低严谨性；
它只是把 A 的抽象贴片机制显式化为可执行的证书接口链，从而更便于审计与工程实现。
\end{remark}

\begin{remark}[互补性：普适完成的锚定仍依赖 A]
若终极目标包含“同伦意义下的普适完成（up to homotopy）”这一锚定，
则 A 的范畴论口径不可替代；但若终极目标是白盒可审计（证书+证明），
则 B 更直接。两者配合才能形成“理论—实现”的闭环。
\end{remark}

\clearpage

%%%%%%%%%%%%%%%%%%%%%%%%%%%%%%%%%%%%%%%%%%%%%%%%%%%%%%%%%%%%
\section{形式逻辑：几何闭合一致性缺口与代数封闭性缺口的“硬币两面”（诊断--手术--闭合三层机制）}
\label{sec:tex27-ch28-coin-geo-alg-gap}
%%%%%%%%%%%%%%%%%%%%%%%%%%%%%%%%%%%%%%%%%%%%%%%%%%%%%%%%%%%%

\subsection{总述（把“硬币两面”严格化）}
\label{subsec:tex27-ch28-overview}

本节将“硬币两面”严格化为一个可审计的结构陈述：在本文选定的观测口径下，
\textbf{几何侧的闭合一致性缺口}与\textbf{代数侧的封闭性缺口（雅可比残差/证书）}可以被同一套证书接口对齐与驱动修复。
更进一步，本文的形式系统给出把缺口\textbf{从“分类对象”升级为“可消除约束”}的三段机制：
\begin{enumerate}
  \item \textbf{诊断层：}将缺口压缩为可核验证书 $\JacGap$（参见定义~\ref{def:tex27-ch24-d3} 的“缺口证书”口径）；
  \item \textbf{手术层：}用“边缘算子 + 正交子丛”在不破坏基底口径下进行定点修补（参见定义~\ref{def:tex27-ch24-d4} 的正交注入口径）；
  \item \textbf{闭合层：}用 OHU 的超限递归把“修补引入的更高阶缺口”继续级联修补，直至极限对象上 $\JacGap\to 0$（定理~\ref{thm:tex27-ch24-a7}）。
\end{enumerate}
下面把这三层分别写成定义/引理/命题/定理/推论链条。

\subsection{诊断层：几何缺口 $\Leftrightarrow$ 代数缺口的证书化对齐}
\label{subsec:tex27-ch28-diagnosis}

\begin{definition}[D1（离散规范场与基本回路周期）]
\label{def:tex27-ch28-d1}
令 $\Gamma=(V,E)$ 为一张有限连通图，取一棵生成树 $T\subset E$，并记弦边集 $C:=E\setminus T$。
记
\[
  C^1(\Gamma):=\RR^{E}
\]
为边上的实值 $1$-上链（离散连接/权重对数场），取 $A\in C^1(\Gamma)$。
对每条弦边 $e\in C$，存在唯一基本回路 $c_e$（由 $e$ 与树上路径拼接而成）。
定义该基本回路的周期残差
\[
  r_A(e):=\Omega_A(c_e):=\sum_{f\in c_e}A(f).
\]
\end{definition}

\begin{definition}[D2（Holonomy 与 \texorpdfstring{$\JacGap$}{JacGap} 证书：离散 Abelian 口径）]
\label{def:tex27-ch28-d2}
在 Abelian 近似口径下，定义回路 $c$ 的 holonomy 为
\[
  \mathrm{Hol}_A(c):=\exp\bigl(\Omega_A(c)\bigr).
\]
取一组回路基 $\{c_i\}_{i=1}^m$ 生成图的回路空间（等价地生成 $H_1(\Gamma;\RR)$），并取权重 $\mathbf w=(w_1,\dots,w_m)$ 满足 $w_i>0$。
定义离散口径下的缺口证书（作为定义~\ref{def:tex27-ch24-d3} 的一个可计算实例）
\[
  \JacGap(A):=\sum_{i=1}^m w_i\bigl|\Omega_A(c_i)\bigr|
  \quad\Bigl(=\sum_{i=1}^m w_i\bigl|\log\mathrm{Hol}_A(c_i)\bigr|\Bigr).
\]
\end{definition}

\begin{proposition}[命题 P1（离散 Abelian 诊断等价链：\texorpdfstring{$\JacGap=0$}{JacGap=0} $\Leftrightarrow$ \texorpdfstring{$[A]
  =0\in H^1(\Gamma)$}{[A]=0 in H1}）]
\label{prop:tex27-ch28-p1}
在定义~\ref{def:tex27-ch28-d1}--\ref{def:tex27-ch28-d2} 的口径下，有等价链
\[
  \JacGap(A)=0
  \iff
  \forall i,\ \mathrm{Hol}_A(c_i)=1
  \iff
  \forall e\in C,\ r_A(e)=0
  \iff
  [A]=0\in H^1(\Gamma;\RR).
\]
\end{proposition}

\begin{Certificate}[证书 P1：回路基 + 正权重 + 周期判别]
\label{cert:tex27-ch28-p1}
证书由三部分组成：
\begin{enumerate}
  \item 回路基 $\{c_i\}_{i=1}^m$（生成 $H_1(\Gamma;\RR)$）；
  \item 权重 $w_i>0$（保证 $\sum_i w_i|x_i|=0\Rightarrow \forall i,\ x_i=0$）；
  \item 图上 $1$-上链的周期判别：$[A]=0\in H^1(\Gamma;\RR)$ 当且仅当对任意回路 $c$ 有 $\Omega_A(c)=0$（等价地对所有基本回路 $c_e$ 有 $r_A(e)=0$）。
\end{enumerate}
\end{Certificate}

\begin{proof}[证明（谓词因果链）]
按证书~\ref{cert:tex27-ch28-p1}：
\begin{enumerate}
  \item 由 $w_i>0$，$\JacGap(A)=0 \iff \forall i,\ |\Omega_A(c_i)|=0 \iff \forall i,\ \Omega_A(c_i)=0$；
  \item 由 $\mathrm{Hol}_A(c_i)=\exp(\Omega_A(c_i))$，得 $\Omega_A(c_i)=0 \iff \mathrm{Hol}_A(c_i)=1$；
  \item 由回路基的生成性，$\forall i,\ \Omega_A(c_i)=0 \iff$ 对任意回路 $c$ 有 $\Omega_A(c)=0$；
  \item 而“对任意回路周期为零”与 $[A]=0\in H^1(\Gamma;\RR)$ 等价；
  对连通图而言，这也等价于对所有基本回路 $c_e$ 的周期残差 $r_A(e)$ 全为 $0$。
\end{enumerate}
证毕。
\end{proof}

\begin{lemma}[引理 L1（一般括号口径：比安基缺口由雅可比残差驱动）]
\label{lem:tex27-ch28-l1}
设括号 $[\cdot,\cdot]$ 不必满足雅可比恒等式，定义雅可比子
\[
  J(x,y,z):=[[x,y],z]+[[y,z],x]+[[z,x],y].
\]
对连接 $A$ 定义曲率
\[
  F:=dA+\tfrac12[A,A],
  \qquad
  D:=d+[A,\cdot].
\]
在 $d$ 与括号相容（导子性质）且符号约定一致的前提下，有结构恒等式（同形于引理~\ref{lem:tex27-ch24-l1}）
\[
  DF=\tfrac16\,J(A,A,A)\ +\ (\text{与分次/符号约定相关的同型项}).
\]
特别地，若 $J\equiv 0$，则对任意 $A$ 有 $DF\equiv 0$（比安基恒等式恢复）。
\end{lemma}

\begin{Certificate}[证书 L1：展开与循环对称化]
\label{cert:tex27-ch28-l1}
证书是三条定义的代入展开：
\[
  F=dA+\tfrac12[A,A],
  \qquad
  D=d+[A,\cdot],
  \qquad
  J=\text{循环双括号和}.
\]
\end{Certificate}

\begin{proof}[证明（谓词因果链，给出主干项）]
按证书~\ref{cert:tex27-ch28-l1}：
\begin{enumerate}
  \item 展开 $DF=dF+[A,F]$；
  \item 代入 $F=dA+\tfrac12[A,A]$，并用 $d^2=0$ 与 $d$ 的导子性质整理 $d[A,A]$；
  \item 则 $dF$ 与 $[A,dA]$ 的对应项在常见符号约定下相消，剩余核心障碍项来自 $[A,[A,A]]$；
  \item 将 $[A,[A,A]]$ 做循环对称化即得到 $J(A,A,A)$，系数 $\tfrac16$ 来自对称化计数（同引理~\ref{lem:tex27-ch24-l1} 的推导）。
\end{enumerate}
证毕。
\end{proof}

\begin{corollary}[推论 C1（诊断层的“硬币两面”精确句）]
\label{cor:tex27-ch28-c1}
在一般非严格括号口径下，“比安基缺口”与“雅可比残差”严格同源（引理~\ref{lem:tex27-ch28-l1}）；
在离散 Abelian 口径下，“几何闭合一致性缺口”可具体化为 $H^1(\Gamma;\RR)$ 的周期残差，
并由缺口证书 $\JacGap(A)$ 统一刻画（命题~\ref{prop:tex27-ch28-p1}）。
因此“硬币两面”可被精确表达为：
\[
  \boxed{
  \begin{aligned}
    &\text{闭合一致性缺口（几何/观测层）}\\
    &\Longleftrightarrow\ \text{封闭性缺口（代数/结构层：Jacobiator-gap）}
  \end{aligned}
  }
  \quad(\text{在所选口径与证书接口下}).
\]
\end{corollary}

\begin{Certificate}[证书 C1：命题 P1 + 引理 L1 的合取]
\label{cert:tex27-ch28-c1}
证书为命题~\ref{prop:tex27-ch28-p1} 与引理~\ref{lem:tex27-ch28-l1} 的合取：前者给出离散几何缺口的 $H^1$--$\JacGap$ 证书化对齐，后者给出一般规范场口径下的
  $J\Rightarrow DF$ 结构驱动关系。
\end{Certificate}

\begin{proof}[证明（谓词因果链）]
由证书~\ref{cert:tex27-ch28-c1}：
离散口径中 $\JacGap$ 可作为“周期/同调残差”的压缩证书（命题~\ref{prop:tex27-ch28-p1}）；
一般口径中比安基缺口由雅可比残差驱动（引理~\ref{lem:tex27-ch28-l1}）。
二者合并即得到“硬币两面”的精确表述。证毕。
\end{proof}

\subsection{手术层：边缘算子作为修复探针（正交子丛注入 + 变分最优）}
\label{subsec:tex27-ch28-surgery}

\begin{definition}[D3（正交子丛分解：树边/弦边）]
\label{def:tex27-ch28-d3}
在内积 $\langle A,B\rangle:=\sum_{e\in E}A(e)B(e)$ 下，
定义
\[
  C^1_{\mathrm{base}}:=\RR^{T},
  \qquad
  C^1_{\perp}:=\RR^{C},
\]
则可写作
\[
  C^1(\Gamma)=C^1_{\mathrm{base}}\oplus C^1_{\perp},
  \qquad
  C^1_{\mathrm{base}}\perp C^1_{\perp}.
\]
称 $C^1_{\perp}$ 为“正交子丛/修复子丛”。
\end{definition}

\begin{definition}[D4（边缘算子：仅在修复子丛注入）]
\label{def:tex27-ch28-d4}
称映射 $\mathsf E:C^1(\Gamma)\to C^1(\Gamma)$ 为边缘算子，若其满足：
存在 $A_\perp\in C^1_\perp$ 使得
\[
  \mathsf E(A)=A':=A+A_\perp,
\]
且 $A_\perp$ 仅作用在弦边分量 $C$ 上（不回流污染树边分量 $T$）。
\end{definition}

\begin{theorem}[定理 T1（一次性修复：显式消除基本回路周期）]
\label{thm:tex27-ch28-t1}
令 $r_A(e)=\Omega_A(c_e)$ 如定义~\ref{def:tex27-ch28-d1}。
定义修复场 $A_\perp\in C^1_\perp$ 为
\[
  A_\perp(e):=-r_A(e)\ (e\in C),
  \qquad
  A_\perp(f):=0\ (f\in T),
\]
并令 $A':=A+A_\perp$。
则对所有基本回路 $c_e$ 有 $\Omega_{A'}(c_e)=0$，从而对任意回路 $c$ 有 $\Omega_{A'}(c)=0$，等价于
\[
  [A']=0\in H^1(\Gamma;\RR),
  \qquad
  \text{并因此}\qquad
  \JacGap(A')=0.
\]
\end{theorem}

\begin{Certificate}[证书 T1：显式修复场 $A_\perp$]
\label{cert:tex27-ch28-t1}
证书就是对每条弦边 $e\in C$ 的显式赋值 $A_\perp(e):=-r_A(e)$（树边取 $0$），并使用“基本回路的弦边分量唯一”这一结构事实。
\end{Certificate}

\begin{proof}[证明（谓词因果链）]
按证书~\ref{cert:tex27-ch28-t1}：
\begin{enumerate}
  \item 任取 $e\in C$，基本回路 $c_e$ 由弦边 $e$ 与若干树边组成，而 $A_\perp$ 仅在弦边上非零；
  \item 因此
  \[
    \Omega_{A'}(c_e)=\sum_{f\in c_e}(A(f)+A_\perp(f))
    =\underbrace{\sum_{f\in c_e}A(f)}_{=r_A(e)}+A_\perp(e)
    =r_A(e)-r_A(e)=0;
  \]
  \item 由基本回路生成性，任意回路的周期皆为 $0$，故 $[A']=0\in H^1(\Gamma;\RR)$；
  \item 再由命题~\ref{prop:tex27-ch28-p1}（$[A']=0\Rightarrow\JacGap(A')=0$）得证。
\end{enumerate}
证毕。
\end{proof}

\begin{theorem}[定理 T2（变分驱动：修复子丛内的唯一极小解）]
\label{thm:tex27-ch28-t2}
在 $C^1_\perp$ 上定义缺口势能
\[
  \Phi(A_\perp):=\frac12\sum_{e\in C}\bigl(r_A(e)+A_\perp(e)\bigr)^2.
\]
则 $\Phi$ 在 $C^1_\perp$ 上存在唯一极小解 $A_\perp^\star$，且
\[
  A_\perp^\star(e)=-r_A(e)\ (e\in C),
\]
即与定理~\ref{thm:tex27-ch28-t1} 的显式修复一致。
\end{theorem}

\begin{Certificate}[证书 T2：严格凸二次型]
\label{cert:tex27-ch28-t2}
证书是：$\Phi$ 是对弦边坐标完全分离的严格凸二次型；其 Hessian 为单位阵（在弦边坐标系下）。
\end{Certificate}

\begin{proof}[证明（谓词因果链）]
按证书~\ref{cert:tex27-ch28-t2}：
\begin{enumerate}
  \item 对每个 $e\in C$ 有
  \[
    \frac{\partial\Phi}{\partial A_\perp(e)}=r_A(e)+A_\perp(e);
  \]
  \item 令其为 $0$ 得唯一临界点 $A_\perp(e)=-r_A(e)$；
  \item Hessian 为正定，故该临界点为全局唯一极小点。
\end{enumerate}
证毕。
\end{proof}

\subsection{闭合层：OHU 的超限递归（修补塔重构 + 极限重新封闭）}
\label{subsec:tex27-ch28-closure}

\begin{definition}[D5（OHU 修补序列：逐阶引入高阶修正算子）]
\label{def:tex27-ch28-d5}
给定初始系统 $L^{(2)}:=L$，构造修补序列
\[
  L^{(2)}\to L^{(3)}\to \cdots \to L^{(n)}\to\cdots,
\]
其中每一步引入更高阶修正算子 $l_n$（$n\ge 3$），以吸收“修补本身引入的更高阶一致性缺口”。
\end{definition}

\begin{definition}[D6（超限正交子丛：避免回流污染）]
\label{def:tex27-ch28-d6}
设修补自由度空间为 $\mathcal H$，允许取（可能超限长的）正交分解
\[
  \mathcal H=\bigoplus_{\alpha<\kappa}\mathcal H_\alpha,
  \qquad
  \mathcal H_\alpha\perp\mathcal H_\beta\ (\alpha\neq\beta),
\]
并约定第 $\alpha$ 步仅在新分量 $\mathcal H_\alpha$ 上注入修补项，从而不回流污染已修补分量。
\end{definition}

\begin{theorem}[定理 T3（OHU 超限闭合：极限对象上 \texorpdfstring{$\JacGap\to 0$}{JacGap->0}）]
\label{thm:tex27-ch28-t3}
在定义~\ref{def:tex27-ch28-d5}--\ref{def:tex27-ch28-d6} 的口径与本文可容许性假设下，
存在一个（可能超限长的）OHU 修补过程及其极限对象 $L^{\mathrm{OHU}}$，使得
\[
  \JacGap(L^{\mathrm{OHU}})=0
  \quad(\text{同伦意义/所选口径下}).
\]
\end{theorem}

\begin{Certificate}[证书 T3：定理 A7 的直接调用]
\label{cert:tex27-ch28-t3}
证书为本文第~\ref{sec:tex27-ch24-jacobiator-bianchi-ohu}~节的定理~\ref{thm:tex27-ch24-a7} 及其“超限归纳 + 极限完备化 + 不回流正交注入”的证书链。
\end{Certificate}

\begin{proof}[证明（谓词因果链）]
由证书~\ref{cert:tex27-ch28-t3}，定理~\ref{thm:tex27-ch24-a7} 断言并构造出一个超限正交修补过程，使缺口证书在极限对象上为 $0$。
将该过程记为 $L\rightsquigarrow L^{\mathrm{OHU}}$ 即得定理结论。证毕。
\end{proof}

\begin{corollary}[推论 C2（从“分类”到“修补”：可检验分界线）]
\label{cor:tex27-ch28-c2}
传统路线常把“缺口”停留为结构标签（输出静态分类信息，例如“$\JacGap\neq 0$”或“上同调非平凡”）；
而在本文的证书接口口径下，缺口被写成可计算证书并可被边缘算子与 OHU 过程驱动压低：
\[
  \boxed{\begin{aligned}
    &\text{静态：记录 }\JacGap\neq 0\\
    &\Longrightarrow\ \text{动态：构造 }\mathsf E,\ l_3,l_4,\dots\ \text{使 }\JacGap\to 0
  \end{aligned}}.
\]
\end{corollary}

\begin{Certificate}[证书 C2：诊断--手术--闭合三段链的合取]
\label{cert:tex27-ch28-c2}
证书由三段构造合取给出：
\begin{enumerate}
  \item 诊断：命题~\ref{prop:tex27-ch28-p1} 将几何缺口压缩为证书 $\JacGap$；
  \item 手术：定理~\ref{thm:tex27-ch28-t1}--\ref{thm:tex27-ch28-t2} 给出在正交子丛上的可执行修补与其最优性；
  \item 闭合：定理~\ref{thm:tex27-ch28-t3} 给出 OHU 超限递归使证书在极限对象上归零。
\end{enumerate}
\end{Certificate}

\begin{proof}[证明（谓词因果链）]
由证书~\ref{cert:tex27-ch28-c2}：
诊断层给出缺口的证书化口径，手术层给出可执行且可最优化的局部修补，闭合层给出对“修补引入的更高阶缺口”的级联吸收。
合并即得“分类 $\Rightarrow$ 可修补”的严格分界陈述。证毕。
\end{proof}

\subsection{备注（把容易误读处先钉死）}
\label{subsec:tex27-ch28-remarks}

\begin{remark}[关于“空洞 $\Rightarrow$ holonomy 非平凡”的边界]
严格说，“拓扑空洞”并不必然推出“holonomy 非平凡”：
空洞提供的是\textbf{允许出现} $[A]\neq 0$ 的全局解的可能性；
而 $\JacGap$（定义~\ref{def:tex27-ch28-d2}）是对该类全局障碍在所选回路基下的证书化测量。
\end{remark}

\begin{remark}[关于“比安基缺口”的严格口径]
在标准李代数规范场里，比安基恒等式通常作为结构恒等式恒成立；
若要出现“比安基缺口”，需要采用引理~\ref{lem:tex27-ch28-l1} 那种更一般口径：
\textbf{当括号不满足雅可比时，比安基残差出现由雅可比子驱动的缺口项}。
因此，“几何空洞”与“比安基缺口”并非同一事物；本文仅在“可审计证书接口”口径下把二者对齐到同一类缺口测量与修复叙事中。
\end{remark}

\begin{remark}[把“真理/秩序恢复”等自然语言落入形式系统的前提]
若要把“主动修复真理/世界秩序恢复”等自然语言进入形式系统，
需先把“真理/秩序”明确为可测约束集合与证书（例如 $\JacGap$、风险泛函、规则一致性证书等）；
否则属于\textbf{未证明断言}。
在本文当前口径下，最硬、可核验的结论是：存在一条可审计的构造链将缺口证书压到 $0$（或压到工程阈值以内）。
\end{remark}

\clearpage

\begin{thebibliography}{99}

\bibitem{Tex20SemanticFunctorGaugeField}
\GFramework 工作笔记：\emph{语义函子场与语义规范场：词包判定--规范化接口（形式系统版）}
（\code{NoteBook/notebook1/tex20_pub/gframework_semantic_functor_field_gauge_field_normalization_interface_formal_system.
  tex}）。

\bibitem{GaoZheng2025GFramework}
Gao, Z. (2025).
\newblock Meta-Mathematical Theory based on Pan-Logic Analysis and
Pan-Iterative Analysis (GaoZheng G-Framework) and the
Principal-Bundle-Based Generalized Noncommutative Lie Algebra
(GaoZheng G-Algebra).
\newblock In \emph{Meta-Mathematical Theory based on Pan-Logic Analysis and
Pan-Iterative Analysis (GaoZheng G-Framework) and the
Principal-Bundle-Based Generalized Noncommutative Lie Algebra
(GaoZheng G-Algebra): An Integrated Construction (v1.0)}.
\newblock Zenodo.
\newblock \url{https://doi.org/10.5281/zenodo.17651584}.


\bibitem{PlotkinPowerdomain}
G.\ D.\ Plotkin,
\emph{A Powerdomain Construction},
SIAM Journal on Computing, 5(3):452--487, 1976.
\\ \url{https://doi.org/10.1137/0205035}

\bibitem{SmythPowerdomain}
M.\ B.\ Smyth,
\emph{Power domains},
Journal of Computer and System Sciences, 16(1):23--36, 1978.
\\ \url{https://doi.org/10.1016/0022-0000(78)90048-X}

\bibitem{Wilson1974}
K.\ G.\ Wilson,
\emph{Confinement of quarks},
Physical Review D, 10(8):2445--2459, 1974.
\\ \url{https://doi.org/10.1103/PhysRevD.10.2445}

\bibitem{HatcherAT}
Allen Hatcher,
\emph{Algebraic Topology},
Cambridge University Press, 2002.
\\ \url{https://pi.math.cornell.edu/~hatcher/AT/AT.pdf}

\end{thebibliography}

\end{CJK*}
\end{document}
